% !TeX program = xelatex
% 数学 LaTeX 排版 Skill v4.0.0；版式保持 v3 金标准. XeLaTeX 编译至少两次. 不分发字体.
% WZ-LOCKED-CLASS-BEGIN
\begin{filecontents*}[overwrite]{elegantbook-original-adapter.cls}
\NeedsTeXFormat{LaTeX2e}
\ProvidesClass{elegantbook-original-adapter}[2026/09/23 v3.0 fixed mathematical book environments]
\LoadClass[10pt,letterpaper,oneside]{book}
\RequirePackage[UTF8,scheme=plain,fontset=fandol]{ctex}
\RequirePackage{fontspec}
\setmainfont{cmunrm.otf}[BoldFont=cmunbx.otf,ItalicFont=cmunti.otf,BoldItalicFont=cmunbi.otf]
\setCJKfamilyfont{sourceheader}[ItalicFont=FandolKai-Regular.otf]{FandolKai-Regular.otf}
\RequirePackage{amsmath,amssymb,mathrsfs}
\RequirePackage{amsthm,etoolbox}
\RequirePackage{xcolor,geometry,setspace,fancyhdr,enumitem,needspace,xurl}
\definecolor{structurecolor}{HTML}{880E4F}
\definecolor{main}{HTML}{9C27B0}
\geometry{paperwidth=612bp,paperheight=792bp,left=20mm,right=20mm,top=95.04bp,bottom=20mm,headheight=12pt,headsep=25pt,footskip=30pt}
\setstretch{1.6}
\setlength{\parindent}{2em}
\setlength{\parskip}{0pt}
\setlist{nosep,leftmargin=2em}
\AtBeginDocument{\setlength{\abovedisplayskip}{4pt plus 1pt minus 1pt}\setlength{\belowdisplayskip}{4pt plus 1pt minus 1pt}\setlength{\abovedisplayshortskip}{2pt}\setlength{\belowdisplayshortskip}{3pt}}
\fancyhf{}
\fancyhead[L]{{\itshape\CJKfamily{sourceheader}\rightmark}}
\fancyhead[R]{{\itshape\CJKfamily{sourceheader}\leftmark}}
\fancyfoot[C]{\thepage}
\renewcommand{\headrulewidth}{0.4pt}
\renewcommand{\headrule}{{\color{structurecolor}\hrule height\headrulewidth width\headwidth\vskip-\headrulewidth}}
\pagestyle{fancy}
\fancypagestyle{cover}{\fancyhf{}\fancyfoot[C]{\thepage}\renewcommand{\headrulewidth}{0pt}\renewcommand{\headrule}{}}
\RequirePackage{hyperref}
\hypersetup{unicode=true,colorlinks=true,linkcolor=structurecolor,urlcolor=structurecolor,citecolor=structurecolor,linktoc=section,bookmarksnumbered=true,bookmarksopen=true,pdfborder={0 0 0},pdfstartview=FitH}
\setcounter{tocdepth}{1}
\renewcommand*\l@section{\@dottedtocline{1}{1.5em}{2.4em}}
\renewcommand{\thesection}{\thechapter.\arabic{section}}
\newcommand{\originalchapter}[1]{\par\addvspace{14pt}\Needspace{5\baselineskip}\refstepcounter{chapter}\setcounter{section}{0}\addcontentsline{toc}{chapter}{\protect\numberline{\thechapter}#1}\markboth{\thechapter\quad #1}{}\begingroup\centering\color{structurecolor}\bfseries\fontsize{14.4}{20}\selectfont\thechapter\quad #1\par\endgroup\nobreak\vspace{8pt}}
\newcommand{\originalsection}[1]{\par\addvspace{9pt}\Needspace{4\baselineskip}\refstepcounter{section}\addcontentsline{toc}{section}{\protect\hspace*{1.5em}\protect\numberline{\protect\textcolor{structurecolor}{\thesection}}#1}\markright{\thesection\quad #1}\noindent{\fontsize{12}{17}\selectfont\color{structurecolor}\thesection\quad}{\bfseries\fontsize{12}{17}\selectfont #1}\par\nobreak\vspace{4pt}}

% Semantic environments are registered once here, not re-designed in manuscripts.
\newtheoremstyle{wzstatement}{6pt}{5pt}
  {\normalfont\kaishu}{0pt}{\normalfont\bfseries\color{structurecolor}}
  {}{.7em}{\thmname{#1}\thmnumber{ #2}\thmnote{\normalfont\ (#3)}}
\newtheoremstyle{wzdefinition}{6pt}{5pt}
  {\normalfont\songti}{0pt}{\normalfont\bfseries\color{structurecolor}}
  {}{.7em}{\thmname{#1}\thmnumber{ #2}\thmnote{\normalfont\ (#3)}}
\newtheoremstyle{wzproblem}{7pt}{5pt}
  {\normalfont\songti}{2em}{\normalfont\bfseries\color{main}}
  {}{.7em}{\thmname{#1}\thmnumber{ #2}}
\newtheoremstyle{wzremark}{4pt}{4pt}
  {\normalfont\songti}{0pt}{\normalfont\bfseries}
  {}{.7em}{\thmname{#1}}
\theoremstyle{wzstatement}
\newtheorem{theorem}{定理}[chapter]
\newtheorem{lemma}[theorem]{引理}
\newtheorem{proposition}[theorem]{命题}
\newtheorem{corollary}[theorem]{推论}
\theoremstyle{wzdefinition}
\newtheorem{definition}[theorem]{定义}
\theoremstyle{wzproblem}
\newtheorem{example}{例题}[chapter]
\newtheorem{exercise}{习题}[chapter]
\theoremstyle{wzremark}
\newtheorem*{remark}{注}
\renewcommand{\proofname}{证明}
% Retain the native amsthm QED stack; only adjust the fixed typography.
\expandafter\patchcmd\csname\string\proof\endcsname{\normalfont \topsep6\p@\@plus6\p@\relax}
  {\normalfont\songti \topsep5\p@\@plus1\p@\relax}
  {}{\ClassError{elegantbook-original-adapter}{Cannot patch proof spacing}{Check the amsthm version.}}
\expandafter\patchcmd\csname\string\proof\endcsname{\itshape}{\normalfont\bfseries}
  {}{\ClassError{elegantbook-original-adapter}{Cannot patch proof heading}{Check the amsthm version.}}
\expandafter\patchcmd\csname\string\proof\endcsname{\ignorespaces}
  {\setlength{\parindent}{2em}\ignorespaces}
  {}{\ClassError{elegantbook-original-adapter}{Cannot patch proof paragraphs}{Check the amsthm version.}}
\AtBeginEnvironment{proof}{\par\Needspace{3\baselineskip}}
\renewcommand{\qedsymbol}{\ensuremath{\square}}
\newenvironment{solution}{\begin{proof}[解答]}{\end{proof}}
\newcommand{\wzcheckchapter}{\ifnum\value{chapter}<1
  \ClassError{elegantbook-original-adapter}{Numbered mathematics before chapter 1}{Move it after the first originalchapter.}\fi}

\newcommand{\wzregisterguard}[1]{\AtBeginEnvironment{#1}{\wzcheckchapter\par\Needspace{3\baselineskip}}}
\forcsvlist{\wzregisterguard}{theorem,lemma,proposition,corollary,definition,example,exercise}
% Front matter and bibliography are not numbered mathematical chapters.
\newcommand{\originalpreface}{\par\addvspace{10pt}\Needspace{5\baselineskip}
  \noindent{\bfseries\fontsize{12}{17}\selectfont 前言}\par\nobreak\vspace{4pt}}
\newcommand{\originalcontents}{\par\addvspace{14pt}
  \begin{center}{\color{structurecolor}\bfseries\fontsize{14.4}{20}\selectfont 目录}\end{center}
  \vspace{-10pt}\@starttoc{toc}}
\renewcommand{\bibname}{参考文献}
\renewenvironment{thebibliography}[1]{
  \par\addvspace{12pt}\Needspace{3\baselineskip}\phantomsection
  \addcontentsline{toc}{chapter}{\bibname}
  \markboth{\bibname}{}
  \noindent{\color{structurecolor}\bfseries\fontsize{14.4}{20}\selectfont\bibname}\par\nobreak\vspace{6pt}
  \list{\@biblabel{\@arabic\c@enumiv}}
    {\settowidth\labelwidth{\@biblabel{#1}}\leftmargin\labelwidth
     \advance\leftmargin\labelsep\usecounter{enumiv}
     \let\p@enumiv\@empty\renewcommand\theenumiv{\@arabic\c@enumiv}
     \setlength{\itemsep}{3pt}\setlength{\parsep}{0pt}}
  \sloppy\clubpenalty4000\widowpenalty4000\sfcode`\.=1000\relax
}{\def\@noitemerr{\@latex@warning{Empty bibliography}}\endlist}

\numberwithin{equation}{chapter}
\renewcommand{\theequation}{\thechapter.\arabic{equation}}
\allowdisplaybreaks[2]
\emergencystretch=2em
\clubpenalty=6000
\widowpenalty=6000
\raggedbottom
\endinput
\end{filecontents*}
% WZ-LOCKED-CLASS-END
\documentclass{elegantbook-original-adapter}
% ===== 元数据内容槽 =====
\newcommand{\BookTitle}{度量空间、拓扑与几何基础}
\newcommand{\BookAuthor}{}
\newcommand{\BookDate}{}
\newcommand{\PDFTitle}{度量空间、拓扑与几何基础: MIT OCW 中文修订版}
\newcommand{\PDFAuthor}{}
\hypersetup{pdftitle={\PDFTitle},pdfauthor={\PDFAuthor}}
% 数学记号可以增加；禁止重定义版式、环境或证毕符号.
\usepackage{tikz}
\newcommand{\dd}{\,\mathrm d}
\newcommand{\R}{\mathbb R}
\newcommand{\norm}[1]{\lVert #1\rVert}
\newcommand{\inner}[2]{\langle #1,#2\rangle}
\begin{document}
\frontmatter
\begin{center}
{\bfseries\BookTitle}\par
MIT OpenCourseWare 中文重编与学习补充 · 修订版
\end{center}
\originalpreface
本书面向已学实数极限、连续函数、一元微积分、基本集合与线性代数的读者. 第1--6章从距离建立开闭集、紧致性、完备性和不动点应用; 第7--8章保留基础拓扑及函数空间的应用; 第9--15章把一般拓扑、基本群与覆盖、组合曲面和流形串成进一步的主线. 群与自由群的必要定义在正文给出, 不预设代数拓扑知识. 第8章Dirichlet动机及第15章微分部分另需基本多元微积分和链式法则.

本修订保留原《度量空间与基础拓扑》的八章、例题和37道完整练习解答, 新增七章与21道完整练习解答; 原第8章只更新流形可度量化的回指, 数学内容不删减. 主源仍为Paige Bright的MIT OCW 18.S190, January IAP 2023, 六讲43个PDF物理页(正文37页), 六份官方TeX及三份习题. 本书按概念依赖重编, 不是逐字全译; 原稿略证、留给读者的论证和条件缺漏均补齐, 修改与补充在各章说明. 练习解答由编者撰写, 不冒称OCW官方答案.

新增部分的直接来源包括James Munkres的18.901 Fall 2004补充讲义, Haynes Miller的18.905 Fall 2016同伦与胞腔讲义(课堂记录Sanath Devalapurkar, 插图Xianglong Ni), Paul Seidel的18.900 Spring 2023组合曲面材料及18.950 Fall 2008长度材料, Tomasz S. Mrowka的18.965 Fall 2004基础坐标讲义. 18.901主体使用商业教材, 不存在本项目已归档的全课公开讲义; 18.905从同调起步, 把基本群与覆盖作为先修. 本书这两部分的完整证明为编者补充, 教学路线参照Andrew Snowden的18.904 Spring 2011公开Lecture Summaries, 不将摘要冒充完整原讲义.

\begin{center}
\begin{tabular}{p{.09\linewidth}p{.38\linewidth}p{.41\linewidth}}
章 & 内容 & 来源及角色\\
1--2 & 度量、序列与连续性 & 18.S190第1--2讲, 证明补全\\
3--4 & 欧氏及一般度量紧致性 & 18.S190第3--4讲, 等价判据\\
5--6 & 完备化与不动点应用 & 18.S190第5讲, Picard补充\\
7--8 & 基础拓扑及函数空间 & 18.S190第6讲, 18.102及增补\\
9--10 & 分离、积与商、紧化 & 18.901安排与Notes D、G, 编者证明\\
11--13 & 同伦、基本群、覆盖与拼接 & 18.905第5、14讲; 18.904路线及增补\\
14 & 组合曲面、Betti数与角亏 & 18.900第29--31、33、40讲; 18.905第18讲\\
15 & 流形、切空间与内在距离 & 18.965第1、2、4讲; 18.950第4章选段
\end{tabular}
\end{center}

本册覆盖本科几何拓扑的基础主线, 不声称覆盖所有拓扑学. 曲面部分证明有限组合模型、显式连通和族和离散Gauss--Bonnet, 不含任意拓扑曲面的三角剖分及完整分类证明; 同调部分限于有限边界矩阵和Euler--Poincar\'e恒等式, 不含完整奇异同调、上同调与对偶. 光滑曲率、测地线全局理论、Hopf--Rinow、Sard、Whitney及Morse理论未收入. 18.950的完整几何范围适合另成一册. 原八章中的分析拓展保留为学习应用, 并非据此判断每项都为考试必考内容.

原版考纲说明仍保留其证据边界: 已核清华大学2027年新领军招生简章及详细考纲发布通知, 但此前未重新获得该通知所链考纲PDF原件; 本次扩充根据MIT官方可下载材料, 不把旧考纲摘要提升为重新逐页核验的证据, 也不声称覆盖整份考纲.

MIT OCW对应材料及本书中文译编、重排、勘误和增补采用CC BY-NC-SA 4.0. 原作者、课程学期、真实文件链接、物理页数及第三方署名见参考文献与\texttt{sources/ATTRIBUTION.md}. 商业教材未复制; 第三方“Used with permission”须结合OCW条款与具体署名判断, 不等于“All rights reserved”. MIT和原作者不为本项目背书. 固定版式沿用用户提供的\texttt{math-latex-typesetting} v4.0.0真实模板, 不分发字体.

\clearpage
% 原课程信息. 在 \frontmatter 内输入, 不增加数学章号或数学环境编号.
\par\addvspace{14pt}\Needspace{5\baselineskip}
\phantomsection\addcontentsline{toc}{chapter}{原课程信息与学习安排}
\markboth{原课程信息与学习安排}{}
\noindent{\color{structurecolor}\bfseries\fontsize{14.4}{20}\selectfont 原课程信息与学习安排}\par\nobreak\vspace{6pt}

以下分别保留本书所用八门 MIT OpenCourseWare 课程的原始身份和教学安排. 18.S190 是度量空间部分的主源; 18.901 支持一般拓扑主线; 18.900 支持组合曲面部分. 其余五门课程用于补充路线或选读材料. 原课先修与评分属于各自学期, 不构成本书的统一先修要求、考纲或评分制度. 原课讲次与本书章次也不相同.

课程信息据官方主页、Syllabus、Calendar 及相关材料页核对. 有些页面只列讲次和作业节点, 没有实际日期或星期; 以下按原有粒度转述. 对未在所用公开页面列出的信息明确说明, 不依据每周课时数倒推课表. 本书的证明补全、中文解答与编者练习另见正文和来源说明.

\par\addvspace{9pt}\Needspace{4\baselineskip}
\phantomsection\addcontentsline{toc}{section}{18.S190: 度量空间入门}
\noindent{\bfseries\fontsize{12}{17}\selectfont 18.S190: 度量空间入门}\par\nobreak\vspace{4pt}

\textbf{作者、学期与角色.} Introduction to Metric Spaces, January IAP 2023, 授课者 Paige Bright, 本科课程. 本书第1--8章以六讲公开讲义为主要起点, 按概念依赖重编并补全论证. 原课介绍度量、开闭集、连续映射、函数空间、完备性与紧致性, 是短期专题课, 并非一般拓扑的全学期课程.

\textbf{上课安排与先修.} 每周2次讲课, 每次1.5小时. 官方先修为18.100A Real Analysis 或18.100P Real Analysis. 这说明原课面向已接触分析和证明写作的学生; 本书另在前言说明自身读者要求.

\textbf{作业、考试与评分.} 原课采用 Pass/No Record, 依据作业与参与情况评定; 共3份 problem sets, 无考试. 官方没有列出作业与参与的百分比, 不应另配权重.

\textbf{官方 Calendar.} 第1讲为动机、直观与例子; 第2讲为一般理论, 交第1份作业; 第3讲为欧氏空间中的紧集, 交第2份作业; 第4讲为紧度量空间; 第5讲为完备度量空间, 交第3份作业; 第6讲为后续方向. Calendar 位于 Syllabus 内, 只列讲次及上述提交节点, 未列实际日期.

\textbf{依据.} \href{https://ocw.mit.edu/courses/18-s190-introduction-to-metric-spaces-january-iap-2023/}{官方课程主页}; \href{https://ocw.mit.edu/courses/18-s190-introduction-to-metric-spaces-january-iap-2023/pages/syllabus/}{Syllabus 与 Calendar}.

\par\addvspace{9pt}\Needspace{4\baselineskip}
\phantomsection\addcontentsline{toc}{section}{18.901: 一般拓扑主线}
\noindent{\bfseries\fontsize{12}{17}\selectfont 18.901: 一般拓扑主线}\par\nobreak\vspace{4pt}

\textbf{作者、学期与角色.} Introduction to Topology, Fall 2004, 授课者 James Munkres, 本科课程. 原课以逻辑与集合论、一般拓扑和专业数学写作为目标. 教材为 Munkres 的 \emph{Topology}, 第2版; OCW 提供补充讲义和按教材定位的练习, 不应把这些文件称作全课完整公开讲义. 本书第9--10章采用其一般拓扑安排及补充 Notes D、G, 完整教学论证由编者组织和补充.

\textbf{上课安排与先修.} 每周2次讲课, 每次1.5小时. 原课要求第一门分析课的水平达到 Rudin 的 \emph{Principles of Mathematical Analysis}, 在 MIT 对应18.100B; 既要求分析知识, 也要求构造和书写证明的经验. 每周在课堂讨论学生希望提问的习题.

\textbf{周练与正式作业.} 每次讲课配套的教材习题用于训练, 不提交评分; 原课鼓励把全部解答写成笔记, 考试几乎全部以这些习题为基础. 另有4份正式 problem sets, 须独立完成、提交并按数学写作要求评分. PS0 是集合论语言的诊断作业, 第二次课在课堂评阅, 成绩仅用于学习建议; PS5 是学期末的可选作业. 不应把编号0--5的六组材料都计作正式评分作业.

\textbf{评分.} 正式作业共400分, 期中考试100分, 期末考试200分, 总计700分. Syllabus 未把它们写成百分制比例; 本书保留原分值.

\textbf{官方 Calendar: 第1--5周.} 第1周为逻辑与基础(第1次课), 列 PS0 提交节点; 第2周为关系、基数与选择公理(第2--3次课); 第3周为拓扑与闭集(第4--5次课); 第4周为连续函数与任意积(第6--7次课), 在该周第二次课交 PS1; 第5周为度量拓扑(第8--9次课).

\textbf{官方 Calendar: 第6--10周.} 第6周为商拓扑(第10次课); 第7周为连通空间与紧空间(第11--12次课); 第8周继续紧致性(第13--14次课), 在第二次课交 PS2; 第9周先讲良序集与极大原理(第15次课), 第16次课为期中考试; 第10周为可数性与分离公理(第17--18次课).

\textbf{官方 Calendar: 第11周以后.} 第11周为 Urysohn 引理与度量化(第19--20次课), 在第二次课交 PS3; 第12周为 Tietze 定理(第21次课); 第13周为 Tychonoff 定理与 Stone--\v{C}ech 紧化(第22--23次课); 第14周为 Baire 空间与维数理论(第24--25次课), 在第二次课交 PS4; 第15周为嵌入欧氏空间(第26次课). 最后一行列期末考试和可选 PS5 的提交, 未给该行周号或实际日期.

\textbf{依据.} \href{https://ocw.mit.edu/courses/18-901-introduction-to-topology-fall-2004/}{官方课程主页}; \href{https://ocw.mit.edu/courses/18-901-introduction-to-topology-fall-2004/pages/syllabus/}{Syllabus}; \href{https://ocw.mit.edu/courses/18-901-introduction-to-topology-fall-2004/pages/calendar/}{Calendar}; \href{https://ocw.mit.edu/courses/18-901-introduction-to-topology-fall-2004/pages/assignments/}{Assignments}.

\par\addvspace{9pt}\Needspace{4\baselineskip}
\phantomsection\addcontentsline{toc}{section}{18.900: 平面几何与拓扑}
\noindent{\bfseries\fontsize{12}{17}\selectfont 18.900: 平面几何与拓扑}\par\nobreak\vspace{4pt}

\textbf{作者、学期与角色.} Geometry and Topology in the Plane, Spring 2023, 授课者 Paul Seidel, 本科课程. 原课选取易于可视化的几何与拓扑专题, 不要求已有微分几何或拓扑学训练. 本书使用第4讲的绕数相关材料及第29--31、33、40讲的组合曲面材料; 这只是原课的一部分, 不包含其台球、代数曲线或双曲几何等完整专题.

\textbf{上课安排与先修.} 每周3次讲课, 每次50分钟. 官方先修为18.03 Differential Equations 或18.06 Linear Algebra, 并预期熟悉18.02 Multivariable Calculus. 原课解释线性代数在第III章短暂出现, 在第VII章作用更大; 微分方程用于第IX章, 复数用于第14讲及第VIII章.

\textbf{作业、考试与评分.} 理解性问题(comprehension questions, 简称 CQ)占30\%, 较有挑战的正式 problem sets 占30\%, 考试占40\%. CQ 在 OCW 公开, 正式 problem sets 和考试不向 OCW 用户提供. 因而不能从网页的 Problem Sets 资源标签推断正式评分题或试卷已公开, 本书所写解答也不冒称这些未公开材料的官方答案.

\textbf{公开讲义顺序与实际授课.} 公开目录有9章、40讲: 第1--5讲为多边形, 第6--8讲为台球, 第9--11讲为主频, 第12--15讲为光滑闭路, 第16--18讲为浸入闭路, 第19--27讲为代数曲线, 第28--33讲为二维复形, 第34--36讲为双曲几何, 第37--40讲为弯曲几何. 这是可用材料的顺序; Syllabus 明确说明2023年春实际略过第20、24--25、38和40讲. 本书第14章采用第40讲 \emph{Geometry of Combinatorial Surfaces} 作为 OCW 公开选读, 不把它标为该学期实际讲授内容.

\textbf{官方 Calendar.} 独立的 Calendar 及带实际日期的授课、作业或考试时间表未在所用公开页面列出. 上述40讲目录与实际略讲说明可以确定材料范围, 不能据此合成逐周课表.

\textbf{依据.} \href{https://ocw.mit.edu/courses/18-900-geometry-and-topology-in-the-plane-spring-2023/}{官方课程主页与讲义目录}; \href{https://ocw.mit.edu/courses/18-900-geometry-and-topology-in-the-plane-spring-2023/pages/syllabus/}{Syllabus}.

\par\addvspace{9pt}\Needspace{4\baselineskip}
\phantomsection\addcontentsline{toc}{section}{18.904: 拓扑研讨课}
\noindent{\bfseries\fontsize{12}{17}\selectfont 18.904: 拓扑研讨课}\par\nobreak\vspace{4pt}

\textbf{作者、学期与角色.} Seminar in Topology, Spring 2011, 授课者 Andrew Snowden, 本科课程. 学生每次课分两人讲授, 每人约25分钟; 目标同时包括基本群、同调、上同调和数学报告能力. 本书第11--13章参照公开 Lecture Summaries 的基本群与覆盖路线, 摘要不等同于完整原讲义.

\textbf{上课安排、先修与评分.} 每周3次, 每次1小时; 先修18.901 Introduction to Topology. 讲授与参与占60\%, 期末论文占30\%, 作业占10\%; 无考试. Syllabus 预计约4份正式作业, 讲者提供的练习集可选且不提交. 须出勤, 每缺3次课降低一个字母等级; 期末论文至少10页, 用 LaTeX 编写.

\textbf{官方 Calendar.} 第1--14次课为基本群与覆盖, PS1 于第9次课提交; 第15--25次课为同调, 第15次课定论文题目, PS2 于第20次课提交; 第26--34次课为上同调与 Poincar\'e 对偶, 第27次课交论文初稿、第32次课交终稿; 第35--39次课为期末论文报告, PS3 于第38次课提交. Calendar 只明确列出这3份作业节点, 不应为预计的第4份补造时间.

\textbf{依据.} \href{https://ocw.mit.edu/courses/18-904-seminar-in-topology-spring-2011/}{官方课程主页}; \href{https://ocw.mit.edu/courses/18-904-seminar-in-topology-spring-2011/pages/syllabus/}{Syllabus 与 Calendar}.

\par\addvspace{9pt}\Needspace{4\baselineskip}
\phantomsection\addcontentsline{toc}{section}{18.905: 代数拓扑选读}
\noindent{\bfseries\fontsize{12}{17}\selectfont 18.905: 代数拓扑选读}\par\nobreak\vspace{4pt}

\textbf{作者、学期与角色.} Algebraic Topology I, Fall 2016, 授课者 Haynes Miller, 研究生课程. 官方说明公开讲义依据 Sanath Devalapurkar 的课堂 LaTeX 记录, 插图由 Xianglong Ni 绘制. 本书选用第5讲同伦、第14讲 CW 复形及第18讲 Euler 示性数材料. 原课从奇异同调起步, 不能当作从零讲授基本群与覆盖的课程.

\textbf{上课安排、先修与评分.} 每周2次, 每次1小时; 先修为18.701 Algebra I 或18.703 Modern Algebra, 加18.901 Introduction to Topology. 须熟悉拓扑空间、覆盖和基本群, 以及主理想整环上有限生成模的结构. 共6份作业, 占75\%; 期末考试周有40分钟口试, 占25\%.

\textbf{官方 Calendar.} 第1--13讲为奇异同调, 第14--25讲为计算方法, 第26--38讲为上同调与对偶. PS1--6 分别于第7、12、16、22、28、35讲提交; 口试安排在期末考试周. 页面按讲次列出节点, 未列实际日期.

\textbf{依据.} \href{https://ocw.mit.edu/courses/18-905-algebraic-topology-i-fall-2016/}{官方课程主页}; \href{https://ocw.mit.edu/courses/18-905-algebraic-topology-i-fall-2016/pages/syllabus/}{Syllabus}; \href{https://ocw.mit.edu/courses/18-905-algebraic-topology-i-fall-2016/pages/calendar/}{Calendar}; \href{https://ocw.mit.edu/courses/18-905-algebraic-topology-i-fall-2016/pages/lecture-notes/}{讲义署名与目录}.

\par\addvspace{9pt}\Needspace{4\baselineskip}
\phantomsection\addcontentsline{toc}{section}{18.950: 微分几何选读}
\noindent{\bfseries\fontsize{12}{17}\selectfont 18.950: 微分几何选读}\par\nobreak\vspace{4pt}

\textbf{作者、学期与角色.} Differential Geometry, Fall 2008, 授课者 Paul Seidel, 本科课程. 原课围绕曲率, 讨论平面曲线、超曲面的局部及整体几何、长度与距离. 本书第15章只选取公开第4章的长度与距离材料, 不据此宣称完整覆盖18.950.

\textbf{上课安排与先修.} 每周上课次数和每次时长未在所用公开页面列出. 官方先修为18.100 Analysis I, 加18.06或18.700 Linear Algebra 或18.701 Algebra I 中的一门; 课程说明同时要求良好的多元微积分和线性代数基础, 并适应定义、定理与证明的叙述.

\textbf{评分与公开作业.} 期中考试占30\%, 作业占30\%, 期末考试占40\%. Assignments 页公开 Homework 1--10. 十份作业的中文题面、独立解答及覆盖记录见本项目独立分册\href{https://github.com/xhc144/mit-ocw-notes/blob/80b9182650861c1a4830fb812032abbc93a429e7/subjects/differential-geometry/differential-geometry.pdf}{《曲线与曲面的微分几何》PDF}, 可下载其\href{https://github.com/xhc144/mit-ocw-notes/blob/80b9182650861c1a4830fb812032abbc93a429e7/subjects/differential-geometry/differential-geometry-source.zip}{源码 ZIP}. 该分册说明32个顶层题中30题的全部要求已有独立解答, PS3.3末尾引用的 Proposition 6.3 与 PS9.4引用的 Lemma 28.3 尚有定位缺口. 本册不重复这些作业解答.

\textbf{官方 Calendar.} 独立 Calendar 及作业、考试实际日期未在所用公开页面列出. Lecture Notes 目录仅给讲次分段: 第1--10讲为平面曲线的局部与整体几何, 第11--23讲为超曲面局部几何, 第24--35讲为超曲面整体几何, 第36--41讲为长度与距离. 这不是带日期的周课表.

\textbf{依据.} \href{https://ocw.mit.edu/courses/18-950-differential-geometry-fall-2008/}{官方课程主页}; \href{https://ocw.mit.edu/courses/18-950-differential-geometry-fall-2008/pages/syllabus/}{Syllabus}; \href{https://ocw.mit.edu/courses/18-950-differential-geometry-fall-2008/pages/lecture-notes/}{讲义目录}; \href{https://ocw.mit.edu/courses/18-950-differential-geometry-fall-2008/pages/assignments/}{10份作业}; \href{https://github.com/xhc144/mit-ocw-notes/blob/80b9182650861c1a4830fb812032abbc93a429e7/subjects/differential-geometry/README.md}{微分几何分册说明}.

\par\addvspace{9pt}\Needspace{4\baselineskip}
\phantomsection\addcontentsline{toc}{section}{18.965: 流形几何选读}
\noindent{\bfseries\fontsize{12}{17}\selectfont 18.965: 流形几何选读}\par\nobreak\vspace{4pt}

\textbf{作者、学期与角色.} Geometry of Manifolds, Fall 2004, 授课者 Tomasz Mrowka, 研究生课程. 本书第15章选用第1、2、4讲的基础坐标材料, 不包含原课的全部微分拓扑理论.

\textbf{上课安排、先修与评分.} 每周2次, 每次1.5小时; 先修18.101 Analysis II 与18.905 Algebraic Topology. 成绩100\%依据作业. 所用 Syllabus 未另列考试权重.

\textbf{官方 Calendar.} 共列35讲, 从流形、光滑映射、导数和逆函数及隐函数定理出发, 进入向量丛、嵌入、Sard 定理、纤维丛、Fredholm 算子与横截性. 第23--28讲为 Morse 理论; 第29--30讲涉及 Lie 导数和 Frobenius 可积性; 第31--34讲为微分形式及 de Rham 定理相关内容; 第35讲为 Smale 浸入定理. Calendar 未列作业提交或考试日期.

\textbf{依据.} \href{https://ocw.mit.edu/courses/18-965-geometry-of-manifolds-fall-2004/}{官方课程主页}; \href{https://ocw.mit.edu/courses/18-965-geometry-of-manifolds-fall-2004/pages/syllabus/}{Syllabus}; \href{https://ocw.mit.edu/courses/18-965-geometry-of-manifolds-fall-2004/pages/calendar/}{Calendar}.

\par\addvspace{9pt}\Needspace{4\baselineskip}
\phantomsection\addcontentsline{toc}{section}{18.102: 泛函分析补充}
\noindent{\bfseries\fontsize{12}{17}\selectfont 18.102: 泛函分析补充}\par\nobreak\vspace{4pt}

\textbf{作者、学期与角色.} Introduction to Functional Analysis, Spring 2021, 授课者 Casey Rodriguez, 本科课程. 官方 Lecture Notes and Readings 页区分 Richard Melrose 的2020年讲义与 Andrew Lin 在 Rodriguez 2021年课堂所作的记录; 本书第8章使用后者第3讲中的 Baire 定理与一致有界原理. 该补充支持分析应用, 不代表本册包含原课的测度论、Hilbert 空间或谱定理全程.

\textbf{上课安排与先修.} 每周2次, 每次1.5小时. 先修为18.06、18.700或18.701之一, 加18.100A、18.100B、18.100P或18.100Q之一; 或获教师许可.

\textbf{作业、考试与评分.} 每周作业占50\%, 去掉最低一次作业成绩; 期中考试占25\%, 为第7周的24小时带回考试; 期末 assignment 占25\%, 为课程结束时的48小时作业. 原页将后者称为 Final Assignment, 本书不改称正式期末考试.

\textbf{官方 Calendar.} 共列23讲: 第1--5讲建立 Banach 空间、算子、Baire 及 Hahn--Banach 等基础; 第6--13讲转入测度、Lebesgue 积分与 $L^p$ 空间; 第14--22讲为 Hilbert 空间、紧算子与谱理论; 第23讲为区间上的 Dirichlet 问题. Assignment 1--10 分别于第2、4、6、8、10、14、16、18、20、22讲提交. 期中考试列在第11讲之后、第12讲之前, 最后一行列 Final Assignment 提交; 页面未列实际日期.

\textbf{依据.} \href{https://ocw.mit.edu/courses/18-102-introduction-to-functional-analysis-spring-2021/}{官方课程主页}; \href{https://ocw.mit.edu/courses/18-102-introduction-to-functional-analysis-spring-2021/pages/syllabus/}{Syllabus}; \href{https://ocw.mit.edu/courses/18-102-introduction-to-functional-analysis-spring-2021/pages/calendar/}{Calendar}; \href{https://ocw.mit.edu/courses/18-102-introduction-to-functional-analysis-spring-2021/pages/lecture-notes-and-readings/}{讲义署名与阅读安排}.

\clearpage
\originalcontents
\clearpage
\mainmatter
\originalchapter{度量与典型空间}\label{ch:metric}
\emph{本章以 Paige Bright 的 MIT OCW 18.S190 IAP 2023 第1讲正文第1--7页为主要来源, 重新组织叙述并补全证明; 范数比较, 球面距离的严格解释及本章练习为编者补充.}

在实数上, ``两项越来越接近''可以写成 $|x-y|$ 越来越小. 如果对象换成向量, 可以量线段的长度; 如果对象换成函数, 则需要先决定要比较图像的最大偏差, 还是整个区间上的累计偏差. 度量空间保留这些比较共同遵守的规则, 使收敛与连续的讨论不再依赖坐标或函数值的具体形式.

本章默认读者已经学过实数的完备性, 连续函数在闭区间上的有界性, 最大值定理与一致连续性, 以及 Riemann 积分的基本性质. 这些一元分析结论会在函数空间的例子中使用. 抽象度量空间的完备性则将在后文重新研究.

\originalsection{度量公理}\label{ch:metric:axioms}
\begin{definition}[度量空间]
设 $X$ 是一个集合. 一个函数 $d:X\times X\to[0,\infty)$ 称为 $X$ 上的\textbf{度量}, 如果对任意 $x,y,z\in X$, 都有
\[
 d(x,y)=0\Longleftrightarrow x=y,\qquad
 d(x,y)=d(y,x),\qquad
 d(x,z)\le d(x,y)+d(y,z).
\]
最后一个条件称为\textbf{三角不等式}. 配备了度量的集合 $(X,d)$ 称为\textbf{度量空间}.
\end{definition}
距离取值有限且非负, 两个不同对象之间的距离严格为正. 三角不等式表达的是: 经过中间点走一段路, 不会比两端之间的直接距离更短. 它也提供了本书最常用的估计方法: 如果目标中的两个点不容易直接比较, 就插入一个已知接近两者的中间点.

集合 $X$ 不必具有加法, 数乘或大小关系. 因而 $x+y$ 在一般度量空间中可能没有意义, 但 $d(x,y)$ 总有意义. 这一区别决定了哪些实分析论证能够原样推广, 哪些需要重新表达.

\begin{lemma}[反三角不等式]\label{ch:metric:reverse}
在任意度量空间中, 对 $x,y,z\in X$, 有
\[
 |d(x,z)-d(y,z)|\le d(x,y).
\]
更一般地, 对 $x,x',y,y'\in X$, 有
\[
 |d(x,y)-d(x',y')|\le d(x,x')+d(y,y').
\]
\end{lemma}
\begin{proof}
由 $d(x,z)\le d(x,y)+d(y,z)$ 可得 $d(x,z)-d(y,z)\le d(x,y)$. 交换 $x,y$, 得到相反方向的差也不超过 $d(x,y)$, 合并即得第一式. 对第二式, 插入 $x'$ 与 $y'$, 有
\[
 d(x,y)\le d(x,x')+d(x',y')+d(y',y).
\]
于是 $d(x,y)-d(x',y')\le d(x,x')+d(y,y')$. 再交换带撇与不带撇的点即可.
\end{proof}
第一式意味着, 当一个端点移动很少时, 它到固定点的距离也只会改变很少. 第二式允许两个端点同时移动. 下一章中距离与极限之间的相容性, 都将从这一估计直接得到.

\begin{definition}[球, 子空间与有界集]\label{ch:metric:balls}
设 $(X,d)$ 是度量空间, $x\in X$ 且 $r>0$. 定义\textbf{开球}与\textbf{闭球}
\[
 B_X(x,r)=\{y\in X:d(x,y)<r\},\qquad
 \overline B_X(x,r)=\{y\in X:d(x,y)\le r\}.
\]
不易混淆时省略下标 $X$. 如果 $A\subseteq X$, 将 $d$ 限制在 $A\times A$ 上仍得到度量, 称为\textbf{子空间度量}. 若存在 $p\in X$ 与 $R\ge0$, 使每个 $a\in A$ 都满足 $d(a,p)\le R$, 则称 $A$ \textbf{有界}. 空集约定为有界.
\end{definition}
子空间度量满足公理, 因为对 $A$ 中任意三个点, 原来的度量公理仍然成立. 子空间中的球由
\[
 B_A(x,r)=A\cap B_X(x,r)\qquad(x\in A)
\]
给出. ``球''指距离小于指定半径的所有对象, 它不必长成欧氏几何中的圆或圆球. 此外, 闭球记号中的横线只是名称记号; 在一般度量空间中, 它不一定等于相应开球的闭包, 这一差异将在下一章出现.

有界性的中心可以更换. 如果 $A$ 被中心 $p$, 半径 $R$ 的闭球包含, 那么对另一点 $q\in X$, 三角不等式给出 $d(a,q)\le R+d(p,q)$. 因而有界性不取决于最初选择哪个中心.

\originalsection{范数与诱导度量}\label{ch:metric:norms}
向量空间中, ``两个点相距多远''通常由它们的差向量有多长决定. 这使我们可以把度量的构造转化为长度函数的构造.

\begin{definition}[范数]
设 $V$ 是实向量空间. 函数 $\lVert\cdot\rVert:V\to[0,\infty)$ 称为\textbf{范数}, 如果对任意 $u,v\in V$ 与 $\lambda\in\mathbb{R}$, 满足
\[
 \lVert v\rVert=0\Longleftrightarrow v=0,\qquad
 \lVert\lambda v\rVert=|\lambda|\lVert v\rVert,\qquad
 \lVert u+v\rVert\le\lVert u\rVert+\lVert v\rVert.
\]
配备了范数的向量空间称为\textbf{赋范空间}.
\end{definition}
\begin{proposition}\label{ch:metric:norm-metric}
若 $\lVert\cdot\rVert$ 是 $V$ 上的范数, 则 $d(u,v)=\lVert u-v\rVert$ 是度量.
\end{proposition}
\begin{proof}
正定性由 $u-v=0\Longleftrightarrow u=v$ 得到. 齐次性给出 $\lVert v-u\rVert=\lVert-(u-v)\rVert=\lVert u-v\rVert$, 所以距离对称. 最后,
\[
 \lVert u-w\rVert=\lVert(u-v)+(v-w)\rVert
 \le\lVert u-v\rVert+\lVert v-w\rVert,
\]
正是度量的三角不等式.
\end{proof}

在 $\mathbb{R}^n$ 上, 对 $v=(v_1,\ldots,v_n)$, 最常用的三个范数是
\[
 \lVert v\rVert_1=\sum_{i=1}^n|v_i|,\qquad
 \lVert v\rVert_2=\left(\sum_{i=1}^nv_i^2\right)^{1/2},\qquad
 \lVert v\rVert_\infty=\max_{1\le i\le n}|v_i|.
\]
它们分别产生 $d_1,d_2,d_\infty$. 其中 $d_2$ 是欧氏距离. 对 $v\ne0$, 至少一个坐标不为零, 所以三个范数均严格为正. 齐次性也都直接来自绝对值的齐次性. 三角不等式值得分别核对, 因为它反映了这三种长度对各坐标的不同组合方式.

\begin{lemma}[有限维 Cauchy--Schwarz 不等式]\label{ch:metric:cs}
对 $u,v\in\mathbb{R}^n$, 记 $u\cdot v=\sum_{i=1}^nu_iv_i$. 则
\[
 |u\cdot v|\le\lVert u\rVert_2\lVert v\rVert_2.
\]
\end{lemma}
\begin{proof}
如果 $v=0$, 结论成立. 如果 $v\ne0$, 则对任意 $t\in\mathbb{R}$, 有
\[
 0\le\sum_{i=1}^n(u_i-tv_i)^2
 =\lVert u\rVert_2^2-2t(u\cdot v)+t^2\lVert v\rVert_2^2.
\]
取 $t=(u\cdot v)/\lVert v\rVert_2^2$, 得
\[
 0\le\lVert u\rVert_2^2-
 \frac{(u\cdot v)^2}{\lVert v\rVert_2^2}.
\]
移项并开平方即得所需不等式. 这里选择 $t$ 是为了尽量消去 $u$ 在 $v$ 方向上的分量, 也就是让平方和最小.
\end{proof}

\begin{proposition}\label{ch:metric:three-norms}
$\lVert\cdot\rVert_1,\lVert\cdot\rVert_2,\lVert\cdot\rVert_\infty$ 都是 $\mathbb{R}^n$ 上的范数. 并且对每个 $v\in\mathbb{R}^n$, 有
\[
 \lVert v\rVert_\infty\le\lVert v\rVert_2
 \le\lVert v\rVert_1\le n\lVert v\rVert_\infty,
 \qquad
 \lVert v\rVert_1\le\sqrt n\lVert v\rVert_2.
\]
\end{proposition}
\begin{proof}
$1$-范数的三角不等式来自对每个坐标使用 $|u_i+v_i|\le|u_i|+|v_i|$, 然后求和. 对最大值范数, 每个坐标都满足
\[
 |u_i+v_i|\le|u_i|+|v_i|
 \le\lVert u\rVert_\infty+\lVert v\rVert_\infty.
\]
对 $i$ 取最大值即可. 欧氏范数则使用上一个引理:
\begin{align*}
 \lVert u+v\rVert_2^2
 &=\lVert u\rVert_2^2+2u\cdot v+\lVert v\rVert_2^2\\
 &\le\lVert u\rVert_2^2+2\lVert u\rVert_2\lVert v\rVert_2
       +\lVert v\rVert_2^2
 =(\lVert u\rVert_2+\lVert v\rVert_2)^2.
\end{align*}
两边均非负, 开平方得到三角不等式.

对于比较不等式, 每一项 $|v_i|$ 都不超过 $\lVert v\rVert_2$, 从而 $\lVert v\rVert_\infty\le\lVert v\rVert_2$. 又因为
\[
 \sum_i|v_i|^2\le\left(\sum_i|v_i|\right)^2,
 \qquad
 \sum_i|v_i|\le n\max_i|v_i|,
\]
得到中间两项. 最后一式对向量 $(|v_1|,\ldots,|v_n|)$ 和 $(1,\ldots,1)$ 使用 Cauchy--Schwarz 不等式即可.
\end{proof}

\begin{example}
在 $\mathbb{R}^2$ 中取 $x=(1,-1)$ 与 $y=(4,3)$. 求它们在 $d_1,d_2,d_\infty$ 下的距离, 并描述三种度量下以原点为中心, 半径为 $1$ 的开球.
\end{example}
\begin{solution}
在 $\mathbb{R}^2$ 中, 取 $x=(1,-1)$ 与 $y=(4,3)$. 差向量为 $(3,4)$, 故
\[
 d_1(x,y)=7,\qquad d_2(x,y)=5,\qquad d_\infty(x,y)=4.
\]
以原点为中心, 半径为 $1$ 的开球, 在 $d_2$ 下是 $x_1^2+x_2^2<1$ 的圆盘, 在 $d_1$ 下是 $|x_1|+|x_2|<1$ 的菱形内部, 在 $d_\infty$ 下是 $\max\{|x_1|,|x_2|\}<1$ 的正方形内部. 距离数值和球形状都发生了变化, 但上述比较不等式表明: 任何一种距离足够小, 另外两种也会足够小. 下一章将把这一点精确表述为它们产生相同的开集.
\end{solution}

\begin{remark}
原讲义还提及一般的 $p$-范数
$\lVert v\rVert_p=(\sum_i|v_i|^p)^{1/p}$, 其中 $1\le p<\infty$. 其三角不等式是 Minkowski 不等式, 不能从绝对值三角不等式直接推出. 本书这一阶段只使用已经完整证明的 $p=1,2,\infty$ 三个版本, 不把一般版本作为隐藏的先修知识.
\end{remark}

\originalsection{离散集合与有限空间}\label{ch:metric:discrete}
\begin{example}\label{ch:metric:discrete-example}
对任意集合 $X$, 定义
\[
 d_{\mathrm{dis}}(x,y)=
 \begin{cases}
 0,&x=y,\\
 1,&x\ne y.
 \end{cases}
\]
证明这是度量, 求所有半径的开球, 并判断空间是否有界.
\end{example}
\begin{solution}
正定性和对称性由定义直接得到. 若 $x=z$, 三角不等式左边为零. 若 $x\ne z$, 则 $y$ 不可能同时等于 $x$ 和 $z$, 所以 $d_{\mathrm{dis}}(x,y)+d_{\mathrm{dis}}(y,z)\ge1=d_{\mathrm{dis}}(x,z)$.

在这个度量下, 当 $0<r\le1$ 时, $B(x,r)=\{x\}$; 当 $r>1$ 时, $B(x,r)=X$. 即使 $X$ 有无限多个点, 它仍然有界, 因为所有距离都不超过 $1$. 有界因而不等于点数有限.
\end{solution}

\begin{proposition}[有限度量空间中的最小间隔]\label{ch:metric:finite}
设 $(X,d)$ 是有限度量空间, 且 $X$ 至少有两个点. 则
\[
 \eta=\min\{d(x,y):x,y\in X,\ x\ne y\}>0,
\]
而且对每个 $x\in X$, 都有 $B(x,\eta)=\{x\}$. 若 $X$ 只有一个点, 每个正半径球也都是这个单点集.
\end{proposition}
\begin{proof}
不同点对的距离组成一个非空有限集, 每个数都严格为正, 故其最小值存在且为正. 当 $y\ne x$ 时, $d(x,y)\ge\eta$, 所以 $y\notin B(x,\eta)$; 而 $x$ 总在自己的正半径球中.
\end{proof}
有限空间中不同点的距离未必全相等, 因而其度量未必是离散度量. 但是它总有一种共同性质: 每个点都可以用足够小的球单独分离出来. 这将使有限度量空间与离散空间在开集, 收敛等性质上相同.

\originalsection{连续函数空间的度量}\label{ch:metric:functions}
设 $a<b$ 为实数, $C([a,b])$ 表示区间 $[a,b]$ 上的实值连续函数全体. 两个函数相等, 指它们在区间的每一点都相等. 函数空间中的``点''是一整个函数; 因而函数空间里的球包含的是许多函数, 并不是定义域上的一个小区间.

\begin{proposition}[上确界范数]\label{ch:metric:sup}
对 $f\in C([a,b])$, 定义
\[
 \lVert f\rVert_\infty=\sup_{t\in[a,b]}|f(t)|.
\]
这是一个范数, 从而
\[
 d_\infty(f,g)=\sup_{t\in[a,b]}|f(t)-g(t)|
\]
是 $C([a,b])$ 上的度量.
\end{proposition}
\begin{proof}
连续函数 $|f|$ 在闭有界区间上有最大值, 所以上确界有限. 若范数为零, 则对所有 $t$, $0\le|f(t)|\le\lVert f\rVert_\infty=0$, 从而 $f$ 是零函数. 反之零函数的范数为零. 齐次性来自 $|\lambda f(t)|=|\lambda||f(t)|$. 最后, 对区间中的每个 $t$, 都有
\[
 |f(t)+g(t)|\le|f(t)|+|g(t)|
 \le\lVert f\rVert_\infty+\lVert g\rVert_\infty.
\]
对左边取上确界, 得到范数三角不等式. 再用命题\ref{ch:metric:norm-metric} 即得度量结论.
\end{proof}
证明三角不等式时不必先找最大值点. 只要右边是对每个 $t$ 都成立的同一个上界, 它自然也是左边上确界的上界. 这种``先逐点估计, 再取上确界''的方法, 在后续一致收敛与算子连续性中十分常用.

\begin{example}
固定 $f\in C([a,b])$ 和 $r>0$, 描述上确界度量下 $B(f,r)$ 中的函数. 再在 $[0,1]$ 上取 $f(t)=t$, $g(t)=t^2$, 求 $d_\infty(f,g)$.
\end{example}
\begin{solution}
固定 $f\in C([a,b])$ 和 $r>0$. 条件 $g\in B(f,r)$ 等价于
\[
 \sup_{t\in[a,b]}|g(t)-f(t)|<r.
\]
由于 $|g-f|$ 在闭区间上取到最大值, 这又等价于每个 $t$ 都满足 $f(t)-r<g(t)<f(t)+r$. 可以把它理解为 $g$ 的整个图像落在 $f$ 上下宽度为 $r$ 的带内. 若定义域不是紧区间, 仅仅逐点严格小于 $r$ 未必能保证上确界也严格小于 $r$.

例如在 $[0,1]$ 上取 $f(t)=t$ 与 $g(t)=t^2$, 则
\[
 d_\infty(f,g)=\max_{0\le t\le1}(t-t^2)=\frac14,
\]
因为 $t-t^2=\frac14-(t-\frac12)^2$. 所以 $g\in B(f,r)$ 当且仅当 $r>1/4$.
\end{solution}

最大偏差并不是唯一合理的函数距离. 如果希望比较整个区间上的累计偏差, 积分可以给出另一种度量.

\begin{proposition}[积分距离]\label{ch:metric:l1}
在 $C([a,b])$ 上,
\[
 d_{L^1}(f,g)=\int_a^b|f(t)-g(t)|\,\mathrm{d}t
\]
是度量. 此外,
\[
 d_{L^1}(f,g)\le(b-a)d_\infty(f,g).
\]
\end{proposition}
\begin{proof}
非负性, 对称性和距离有限都由连续函数积分的基本性质得到. 要证明正定性, 不能只说``积分为零则函数为零'', 因为这一断言对任意函数并不成立. 此处连续性是关键.

设 $f(t_0)\ne g(t_0)$, 并记 $c=|f(t_0)-g(t_0)|>0$. 连续性给出 $t_0$ 在 $[a,b]$ 中的一个相对邻域, 在其中 $|f(t)-g(t)|>c/2$. 即使 $t_0$ 是端点, 该邻域也包含一个正长度区间 $J\subseteq[a,b]$. 因而
\[
 \int_a^b|f-g|\,\mathrm{d}t
 \ge\int_J|f-g|\,\mathrm{d}t
 \ge\frac c2\,|J|>0.
\]
所以积分距离为零时, $f=g$ 在每一点都成立. 反方向当然成立. 三角不等式由逐点的 $|f-h|\le|f-g|+|g-h|$ 积分得到. 最后, 对所有 $t$ 有 $|f(t)-g(t)|\le d_\infty(f,g)$, 对这个常数上界积分即得比较式.
\end{proof}

\begin{example}
\label{ch:metric:spikes}在 $[0,1]$ 上定义 $f_n(t)=\max\{1-nt,0\}$. 求 $d_\infty(f_n,0)$ 与 $d_{L^1}(f_n,0)$, 并判断积分距离能否用一个固定常数反向控制上确界距离.
\end{example}
\begin{solution}

在 $[0,1]$ 上定义 $f_n(t)=\max\{1-nt,0\}$. 它是一条高度为 $1$, 底边长为 $1/n$ 的三角形尖峰. 于是
\[
 d_\infty(f_n,0)=1,\qquad
 d_{L^1}(f_n,0)=\int_0^{1/n}(1-nt)\,\mathrm{d}t=\frac1{2n}.
\]
积分距离趋于零, 最大偏差却始终是 $1$. 因而上一个命题中的估计不能在全体连续函数上反向成立: 不存在与函数无关的常数 $C$, 使 $d_\infty(f,g)\le C d_{L^1}(f,g)$ 总成立. 同一个集合上的不同度量, 确实可能表达不同的接近方式.
\end{solution}

\originalsection{C¹空间与微分算子}\label{ch:metric:c1}
记 $C^1([a,b])$ 为连续可微实函数的空间, 其中区间内部有导数, 且导数连续延拓到端点; 端点处也可等价地采用相应的单侧导数. 如果只用上确界距离比较函数值, 便可能遗漏其变化速度.

\begin{proposition}[$C^1$ 范数与微分算子]\label{ch:metric:derivative}
在 $C^1([a,b])$ 上定义
\[
 \lVert f\rVert_{C^1}=\lVert f\rVert_\infty+\lVert f'\rVert_\infty,
 \qquad
 d_{C^1}(f,g)=\lVert f-g\rVert_{C^1}.
\]
这是范数及其诱导度量. 微分算子
\[
 D:C^1([a,b])\to C([a,b]),\qquad Df=f'
\]
满足
\[
 d_\infty(Df,Dg)\le d_{C^1}(f,g).
\]
因此, 对任意 $\varepsilon>0$, 若 $d_{C^1}(f,g)<\varepsilon$, 就有 $d_\infty(Df,Dg)<\varepsilon$; 特别地, 微分算子连续.
\end{proposition}
\begin{proof}
两个上确界均有限. 若 $\lVert f\rVert_{C^1}=0$, 第一项已保证 $f=0$. 齐次性由微分的线性和上确界范数的齐次性得到. 对三角不等式, 分别应用
\[
 \lVert f+g\rVert_\infty\le\lVert f\rVert_\infty+\lVert g\rVert_\infty,
 \qquad
 \lVert f'+g'\rVert_\infty\le\lVert f'\rVert_\infty+\lVert g'\rVert_\infty,
\]
再相加即可. 算子估计则来自
\[
 d_\infty(Df,Dg)=\lVert f'-g'\rVert_\infty
 \le\lVert f-g\rVert_\infty+\lVert f'-g'\rVert_\infty.
\]
连续性的精确定义将在下一章给出; 此处估计已经给出了对任意精度都可使用的控制.
\end{proof}

\begin{example}
在 $C^1([0,1])$ 上只用 $d_\infty$ 比较函数, 并在 $C([0,1])$ 上同样使用 $d_\infty$. 判断微分算子 $D:C^1([0,1])\to C([0,1])$ 是否连续.
\end{example}
\begin{solution}
在 $C^1([0,1])$ 上, 若只用 $d_\infty$ 比较函数, 微分算子便不连续. 取
\[
 f_n(t)=\frac{\sin(nt)}n.
\]
则 $\lVert f_n\rVert_\infty\le1/n\to0$, 但 $f_n'(t)=\cos(nt)$ 且 $\lVert f_n'\rVert_\infty=1$, 因为 $t=0$ 时导数等于 $1$. 这说明连续性需要同时说明定义域与值域各使用什么度量. $C^1$ 范数中加入导数项, 正是为了排除``函数很小但振荡很快''的现象.
\end{solution}

只比较导数的表达式 $\lVert f'-g'\rVert_\infty$ 也不是整个 $C^1([a,b])$ 上的度量. 不同的常数函数会使它等于零, 违反正定性. 若额外固定一个函数值, 例如限制在满足 $f(a)=0$ 的子空间, 这一缺陷就消失了. 章末练习将核对这个区别.

\originalsection{球面上的弦长与测地距离}\label{ch:metric:sphere}
还有一种差异与函数空间无关: 我们可以用环境空间中的直线距离, 也可以用被限制在对象内部的路径长度. 必须先说清楚对象究竟是哪个集合.

\begin{proposition}[单位球面上的两种距离]\label{ch:metric:sphere-example}
设
\[
 S^2=\{x\in\mathbb{R}^3:\lVert x\rVert_2=1\}.
\]
这是单位\textbf{球面}, 不包含球体内部. 欧氏距离限制在 $S^2$ 上给出\textbf{弦长距离} $d_{\mathrm{ch}}(x,y)=\lVert x-y\rVert_2$. 另一种距离为
\[
 d_{\mathrm{geo}}(x,y)=\arccos(x\cdot y)\in[0,\pi],
\]
称为单位球面的\textbf{测地距离}. 它等于连接两点的较短大圆弧的长度; 对互为对跖点的两点, 任意相应半大圆的长度都为 $\pi$.
\end{proposition}
\begin{proof}
先证明角度公式确实满足度量公理. Cauchy--Schwarz 不等式保证 $x\cdot y\in[-1,1]$, 所以公式有定义. 又由
$\lVert x-y\rVert_2^2=2-2x\cdot y$, 得到角度为零当且仅当 $x=y$. 对称性也直接成立.

对三角不等式, 记 $\alpha=d_{\mathrm{geo}}(x,z)$, $\beta=d_{\mathrm{geo}}(z,y)$. 若 $\alpha+\beta\ge\pi$, 则 $d_{\mathrm{geo}}(x,y)\le\pi\le\alpha+\beta$. 若 $\alpha+\beta<\pi$, 将向量分解为沿 $z$ 的分量和与 $z$ 垂直的分量:
\[
 x=(\cos\alpha)z+u,\qquad
 y=(\cos\beta)z+v,
\]
其中 $u\cdot z=v\cdot z=0$, $\lVert u\rVert_2=\sin\alpha$, $\lVert v\rVert_2=\sin\beta$. Cauchy--Schwarz 给出
\[
 x\cdot y=\cos\alpha\cos\beta+u\cdot v
 \ge\cos\alpha\cos\beta-\sin\alpha\sin\beta
 =\cos(\alpha+\beta).
\]
由于 $\arccos$ 在 $[-1,1]$ 上递减, 有 $d_{\mathrm{geo}}(x,y)\le\alpha+\beta$.

为说明``最短路径''确有数学含义, 可把连续路径 $\gamma:[0,1]\to S^2$ 的长度定义为各分割下弦长和的上确界:
\[
 L(\gamma)=\sup_{0=t_0<\cdots<t_m=1}
 \sum_{j=1}^m\lVert\gamma(t_j)-\gamma(t_{j-1})\rVert_2.
\]
对夹角 $\theta\in[0,\pi]$, 相应弦长 $c$ 满足 $c=2\sin(\theta/2)$, 故当 $c\to0$ 时 $\theta/c\to1$. 给定 $\varepsilon>0$, 充分短的弦满足 $\theta\le(1+\varepsilon)c$. 路径的三个坐标函数都在 $[0,1]$ 上连续, 由本章声明的一元分析先修, 它们分别一致连续. 对给定弦长精度, 分别把三个坐标变化控制在该精度的 $1/\sqrt3$ 以下, 再取三个有效参数间隔的最小值, 就得到整条路径关于欧氏距离的一致连续性. 因而可以把参数区间分得足够细, 使每一小段两端的弦都满足这一估计. 用已经证明的角度三角不等式, 得
\[
 d_{\mathrm{geo}}(\gamma(0),\gamma(1))
 \le\sum_jd_{\mathrm{geo}}(\gamma(t_{j-1}),\gamma(t_j))
 \le(1+\varepsilon)L(\gamma).
\]
若路径长度有限, 令 $\varepsilon\downarrow0$ 得角度距离不超过路径长度; 长度无限时结论也成立. 较短大圆弧的长度恰为夹角: 分割弦长和不超过总圆弧长, 而等分后的弦长和 $2m\sin(\theta/(2m))$ 趋于 $\theta$. 所以大圆弧达到上述下界, 测地距离确实是所有球面内连续路径长度的最小值.
\end{proof}

\begin{remark}
原讲义第1讲 Example 18 用了 ``unit ball'' 一词, 但随后描述的是球面上两点之间的球面路径. 本章明确改为 $S^2$. 若对象是实心球 $\{x:\lVert x\rVert_2\le1\}$, 两点间的直线段留在球内, 其内部最短路径长度就是欧氏距离, 不能与球面大圆弧混为一谈.
\end{remark}

例如北极与赤道上的一点, 弦长为 $\sqrt2$, 测地距离为 $\pi/2$; 对跖点的弦长为 $2$, 测地距离为 $\pi$. 两种距离满足
\[
 d_{\mathrm{ch}}(x,y)\le d_{\mathrm{geo}}(x,y)
 \le\frac\pi2 d_{\mathrm{ch}}(x,y).
\]
这里第一式来自 $\sin u\le u$, 第二式来自 $\sin u\ge2u/\pi$ 对 $0\le u\le\pi/2$ 的成立; 后一不等式可由正弦函数在该区间的凹性与连接两端的弦线得到. 因而这两种距离虽有不同数值, 对``趋于某一点''给出的判断相同.

\originalsection{练习与解答}\label{ch:metric:exercises}
以下各题是围绕本章内容设计的编者练习, 不标作 MIT 原习题.

\begin{exercise}\label{ch:metric:ex-bounded}
设 $d$ 是 $X$ 上的度量. 证明
$\rho(x,y)=d(x,y)/(1+d(x,y))$ 也是度量, 且其所有距离都小于 $1$. 对 $0<r<1$, 求 $\rho$-球与 $d$-球之间的关系.
\end{exercise}
\begin{solution}
令 $h(t)=t/(1+t)$. 它在 $[0,\infty)$ 上严格递增, 且只在 $t=0$ 取零. 所以正定性, 对称性成立. 对 $a,b\ge0$, 有
\[
 h(a+b)=\frac{a}{1+a+b}+\frac{b}{1+a+b}
 \le\frac{a}{1+a}+\frac{b}{1+b}=h(a)+h(b).
\]
由三角不等式与 $h$ 的单调性,
$\rho(x,z)\le h(d(x,y)+d(y,z))\le\rho(x,y)+\rho(y,z)$.
另外 $h(t)<1$. 最后,
\[
 \frac{d(x,y)}{1+d(x,y)}<r
 \Longleftrightarrow d(x,y)<\frac r{1-r},
\]
故 $B_\rho(x,r)=B_d(x,r/(1-r))$. 若 $r\ge1$, 则 $B_\rho(x,r)=X$.
\end{solution}

\begin{exercise}\label{ch:metric:ex-c1}
在 $V=\{f\in C^1([a,b]):f(a)=0\}$ 上, 证明 $\lVert f\rVert_D=\lVert f'\rVert_\infty$ 是范数, 并证明
\[
 \lVert f\rVert_D\le\lVert f\rVert_{C^1}
 \le(1+b-a)\lVert f\rVert_D.
\]
说明为何在整个 $C^1([a,b])$ 上第一项不是范数.
\end{exercise}
\begin{solution}
若 $\lVert f'\rVert_\infty=0$, 则 $f'=0$. 由微积分基本定理与 $f(a)=0$,
$f(t)=\int_a^tf'(s)\,\mathrm{d}s=0$. 因而在 $V$ 上正定性成立. 齐次性与三角不等式由导数的线性及上确界范数的相应性质得到. 同一积分表达式给出
\[
 |f(t)|\le\int_a^t|f'(s)|\,\mathrm{d}s
 \le(b-a)\lVert f'\rVert_\infty,
\]
故 $\lVert f\rVert_\infty\le(b-a)\lVert f\rVert_D$, 从而得到第二个估计. 第一个估计来自 $C^1$ 范数包含导数项. 在整个空间中, 非零常数函数的导数为零, 所以正定性失败.
\end{solution}

\begin{exercise}\label{ch:metric:ex-integration}
对 $f\in C([a,b])$, 定义 $Tf(t)=\int_a^tf(s)\,\mathrm{d}s$. 证明 $Tf\in C^1([a,b])$, 且
\[
 \lVert Tf-Tg\rVert_{C^1}\le(1+b-a)\lVert f-g\rVert_\infty.
\]
\end{exercise}
\begin{solution}
微积分基本定理给出 $(Tf)'=f$, 而 $f$ 连续, 所以 $Tf\in C^1([a,b])$. 对每个 $t$,
\[
 |Tf(t)-Tg(t)|
 \le\int_a^t|f(s)-g(s)|\,\mathrm{d}s
 \le(b-a)\lVert f-g\rVert_\infty.
\]
取上确界, 再与 $\lVert(Tf-Tg)'\rVert_\infty=\lVert f-g\rVert_\infty$ 相加, 即得估计. 它说明积分算子不仅控制函数值, 也控制所得函数的一阶导数.
\end{solution}

\begin{exercise}\label{ch:metric:ex-nonmetric}
判断 $\mathbb{R}$ 上的 $d(x,y)=|x-y|^2$ 是否为度量. 再证明对固定 $0<\alpha\le1$, $d_\alpha(x,y)=|x-y|^\alpha$ 是度量.
\end{exercise}
\begin{solution}
平方距离不是度量. 取 $x=0,y=1,z=2$, 有 $d(0,2)=4>1+1=d(0,1)+d(1,2)$.

对于 $0<\alpha\le1$, 先证 $(a+b)^\alpha\le a^\alpha+b^\alpha$. 当 $a+b=0$ 时成立; 否则令 $s=a/(a+b)\in[0,1]$. 因 $s^\alpha\ge s$ 且 $(1-s)^\alpha\ge1-s$, 两式相加并乘以 $(a+b)^\alpha$ 即得. 于是
\[
 |x-z|^\alpha\le(|x-y|+|y-z|)^\alpha
 \le|x-y|^\alpha+|y-z|^\alpha.
\]
正定性与对称性由绝对值得到, 所以 $d_\alpha$ 是度量. 本题也说明: 把已有距离代入递增函数, 还必须核对该函数是否保留三角不等式.
\end{solution}

\begin{exercise}\label{ch:metric:ex-finite}
设 $X=\{a,b,c\}$, 指定 $d(a,b)=2$, $d(b,c)=3$, $d(a,c)=s$, 并用对称性及 $d(x,x)=0$ 补齐距离. 求所有使 $d$ 成为度量的 $s$, 并在 $s=1$ 时写出 $B(a,1)$, $B(a,2)$, $B(a,3)$.
\end{exercise}
\begin{solution}
不同点距离为正要求 $s>0$. 三个可能有内容的三角不等式分别是 $s\le2+3$, $2\le3+s$, $3\le2+s$. 合并得 $1\le s\le5$, 而这一范围也保证全部公理. 当 $s=1$ 时, 开球的严格不等号不能遗漏, 因而
\[
 B(a,1)=\{a\},\qquad B(a,2)=\{a,c\},\qquad B(a,3)=X.\qedhere
\]
\end{solution}

\originalchapter{序列, 开闭集与连续映射}\label{ch:seq}
\emph{本章以 Paige Bright 的 MIT OCW 18.S190 IAP 2023 第2讲正文第1--7页为主要来源, 并衔接第1讲第3--4页的定义; 子空间, 闭包与边界, 距离函数以及度量与拓扑的比较为编者补充.}

有了距离, 就能把实数上的极限定义推广到函数, 向量或更一般的对象. 但度量不仅给出一个数值: 它还说明哪些集合可以作为某一点周围的活动范围. 本章先研究序列如何接近一点, 再把这种接近与开集联系起来, 最后得到连续性的三种等价表达. 这条联系也解释了为何基础拓扑能够从距离的具体公式中抽离出来.

\originalsection{收敛与 Cauchy 条件}\label{ch:seq:convergence}
\begin{definition}[序列收敛]\label{ch:seq:def-convergence}
度量空间 $(X,d)$ 中的\textbf{序列}是一个映射 $\mathbb{N}\to X$, 通常记为 $(x_n)_{n\ge1}$. 对 $x\in X$, 若
\[
 \forall\varepsilon>0,\ \exists N\in\mathbb{N},\
 \forall n\ge N,\quad d(x_n,x)<\varepsilon,
\]
则称 $x_n$ \textbf{收敛于} $x$, 写作 $x_n\to x$; $x$ 称为序列的极限.
\end{definition}
这里 $N$ 可以依赖于 $\varepsilon$, 但在 $N$ 之后必须控制所有项. 极限还必须属于所讨论的空间 $X$. 比如序列 $1/(n+1)$ 在 $\mathbb{R}$ 中趋于 $0$, 但在子空间 $(0,1)$ 中没有极限, 因为 $0$ 不是该空间的点. 定义中使用严格不等号只是统一记号; 如果把它改成 $d(x_n,x)\le\varepsilon$, 仍得到相同概念, 因为可以先对 $\varepsilon/2$ 应用该版本.

\begin{theorem}[极限的基本性质]\label{ch:seq:basic-limits}
在度量空间中, 一个收敛序列的极限唯一, 且该序列有界. 若 $x_n\to x$, 则任何子序列 $x_{n_k}$ 都收敛于同一个 $x$. 若另有 $y_n\to y$, 则
\[
 d(x_n,y_n)\longrightarrow d(x,y).
\]
\end{theorem}
\begin{proof}
若 $x_n$ 同时趋于 $x$ 和 $y$, 给定 $\varepsilon>0$, 可取足够大的 $n$, 使 $d(x,x_n)<\varepsilon/2$ 且 $d(x_n,y)<\varepsilon/2$. 三角不等式给出 $0\le d(x,y)<\varepsilon$. 由于这对每个正 $\varepsilon$ 成立, 必有 $d(x,y)=0$, 即 $x=y$.

为了证明有界, 先对 $\varepsilon=1$ 取 $N$, 使 $n\ge N$ 时 $d(x_n,x)<1$. 余下只有有限多个项. 因而
\[
 R=\max\bigl(\{1\}\cup\{d(x_n,x):1\le n<N\}\bigr)
\]
是有限数, 且每一项都满足 $d(x_n,x)\le R$. 使用集合中额外的 $1$ 可以同时处理 $N=1$, 没有前面例外项的情况.

子序列的下标 $n_1<n_2<\cdots$ 严格递增, 所以 $n_k\ge k$. 对给定 $\varepsilon$, 原序列在下标 $N$ 以后距 $x$ 小于 $\varepsilon$. 当 $k\ge N$ 时也有 $n_k\ge N$, 故 $d(x_{n_k},x)<\varepsilon$.

最后, 引理\ref{ch:metric:reverse} 给出
\[
 |d(x_n,y_n)-d(x,y)|\le d(x_n,x)+d(y_n,y).
\]
右边趋于零, 从而得到距离的极限公式. 固定 $y_n=y$ 时, 特别得到 $d(x_n,y)\to d(x,y)$.
\end{proof}
有界性只能保证各项留在一个大球中, 不能保证它们向同一点靠近. 在实数中, $(-1)^n$ 有界但不收敛. 更强的尾部控制是 Cauchy 条件: 它直接比较后面的项彼此之间的距离, 而无需预先知道极限应在哪里.

\begin{definition}[Cauchy 序列与完备性]\label{ch:seq:def-cauchy}
若序列 $(x_n)$ 满足
\[
 \forall\varepsilon>0,\ \exists N\in\mathbb{N},\
 \forall m,n\ge N,\quad d(x_m,x_n)<\varepsilon,
\]
则称其为\textbf{Cauchy 序列}. 若 $X$ 中每个 Cauchy 序列都收敛于 $X$ 中某一点, 则称 $(X,d)$ \textbf{完备}.
\end{definition}

\begin{proposition}[Cauchy 序列的基本性质]\label{ch:seq:cauchy-properties}
每个收敛序列都是 Cauchy 序列. 每个 Cauchy 序列都有界. 若一个 Cauchy 序列有子序列收敛于 $x\in X$, 则原序列也收敛于 $x$.
\end{proposition}
\begin{proof}
若 $x_n\to x$, 对给定 $\varepsilon>0$, 取 $N$ 使 $n\ge N$ 时 $d(x_n,x)<\varepsilon/2$. 对 $m,n\ge N$, 有
\[
 d(x_m,x_n)\le d(x_m,x)+d(x,x_n)<\varepsilon.
\]

若 $(x_n)$ 是 Cauchy 序列, 对精度 $1$ 取 $N$. 则对 $n\ge N$, 有 $d(x_n,x_N)<1$. 将前 $N-1$ 项到 $x_N$ 的距离与 $1$ 一起取有限最大值, 就得到整个序列的统一界.

最后设 $x_{n_k}\to x$. 对精度 $\varepsilon/2$ 取 Cauchy 下标 $N$. 再取 $k$ 足够大, 使 $n_k\ge N$ 且 $d(x_{n_k},x)<\varepsilon/2$. 对每个 $n\ge N$, 都有
\[
 d(x_n,x)\le d(x_n,x_{n_k})+d(x_{n_k},x)<\varepsilon.
\]
这里固定一个足够靠后的子序列项作为中间点, 就把子序列的信息传给了整个尾部.
\end{proof}

\begin{proposition}[两个 Cauchy 序列之间的距离]\label{ch:seq:cauchy-distance}
若 $(x_n)$ 和 $(y_n)$ 都是 $(X,d)$ 中的 Cauchy 序列, 则实数序列 $(d(x_n,y_n))$ 收敛, 即使 $X$ 本身不完备.
\end{proposition}
\begin{proof}
给定 $\varepsilon>0$, 取共同下标 $N$, 使 $m,n\ge N$ 时两序列各自项间距离均小于 $\varepsilon/2$. 由反三角不等式的双端点形式,
\[
 |d(x_n,y_n)-d(x_m,y_m)|
 \le d(x_n,x_m)+d(y_n,y_m)<\varepsilon.
\]
因此 $(d(x_n,y_n))$ 是实数中的 Cauchy 序列. 实数的完备性保证它收敛. 此处没有假设 $x_n$ 或 $y_n$ 在 $X$ 中有极限, 也没有把不存在的极限代入距离.
\end{proof}

\begin{example}
证明 $(0,1)$ 的通常距离不完备, 并证明每个离散度量空间都完备.
\end{example}
\begin{solution}
在子空间 $(0,1)$ 的通常距离下, 令 $x_n=1/(n+1)$. 给定 $\varepsilon>0$, 取 $N$ 使 $1/(N+1)<\varepsilon/2$. 对 $m,n\ge N$, 有 $|x_m-x_n|\le x_m+x_n<\varepsilon$, 所以它是 Cauchy 序列. 但它在 $\mathbb{R}$ 中的极限为 $0$, 而 $0\notin(0,1)$. 若它在 $(0,1)$ 中还有极限, 由同一距离下的极限唯一性, 该极限也只能是 $0$, 矛盾. 因而 $(0,1)$ 不完备.

在离散度量下, 对 $\varepsilon=1/2$ 应用 Cauchy 条件, 得到尾部任意两项的距离小于 $1/2$. 由于离散距离只能是 $0$ 或 $1$, 这些项必须全部相同. 因而离散度量空间中的 Cauchy 序列最终恒定, 所以每个离散度量空间都完备.
\end{solution}
完备性讨论的是空间是否容纳全部 Cauchy 极限. 后面的完备性与完备化章节会系统发展这个主题. 当前先记住: ``收敛 $\Rightarrow$ Cauchy $\Rightarrow$ 有界''在任何度量空间中成立, 反向则需要额外条件.

\originalsection{开集, 闭集及其运算}\label{ch:seq:open-closed}
\begin{definition}[开集与闭集]\label{ch:seq:def-open}
设 $A\subseteq X$. 若对每个 $x\in A$, 都存在 $r>0$, 使 $B_X(x,r)\subseteq A$, 则称 $A$ 在 $X$ 中\textbf{开}. 若 $X\setminus A$ 在 $X$ 中开, 则称 $A$ 在 $X$ 中\textbf{闭}. 若 $U$ 是包含 $x$ 的开集, 则称 $U$ 为 $x$ 的\textbf{开邻域}.
\end{definition}
本文用``开邻域''明确表示邻域本身开. 有些教材允许``邻域''是任何包含一个开邻域的集合; 两种约定对这里的收敛与连续性结论没有影响.

\begin{proposition}[开球与闭球]\label{ch:seq:balls-open}
每个开球 $B(x,r)$ 都是开集, 每个闭球 $\overline B(x,r)$ 都是闭集. 每个单点集 $\{x\}$ 都是闭集.
\end{proposition}
\begin{proof}
若 $y\in B(x,r)$, 则 $s=r-d(x,y)>0$. 对 $z\in B(y,s)$, 三角不等式给出
\[
 d(x,z)\le d(x,y)+d(y,z)<d(x,y)+s=r.
\]
所以 $B(y,s)\subseteq B(x,r)$, 证明了开球是开集. 这一论证对 $y=x$ 同样成立, 不必另设特殊情况.

若 $y\notin\overline B(x,r)$, 则 $s=d(x,y)-r>0$. 对 $z\in B(y,s)$, 反三角不等式给出
\[
 d(x,z)\ge d(x,y)-d(y,z)>d(x,y)-s=r.
\]
所以这个球留在闭球的补集中, 补集开, 闭球闭.

若 $y\ne x$, 则 $B(y,d(x,y))$ 不含 $x$. 因而单点集补集中每一点都有一个仍在补集内的球, 证明单点集闭.
\end{proof}

\begin{theorem}[开闭集的集合运算]\label{ch:seq:operations}
在度量空间 $X$ 中, $\varnothing$ 和 $X$ 都既开又闭. 任意族开集的并为开集, 有限多个开集的交为开集. 任意族闭集的交为闭集, 有限多个闭集的并为闭集.
\end{theorem}
\begin{proof}
空集开, 因为其中没有需要核对的小球条件的点. $X$ 开, 因为任何球按定义都在 $X$ 中. 二者互为补集, 所以也都闭.

设 $U=\bigcup_{i\in I}U_i$, 且每个 $U_i$ 开. 若 $x\in U$, 则 $x\in U_{i_0}$ 对某个下标 $i_0$ 成立. 取 $r>0$, 使 $B(x,r)\subseteq U_{i_0}$, 自然也有 $B(x,r)\subseteq U$. 所以任意并开.

设 $U_1,\ldots,U_m$ 开且 $x$ 属于它们的交. 分别取 $r_i>0$, 使 $B(x,r_i)\subseteq U_i$. 因为只有有限多个正数, $r=\min_i r_i$ 仍然严格为正. 因而 $B(x,r)\subseteq\bigcap_iU_i$, 证明有限交开. 零个集合的交约定为 $X$, 也与结论一致.

闭集结论由补集与 De Morgan 律得到. 具体地, 如果每个 $F_i$ 闭, 则
\[
 X\setminus\bigcap_{i\in I}F_i
 =\bigcup_{i\in I}(X\setminus F_i)
\]
是开集, 所以交闭. 对有限并, 其补集是有限多个开集的交, 因而同样开, 证明并闭. De Morgan 律本身来自逐点判断: 一点不属于并集, 当且仅当它不属于任何一个集合; 一点不属于交集, 当且仅当它不属于其中至少一个集合.
\end{proof}
``有限''这个条件不能删去. 在 $\mathbb{R}$ 中,
\[
 \bigcap_{n=1}^\infty(-1/n,1/n)=\{0\}
\]
不是开集. 与之对应,
$\bigcup_{n=1}^\infty[1/n,1]=(0,1]$ 不是闭集. 第一个例子中, 每个邻域半径都可以取正, 但无限多个半径的下确界可能为零.

有限集闭, 因为它是有限多个闭单点集的并. 开与闭也不是互相否定的两个形容词. 例如 $[0,1)$ 在 $\mathbb{R}$ 中既不开也不闭: 点 $0$ 的任何小球都会越过左端点, 而补集在点 $1$ 的任何小球都会遇到 $[0,1)$. 离散空间的任何子集则既开又闭, 因为每个点都可用单点球隔离. 开闭集合与空间能否被分成两个部分的关系, 将在连通性章节中专门讨论.

\begin{proposition}[开集是球的并]\label{ch:seq:union-balls}
$U\subseteq X$ 是开集当且仅当它可以写成 $X$ 中某些开球的并.
\end{proposition}
\begin{proof}
若 $U$ 开, 则对每个 $x\in U$, 取一个 $r_x>0$ 使 $B(x,r_x)\subseteq U$. 每个 $x$ 都在相应球中, 所以 $U=\bigcup_{x\in U}B(x,r_x)$. 若 $U=\varnothing$, 用空族的并即可. 反方向由每个球开及开集任意并仍开得到.
\end{proof}

\originalsection{相对开闭性}\label{ch:seq:subspaces}
``这个集合开不开''必须说明它被放在哪个空间中. 子空间中的球只搜索子空间里的点, 因此同一个集合在不同环境里可能有不同的开闭性质.

\begin{theorem}[子空间中的开闭集]\label{ch:seq:subspace-open}
设 $A\subseteq X$, 并在 $A$ 上使用子空间度量. 对 $U\subseteq A$, $U$ 在 $A$ 中开当且仅当存在 $X$ 中的开集 $O$, 使 $U=A\cap O$. 对 $F\subseteq A$, $F$ 在 $A$ 中闭当且仅当存在 $X$ 中的闭集 $C$, 使 $F=A\cap C$.
\end{theorem}
\begin{proof}
若 $U=A\cap O$, 对每个 $x\in U$, 由 $O$ 开可取 $r_x>0$ 使 $B_X(x,r_x)\subseteq O$. 于是
$B_A(x,r_x)=A\cap B_X(x,r_x)\subseteq U$, 所以 $U$ 在 $A$ 中开.

反过来, 若 $U$ 在 $A$ 中开, 对每个 $x\in U$ 取 $r_x>0$, 使 $A\cap B_X(x,r_x)\subseteq U$. 令 $O=\bigcup_{x\in U}B_X(x,r_x)$, 则 $O$ 在 $X$ 中开. 每个 $x\in U$ 都属于 $A\cap O$; 另一方面, $A\cap O$ 中每点落在某个 $A\cap B_X(x,r_x)$ 中, 因而属于 $U$. 所以 $A\cap O=U$.

闭集版本用 $A$ 中的补集转化为开集版本. 若 $A\setminus F=A\cap O$, 则
\[
 F=A\setminus(A\cap O)=A\cap(X\setminus O).
\]
右边是 $A$ 与 $X$ 中闭集的交. 相反方向使用同一补集关系即可.
\end{proof}

\begin{example}
说明 $[0,1/2)$ 在 $[0,1]$ 中相对开, 在 $\mathbb{R}$ 中却不开, 以及 $(0,1/2]$ 在 $(0,1)$ 中相对闭, 在 $\mathbb{R}$ 中却不闭.
\end{example}
\begin{solution}
在 $X=\mathbb{R}$ 中取 $A=[0,1]$. 集合 $U=[0,1/2)$ 是 $A$ 中的开集, 因为 $U=A\cap(-1,1/2)$, 但它不是 $\mathbb{R}$ 中的开集. 在另一个子空间 $A=(0,1)$ 中, $F=(0,1/2]$ 是相对闭集, 因为 $F=A\cap(-\infty,1/2]$, 但它不是 $\mathbb{R}$ 中的闭集.

任一子空间 $A$ 在自身中总是既开又闭. 这并不表明 $A$ 在环境空间中也开或闭. 例如 $(0,1)$ 作为自己的全集是闭的, 但在 $\mathbb{R}$ 中不是闭的.
\end{solution}

\originalsection{内部, 闭包与边界}\label{ch:seq:closure}
开集描述一个集合中的点能否容许小幅移动而仍留在集合内. 闭包则描述一个点能否被集合中的点任意逼近. 这两种描述合起来, 能够分清集合内部, 外部及两者之间的边界.

\begin{definition}[内部, 闭包与边界]\label{ch:seq:def-closure}
对 $A\subseteq X$, 定义
\begin{align*}
 A^\circ&=\{x\in X:\exists r>0,\ B(x,r)\subseteq A\},\\
 \overline A&=\{x\in X:\forall r>0,\ B(x,r)\cap A\ne\varnothing\},\\
 \partial A&=\{x\in X:\forall r>0,\ B(x,r)\cap A\ne\varnothing
             \text{ 且 }B(x,r)\cap(X\setminus A)\ne\varnothing\}.
\end{align*}
它们分别称为 $A$ 在 $X$ 中的\textbf{内部}, \textbf{闭包}和\textbf{边界}. 闭包允许一个点用它自己被``逼近'', 所以 $A\subseteq\overline A$.
\end{definition}

\begin{theorem}[内部与闭包的基本性质]\label{ch:seq:closure-properties}
$A^\circ$ 是包含在 $A$ 中的最大开集; $\overline A$ 是包含 $A$ 的最小闭集. 因而
\[
 A\text{ 开}\Longleftrightarrow A=A^\circ,
 \qquad
 A\text{ 闭}\Longleftrightarrow A=\overline A.
\]
此外,
\[
 \overline A=X\setminus(X\setminus A)^\circ,
 \qquad
 \partial A=\overline A\cap\overline{X\setminus A}
           =\overline A\setminus A^\circ.
\]
特别地, 边界是闭集, 且 $X$ 可分成三个互不相交的集合
\[
 X=A^\circ\ \cup\ \partial A\ \cup\ (X\setminus A)^\circ.
\]
\end{theorem}
\begin{proof}
若 $x\in A^\circ$, 取 $r>0$ 使 $B(x,r)\subseteq A$. 开球本身开, 且它的每个点都是 $A$ 的内点, 故 $B(x,r)\subseteq A^\circ$. 所以 $A^\circ$ 开. 任意开集 $O\subseteq A$ 中的点都有一个留在 $O$, 从而留在 $A$ 中的球, 所以 $O\subseteq A^\circ$. 这证明了最大性及开集判据.

按定义, $x\notin\overline A$ 当且仅当存在 $r>0$ 使 $B(x,r)\cap A=\varnothing$, 即 $B(x,r)\subseteq X\setminus A$. 这又等价于 $x\in(X\setminus A)^\circ$. 得到第一条恒等式, 并由内部开得到 $\overline A$ 闭. 它包含 $A$, 因为每个以 $a\in A$ 为中心的球都包含 $a$.

若 $C$ 闭且包含 $A$, 对 $x\notin C$, 开集 $X\setminus C$ 包含一个以 $x$ 为中心的球; 该球不与 $A$ 相交. 所以 $x\notin\overline A$, 即 $\overline A\subseteq C$. 这证明了最小性. 若 $A$ 已闭, 最小性给出 $\overline A\subseteq A$, 结合反向包含即得闭集判据.

边界定义要求每个球同时与两边相交, 因而正是 $\overline A\cap\overline{X\setminus A}$. 对补集应用第一条恒等式, 得 $\overline{X\setminus A}=X\setminus A^\circ$, 所以边界也等于 $\overline A\setminus A^\circ$. 两个闭集的交闭. 最后, 一个点若不在边界, 则至少有某个球完全落在 $A$ 或完全落在其补集中; 它相应属于其中一个内部. 两个内部不相交, 并且都不与边界相交, 得到所述分解.
\end{proof}

\begin{example}
分别求 $\mathbb{R}$ 中 $[0,1)$ 与 $\mathbb{Q}$ 的内部, 闭包和边界. 再在离散空间中求任意子集的内部, 闭包和边界, 并比较开球的闭包与同半径闭球.
\end{example}
\begin{solution}
在 $\mathbb{R}$ 中, 若 $A=[0,1)$, 则 $A^\circ=(0,1)$, $\overline A=[0,1]$, $\partial A=\{0,1\}$. 边界点可能属于 $A$, 也可能不属于 $A$.

有理数与无理数在 $\mathbb{R}$ 中都稠密, 因为每个非空实区间都包含两种数. 例如有理数的存在可从 Archimedes 性质先选足够大分母, 再选合适整数分子得到; 在更小区间内找到有理数后加上足够小的非零有理数倍 $\sqrt2$, 便得到区间内的无理数. 因而对 $A=\mathbb{Q}$,
\[
 A^\circ=\varnothing,\qquad \overline A=\mathbb{R},\qquad
 \partial A=\mathbb{R}.
\]
这提醒我们, 边界并不一定是一条线或少数端点, 它可以是整个空间.

在离散空间中, 每个集合都有单点球留在自己内部, 因而 $A^\circ=A=\overline A$, $\partial A=\varnothing$. 在具有至少两个点的离散空间里, $B(x,1)=\{x\}$, 所以其闭包仍是 $\{x\}$; 但闭球 $\overline B(x,1)=X$. 因此``闭球等于开球闭包''在一般度量空间中不成立.
\end{solution}

\begin{definition}[稠密, 聚点与孤立点]
若 $\overline A=X$, 称 $A$ 在 $X$ 中\textbf{稠密}. 若 $x\in X$ 的每个球都与 $A\setminus\{x\}$ 相交, 称 $x$ 为 $A$ 的\textbf{聚点}. 若 $x\in A$ 且存在球 $B(x,r)$ 满足 $B(x,r)\cap A=\{x\}$, 称 $x$ 为 $A$ 的\textbf{孤立点}.
\end{definition}
聚点的定义排除了只由点本身贡献的相交. 因此每个孤立点都属于闭包, 但不是聚点. 例如 $A=\{0\}\cup\{1/n:n\ge1\}\subset\mathbb{R}$ 中, $0$ 是聚点, 而每个 $1/n$ 都是孤立点. 对固定 $n$, 相邻数值与 $1/n$ 的距离严格为正, 取小于这些间隔的半径即可隔离它; $n=1$ 时只需核对右侧空白与下一项的距离.

\begin{proposition}[闭包的集合运算]\label{ch:seq:closure-operations}
若 $A\subseteq B$, 则 $\overline A\subseteq\overline B$ 且 $A^\circ\subseteq B^\circ$. 对任意 $A,B\subseteq X$, 有
\[
 \overline{A\cup B}=\overline A\cup\overline B,
 \qquad
 (A\cap B)^\circ=A^\circ\cap B^\circ,
 \qquad
 \overline{A\cap B}\subseteq\overline A\cap\overline B.
\]
\end{proposition}
\begin{proof}
包含关系的结论直接来自球的定义: 一个球与 $A$ 相交就与 $B$ 相交; 一个球包含在 $A$ 中就包含在 $B$ 中.

单调性给出 $\overline A\cup\overline B\subseteq\overline{A\cup B}$. 另一方面, $\overline A\cup\overline B$ 是包含 $A\cup B$ 的闭集, 由闭包最小性得到反向包含. 对内部, 若球包含在 $A\cap B$ 中, 它也分别包含在两者中. 反之, 若一点分别为两者的内点, 取两个有效半径的较小者, 得到留在交集内的球. 最后一式由 $A\cap B$ 分别包含在 $A$ 与 $B$ 中, 再用单调性得到.
\end{proof}
最后一式可能严格. 取 $A=(0,1)$, $B=(1,2)$, 二者的交为空, 但闭包的交为 $\{1\}$. 任意无限并的闭包也未必等于各自闭包的并: 取单点集 $A_n=\{1/n\}$, 各自都闭, 但并集的闭包还包含 $0$.

\originalsection{用序列识别拓扑性质}\label{ch:seq:sequence-criteria}
度量空间的小球允许我们用半径 $1/n$ 构造一串逐渐精确的逼近点. 这一可数构造, 是序列能完整识别闭包和连续性的原因.

\begin{theorem}[邻域中的最终包含]\label{ch:seq:neighborhood-limit}
对 $(X,d)$ 中的序列 $(x_n)$ 与点 $x$, 以下两条等价: $x_n\to x$; 对每个 $x$ 的开邻域 $U$, 除有限多个下标外, 所有 $x_n$ 都属于 $U$.
\end{theorem}
\begin{proof}
若 $x_n\to x$ 且 $U$ 是 $x$ 的开邻域, 取 $r>0$ 使 $B(x,r)\subseteq U$. 收敛定义保证某个 $N$ 之后各项都在这个球内, 从而都在 $U$ 中.

反过来, 对每个 $\varepsilon>0$, $B(x,\varepsilon)$ 是 $x$ 的开邻域. 只有有限多个下标的项不在这个球中, 所以取一个大于这些例外下标的 $N$, 就有 $n\ge N$ 时 $d(x_n,x)<\varepsilon$. 这正是收敛定义.
\end{proof}
\begin{remark}
这里正确的结论是序列最终\textbf{在}每个邻域中. 原讲义第2讲 Theorem 27 的陈述多出否定词 ``not'', 与前文实数例子及证明意图不一致. 本章已按收敛定义修正. ``有限多个例外''指下标个数有限, 而不只是例外值的种类有限.
\end{remark}

\begin{theorem}[闭包与闭集的序列判据]\label{ch:seq:closure-sequences}
对 $A\subseteq X$ 与 $x\in X$, 有 $x\in\overline A$ 当且仅当存在 $A$ 中的序列 $(a_n)$ 使 $a_n\to x$. 因而 $A$ 闭当且仅当: 每个在 $X$ 中收敛, 且各项都在 $A$ 中的序列, 其极限仍属于 $A$.
\end{theorem}
\begin{proof}
若 $x\in\overline A$, 对每个 $n\ge1$, 球 $B(x,1/n)$ 与 $A$ 相交, 选 $a_n$ 为交集中一点. 有 $d(a_n,x)<1/n$, 因而 $a_n\to x$. 若 $A$ 为空, 没有这样的 $x$, 两边同时不成立.

反过来, 若 $a_n\in A$ 且 $a_n\to x$, 则对任意 $r>0$, 足够靠后的 $a_n$ 落在 $B(x,r)$ 中, 所以该球与 $A$ 相交, 即 $x\in\overline A$.

若 $A$ 闭, 就有 $\overline A=A$, 因而上一结论表明所有这样的极限都在 $A$ 中. 若所有这样的极限都在 $A$ 中, 对任意 $x\in\overline A$, 用第一部分构造序列并应用假设, 得 $x\in A$. 所以 $\overline A\subseteq A$, 从而 $A$ 闭.
\end{proof}
注意判据所说的是``若序列在 $X$ 中收敛'', 并不要求每个 $A$ 中序列都收敛. 例如 $\mathbb{R}$ 是闭集, 其中序列 $n$ 却不收敛.

\begin{proposition}[聚点的序列判据]\label{ch:seq:accumulation-sequences}
$x$ 是 $A$ 的聚点当且仅当存在由 $A\setminus\{x\}$ 中互不相同的点组成的序列, 收敛于 $x$.
\end{proposition}
\begin{proof}
若 $x$ 是聚点, 先在 $B(x,1)\cap(A\setminus\{x\})$ 中取 $a_1$. 已取互不相同的 $a_1,\ldots,a_{n-1}$ 后, 取
\[
 r_n=\min\left\{\frac1n,\frac12d(x,a_1),\ldots,
 \frac12d(x,a_{n-1})\right\}>0.
\]
聚点条件允许在 $B(x,r_n)\cap(A\setminus\{x\})$ 中选 $a_n$. 此球排除了所有此前的 $a_i$, 且 $d(a_n,x)<1/n$, 因而得到互不相同且趋于 $x$ 的序列.

反过来, 若这些点趋于 $x$, 则每个以 $x$ 为中心的球都包含一个序列项, 该项属于 $A\setminus\{x\}$, 所以 $x$ 是聚点. 互不相同的条件还体现了聚点周围有无限多个不同的点; 仅有常值序列不能证明聚点性质.
\end{proof}

\originalsection{连续性的等价表达}\label{ch:seq:continuity}
函数连续意味着输入的小变化只能引起输出的小变化. 为了准确表达, 必须同时指定输入与输出所用的距离, 并区分在一点连续与在整个空间连续.

\begin{definition}[连续与 Lipschitz 连续]\label{ch:seq:def-continuity}
设 $(X,d_X),(Y,d_Y)$ 为度量空间, $f:X\to Y$, $c\in X$. 若
\[
 \forall\varepsilon>0,\ \exists\delta>0,\
 \forall x\in X,\quad
 d_X(x,c)<\delta\Longrightarrow d_Y(f(x),f(c))<\varepsilon,
\]
则称 $f$ \textbf{在 $c$ 连续}. 若 $f$ 在 $X$ 中每一点连续, 称 $f$ \textbf{连续}. 如果存在一个与 $x,x'$ 无关的常数 $L\ge0$, 使
\[
 d_Y(f(x),f(x'))\le Ld_X(x,x')\qquad(x,x'\in X),
\]
则称 $f$ 为\textbf{Lipschitz 连续映射}.
\end{definition}
一般连续性的 $\delta$ 可以依赖于 $c$ 与 $\varepsilon$. Lipschitz 估计则用一个常数控制整个空间. 若 $L>0$, 取 $\delta=\varepsilon/L$ 即得连续; 若 $L=0$, $f$ 是常值映射, 也连续. 第1章微分算子从 $C^1$ 范数到上确界范数的估计, 就是 $L=1$ 的例子. 若函数原本定义在 $A\subseteq X$ 上, 应把 $A$ 当作定义域并使用子空间度量, 量词中的 $x$ 只在 $A$ 中变化.

\begin{theorem}[点处连续的三个判据]\label{ch:seq:continuity-equivalence}
对 $f:X\to Y$ 和 $c\in X$, 下列条件等价:
\begin{enumerate}
 \item $f$ 在 $c$ 满足 $\varepsilon$--$\delta$ 连续性定义;
 \item 对 $X$ 中每个满足 $x_n\to c$ 的序列, 都有 $f(x_n)\to f(c)$;
 \item 对 $f(c)$ 的每个开邻域 $U\subseteq Y$, 逆像 $f^{-1}(U)$ 包含 $c$ 的一个开邻域.
\end{enumerate}
\end{theorem}
\begin{proof}
先由第一条推出第二条. 对给定 $\varepsilon>0$, 连续性给出 $\delta>0$. 因为 $x_n\to c$, 某个 $N$ 之后都有 $d_X(x_n,c)<\delta$. 故相同尾部满足 $d_Y(f(x_n),f(c))<\varepsilon$, 得到像序列收敛.

反过来, 假设 $f$ 在 $c$ 不连续. 否定定义中的量词, 得到一个固定的 $\varepsilon_0>0$, 使对每个 $\delta>0$, 都能找到 $x\in X$ 满足
\[
 d_X(x,c)<\delta,\qquad d_Y(f(x),f(c))\ge\varepsilon_0.
\]
对 $\delta=1/n$ 逐次选择 $x_n$. 这样 $x_n\to c$, 但所有像点到 $f(c)$ 的距离都至少为 $\varepsilon_0$, 所以像序列不趋于 $f(c)$. 这与第二条矛盾. 反证中 $\varepsilon_0$ 必须是固定的, 不能随 $n$ 变化.

再证第一条推出第三条. 若 $U$ 是 $f(c)$ 的开邻域, 取 $\varepsilon>0$ 使 $B_Y(f(c),\varepsilon)\subseteq U$. 连续性给出 $\delta>0$, 从而
\[
 B_X(c,\delta)\subseteq f^{-1}(B_Y(f(c),\varepsilon))
 \subseteq f^{-1}(U).
\]
左边是所需开邻域. 最后, 若第三条成立, 对 $U=B_Y(f(c),\varepsilon)$ 应用它. 所得到的 $c$ 的开邻域包含某个球 $B_X(c,\delta)$, 故这个球的像落在 $U$ 中, 正是第一条.
\end{proof}
注意第三条只要求逆像在点 $c$ 附近包含一个开邻域. 若 $f$ 只在 $c$ 连续, 整个逆像未必开. 当连续性在每一点成立时, 才能得到如下全局结论.

\begin{theorem}[连续映射的逆像判据]\label{ch:seq:preimage}
对度量空间间的映射 $f:X\to Y$, 下列条件等价: $f$ 连续; $Y$ 中每个开集的逆像在 $X$ 中开; $Y$ 中每个闭集的逆像在 $X$ 中闭.
\end{theorem}
\begin{proof}
若 $f$ 连续且 $U\subseteq Y$ 开, 对每个 $x\in f^{-1}(U)$, $U$ 是 $f(x)$ 的开邻域. 上一定理给出包含在 $f^{-1}(U)$ 中的 $x$ 的开邻域, 进而给出一个包含在逆像中的小球. 因而逆像开.

若每个开集的逆像开, 固定 $c\in X$ 与 $\varepsilon>0$. 集合 $f^{-1}(B_Y(f(c),\varepsilon))$ 是包含 $c$ 的开集, 所以包含某个 $B_X(c,\delta)$. 这证明 $f$ 在 $c$ 连续. 因为 $c$ 任意, $f$ 连续.

最后, 逆像与补集可交换:
\[
 f^{-1}(Y\setminus A)=X\setminus f^{-1}(A).
\]
因此开集逆像开与闭集逆像闭互相等价.
\end{proof}

\begin{example}
给出连续映射不保持开集的像为开集, 也不保持闭集的像为闭集的例子.
\end{example}
\begin{solution}
常值映射 $f:\mathbb{R}\to\mathbb{R}$, $f(x)=0$, 是连续的, 但非空开区间的像为 $\{0\}$, 并非开集. 连续性因而保证的是\emph{开集的逆像}开, 不是开集的像开.

同样, 闭集的像不一定闭. 连续函数 $f(x)=e^x$ 把闭集 $\mathbb{R}$ 映成 $(0,\infty)$, 该集合在 $\mathbb{R}$ 中不闭. 后文将看到, 若定义域集合紧, 连续像具有更强的性质; 当前不能把这一额外结论提前套用到任意闭集.
\end{solution}

\begin{proposition}[复合与限制]\label{ch:seq:composition}
若 $f:X\to Y$, $g:Y\to Z$ 连续, 则 $g\circ f$ 连续. 若 $A\subseteq X$, 则 $f|_A:A\to Y$ 在子空间度量下连续.
\end{proposition}
\begin{proof}
对 $Z$ 中任意开集 $U$, 有
$(g\circ f)^{-1}(U)=f^{-1}(g^{-1}(U))$. 先由 $g$ 连续知内部逆像在 $Y$ 中开, 再由 $f$ 连续知外部逆像在 $X$ 中开. 应用逆像判据, 得到复合连续.

对限制映射, $(f|_A)^{-1}(U)=A\cap f^{-1}(U)$. 右边是 $A$ 与 $X$ 中开集的交, 根据子空间定理在 $A$ 中开. 所以限制也连续.
\end{proof}

\originalsection{到集合的距离}\label{ch:seq:distance-functions}
反三角不等式已经告诉我们, 到固定点的距离是连续函数. 现在将固定点换成一个非空集合, 就得到一个能够识别闭包的重要工具.

\begin{definition}[距离函数]
设 $A\subseteq X$ 非空. 对 $x\in X$, 定义
\[
 \operatorname{dist}(x,A)=\inf_{a\in A}d(x,a).
\]
它是 $x$ 到 $A$ 的\textbf{距离}. 因为集合非空, 所取下确界是一组有限非负数的下确界, 因而总是一个有限非负实数. 这里不要求存在某个 $a$ 达到下确界.
\end{definition}

\begin{theorem}[距离函数的连续性与零点]\label{ch:seq:distance-set}
对非空 $A\subseteq X$, 有
\[
 |\operatorname{dist}(x,A)-\operatorname{dist}(y,A)|\le d(x,y).
\]
因此距离函数为 $1$-Lipschitz 连续函数. 此外,
\[
 \operatorname{dist}(x,A)=0\Longleftrightarrow x\in\overline A.
\]
若 $A$ 闭, 零点集恰好是 $A$.
\end{theorem}
\begin{proof}
对每个 $a\in A$, 有 $d(x,a)\le d(x,y)+d(y,a)$. 对 $a$ 取下确界, 得
\[
 \operatorname{dist}(x,A)\le d(x,y)+\operatorname{dist}(y,A).
\]
交换 $x,y$, 就得到绝对值估计.

若下确界等于零, 对每个 $r>0$ 都能找到 $a\in A$ 使 $d(x,a)<r$; 否则 $r$ 会是一个正下界, 与下确界为零矛盾. 所以每个球都与 $A$ 相交, 即 $x\in\overline A$. 反过来, 若 $x\in\overline A$, 每个球都与 $A$ 相交, 则下确界不超过每个正 $r$, 只能为零. 闭集情形由 $\overline A=A$ 得到.
\end{proof}

\begin{example}
举例说明到一个非空集合的距离可以为零却没有最近点, 以及两个非空, 不相交闭集之间的距离可以为零.
\end{example}
\begin{solution}
在 $\mathbb{R}$ 中取 $A=(0,1)$, 有 $\operatorname{dist}(0,A)=0$, 但 $0\notin A$, 所以没有 $a\in A$ 能使 $|0-a|$ 等于零. 即使 $A$ 闭, 一般度量空间中也不能自动断言最近点存在. 例如取 $X=\{p\}\cup\{a_n:n\ge1\}$, 对不同点指定
\[
 d(p,a_n)=1+\frac1n,\qquad d(a_n,a_m)=2\quad(n\ne m),
\]
并用对称性及对角线上取零补齐定义. 所有非零距离都在 $(1,2]$ 中, 所以三个互不相同的点满足三角不等式: 左边不超过 $2$, 右边大于 $2$; 有点重复的情形也直接成立. 因而这是度量. 集合 $A=\{a_n:n\ge1\}$ 闭, 因为它的补集 $\{p\}=B(p,1)$ 是开集. 但是 $\operatorname{dist}(p,A)=1$, 而每个 $d(p,a_n)>1$, 所以没有最近点. 最近点问题需要额外紧性或几何条件.

两个不相交闭集之间的距离也可能为零. 在 $\mathbb{R}^2$ 的欧氏距离下, 取
\[
 F=\{(t,0):t\ge1\},\qquad
 G=\{(t,1/t):t\ge1\}.
\]
它们不相交. $F$ 显然闭; 若 $G$ 中序列 $(t_n,1/t_n)$ 在 $\mathbb{R}^2$ 中趋于 $(t,s)$, 则 $t_n\to t\ge1$, 且 $1/t_n\to1/t$, 所以 $s=1/t$, 极限仍在 $G$ 中, 由序列判据知 $G$ 闭. 但两点 $(n,0)$ 与 $(n,1/n)$ 的距离为 $1/n$, 因而
$\inf\{d(u,v):u\in F,v\in G\}=0$. 对一个固定外点到闭集的距离严格为正, 与两个无限闭集之间缺少统一正间隔, 是不同的命题.
\end{solution}

\originalsection{度量可以不同而拓扑相同}\label{ch:seq:equivalent-metrics}
开集保留了每一点附近有哪些可用的活动范围, 但不保留距离的具体数值. 将一个度量所产生的全部开集记为 $\mathcal{T}_d$, 称为它诱导的\textbf{度量拓扑}. 定理\ref{ch:seq:operations} 证明了它满足空集和全集属于其中, 任意并封闭及有限交封闭三个性质. 后文的基础拓扑将直接把这三条作为开集系统的公理.

\begin{definition}[同拓扑的度量与同胚]
同一集合 $X$ 上的度量 $d$ 与 $\rho$ 若满足 $\mathcal{T}_d=\mathcal{T}_\rho$, 称为\textbf{产生相同拓扑的度量}. 若双射 $h:X\to Y$ 及其逆映射 $h^{-1}$ 都连续, 称 $h$ 为\textbf{同胚}. 同胚保留开集系统, 因而也保留由开集表述的性质.
\end{definition}

\begin{proposition}[比较小球]\label{ch:seq:same-topology}
$d$ 与 $\rho$ 产生相同拓扑, 当且仅当对每个 $x\in X$ 与每个 $r>0$, 都能找到 $s,t>0$ 使
\[
 B_\rho(x,s)\subseteq B_d(x,r),\qquad
 B_d(x,t)\subseteq B_\rho(x,r).
\]
它们产生相同拓扑时, 有完全相同的收敛序列及极限, 并且用二者之一作为定义域或值域的度量, 不改变映射的连续性.
\end{proposition}
\begin{proof}
若拓扑相同, $d$-球是包含 $x$ 的 $\rho$-开集, 因而包含某个 $\rho$-球. 交换角色得到另一包含关系.

反过来, 若小球互相包含的条件成立, 对 $d$-开集 $U$ 中每一点 $x$, 先取留在 $U$ 中的 $d$-球, 再在该球中取一个 $\rho$-球, 便证明 $U$ 为 $\rho$-开集. 交换角色得到所有开集相同. 收敛序列相同由邻域判据得到, 连续性相同则由开集逆像判据得到. 此时恒等映射在两个度量空间之间是同胚.
\end{proof}

例如 $\mathbb{R}^n$ 上的 $d_1,d_2,d_\infty$ 产生相同拓扑, 因为第1章的范数比较给出了互相包含的小球. 球面弦长距离与测地距离也由同样的理由产生相同拓扑. 练习\ref{ch:metric:ex-bounded} 中 $d/(1+d)$ 的球可明确换算, 所以它与 $d$ 也产生相同拓扑.

\begin{proposition}[双边常数估计的额外作用]\label{ch:seq:bilipschitz-metrics}
若存在 $c,C>0$, 使 $cd(x,y)\le\rho(x,y)\le Cd(x,y)$ 对所有 $x,y\in X$ 成立, 则两种度量产生相同拓扑, 具有相同的 Cauchy 序列, 并且其中一个完备当且仅当另一个完备.
\end{proposition}
\begin{proof}
由 $\rho\le Cd$ 得 $B_d(x,r/C)\subseteq B_\rho(x,r)$; 由 $cd\le\rho$ 得 $B_\rho(x,cr)\subseteq B_d(x,r)$. 上一命题给出拓扑相同.

若序列为 $d$-Cauchy, 对精度 $\varepsilon/C$ 应用它, 便得到 $\rho$-Cauchy 条件; 反方向使用精度 $c\varepsilon$. 因而 Cauchy 序列相同. 若 $d$ 完备, 则任何 $\rho$-Cauchy 序列先是 $d$-Cauchy, 在 $d$ 中收敛, 再由拓扑相同知也在 $\rho$ 中收敛. 反方向相同.
\end{proof}
这里只说这一\emph{较强的双边估计}能够保留完备性. 单有拓扑相同, 结论并不成立.

\begin{example}
\label{ch:seq:not-topological-complete}在 $\mathbb{R}$ 上考虑 $d(x,y)=|x-y|$ 和 $\rho(x,y)=|\arctan x-\arctan y|$. 证明二者都是度量, 产生相同拓扑, 但一个完备, 另一个不完备. 同时比较二者的有界性.
\end{example}
\begin{solution}

在同一集合 $\mathbb{R}$ 上, 考虑
\[
 d(x,y)=|x-y|,\qquad
 \rho(x,y)=|\arctan x-\arctan y|.
\]
$\arctan$ 是单射, 因而 $\rho$ 的正定性成立; 对称性与三角不等式来自实数绝对值, 所以 $\rho$ 是度量.

它与通常距离产生相同拓扑. 一方面, 中值定理与 $0<(\arctan)'(t)=1/(1+t^2)\le1$ 给出 $\rho(x,y)\le|x-y|$, 因此通常距离足够小便使 $\rho$ 足够小. 另一方面, 固定 $x$, $\tan$ 在 $\arctan x$ 处连续, 所以对于任意 $\varepsilon>0$, 存在 $\delta>0$, 使
\[
 |\arctan y-\arctan x|<\delta
 \Longrightarrow |y-x|<\varepsilon.
\]
这就给出了反方向的小球包含.

通常实数空间完备, 但 $(\mathbb{R},\rho)$ 不完备. 序列 $x_n=n$ 的反正切值趋于 $\pi/2$, 所以这些值构成实数 Cauchy 序列, 从而 $(n)$ 为 $\rho$-Cauchy. 若它在 $\rho$ 中趋于某个 $x\in\mathbb{R}$, 则 $|\arctan n-\arctan x|\to0$, 因而 $\arctan x=\pi/2$, 不可能. 所以该序列没有 $\rho$-极限.

同一开集系统保留了收敛与连续性, 却没有保留 Cauchy 条件. 原因是 Cauchy 条件要求尾部任意两点在一个\emph{统一的距离尺度}下接近, 而拓扑只指定每一点局部有哪些邻域. 本例也说明有界性不是拓扑不变量: $\mathbb{R}$ 在通常距离下无界, 但所有 $\rho$-距离都小于 $\pi$.
\end{solution}

\originalsection{练习与解答}\label{ch:seq:exercises}
以下题目均为编者练习. 前几题训练量词与相对拓扑, 后几题说明常见结论的边界.

\begin{exercise}\label{ch:seq:ex-subsequence}
设 $(x_n)$ 是度量空间中的序列. 证明 $x_n\to x$ 当且仅当每个子序列都有一个进一步的子序列收敛于 $x$.
\end{exercise}
\begin{solution}
正方向由子序列保持同一极限得到, 每个子序列本身就已经趋于 $x$.

反方向用反证法. 若 $x_n$ 不趋于 $x$, 则存在固定 $\varepsilon_0>0$, 使对每个 $N$, 都存在 $n\ge N$ 满足 $d(x_n,x)\ge\varepsilon_0$. 递归选取严格递增的下标 $n_k$, 使这些项全部满足该不等式. 这个子序列的任何进一步子序列仍然距 $x$ 至少为 $\varepsilon_0$, 不可能趋于 $x$, 与题设矛盾.
\end{solution}

\begin{exercise}\label{ch:seq:ex-relative-closure}
设 $A\subseteq X$, $E\subseteq A$. 证明 $E$ 在子空间 $A$ 中的闭包为 $A\cap\overline E^{\,X}$, 其中上标说明闭包在 $X$ 中取. 再取 $X=\mathbb{R}$, $A=(0,1]$, $E=(0,1)$, 求 $E$ 在 $A$ 中的内部, 闭包与边界.
\end{exercise}
\begin{solution}
对 $x\in A$, 有
\[
 B_A(x,r)\cap E=(A\cap B_X(x,r))\cap E=B_X(x,r)\cap E,
\]
因为 $E\subseteq A$. 所以 $x$ 的每个 $A$-球与 $E$ 相交, 当且仅当每个 $X$-球与 $E$ 相交. 将 $x$ 限制在 $A$ 中, 就得到闭包公式.

在具体例子中, $E=A\cap(-1,1)$, 所以 $E$ 在 $A$ 中开, 内部为 $(0,1)$. 它在 $X$ 中的闭包是 $[0,1]$, 与 $A$ 相交得 $(0,1]$. 因而相对边界为 $\{1\}$. 点 $0$ 不在环境空间 $A$ 中, 不能出现在相对边界里.
\end{solution}

\begin{exercise}\label{ch:seq:ex-separation}
设 $F,G$ 是度量空间中两个非空, 不相交的闭集. 证明
\[
 u(x)=\frac{\operatorname{dist}(x,F)}
 {\operatorname{dist}(x,F)+\operatorname{dist}(x,G)}
\]
是连续函数, 满足 $u|_F=0$ 与 $u|_G=1$. 不要假定两个集合之间的距离有正下界.
\end{exercise}
\begin{solution}
两个距离函数连续且非负. 若某点分母为零, 则它到 $F$ 与 $G$ 的距离都为零. 由于两集合闭, 该点同时属于 $F$ 与 $G$, 与不相交矛盾. 因而分母在每一点严格为正; 通过连续实函数的加法与除法, 得到 $u$ 连续.

若 $x\in F$, 分子为零, 而 $x\notin G$ 保证到 $G$ 的距离严格为正, 故 $u(x)=0$. 若 $x\in G$, 则到 $G$ 的距离为零, 到 $F$ 的距离为正, 故 $u(x)=1$. 证明只使用了每个固定点的分母为正, 并未要求所有点的分母具有统一正下界.
\end{solution}

\begin{exercise}\label{ch:seq:ex-open-union}
设 $X=U\cup V$, 其中 $U,V$ 是 $X$ 中的开集. 若 $f:X\to Y$ 的限制 $f|_U$ 和 $f|_V$ 均连续, 证明 $f$ 连续. 再证明将有限个开集换成任意族开集, 结论仍成立.
\end{exercise}
\begin{solution}
对 $Y$ 中开集 $O$, 限制连续意味着 $f^{-1}(O)\cap U$ 在 $U$ 中开. 因为 $U$ 本身在 $X$ 中开, 子空间中相对开集也是 $X$ 中开集: 它可写成 $U\cap W$ 且 $W$ 在 $X$ 中开. $V$ 同理. 因而
\[
 f^{-1}(O)=(f^{-1}(O)\cap U)\cup(f^{-1}(O)\cap V)
\]
在 $X$ 中开, 逆像判据给出连续. 对任意开覆盖 $(U_i)_{i\in I}$, 使用 $f^{-1}(O)=\bigcup_i(f^{-1}(O)\cap U_i)$, 并应用任意开集并仍开即可. 这里不需要覆盖有限.
\end{solution}

\begin{exercise}\label{ch:seq:ex-graph}
在 $X\times Y$ 上使用度量
$D((x,y),(x',y'))=\max\{d_X(x,x'),d_Y(y,y')\}$.
证明这个公式给出度量, 且连续映射 $f:X\to Y$ 的图像
$\Gamma_f=\{(x,f(x)):x\in X\}$ 是闭集. 说明反过来不一定成立.
\end{exercise}
\begin{solution}
正定性与对称性由各个分量得到. 每个分量距离都满足三角不等式, 因而两分量中的任一个都不超过相邻两段的 $D$ 距离之和; 取最大值即得 $D$ 的三角不等式. 此外, 在 $D$ 中收敛等价于两个分量分别收敛.

设 $\Gamma_f$ 中序列 $(x_n,f(x_n))$ 在积空间中趋于 $(x,y)$. 则 $x_n\to x$ 且 $f(x_n)\to y$. 连续性给出 $f(x_n)\to f(x)$, 极限唯一性给出 $y=f(x)$. 所以极限仍在图像中, 由闭集序列判据知图像闭.

反例取 $X=Y=\mathbb{R}$, 并定义 $f(0)=0$, $f(x)=1/x$ 对 $x\ne0$. 它在 $0$ 不连续, 因为 $1/n\to0$, 但 $f(1/n)=n$ 不趋于 $0$. 图像却闭. 若图像中序列 $(x_n,f(x_n))$ 趋于有限点 $(x,y)$, 当 $x\ne0$ 时, 最终 $x_n\ne0$, 且 $f(x_n)\to1/x$, 故 $y=1/x$. 当 $x=0$ 时, 如果有无限多个 $x_n\ne0$, 沿这些下标 $x_n\to0$, 则 $|f(x_n)|=1/|x_n|\to\infty$, 与像序列收敛从而有界矛盾. 因此最终 $x_n=0$, 得 $y=0$. 两种情形的极限都在图像中, 图像闭.
\end{solution}

\begin{exercise}\label{ch:seq:ex-truncated-cauchy}
对于练习\ref{ch:metric:ex-bounded} 的 $\rho=d/(1+d)$, 证明 $d$ 与 $\rho$ 具有相同的 Cauchy 序列, 因而完备性相同. 解释这为何不与例\ref{ch:seq:not-topological-complete} 矛盾.
\end{exercise}
\begin{solution}
因 $\rho(x,y)\le d(x,y)$, $d$-Cauchy 立即推出 $\rho$-Cauchy. 反方向, 给定 $\varepsilon>0$, 令 $\eta=\varepsilon/(1+\varepsilon)\in(0,1)$. 由于 $t/(1+t)$ 严格递增,
\[
 \rho(x_n,x_m)<\eta\Longrightarrow d(x_n,x_m)<\varepsilon.
\]
对 $\rho$-Cauchy 条件使用精度 $\eta$, 就得到 $d$-Cauchy 条件. 两者已有相同拓扑, 所以收敛序列相同; Cauchy 序列也相同, 完备性因而相同.

例\ref{ch:seq:not-topological-complete} 只说明同拓扑未必足以保留完备性, 并没有说每一对同拓扑度量都改变完备性. 本题的换算对所有点对使用同一个精度 $\eta$, 而反正切例子中从 $\rho$ 小推出 $d$ 小, 只能在一个固定有限点附近使用依赖该点的局部控制.
\end{solution}

\originalchapter{欧氏空间的紧性}
\label{ch:eucompact}

本章以 Paige Bright 的 MIT OCW 18.S190 IAP 2023 第 3 讲为主线重编, 讨论范数、开覆盖与欧氏空间的紧集. 为使证明可以独立追踪, 补写原讲义中的区间端点论证, 并补充嵌套方盒、有限维范数等价与紧支撑的反例. 这些补充是编者组织的学习内容, 不作为原课程的逐字翻译或试题题源.
\emph{来源定位: 第 3 讲正文第 1--3 页为范数与支撑, 第 4--7 页为紧性、Heine--Borel 与 Bolzano--Weierstrass. 页码采用原讲义印刷页码, 不含末尾 OCW 许可页.}

\originalsection{范数与线性结构}
\label{ch:eucompact:norms}

度量比较两个点之间的距离. 向量空间还允许作差, 因而可以先衡量一个向量的长度, 再把两点之间的距离定义为它们之差的长度. 范数正是与加法和数乘相容的长度概念. 本节以实向量空间为例; 向量空间的代数公理与有限维线性代数作为先修知识.

\begin{definition}[范数]
设 $V$ 是实向量空间. 函数 $\|\cdot\|:V\to[0,\infty)$ 称为范数, 如果对任意 $u,v\in V$ 与 $\lambda\in\mathbb{R}$ 都有
\[
 \|v\|=0\Longleftrightarrow v=0,\qquad
 \|\lambda v\|=|\lambda|\|v\|,\qquad
 \|u+v\|\leq\|u\|+\|v\|.
\]
配备范数的向量空间称为赋范空间.
\end{definition}

\begin{proposition}[范数诱导的度量]
若 $\|\cdot\|$ 是 $V$ 上的范数, 则 $d(u,v)=\|u-v\|$ 是度量. 此外,
\[
 \bigl|\|u\|-\|v\|\bigr|\leq\|u-v\|.
\]
\end{proposition}
\begin{proof}
范数的正定性给出 $d(u,v)\geq0$ 及 $d(u,v)=0$ 当且仅当 $u=v$. 由齐次性,
$\|u-v\|=\|-(v-u)\|=\|v-u\|$, 所以 $d$ 对称. 把 $u-w$ 写成 $(u-v)+(v-w)$, 再使用范数的三角不等式, 得到
\[
 d(u,w)\leq d(u,v)+d(v,w).
\]
最后, $\|u\|\leq\|u-v\|+\|v\|$ 给出 $\|u\|-\|v\|\leq\|u-v\|$. 交换 $u,v$ 后得到反方向的估计, 合在一起即为所求.
\end{proof}

\begin{example}
对 $x=(x_1,\ldots,x_n)$, 定义
\[
 \|x\|_1=\sum_{j=1}^n|x_j|,\qquad
 \|x\|_2=\left(\sum_{j=1}^n x_j^2\right)^{1/2},\qquad
 \|x\|_\infty=\max_{1\leq j\leq n}|x_j|.
\]
证明这三个函数都是范数.
\end{example}
\begin{solution}
对 $\|\cdot\|_1$, 三角不等式由各坐标上的 $|x_j+y_j|\leq|x_j|+|y_j|$ 相加得到. 对 $\|\cdot\|_\infty$, 各坐标都有 $|x_j+y_j|\leq\|x\|_\infty+\|y\|_\infty$, 取最大值得到三角不等式. 欧氏范数的三角不等式则由有限维 Cauchy--Schwarz 不等式给出:
\[
 \|x+y\|_2^2
 =\|x\|_2^2+2\sum_j x_jy_j+\|y\|_2^2
 \leq(\|x\|_2+\|y\|_2)^2.
\]
其余两条范数公理均可直接由坐标定义核对.
\end{solution}

三种长度的数值不相同, 但有不依赖于向量 $x$ 的比较常数:
\begin{equation}
 \|x\|_\infty\leq\|x\|_2\leq\|x\|_1
 \leq\sqrt n\,\|x\|_2\leq n\|x\|_\infty.
 \label{ch:eucompact:coordinate-norms}
\end{equation}
其中 $\|x\|_1\leq\sqrt n\,\|x\|_2$ 来自对 $(|x_j|)$ 与 $(1)$ 使用 Cauchy--Schwarz; 其余估计来自最大项与求和的比较. 因此, 三种度量下的收敛完全相同: 一种长度趋于零, 另外两种也趋于零. 后面将证明这不是三个特殊公式之间的巧合, 而是所有有限维范数共有的性质. 为避免用尚未证明的紧性, 我们先建立欧氏空间的紧性理论, 再完成该结论.

\originalsection{开覆盖与紧性}
\label{ch:eucompact:open-covers}

有限集合有一个简单性质: 不论怎样给每个点安排一个包含它的开集, 总能从这批开集中挑出有限多个, 仍然覆盖全部点. 紧性把这个性质保留下来, 却不要求集合本身只有有限个点.

\begin{definition}[开覆盖与紧集]
设 $A\subseteq X$, 其中 $(X,d)$ 是度量空间. 若每个 $U_i$ 都在 $X$ 中开, 且
\[
 A\subseteq\bigcup_{i\in I}U_i,
\]
则称 $\{U_i\}_{i\in I}$ 为 $A$ 的开覆盖. 从中选出的子族若仍覆盖 $A$, 称为子覆盖; 只有有限个成员的子覆盖称为有限子覆盖. 若 $A$ 的每个开覆盖都有有限子覆盖, 就称 $A$ 是紧集.
\end{definition}

覆盖要求的是包含关系, 而不要求开集的并恰好等于 $A$. 例如一个包含 $[0,1]$ 的开区间就构成 $[0,1]$ 的开覆盖. 当我们把 $A$ 本身作为度量空间时, 开集是相对于 $A$ 开的. 这两个说法与紧性相容, 但需要说明理由.

\begin{proposition}[紧性与子空间]
\label{ch:eucompact:subspace}
设 $A\subseteq Y\subseteq X$, 均使用从 $X$ 限制而来的度量. 则 $A$ 作为 $X$ 的子集是紧集, 当且仅当 $A$ 作为 $Y$ 的子集是紧集, 也当且仅当度量空间 $A$ 是紧的.
\end{proposition}
\begin{proof}
先比较 $X$ 与 $Y$. 若 $A$ 在 $X$ 中紧, 而 $\{V_i\}$ 是由 $Y$ 中开集组成的 $A$ 的覆盖, 则存在 $X$ 中的开集 $U_i$ 使 $V_i=Y\cap U_i$. 因为 $A\subseteq Y$, $\{U_i\}$ 也覆盖 $A$. 取其中有限子覆盖后, 相应的 $V_i$ 仍覆盖 $A$.

反过来, 若 $A$ 在 $Y$ 中紧, 对任意 $X$ 中开集给出的覆盖 $\{U_i\}$, 子空间开集 $Y\cap U_i$ 覆盖 $A$. 取有限多个后, 原来的相应 $U_i$ 也覆盖 $A$. 再取 $Y=A$, 就得到最后一个等价关系.
\end{proof}

\begin{example}
证明度量空间中的任意有限子集与空集都是紧集.
\end{example}
\begin{solution}
对有限子集 $\{a_1,\ldots,a_m\}$ 的任意开覆盖, 对每个 $a_j$ 从覆盖中选一个包含它的成员, 最多选出 $m$ 个就足够. 空集也紧, 因为空子族已经覆盖空集.
\end{solution}
以后关于取得最大值、选定一个点等陈述则要另外要求集合非空.

\begin{example}
证明 $\mathbb{R}$ 与 $(0,1]$ 在通常实数度量下均不紧.
\end{example}
\begin{solution}
$\mathbb{R}$ 上的开区间 $U_m=(-m,m)$, $m\geq1$, 覆盖整个 $\mathbb{R}$. 有限多个的并只等于 $(-M,M)$, 其中 $M$ 是所选指标的最大值, 因而不能覆盖 $\mathbb{R}$. 这说明 $\mathbb{R}$ 不紧.

集合 $(0,1]$ 虽然有界, 却也不紧. 取 $U_m=(1/m,2)$, $m\geq1$. 对每个 $x\in(0,1]$, 可取 $m>1/x$ 使 $x\in U_m$, 所以这是一族开覆盖. 有限子族的并为 $(1/M,2)$, 未覆盖点 $1/(2M)$. 前一个例子中的点可以向无穷远逃逸, 后一个例子中的点可以向缺失的端点 $0$ 逃逸. 这正是接下来闭性与有界性分别要阻止的现象.
\end{solution}

\begin{theorem}[度量空间中的紧集闭且有界]
\label{ch:eucompact:compact-closed-bounded}
度量空间 $X$ 的每个紧子集 $K$ 都在 $X$ 中闭且有界.
\end{theorem}
\begin{proof}
若 $K=\varnothing$, 两个结论均成立. 以下设 $K\ne\varnothing$, 并固定 $p\in K$. 正整数半径的球 $B_X(p,m)$ 覆盖 $X$, 当然也覆盖 $K$. 由紧性, 有有限多个球覆盖 $K$. 若所选半径的最大值为 $M$, 则 $K\subseteq B_X(p,M)$, 所以 $K$ 有界.

为了证明闭性, 固定 $z\in X\setminus K$, 并对每个 $q\in K$ 令
\[
 r_q=\frac{d(z,q)}{3}>0,\qquad W_q=B_X(q,r_q).
\]
这些 $W_q$ 覆盖 $K$, 因而存在 $q_1,\ldots,q_s$ 使 $K\subseteq\bigcup_{j=1}^sW_{q_j}$. 令 $r=\min_jr_{q_j}>0$. 若 $x\in B_X(z,r)\cap K$, 则 $x$ 落在某个 $W_{q_j}$ 内, 于是
\[
 d(z,q_j)\leq d(z,x)+d(x,q_j)
 <r+r_{q_j}\leq2r_{q_j}=\frac23d(z,q_j),
\]
矛盾. 因此 $B_X(z,r)$ 与 $K$ 不相交. 每个 $z\notin K$ 都有这样的邻域, 所以 $X\setminus K$ 开, 即 $K$ 在 $X$ 中闭.
\end{proof}

上述闭性是相对于指定的环境空间 $X$ 而言的. 每个空间在自身中都是闭集, 却未必紧. 例如 $(0,1]$ 作为自身的子空间是闭且有界的, 仍不紧. 因而不能把后面的欧氏空间定理误读为任意子空间中闭且有界都紧.

\begin{proposition}[紧空间的闭子集]
\label{ch:eucompact:closed-subset}
若 $K$ 是紧度量空间, 且 $F\subseteq K$ 在 $K$ 中闭, 则 $F$ 紧. 特别地, 若 $K\subseteq X$ 紧而 $F\subseteq K$ 在 $X$ 中闭, 也有 $F$ 紧.
\end{proposition}
\begin{proof}
设 $\{V_i\}_{i\in I}$ 是由 $K$ 中开集组成的 $F$ 的覆盖. 因为 $F$ 在 $K$ 中闭, $K\setminus F$ 在 $K$ 中开, 且
\[
 K=(K\setminus F)\cup\bigcup_{i\in I}V_i.
\]
由 $K$ 的紧性选出有限子覆盖. 从这有限子族中去掉 $K\setminus F$ 后, 所剩的 $V_i$ 仍覆盖 $F$, 因为 $K\setminus F$ 不含 $F$ 的任何点. 因此 $F$ 在 $K$ 中紧, 再由命题 \ref{ch:eucompact:subspace} 得到所需的子集紧性. 若 $F$ 在 $X$ 中闭, 它也在 $K$ 中闭, 故属于前一种情形.
\end{proof}

\originalsection{闭区间与闭方盒}
\label{ch:eucompact:boxes}

要证明 $\mathbb{R}^n$ 中的闭有界集紧, 不能只重复紧集闭且有界的论证. 我们需要从实数的完备性出发, 真正构造有限子覆盖. 一维情况下, 上确界能描述已经用有限多个开集覆盖到的最远位置.

\begin{theorem}[闭区间的紧性]
\label{ch:eucompact:interval}
对实数 $a\leq b$, 闭区间 $[a,b]$ 在通常度量下紧.
\end{theorem}
\begin{proof}
若 $a=b$, 这是单点集, 已知紧. 以下设 $a<b$. 设 $\{U_i\}_{i\in I}$ 是 $[a,b]$ 的开覆盖, 并令
\[
 S=\{t\in[a,b]:[a,t]\text{ 可由这族覆盖的有限多个成员覆盖}\}.
\]
$a\in S$, 因为某个 $U_i$ 包含 $a$. 集合 $S$ 非空且上界为 $b$, 故由实数的上确界性质存在 $c=\sup S\in[a,b]$.

取覆盖中包含 $c$ 的成员 $U_{i_0}$. 由开放性, 存在 $\delta>0$ 使 $(c-\delta,c+\delta)\subseteq U_{i_0}$. 上确界的性质保证可以取 $s\in S$ 且 $s>c-\delta/2$: 当 $c=a$ 时取 $s=a$, 当 $c>a$ 时用上确界的逼近性质. 已有有限多个成员覆盖 $[a,s]$, 再加入 $U_{i_0}$, 就覆盖 $[a,c]$. 因此 $c\in S$. 这一步说明上确界本身也能被有限覆盖, 不能只由 $c=\sup S$ 直接断言这一点.

若 $c<b$, 令 $t=\min\{b,c+\delta/2\}$, 则 $t>c$. 仍用覆盖 $[a,s]$ 的有限子族并加入 $U_{i_0}$, 就覆盖 $[a,t]$, 所以 $t\in S$, 与 $c$ 是上界矛盾. 因而 $c=b$. 前一段已证明 $c\in S$, 故 $b\in S$, 即整个 $[a,b]$ 有有限子覆盖.
\end{proof}

高维时, 不再有单一的最右端点可供推进. 但可以把方盒不断二分. 如果原覆盖没有有限子覆盖, 那么每一层总有一个小方盒也没有有限子覆盖. 小方盒不断缩小, 最后集中到一个点; 包含这个点的开集就会覆盖整个足够小的方盒, 产生矛盾. 为了让这个论证完整, 先证明其中的实数存在性.

\begin{lemma}[嵌套方盒]
\label{ch:eucompact:nested-boxes}
设
\[
 Q_k=\prod_{j=1}^n[a_{k,j},b_{k,j}],\qquad k\geq1,
\]
都是非空闭方盒, 且 $Q_{k+1}\subseteq Q_k$. 则 $\bigcap_{k\geq1}Q_k\ne\varnothing$. 若还满足
$\max_j(b_{k,j}-a_{k,j})\to0$, 则这个交集恰好只有一个点.
\end{lemma}
\begin{proof}
对每个坐标 $j$, 嵌套关系给出 $a_{k,j}$ 单调不减、$b_{k,j}$ 单调不增, 且 $a_{k,j}\leq b_{1,j}$. 定义 $x_j=\sup_{k\geq1}a_{k,j}$. 对任意固定 $k$, 有 $a_{k,j}\leq x_j$. 又对所有 $\ell\geq k$ 都有 $a_{\ell,j}\leq b_{k,j}$; 对 $\ell<k$ 则有 $a_{\ell,j}\leq a_{k,j}\leq b_{k,j}$. 所以 $b_{k,j}$ 是所有 $a_{\ell,j}$ 的上界, 从而 $x_j\leq b_{k,j}$. 因此 $x=(x_1,\ldots,x_n)$ 属于每个 $Q_k$.

若 $x,y$ 均属于交集, 对每个 $k,j$ 都有 $|x_j-y_j|\leq b_{k,j}-a_{k,j}$. 当各边长趋于零时, 令 $k\to\infty$ 得到 $x_j=y_j$ 对所有 $j$ 成立, 故 $x=y$.
\end{proof}

\begin{theorem}[闭方盒的紧性]
\label{ch:eucompact:box-compact}
每个非空闭方盒 $Q=\prod_{j=1}^n[a_j,b_j]$ 在 $\mathbb{R}^n$ 的欧氏度量下紧.
\end{theorem}
\begin{proof}
设 $\{U_i\}_{i\in I}$ 是 $Q$ 的开覆盖, 却没有有限子覆盖. 把每个坐标区间从中点二分, 就把 $Q$ 表为至多 $2^n$ 个闭子方盒的并. 至少一个子方盒没有有限子覆盖; 否则把各子方盒的有限子覆盖合在一起, 就会得到 $Q$ 的有限子覆盖. 选定这样的子方盒为 $Q_1$, 再将 $Q_1$ 二分并重复此过程. 因此得到
\[
 Q\supseteq Q_1\supseteq Q_2\supseteq\cdots,
\]
其中每个 $Q_k$ 都没有有限子覆盖, 第 $j$ 个坐标的边长为 $(b_j-a_j)2^{-k}$. 允许原来某些边长为零, 这不影响二分和论证.

由嵌套方盒引理, 存在唯一的 $x\in\bigcap_{k\geq1}Q_k$. 取覆盖中包含 $x$ 的 $U_{i_0}$, 并取 $r>0$ 使欧氏球 $B_2(x,r)\subseteq U_{i_0}$. 对每个 $y\in Q_k$, 因为 $x$ 也在 $Q_k$ 中,
\[
 \|y-x\|_2\leq
 2^{-k}\left(\sum_{j=1}^n(b_j-a_j)^2\right)^{1/2}.
\]
右端随 $k$ 趋于零. 因而当 $k$ 足够大时, 整个 $Q_k$ 都包含在 $B_2(x,r)\subseteq U_{i_0}$ 中. 此时单个 $U_{i_0}$ 就是 $Q_k$ 的有限子覆盖, 矛盾. 所以原开覆盖必有有限子覆盖.
\end{proof}

\begin{theorem}[Heine--Borel 定理]
\label{ch:eucompact:heine-borel}
设 $K\subseteq\mathbb{R}^n$, 其中使用通常欧氏度量. 则 $K$ 紧当且仅当 $K$ 在 $\mathbb{R}^n$ 中闭且有界.
\end{theorem}
\begin{proof}
紧性推出闭且有界已由定理 \ref{ch:eucompact:compact-closed-bounded} 证明. 反过来, 空集紧. 对非空闭有界集 $K$, 可取 $R>0$ 使 $K\subseteq[-R,R]^n$. 该方盒由定理 \ref{ch:eucompact:box-compact} 紧, 而 $K$ 在方盒中闭. 命题 \ref{ch:eucompact:closed-subset} 遂给出 $K$ 紧.
\end{proof}

\begin{example}
在 $\mathbb{R}^n$ 中, 判断单位闭球、单位球面, 以及从单位闭球中删去原点后的集合是否紧.
\end{example}
\begin{solution}
闭球 $\{x:\|x\|_2\leq1\}$ 与单位球面 $\{x:\|x\|_2=1\}$ 都是闭有界集, 因而紧. 删去原点后的集合仍有界, 但不再闭: 点列 $(1/k,0,\ldots,0)$ 在该集合中而趋于原点. 由紧集必闭, 删点后的集合不紧. 集合名字中的 ``闭球'' 指原来的球, 不能自动把后来删点得到的集合也视为闭集.
\end{solution}

\originalsection{有界点列与有限维范数等价}
\label{ch:eucompact:bolzano-weierstrass}

开覆盖紧性处理的是任意多的开集, 而点列常常更容易直接构造. 二分方盒不仅能从开覆盖中导出有限子覆盖, 也能从有界点列中抽出收敛子列. 两个用途依赖相同的实数完备性, 不必让一个结论的证明预先假设另一个结论.

\begin{theorem}[Bolzano--Weierstrass 定理]
\label{ch:eucompact:bw}
$\mathbb{R}^n$ 中的每个有界点列都有在 $\mathbb{R}^n$ 中收敛的子列.
\end{theorem}
\begin{proof}
设 $(x_m)_{m\geq1}$ 有界, 选闭方盒 $Q_0=[-R,R]^n$ 包含全部点列项. 二分 $Q_0$ 后得到有限多个子方盒. 至少一个子方盒含有无穷多个点列项的指标, 否则有限个有限指标集合的并不可能包含全部正整数. 选这样的方盒为 $Q_1$. 再二分 $Q_1$, 并在仍有无穷多个指标的子方盒中选择 $Q_2$. 递归得到嵌套方盒 $Q_k$, 其各边长为 $2R\,2^{-k}$, 而每个指标集合
\[
 I_k=\{m\geq1:x_m\in Q_k\}
\]
都无限.

由嵌套方盒引理, 交集恰好为一个点 $x$. 先取 $m_1\in I_1$, 再逐次取 $m_k\in I_k$ 且 $m_k>m_{k-1}$. 这是可行的, 因为无限的正整数集合不可能有上界. 严格递增保证所选的是原点列的子列. 因为 $x_{m_k}$ 与 $x$ 都在 $Q_k$ 中,
\[
 \|x_{m_k}-x\|_2\leq2R\sqrt n\,2^{-k}\longrightarrow0.
\]
所以 $x_{m_k}\to x$.
\end{proof}

\begin{definition}[序列紧]
度量空间的子集 $A$ 称为序列紧, 如果每个点列 $(x_m)$, $x_m\in A$, 都有一个子列收敛到 $A$ 中的某个点. 这里极限属于 $A$ 是定义的一部分.
\end{definition}

\begin{theorem}[欧氏空间中的序列紧性]
\label{ch:eucompact:sequential-hb}
对 $A\subseteq\mathbb{R}^n$, 以下三项等价: $A$ 紧; $A$ 在 $\mathbb{R}^n$ 中闭且有界; $A$ 序列紧.
\end{theorem}
\begin{proof}
前两项由 Heine--Borel 定理等价. 若 $A$ 闭且有界, 对 $A$ 中的任意点列使用 Bolzano--Weierstrass, 得到在 $\mathbb{R}^n$ 中收敛的子列. 闭性保证这个极限仍在 $A$ 中, 所以 $A$ 序列紧.

反过来, 设 $A$ 序列紧. 若 $x\in\overline A$, 对每个 $m$ 选 $a_m\in A\cap B_2(x,1/m)$. 于是 $a_m\to x$. 序列紧性给出一个子列收敛到某个 $a\in A$. 这个子列也收敛到 $x$, 由度量空间中极限的唯一性得 $x=a\in A$. 因此 $A$ 包含所有闭包点, 所以闭. 若 $A$ 无界, 可以选 $a_m\in A$ 使 $\|a_m\|_2>m$. 对任意严格递增的子列指标 $m_k$, 都有 $\|a_{m_k}\|_2>m_k\to\infty$, 所以该子列不可能收敛: 收敛于 $a$ 会使尾项的长度不超过 $\|a\|_2+1$. 这与序列紧性矛盾. 因而 $A$ 也有界.
\end{proof}

本章的依赖关系至此已经清楚: 实数上确界性质给出嵌套方盒; 嵌套方盒分别给出闭方盒紧性与 Bolzano--Weierstrass; 前者结合闭子集命题给出 Heine--Borel, 后者结合闭性给出序列紧. 在一般度量空间中, ``闭且有界'' 无法替代紧性, 但紧与序列紧仍等价. 下一章将另作证明, 而不把 Heine--Borel 当成一般空间的工具.

\begin{definition}[等价范数]
同一向量空间上的范数 $\|\cdot\|_a$ 与 $\|\cdot\|_b$ 称为等价, 如果存在与 $v$ 无关的常数 $c,C>0$, 使
\[
 c\|v\|_a\leq\|v\|_b\leq C\|v\|_a
 \qquad(v\in V).
\]
\end{definition}

\begin{theorem}[有限维范数等价]
\label{ch:eucompact:norm-equivalence}
$\mathbb{R}^n$ 上的任意范数 $\|\cdot\|$ 都与欧氏范数等价. 因而同一有限维实向量空间上的任意两个范数等价.
\end{theorem}
\begin{proof}
设 $n\geq1$, 并令 $e_1,\ldots,e_n$ 是标准基. 对任意 $x=\sum_jx_je_j$, 三角不等式、齐次性与 Cauchy--Schwarz 给出
\[
 \|x\|\leq\sum_{j=1}^n|x_j|\|e_j\|
 \leq C\|x\|_2,\qquad
 C=\left(\sum_{j=1}^n\|e_j\|^2\right)^{1/2}>0.
\]
因此
\[
 \bigl|\|x\|-\|y\|\bigr|\leq\|x-y\|\leq C\|x-y\|_2,
\]
即 $x\mapsto\|x\|$ 对欧氏度量连续. 上界容易取得, 困难在于不能让单位向量的范数任意接近零.

若没有正数 $c$ 使所有欧氏单位向量 $u$ 满足 $\|u\|\geq c$, 则对每个 $m$ 可选 $u_m$ 使 $\|u_m\|_2=1$ 而 $\|u_m\|<1/m$. 单位球面是欧氏空间中的闭有界集, 由刚证明的序列紧性, 有子列 $u_{m_k}\to u$ 且 $\|u\|_2=1$. 上一段的连续性又给出 $\|u\|=\lim_k\|u_{m_k}\|=0$, 与范数的正定性矛盾. 故这样的 $c>0$ 必定存在.

对 $x\ne0$, 向量 $u=x/\|x\|_2$ 在单位球面上. 齐次性给出
$\|x\|=\|x\|_2\|u\|\geq c\|x\|_2$; 对 $x=0$ 也成立. 所以
\[
 c\|x\|_2\leq\|x\|\leq C\|x\|_2.
\]
若 $V$ 为 $n$ 维, 选定一组基把 $V$ 与 $\mathbb{R}^n$ 线性对应. $V$ 上的范数通过这一对应成为 $\mathbb{R}^n$ 上的范数. 两个范数各自与欧氏范数比较, 再合并常数, 就得到它们互相等价. 若维数为零, $V=\{0\}$, 任取 $c=C=1$ 即可.
\end{proof}

\begin{corollary}[有限维空间中的拓扑与紧性]
\label{ch:eucompact:norm-topology}
有限维向量空间上的等价范数诱导相同的开集、相同的收敛点列与 Cauchy 点列, 也给出相同的有界集与紧集. 对其中任一范数的度量, 闭有界集均紧.
\end{corollary}
\begin{proof}
若 $c\|v\|_a\leq\|v\|_b\leq C\|v\|_a$, 则对每个 $x$ 与 $r>0$,
\[
 B_a(x,r/C)\subseteq B_b(x,r)\subseteq B_a(x,r/c).
\]
这些包含关系使任一度量的每个开邻域都包含另一度量的某个开球, 所以开集相同. 范数比较应用于 $x_m-x$、$x_m-x_\ell$ 以及 $x-p$, 分别给出收敛、Cauchy 性与有界性的等价. 紧性只由开集和有限子覆盖定义, 因而也相同. 最后选定基, 把闭有界集转为欧氏空间中的闭有界集, 再用 Heine--Borel 即得结论.
\end{proof}

\originalsection{函数范数与紧支撑}
\label{ch:eucompact:support}

范数不只衡量有限维向量. 连续函数也组成向量空间, 不同范数可以突出函数的不同特征. 这时首先要检查定义的数值是否有限, 而不能仅验证三角不等式.

\begin{example}
在 $C([0,1])$ 上, 定义
\[
 \|f\|_\infty=\sup_{x\in[0,1]}|f(x)|,
 \qquad \|f\|_1=\int_0^1|f(x)|\,\mathrm{d}x.
\]
证明两者都是范数. 本例假设读者掌握闭区间上连续函数的 Riemann 积分.
\end{example}
\begin{solution}
先说明数值有限. 连续函数 $f$ 的集合
$U_m=\{x\in[0,1]:|f(x)|<m\}$ 在 $[0,1]$ 中开, 且随 $m\geq1$ 覆盖整个区间. 闭区间紧, 所以有限多个 $U_m$ 已覆盖它. 由于 $U_m$ 随 $m$ 增大, 其中指标最大的 $U_M$ 就包含整个区间, 即 $|f(x)|<M$ 对所有 $x$ 成立. 因此上确界范数有限. 连续函数 $|f|$ Riemann 可积, 且积分也有限.

对上确界范数, 正定性来自 $\sup|f|=0$ 当且仅当各点上 $f=0$. 齐次性由 $\sup|\lambda f|=|\lambda|\sup|f|$ 得到. 点态不等式 $|f(x)+g(x)|\leq|f(x)|+|g(x)|$ 在取上确界后给出
\[
 \|f+g\|_\infty\leq\|f\|_\infty+\|g\|_\infty.
\]

对积分范数, 非负性、齐次性与三角不等式分别来自积分的非负性、常数因子法则与点态三角不等式. 唯一需要多解释的是正定性. 若 $f(x_0)\ne0$, 由连续性, 在某个正长度的子区间 $J\subseteq[0,1]$ 上有 $|f(x)|\geq|f(x_0)|/2$. 若 $x_0$ 是端点, 就取朝区间内部的单侧子区间. 因而
\[
 \int_0^1|f(x)|\,\mathrm{d}x
 \geq\frac{|f(x_0)|}{2}\,|J|>0.
\]
所以积分等于零只能发生在 $f$ 恒为零时.
\end{solution}

在 $\mathbb{R}$ 上直接写 $\sup_{x\in\mathbb{R}}|f(x)|$ 或 $\int_{\mathbb{R}}|f(x)|\,\mathrm{d}x$, 不一定得到有限数. 例如 $f(x)=x$ 无界, 而常值函数 $1$ 的绝对值积分发散. 选择合适的函数空间因此属于定义的一部分. 对连续实值函数, 紧支撑是保证这两个数值有限的一种充分条件.

\begin{definition}[支撑与紧支撑]
设 $f:X\to\mathbb{R}$, 其中 $X$ 是度量空间. 它的支撑定义为
\[
 \operatorname{supp}f=
 \overline{\{x\in X:f(x)\ne0\}}^{\,X}.
\]
若 $\operatorname{supp}f$ 紧, 则称 $f$ 有紧支撑. 全体有紧支撑的实值连续函数记作 $C_c(X)$.
\end{definition}

取闭包是必要的. 函数在某个点恰好为零, 该点仍可能被非零点逼近, 于是仍属于支撑. 对连续函数, 非零点集 $f^{-1}(\mathbb{R}\setminus\{0\})$ 是开集; 支撑则始终是闭集, 两者通常不同.

\begin{proposition}[紧支撑的等价表述]
\label{ch:eucompact:compact-support}
函数 $f:X\to\mathbb{R}$ 有紧支撑, 当且仅当存在紧集 $K\subseteq X$ 使 $f$ 在 $X\setminus K$ 上恒为零. 特别地, 在 $\mathbb{R}^n$ 上, 这等价于存在 $R>0$ 使 $f$ 在 $[-R,R]^n$ 外恒为零.
\end{proposition}
\begin{proof}
若支撑紧, 就取 $K=\operatorname{supp}f$. 不在支撑中的点当然不在非零点集中, 所以函数值为零. 反过来, 若非零点集包含在紧集 $K$ 中, 由紧集闭, 它的闭包仍包含在 $K$ 中. 支撑作为 $K$ 的闭子集紧.

在 $\mathbb{R}^n$ 中, 紧集有界, 因而包含在某个 $[-R,R]^n$ 中. 反过来, 若函数在此方盒外为零, 该方盒紧, 再用前一结论即可.
\end{proof}

\begin{example}
在 $\mathbb{R}$ 上定义 $f(x)=\max\{1-|x|,0\}$. 求它的支撑, 并判断它与零函数是否有紧支撑.
\end{example}
\begin{solution}
$f$ 的非零点集是 $(-1,1)$, 支撑为 $[-1,1]$, 所以它有紧支撑. 尽管 $f(-1)=f(1)=0$, 这两个点仍在支撑中. 零函数的支撑是空集, 所以零函数也有紧支撑.
\end{solution}

\begin{example}
判断函数 $g(x)=e^{-|x|}$ 是否有紧支撑, 并计算其绝对值在 $\mathbb{R}$ 上的积分.
\end{example}
\begin{solution}
$g$ 在整个 $\mathbb{R}$ 上均非零, 因而支撑是 $\mathbb{R}$, 不紧. 但
\[
 \int_{-\infty}^{\infty}|g(x)|\,\mathrm{d}x
 =2\int_0^\infty e^{-x}\,\mathrm{d}x=2.
\]
所以对连续实值函数, 紧支撑是绝对可积的充分条件, 不是必要条件. 对 $f\in C_c(\mathbb{R})$, 其支撑包含在某个闭区间中, 连续性保证区间上有界, 区间外又为零, 因而上确界与绝对值积分都有限. 任意函数的支撑仍可定义, 但若没有连续性或其他可积性条件, 单有紧支撑并不能保证函数有界或可积.
\end{solution}

\originalsection{编者练习与解答}
\label{ch:eucompact:exercises}

本节各题由编者为本章知识组织, 不声称来自 MIT 原习题或新领军真题. 做题时先确定环境空间, 再判断闭性、有界性与极限是否留在集合中.

\begin{exercise}
\label{ch:eucompact:ex-finite-union}
证明度量空间中有限多个紧集的并紧. 再说明无限多个紧集的并未必紧.
\end{exercise}
\begin{solution}
设 $K=K_1\cup\cdots\cup K_m$, 各 $K_j$ 紧. $K$ 的任意开覆盖也是每个 $K_j$ 的开覆盖, 所以对每个 $j$ 可取有限子族覆盖 $K_j$. 把这 $m$ 个有限子族合并, 仍只有有限多个成员, 并覆盖 $K$. 故 $K$ 紧. 对无限并, $\mathbb{R}=\bigcup_{m\geq1}[-m,m]$, 各闭区间紧, 但 $\mathbb{R}$ 不紧.
\end{solution}

\begin{exercise}
\label{ch:eucompact:ex-rational}
判断下列集合在 $\mathbb{R}$ 中是否紧, 并说明理由:
\[
 A=\{1/m:m\geq1\},\qquad
 B=A\cup\{0\},\qquad
 C=\mathbb{Q}\cap[0,1].
\]
\end{exercise}
\begin{solution}
$A$ 有界, 但点列 $1/m\to0$ 且 $0\notin A$, 所以 $A$ 在 $\mathbb{R}$ 中不闭, 因而不紧. $B$ 有界且闭. 为核对闭性, 若 $x\notin B$, 则 $x\ne0$. 在不含 $0$ 的一个小邻域中, 只有有限多个 $1/m$ 可能出现; 再缩小半径避开这有限个点, 就得到与 $B$ 不相交的邻域. 故补集开, 从而 $B$ 紧. $C$ 有界但不闭, 因为可取有理数列趋于 $[0,1]$ 内的无理数, 例如 $\sqrt2/2$. 极限不在 $C$ 中, 所以 $C$ 不紧.
\end{solution}

\begin{exercise}
\label{ch:eucompact:ex-norms-infinite}
在 $C([0,1])$ 上证明 $\|f\|_1\leq\|f\|_\infty$. 是否存在常数 $c>0$ 使所有连续函数都满足 $c\|f\|_\infty\leq\|f\|_1$?
\end{exercise}
\begin{solution}
因为区间长度为 $1$, 积分估计给出
$\int_0^1|f(x)|\,\mathrm{d}x\leq\int_0^1\|f\|_\infty\,\mathrm{d}x=\|f\|_\infty$.
对 $m\geq1$, 令 $f_m(x)=\max\{1-mx,0\}$. 则 $f_m$ 连续, $\|f_m\|_\infty=1$, 而
\[
 \|f_m\|_1=\int_0^{1/m}(1-mx)\,\mathrm{d}x=\frac1{2m}.
\]
若存在这样的 $c$, 就有 $c\leq1/(2m)$ 对所有 $m$ 成立, 矛盾. 两个范数因此不等价. 有限维范数等价定理中的维数条件不能删去.
\end{solution}

\begin{exercise}
\label{ch:eucompact:ex-subspace}
在度量空间 $X=(0,1)$ 中, 判断 $A=[1/3,2/3]$ 与 $D=(0,1/2]$ 是否紧. 说明 ``在 $X$ 中闭且有界'' 为什么不能直接决定答案.
\end{exercise}
\begin{solution}
$A$ 在 $\mathbb{R}$ 中闭且有界, 因而紧; 子空间命题说明它作为 $X$ 的子集仍紧. $D$ 在 $X$ 中闭, 因为 $X\setminus D=(1/2,1)$ 在 $X$ 中开; 它也有界. 但取 $x_m=1/(m+2)\in D$, 任何子列在 $\mathbb{R}$ 中都趋于 $0$, 因而不可能在 $X$ 中收敛. 或直接取 $X$ 中开集 $V_m=X\cap(1/m,1)$ 覆盖 $D$, 有限个的并漏掉充分小的正数. 所以 $D$ 不紧. Heine--Borel 的闭性要求是在整个欧氏空间中闭, 不是仅在任意子空间中闭.
\end{solution}

\begin{exercise}
\label{ch:eucompact:ex-support}
设 $f,g\in C_c(X)$, 其中 $X$ 是度量空间. 证明 $f+g$ 与 $fg$ 均有紧支撑, 并给出它们的支撑与原支撑之间的包含关系.
\end{exercise}
\begin{solution}
若 $f(x)+g(x)\ne0$, 则至少一个函数在 $x$ 非零, 所以
\[
 \{f+g\ne0\}\subseteq\operatorname{supp}f\cup\operatorname{supp}g.
\]
右端是两个闭集的并, 因而闭. 取闭包后得
$\operatorname{supp}(f+g)\subseteq\operatorname{supp}f\cup\operatorname{supp}g$.
右端由有限并练习紧, 所以左端作为其闭子集紧. 同理, $fg\ne0$ 要求 $f\ne0$ 且 $g\ne0$, 故
\[
 \operatorname{supp}(fg)\subseteq
 \operatorname{supp}f\cap\operatorname{supp}g.
\]
右端为 $\operatorname{supp}f$ 的闭子集, 因而紧, 左端亦紧. 加法时可能有抵消, 所以第一条通常只能写包含关系, 不能写等号.
\end{solution}

\originalchapter{一般度量空间的紧性}
\label{ch:compact}

本章以 Paige Bright 的 MIT OCW 18.S190 IAP 2023 第 4 讲为主线重编. 核心是把开覆盖、收敛子列和 Cauchy 点列三种语言联系起来, 并用紧性证明连续像、极值与非空交结论. 编者补写递归抽取的指标选择、空集约定及紧集间的距离, 修正原稿中同一半径球内两点的距离估计. 一致连续与完备化分别留待后续章节.
\emph{来源定位: 第 4 讲正文第 1--2 页为 Lebesgue 数、全有界与紧性的等价, 第 3--4 页为连续像、极值、有限交性质及完备性; 本章补充一般度量空间版本与完整证明. 页码采用原讲义印刷页码.}

\originalsection{紧、序列紧与全有界}
\label{ch:compact:three-languages}

上一章中, 闭性使有界点列的极限留在集合中, 而有界性配合有限维坐标保证能抽出收敛子列. 在一般度量空间中, 坐标方盒未必存在, 有界性也可能太弱. 我们需要一种更细的控制: 不仅把全部点装进一个大球, 而且能在每一个指定尺度上, 用有限多个小球覆盖全部点.

本章把 $X$ 本身作为度量空间讨论. 对子集 $K\subseteq X$, 所有结论都可应用于限制度量后的空间 $K$; 其中的球就是 $B_K(x,r)=K\cap B_X(x,r)$.

\begin{definition}[序列紧、全有界与完备]
度量空间 $(X,d)$ 称为序列紧, 如果其中的每个点列都有收敛到 $X$ 中某点的子列. 称 $X$ 全有界, 如果对每个 $\varepsilon>0$, 存在有限多个点 $p_1,\ldots,p_m\in X$, 使
\[
 X=\bigcup_{j=1}^mB_X(p_j,\varepsilon).
\]
这样的有限点集称为一个有限 $\varepsilon$-网, 这里使用严格小于 $\varepsilon$ 的开球约定.

点列 $(x_n)$ 称为 Cauchy 点列, 如果对每个 $\varepsilon>0$, 存在 $N$, 使 $m,n\geq N$ 时 $d(x_m,x_n)<\varepsilon$. 若 $X$ 中每个 Cauchy 点列都收敛到 $X$ 内的点, 则称 $X$ 完备.
\end{definition}

空空间按定义紧、序列紧、全有界且完备. 全有界性在空空间中用空的有限网即可. 下面在需要选取点时可直接排除这一平凡情况.

\begin{example}
在无限集合 $X$ 上定义离散度量
\[
 d(x,y)=\begin{cases}0,&x=y,\\1,&x\ne y.\end{cases}
\]
判断 $X$ 是否有界、全有界、完备以及紧.
\end{example}
\begin{solution}
所有距离都不超过 $1$, 所以 $X$ 有界. 但当 $0<\varepsilon<1$ 时, 每个 $\varepsilon$-球都只是一个单点, 有限多个球不能覆盖无限集合. 因而 $X$ 不全有界, 也不紧: 全部单点组成开覆盖, 却没有有限子覆盖. 此空间完备, 因为 Cauchy 点列取 $\varepsilon=1/2$ 后必须最终恒定. 所以完备与有界合起来仍不足以得到紧性.
\end{solution}

\begin{proposition}[全有界的基本性质]
\label{ch:compact:total-basic}
全有界空间必有界, 它的任意子空间也全有界. 有限维欧氏空间中的每个有界子集全有界.
\end{proposition}
\begin{proof}
若 $X$ 非空且全有界, 取有限 $1$-网 $p_1,\ldots,p_m$. 对任意 $x\in X$, 选 $j$ 使 $d(x,p_j)<1$. 则
\[
 d(x,p_1)<1+\max_{1\leq j\leq m}d(p_j,p_1),
\]
所以 $X$ 包含在某个以 $p_1$ 为中心的有限半径球中.

设 $A\subseteq X$, 要在 $A$ 中选有限 $\varepsilon$-网. 先取 $X$ 的有限 $\varepsilon/2$-网. 对每个与 $A$ 相交的球 $B_X(p_j,\varepsilon/2)$, 从交集中选 $a_j\in A$. 若 $a\in A$ 落在这个球中, 则
$d(a,a_j)\leq d(a,p_j)+d(p_j,a_j)<\varepsilon$.
因此有限多个 $B_A(a_j,\varepsilon)$ 覆盖 $A$. 这一步需要移动球心, 因为原来网的球心未必属于 $A$.

最后, 设 $A\subseteq\mathbb{R}^n$ 有界, 把它装进某个方盒. 对给定 $\varepsilon>0$, 把方盒分成有限个边长小于 $\varepsilon/\sqrt n$ 的小方盒. 从每个与 $A$ 相交的小方盒中选一个 $A$ 中的点. 同一小方盒内任意两点的欧氏距离小于 $\varepsilon$, 所选点便构成有限 $\varepsilon$-网. 空集的情形均用空网处理.
\end{proof}

\begin{lemma}[Cauchy 点列中的收敛子列]
\label{ch:compact:cauchy-subsequence}
收敛点列是 Cauchy 点列. 若一个 Cauchy 点列 $(x_n)$ 有子列 $x_{n_k}\to x\in X$, 则整个点列都收敛到 $x$.
\end{lemma}
\begin{proof}
若 $x_n\to x$, 对给定 $\varepsilon>0$, 当 $m,n$ 足够大时 $d(x_m,x)<\varepsilon/2$ 且 $d(x_n,x)<\varepsilon/2$, 从而 $d(x_m,x_n)<\varepsilon$.

再设 $(x_n)$ 为 Cauchy 点列且 $x_{n_k}\to x$. 取 $N$ 使 $m,n\geq N$ 时 $d(x_m,x_n)<\varepsilon/2$. 选一个固定的子列项 $x_{n_k}$, 使 $n_k\geq N$ 且 $d(x_{n_k},x)<\varepsilon/2$. 对每个 $n\geq N$, 都有
\[
 d(x_n,x)\leq d(x_n,x_{n_k})+d(x_{n_k},x)<\varepsilon.
\]
这对所有尾项成立, 因而 $x_n\to x$. 仅估计某一个 $x_N$ 到 $x$ 的距离不能代替这一步.
\end{proof}

\originalsection{紧性等价定理的证明}
\label{ch:compact:equivalence-proofs}

我们将先证明开覆盖紧性推出序列紧, 再证明序列紧推出全有界与完备. 反向的抽取依靠全有界, 让子列先成为 Cauchy 点列, 再由完备性得到极限. 最后使用 Lebesgue 数, 把序列紧重新转回有限子覆盖. 这样每一步只使用已经证明的结果.

\begin{theorem}[紧性推出序列紧]
\label{ch:compact:compact-to-sequential}
每个紧度量空间都序列紧.
\end{theorem}
\begin{proof}
设 $X$ 紧, 假设存在点列 $(x_n)$ 没有任何在 $X$ 中收敛的子列. 对每个 $x\in X$, 必存在 $r_x>0$, 使 $B_X(x,r_x)$ 只包含有限多个点列项的指标. 否则, 对每个 $k$ 球 $B_X(x,1/k)$ 都包含无穷多个指标, 我们便可逐次选取
\[
 n_1<n_2<\cdots,\qquad x_{n_k}\in B_X(x,1/k),
\]
于是 $x_{n_k}\to x$, 与假设矛盾.

全部 $B_X(x,r_x)$ 构成 $X$ 的开覆盖. 由紧性, 有有限多个这样的球覆盖 $X$. 每个所选球只包含有限多个点列项的指标, 所以它们的有限并只包含有限多个指标. 但它们覆盖 $X$, 就必须包含每个 $x_n$, 即必须包含全部正整数指标, 矛盾. 因而每个点列必有收敛子列.
\end{proof}

这里数的是指标, 不是不同点的个数. 如果某个点在点列中反复出现无穷次, 本来就有恒定的收敛子列. 上述证明统一处理了重复出现与各项不同两种情形, 不需要预先把点列改成两两不同.

\begin{proposition}[序列紧性推出全有界与完备]
\label{ch:compact:sequential-to-complete-total}
每个序列紧的度量空间都全有界且完备.
\end{proposition}
\begin{proof}
先证全有界. 若不全有界, 存在 $\varepsilon_0>0$, 使任何有限多个半径为 $\varepsilon_0$ 的球都不能覆盖 $X$. 任选 $x_1\in X$, 再递归选择
\[
 x_{n+1}\in X\setminus\bigcup_{j=1}^nB_X(x_j,\varepsilon_0).
\]
每一步的补集非空, 所以选择可行. 于是当 $i\ne j$ 时 $d(x_i,x_j)\geq\varepsilon_0$. 这个点列的任何子列都不是 Cauchy 点列, 因为不同两项之间的距离始终至少为 $\varepsilon_0$. 由引理 \ref{ch:compact:cauchy-subsequence}, 它也没有收敛子列, 与序列紧性矛盾.

再证完备. 任取 $X$ 中的 Cauchy 点列, 序列紧性保证它有收敛于 $X$ 中某点的子列. 同一引理说明原点列收敛到这个点. 因此每个 Cauchy 点列都在 $X$ 中收敛, 即 $X$ 完备.
\end{proof}

\begin{lemma}[全有界空间中的子列抽取]
\label{ch:compact:total-to-cauchy-subsequence}
全有界度量空间中的每个点列都有 Cauchy 子列.
\end{lemma}
\begin{proof}
设 $(x_n)$ 是 $X$ 中的点列. 对每个整数 $k\geq1$, 取一个有限的半径 $1/k$ 的球覆盖. 第一层覆盖中, 至少有一个球 $B_X(p_1,1)$ 含有无穷多个点列项的指标. 记这些指标组成的集合为 $I_1$.

已构造无限指标集合 $I_{k-1}$ 后, 用第 $k$ 层的有限球覆盖这些指标对应的点. 至少一个球 $B_X(p_k,1/k)$ 含有无穷多个来自 $I_{k-1}$ 的指标. 令
\[
 I_k=\{n\in I_{k-1}:x_n\in B_X(p_k,1/k)\}.
\]
因此 $I_1\supseteq I_2\supseteq\cdots$, 且每个 $I_k$ 都无限. 这里 $p_k$ 是第 $k$ 层有限覆盖中选出的球心, 不要求各层球彼此包含.

从 $I_k$ 中依次选取 $n_k>n_{k-1}$, 得到真正的子列 $(x_{n_k})$. 对任意固定的 $N$, 若 $k,\ell\geq N$, 则 $n_k\in I_k\subseteq I_N$ 且 $n_\ell\in I_\ell\subseteq I_N$, 所以两点都在 $B_X(p_N,1/N)$ 内. 故
\begin{equation}
 d(x_{n_k},x_{n_\ell})
 \leq d(x_{n_k},p_N)+d(p_N,x_{n_\ell})<\frac2N.
 \label{ch:compact:two-over-n}
\end{equation}
给定 $\varepsilon>0$, 取 $N$ 使 $2/N<\varepsilon$, 就证明该子列为 Cauchy 点列. 无穷多个指标保证每次能选比前一次更大的指标; 单纯在各个球中随便取点, 不足以得到原点列的子列.
\end{proof}

\begin{corollary}[完备且全有界推出序列紧]
\label{ch:compact:complete-total-to-sequential}
若 $X$ 完备且全有界, 则 $X$ 序列紧.
\end{corollary}
\begin{proof}
任取点列. 由上一引理, 它有 Cauchy 子列; 由完备性, 这个子列收敛到 $X$ 中的某点. 这正是序列紧的定义.
\end{proof}

剩下的问题是怎样从收敛子列推出有限子覆盖. 全有界性会给出有限个同半径的球, 但原来的覆盖成员可能大小不一. 我们需要证明: 对每个固定的开覆盖, 总能找到一个统一的正半径, 使以任意点为中心的这个半径的球都落进某个覆盖成员中. 统一的是半径, 而允许使用的覆盖成员随中心改变.

\begin{lemma}[Lebesgue 数引理]
\label{ch:compact:lebesgue}
设 $X$ 是序列紧度量空间, $\mathcal U=\{U_i\}_{i\in I}$ 是 $X$ 的开覆盖. 则存在 $r>0$, 使对每个 $x\in X$, 都有某个 $i\in I$ 满足
\[
 B_X(x,r)\subseteq U_i.
\]
这样的 $r$ 称为这个覆盖的一个 Lebesgue 数, 这里采用球包含版本的定义.
\end{lemma}
\begin{proof}
空空间中取任意 $r>0$ 即可. 设 $X\ne\varnothing$. 若结论不成立, 则对每个 $r>0$, 都能找到一个中心 $x\in X$, 使这个半径的球不包含于任何 $U_i$. 对 $r=1/n$, 选取这样的 $x_n$, 即
\[
 B_X(x_n,1/n)\not\subseteq U_i
 \qquad\text{对每个 }i\in I.
\]
由序列紧性, 有子列 $x_{n_k}\to x\in X$. 取 $i_0$ 使 $x\in U_{i_0}$; 由开放性, 存在 $\rho>0$ 使 $B_X(x,\rho)\subseteq U_{i_0}$. 当 $k$ 足够大时, 同时有
\[
 d(x_{n_k},x)<\rho/2,\qquad 1/n_k<\rho/2.
\]
若 $y\in B_X(x_{n_k},1/n_k)$, 则
\[
 d(y,x)\leq d(y,x_{n_k})+d(x_{n_k},x)
 <1/n_k+\rho/2<\rho.
\]
因此 $B_X(x_{n_k},1/n_k)\subseteq B_X(x,\rho)\subseteq U_{i_0}$, 矛盾. 故所需的 $r$ 存在.
\end{proof}

\begin{remark}[另一种 Lebesgue 数约定]
有的书把 Lebesgue 数定义为正数 $\lambda$, 使每个直径小于 $\lambda$ 的非空子集都包含在某个覆盖成员中. 球包含版本给出这个版本: 若 $\operatorname{diam}A<r$ 且 $a\in A$, 则 $A\subseteq B_X(a,r)$. 反过来, 若使用直径版本的 $\lambda$, 则 $B_X(x,\lambda/3)$ 的直径至多 $2\lambda/3<\lambda$, 从而得到球包含版本. 两种版本表达相同的统一尺度思想, 使用时要注意半径与直径的因子.
\end{remark}

\begin{theorem}[序列紧性推出紧性]
\label{ch:compact:sequential-to-compact}
每个序列紧度量空间都紧.
\end{theorem}
\begin{proof}
设 $\{U_i\}_{i\in I}$ 是 $X$ 的任意开覆盖. 由 Lebesgue 数引理取 $r>0$. 序列紧性已经推出全有界, 因而可取有限 $r$-网 $p_1,\ldots,p_m$. 对每个球 $B_X(p_j,r)$, 由 Lebesgue 数性质选 $i_j$ 使该球包含在 $U_{i_j}$ 中. 所以
\[
 X=\bigcup_{j=1}^mB_X(p_j,r)
 \subseteq\bigcup_{j=1}^mU_{i_j}.
\]
这就是原覆盖的有限子覆盖. 空空间直接用空子覆盖.
\end{proof}

\begin{theorem}[度量紧性的三个等价条件]
\label{ch:compact:equivalence}
对任意度量空间 $X$, 下列条件等价:
\begin{enumerate}
 \item $X$ 紧, 即每个开覆盖都有有限子覆盖.
 \item $X$ 序列紧, 即每个点列都有在 $X$ 中收敛的子列.
 \item $X$ 完备且全有界.
\end{enumerate}
\end{theorem}
\begin{proof}
定理 \ref{ch:compact:compact-to-sequential} 给出第一项推出第二项. 命题 \ref{ch:compact:sequential-to-complete-total} 给出第二项推出第三项. 推论 \ref{ch:compact:complete-total-to-sequential} 给出第三项推出第二项. 定理 \ref{ch:compact:sequential-to-compact} 给出第二项推出第一项. 因此三者等价, 且这些方向的证明没有使用一般空间中不成立的 Heine--Borel 性质.
\end{proof}

\begin{example}
比较区间 $(0,1)$ 与 $\mathbb{R}$ 的全有界性、完备性和紧性.
\end{example}
\begin{solution}
区间 $(0,1)$ 全有界, 因为它是有界的欧氏子集, 但不完备: Cauchy 点列 $1/(n+1)$ 的唯一实数极限为 $0$, 不属于此空间. 它因而不紧. $\mathbb{R}$ 完备, 但不全有界: 有限多个半径为 $1$ 的球的并仍有界, 不能覆盖 $\mathbb{R}$. 前面的无限离散空间则表明, 把全有界放宽成有界, 即使保留完备也不能得到紧性.
\end{solution}

\originalsection{有限交性质与嵌套紧集}
\label{ch:compact:finite-intersections}

开覆盖使用并集来表达 ``所有点都被照顾到''. 取补集后, 同一个事实可以改写成闭集交为空. 因而有限子覆盖与有限交之间有一个对偶关系: 在紧空间中, 若每次有限地施加闭条件都还有解, 就不能在施加全部条件时突然失去所有解.

\begin{definition}[有限交性质]
空间 $X$ 的子集族 $\{F_i\}_{i\in I}$ 称为具有有限交性质, 如果对每个有限指标集 $J\subseteq I$, 都有
\[
 \bigcap_{j\in J}F_j\ne\varnothing.
\]
约定没有任何集合参与时的交集为 $X$, 即 $\bigcap_{j\in\varnothing}F_j=X$. 因此这个定义也包含空的有限子族; 在非空空间中, 它不额外增加条件.
\end{definition}

\begin{theorem}[紧性的有限交刻画]
\label{ch:compact:fip}
度量空间 $X$ 紧, 当且仅当每一族具有有限交性质的、在 $X$ 中闭的子集 $\{F_i\}_{i\in I}$ 都满足
\[
 \bigcap_{i\in I}F_i\ne\varnothing.
\]
\end{theorem}
\begin{proof}
先设 $X$ 紧. 如果闭集族的总交为空, 则其开补集 $U_i=X\setminus F_i$ 覆盖 $X$. 紧性给出有限子覆盖 $U_{i_1},\ldots,U_{i_m}$, 从而
\[
 \bigcap_{j=1}^mF_{i_j}
 =X\setminus\bigcup_{j=1}^mU_{i_j}=\varnothing,
\]
与有限交性质矛盾.

反过来, 假设每个这样的闭集族均有非空总交. 若有开覆盖 $\{U_i\}$ 没有有限子覆盖, 则对每个有限 $J\subseteq I$,
\[
 \bigcap_{j\in J}(X\setminus U_j)
 =X\setminus\bigcup_{j\in J}U_j\ne\varnothing.
\]
所以闭集族 $F_i=X\setminus U_i$ 具有有限交性质, 按假设其总交非空. 但 $\{U_i\}$ 覆盖 $X$ 意味着这个总交为空, 矛盾. 故每个开覆盖都有有限子覆盖, 即 $X$ 紧.

若 $X=\varnothing$, 它紧, 且没有任何子集族满足本定义的有限交性质, 因为空的有限子族的交已经是空集. 所以右侧的全称条件也成立. 上述约定保证空空间不会成为例外.
\end{proof}

有限交刻画与定理 \ref{ch:compact:equivalence} 的三项一起, 就得到原课程中紧性的四种等价描述. 有限交性质必须配合闭性. 例如在紧空间 $[0,1]$ 中, 集合 $F_n=(0,1/n)$ 具有有限交性质, 却有空的总交; 它们都不在 $[0,1]$ 中闭.

\begin{theorem}[嵌套非空紧集的交]
\label{ch:compact:nested}
设 $X$ 是度量空间, 且
\[
 K_1\supseteq K_2\supseteq K_3\supseteq\cdots
\]
是一列非空紧子集. 则 $\bigcap_{n\geq1}K_n\ne\varnothing$.
\end{theorem}
\begin{proof}
每个 $K_n$ 都在 $X$ 中闭, 因而也在 $K_1$ 中闭. 任意有限个 $K_n$ 的交就是其中指标最大的那个集合, 所以非空. 把它们看成紧空间 $K_1$ 的闭子集族, 应用有限交刻画, 就得到总交非空.

也可以直接用序列紧性理解这个结论. 对每个 $n$ 取 $a_n\in K_n$. 因为 $a_n\in K_1$ 且 $K_1$ 紧, 存在同一个收敛子列 $a_{n_k}\to a\in K_1$. 对任意固定的 $m$, 当 $k$ 足够大时 $n_k\geq m$, 故 $a_{n_k}\in K_{n_k}\subseteq K_m$. $K_m$ 闭, 所以 $a\in K_m$. 这对每个 $m$ 都成立, 于是 $a$ 属于全部 $K_m$. 不需要为不同的 $m$ 另选不同的收敛子列, 也不能仅由另一个子列存在就断言它有同一个极限.
\end{proof}

\begin{remark}[直径条件与紧性条件]
此定理不要求 $\operatorname{diam}K_n\to0$. 若另有这个条件, 交集至多含一个点, 因为交集中的任意 $x,y$ 都满足 $d(x,y)\leq\operatorname{diam}K_n$ 对所有 $n$ 成立. 结合非空性, 此时交集恰好是单点. 后面的完备性章节会讨论另一种 Cantor 交集定理: 用完备空间、非空闭集和直径趋于零来替代这里各集合的紧性.
\end{remark}

\begin{example}
在 $\mathbb{R}$ 中取 $K_n=[n,\infty)$ 与 $H_n=(0,1/n)$, $n\geq1$. 分别求它们的总交, 并说明为什么不违背嵌套紧集的交定理.
\end{example}
\begin{solution}
$K_n$ 非空、闭且嵌套, 但不紧. 任意实数 $x$ 都不属于某个 $n>x$ 对应的 $K_n$, 所以总交为空. $H_n$ 非空、有界且嵌套, 直径甚至趋于零, 但不闭. 若 $x$ 在全部 $H_n$ 中, 则 $0<x<1/n$ 对所有 $n$ 成立, 不可能. 所以总交也为空. 两者都没有满足紧性假设, 这两个例子分别说明不能任意删去控制逃逸或保留极限的条件.
\end{solution}

\originalsection{连续像、最值与最小距离}
\label{ch:compact:continuous-images}

连续性用开集的逆像来描述, 因而特别适合与开覆盖定义相接. 若目标空间中的开集覆盖了函数像, 把它们拉回原空间就会覆盖定义域. 定义域紧时, 选出的有限子覆盖又可以送回目标空间.

\begin{theorem}[紧集的连续像]
\label{ch:compact:image}
设 $K$ 是紧度量空间, $Y$ 是度量空间, $f:K\to Y$ 连续. 则 $f(K)$ 是 $Y$ 的紧子集. 特别地, 若 $f:X\to Y$ 连续而 $K\subseteq X$ 紧, 则 $f(K)$ 紧.
\end{theorem}
\begin{proof}
设 $\{U_i\}_{i\in I}$ 是 $f(K)$ 的开覆盖, 各 $U_i$ 在 $Y$ 中开. 连续性保证 $V_i=f^{-1}(U_i)$ 在 $K$ 中开. 因为每个 $x\in K$ 的像落在某个 $U_i$ 中, 所以 $\{V_i\}$ 覆盖 $K$. 由紧性, 可选 $V_{i_1},\ldots,V_{i_m}$ 覆盖 $K$.

任取 $y\in f(K)$, 写为 $y=f(x)$, $x\in K$. 某个 $V_{i_j}$ 包含 $x$, 即 $f(x)\in U_{i_j}$. 因而 $U_{i_1},\ldots,U_{i_m}$ 覆盖 $f(K)$, 得到有限子覆盖. 注意证明只使用 $f(V_i)\subseteq U_i$, 并不需要、也一般没有 $f(V_i)=U_i$. 最后一项应用于限制映射 $f|_K$ 即可.
\end{proof}

\begin{theorem}[紧集上的极值定理]
\label{ch:compact:extreme-values}
设 $K$ 是非空紧度量空间, $f:K\to\mathbb{R}$ 连续. 则 $f$ 有界, 且存在 $x_{\min},x_{\max}\in K$ 使
\[
 f(x_{\min})\leq f(x)\leq f(x_{\max})
 \qquad(x\in K).
\]
\end{theorem}
\begin{proof}
连续像定理说明 $C=f(K)$ 是 $\mathbb{R}$ 中的非空紧集, 因而闭且有界. 设 $M=\sup C$ 与 $m=\inf C$, 它们都是有限实数. 对每个正整数 $n$, 由上确界定义选 $c_n\in C$ 使 $M-1/n<c_n\leq M$. 所以 $c_n\to M$, 由 $C$ 的闭性得 $M\in C$. 同理, 可选 $b_n\in C$ 使 $m\leq b_n<m+1/n$, 得 $m\in C$. 由像集的定义, 有 $x_{\max},x_{\min}\in K$ 分别满足 $f(x_{\max})=M$, $f(x_{\min})=m$. 这就证明最值被实际取得, 而不仅是上确界与下确界存在.
\end{proof}

\begin{example}
在 $(0,1)$ 上取连续函数 $f(x)=x$ 与 $g(x)=1/x$. 判断它们是否有界以及是否取得最值.
\end{example}
\begin{solution}
$f$ 有界, 上确界为 $1$, 下确界为 $0$, 但两个端点值都不在函数像中, 所以最大值与最小值均不存在. $g$ 无界, 因为 $g(1/n)=n$ 对 $n\geq2$ 成立. 它的下确界为 $1$, 不取得, 也没有最大值. 非紧并不表示每个连续函数都会失败: 常值函数仍有最大值与最小值. 极值定理给出的是非空紧定义域上所有连续实值函数均可取得最值的保证.
\end{solution}

\begin{definition}[点到集合、集合到集合的距离]
对非空子集 $A\subseteq X$, 定义
\[
 d(x,A)=\inf_{a\in A}d(x,a).
\]
若 $A,B$ 均非空, 定义
\[
 d(A,B)=\inf\{d(a,b):a\in A,\ b\in B\}.
\]
两者都是有限的非负实数. 这里不把空集的距离纳入定义, 以免在讨论最小值时混入扩展实数 $+\infty$.
\end{definition}

\begin{lemma}[距离函数的连续性]
\label{ch:compact:distance-function}
对非空 $A\subseteq X$, 有
\[
 |d(x,A)-d(y,A)|\leq d(x,y)
 \qquad(x,y\in X).
\]
因此 $x\mapsto d(x,A)$ 连续. 此外, $d(x,A)=0$ 当且仅当 $x\in\overline A$.
\end{lemma}
\begin{proof}
对任意 $a\in A$, 三角不等式给出 $d(x,A)\leq d(x,a)\leq d(x,y)+d(y,a)$. 对 $a$ 取下确界, 得 $d(x,A)\leq d(x,y)+d(y,A)$. 交换 $x,y$, 得到绝对值估计. 若 $d(x,A)=0$, 则对每个 $r>0$, 下确界定义保证某个 $a\in A$ 满足 $d(x,a)<r$, 所以每个球 $B_X(x,r)$ 都与 $A$ 相交, 即 $x\in\overline A$. 反过来, 若 $x\in\overline A$, 每个半径 $r$ 的球都包含 $A$ 中某点, 从而 $0\leq d(x,A)<r$ 对所有 $r>0$ 成立, 故距离为零.
\end{proof}

\begin{theorem}[紧集与闭集的正距离]
\label{ch:compact:positive-distance}
设 $A,B\subseteq X$ 非空且不相交. 若 $A$ 紧而 $B$ 在 $X$ 中闭, 则 $d(A,B)>0$. 特别地, 两个非空不相交紧集之间的距离为正. 若 $A,B$ 均紧, 还存在 $a_0\in A,b_0\in B$ 使
\[
 d(A,B)=d(a_0,b_0)>0.
\]
\end{theorem}
\begin{proof}
因为 $B$ 闭且 $A\cap B=\varnothing$, 对每个 $a\in A$ 都有 $a\notin\overline B$, 所以上一引理给出 $d(a,B)>0$. 连续函数 $a\mapsto d(a,B)$ 在非空紧集 $A$ 上取得最小值, 设取得点为 $a_0$. 因为函数在每个点都严格为正, 特别是在最小值取得点严格为正, 故
\[
 d(A,B)=\inf_{a\in A}d(a,B)=d(a_0,B)>0.
\]
第一处等号来自对全部配对 $(a,b)$ 取下确界与逐次取下确界的等价. 具体地, 任意全体配对距离的下界也是每个 $d(a,B)$ 的下界; 反过来, 任意全部 $d(a,B)$ 的下界又是每个配对距离 $d(a,b)$ 的下界. 因而两边有相同的下界集合, 下确界相同.

若 $B$ 也紧, 连续函数 $b\mapsto d(a_0,b)$ 在 $B$ 上取得最小值, 设取得点为 $b_0$. 则 $d(a_0,b_0)=d(a_0,B)=d(A,B)$, 从而同时证明距离被一对点取得.
\end{proof}

这个证明揭示两个不同作用: $B$ 的闭性使每个固定点 $a\notin B$ 与 $B$ 有正距离; $A$ 的紧性把这些各自可能很小的正数变成一个统一的正下界. 若两者都只闭, 点对可以一起向无穷远移动, 让距离趋于零.

\begin{example}
在 $\mathbb{R}^2$ 中取
\[
 A=\{(x,0):x\geq1\},\qquad
 B=\{(x,1/x):x\geq1\}.
\]
证明 $A,B$ 均闭且不相交, 但 $d(A,B)=0$. 再比较 $\mathbb{R}$ 中的 $\{0\}$ 与 $(0,1)$.
\end{example}
\begin{solution}
两者不相交. $A$ 闭; $B$ 也闭, 因为若 $(x_n,1/x_n)\to(x,y)$, 则 $x_n\geq1$ 给出 $x\geq1$, 连续性给出 $y=1/x$, 所以极限仍在 $B$. 但
\[
 d((n,0),(n,1/n))=1/n\longrightarrow0,
\]
所以 $d(A,B)=0$. 两集合均不有界, 因而均不紧. 另一个反例是 $A=\{0\}$ 与 $B=(0,1)$: $A$ 紧, 但 $B$ 不闭, 仍有 $d(A,B)=0$. 因而在正距离定理中, 不能只保留 ``不相交'' 而删去相应的紧性或闭性条件.
\end{solution}

\originalsection{编者练习与解答}
\label{ch:compact:exercises}

以下练习由编者组织, 不标作 MIT 原习题或新领军真题. 它们分别训练有限网的运用、子列抽取、有限交与最值条件, 每题均附完整解答.

\begin{exercise}
\label{ch:compact:ex-separable}
称空间 $X$ 可分, 如果存在至多可数的子集 $D\subseteq X$ 使 $\overline D=X$. 证明每个全有界度量空间可分, 从而每个紧度量空间可分. 可分是否推出全有界?
\end{exercise}
\begin{solution}
对每个 $n\geq1$, 取一个有限 $1/n$-网 $D_n$, 并令 $D=\bigcup_{n\geq1}D_n$. 可数多个有限集合的并至多可数. 对任意 $x\in X$ 与 $r>0$, 取 $n$ 使 $1/n<r$. 有 $p\in D_n$ 满足 $d(x,p)<1/n<r$, 所以 $B_X(x,r)$ 与 $D$ 相交. 因而每个 $x$ 都在 $\overline D$ 中, 即 $D$ 稠密. 空空间取 $D=\varnothing$. 紧空间全有界, 所以也可分. 反方向不成立: $\mathbb{Q}$ 是 $\mathbb{R}$ 的可数稠密子集, 所以 $\mathbb{R}$ 可分, 但 $\mathbb{R}$ 不全有界.
\end{solution}

\begin{exercise}
\label{ch:compact:ex-closure}
设 $A$ 是度量空间 $X$ 的全有界子集. 证明其在 $X$ 中的闭包 $\overline A$ 也全有界. 若 $X$ 完备, 证明 $\overline A$ 紧.
\end{exercise}
\begin{solution}
若 $A=\varnothing$, 闭包为空, 结论成立. 否则对给定 $\varepsilon>0$, 取 $A$ 中的有限 $\varepsilon/3$-网 $p_1,\ldots,p_m$. 对任意 $x\in\overline A$, 可选 $a\in A$ 使 $d(x,a)<\varepsilon/3$, 再选 $j$ 使 $d(a,p_j)<\varepsilon/3$. 则 $d(x,p_j)<2\varepsilon/3<\varepsilon$. 球心 $p_j$ 均属于 $\overline A$, 所以它们构成闭包的有限 $\varepsilon$-网.

若 $X$ 完备, 则 $\overline A$ 也完备: 其中任意 Cauchy 点列在 $X$ 中收敛, 而闭性使其极限留在 $\overline A$. 闭包因此完备且全有界, 由紧性的等价条件得它紧. 这也说明在完备度量空间中, 全有界性恰好保证闭包紧; 单纯有界性没有这样的保证.
\end{solution}

\begin{exercise}
\label{ch:compact:ex-product}
设 $(X,d_X)$ 与 $(Y,d_Y)$ 是紧度量空间. 在 $X\times Y$ 上定义
\[
 d((x,y),(x',y'))=\max\{d_X(x,x'),d_Y(y,y')\}.
\]
证明这是度量, 且 $X\times Y$ 紧. 说明证明中为什么要再次抽取子列.
\end{exercise}
\begin{solution}
非负性、对称性与正定性来自两个坐标的度量. 对三个点 $u,v,w$, 每个坐标距离都满足三角不等式, 且分别不超过 $d(u,v)+d(v,w)$, 所以取最大值后有 $d(u,w)\leq d(u,v)+d(v,w)$. 因而 $d$ 是度量. 若某个因子为空, 积也为空, 已知紧.

对非空情形, 取任意点列 $((x_n,y_n))$. $X$ 紧, 所以有子列指标 $n_k$ 使 $x_{n_k}\to x\in X$. 再对 $Y$ 中的点列 $(y_{n_k})$ 使用序列紧性, 得到严格递增的 $k_\ell$ 使 $y_{n_{k_\ell}}\to y\in Y$. 因为收敛点列的子列仍收敛到同一极限, 仍有 $x_{n_{k_\ell}}\to x$. 因此
\[
 d((x_{n_{k_\ell}},y_{n_{k_\ell}}),(x,y))
 =\max\{d_X(x_{n_{k_\ell}},x),d_Y(y_{n_{k_\ell}},y)\}
 \longrightarrow0.
\]
所以积空间序列紧, 从而紧. 第二次抽取是在第一次子列内部进行的, 才能同时保留两个坐标的收敛性. 分别独立选两个子列, 未必得到共同的指标. 同一方法可逐次用于任意有限多个因子.
\end{solution}

\begin{exercise}
\label{ch:compact:ex-finite-constraints}
设 $K$ 为非空紧度量空间, $\{f_i\}_{i\in I}$ 是一族连续实值函数. 假设对每个有限指标集 $J\subseteq I$, 都存在 $x_J\in K$ 使 $f_j(x_J)\geq0$ 对所有 $j\in J$ 成立. 证明存在同一个 $x\in K$ 使 $f_i(x)\geq0$ 对所有 $i\in I$ 成立.
\end{exercise}
\begin{solution}
令 $F_i=f_i^{-1}([0,\infty))$. 连续函数对闭集的逆像闭, 所以每个 $F_i$ 都在 $K$ 中闭. 题设恰好说明每个有限交 $\bigcap_{j\in J}F_j$ 非空. 因而这族闭集具有有限交性质. 由 $K$ 的紧性及有限交刻画, 总交非空. 取 $x\in\bigcap_{i\in I}F_i$, 就有所有 $f_i(x)\geq0$. 当 $I$ 为空时, 总交按约定为非空的 $K$, 结论也成立. 这里有限问题的解 $x_J$ 可以依赖 $J$, 而紧性保证最后能找到适合全部条件的同一个点.
\end{solution}

\begin{exercise}
\label{ch:compact:ex-distance-attainment}
构造度量空间中的一个非空紧集 $A$ 和一个非空闭集 $B$, 使 $A\cap B=\varnothing$, $d(A,B)>0$, 但没有任何 $a\in A,b\in B$ 满足 $d(a,b)=d(A,B)$.
\end{exercise}
\begin{solution}
取 $X=\{p\}\cup\{b_n:n\geq1\}$, 这些点彼此不同. 令 $r_n=1+1/n$, 定义
\[
 d(p,b_n)=d(b_n,p)=r_n,\qquad
 d(b_n,b_m)=r_n+r_m\quad(n\ne m),
\]
并令 $d(x,x)=0$. 这是度量: 正定性和对称性直接成立. 对三个不同点, 若中间点是 $p$, 两个 $b$ 点之间的三角不等式取等号; 若三个点均为 $b$ 点, 右端比左端多出中间点的 $2r_j>0$; 若一个端点是 $p$ 而中间点是某个 $b_j$, 右端比左端多出 $2r_j$. 含重复点的三角不等式则直接成立.

令 $A=\{p\}$, $B=\{b_n:n\geq1\}$. 单点集 $A$ 紧. $B$ 在 $X$ 中闭, 因为它的补集 $\{p\}=B_X(p,1)$ 开. 两者不相交, 且
\[
 d(A,B)=\inf_{n\geq1}(1+1/n)=1.
\]
但每个 $d(p,b_n)>1$, 所以下确界不被任何点对取得. 此例与正距离定理相容: $B$ 只有闭性而无紧性. 若 $B$ 也紧, 才能保证这个最后的下确界也取得.
\end{solution}

\begin{exercise}
\label{ch:compact:ex-positive-function}
设 $K$ 是非空紧度量空间, $f,h:K\to\mathbb{R}$ 连续, 且 $f(x)>0$ 对每个 $x\in K$ 成立. 证明存在 $c>0$ 使 $f(x)\geq c$ 对所有 $x$ 成立, 并证明 $h/f$ 有界. 若删去紧性, 第一结论是否仍成立?
\end{exercise}
\begin{solution}
极值定理保证 $f$ 在某个 $x_0\in K$ 取得最小值. 令 $c=f(x_0)$, 题设逐点严格为正, 所以这个已经取得的最小值满足 $c>0$. 同时 $h$ 有界, 可取 $M\geq0$ 使 $|h(x)|\leq M$. 因而
\[
 \left|\frac{h(x)}{f(x)}\right|\leq\frac Mc
 \qquad(x\in K).
\]
删去紧性后不成立: 在 $(0,1)$ 上取 $f(x)=x$, 每个函数值严格为正, 但 $\inf f=0$. 要把逐点的正性变成统一正下界, 需要保证下确界实际取得, 而不能仅写出 ``所有值都正''.
\end{solution}

\originalchapter{完备空间与完备化}\label{ch:complete}

\emph{来源: 本章重编自 Paige Bright, MIT OCW 18.S190, IAP 2023, Lecture 5, pp.\ 5--6, 并承接 Lecture 1, pp.\ 4--7 的函数空间例子. 原课留作练习的完备性证明在此补全; 嵌套闭集判据、Cauchy 序列构造及章末练习为编者补充. 来源和许可详见书末说明.}

\originalsection{Cauchy序列与完备性}

上一章研究紧性时, 我们从任意序列中寻找收敛子列. 完备性回答另一个问题: 如果一个序列的后段已经彼此任意接近, 空间中是否有一个点承接这些越来越集中的项? 这里不要求从所有序列中找到好子列, 而只考察已经满足 Cauchy 条件的序列.

\begin{definition}\label{def:complete}
度量空间 $(X,d)$ 称为\textbf{完备空间}, 如果其中每个 Cauchy 序列都收敛于 $X$ 中某点. 具体地, 对每个满足
\[
\forall\varepsilon>0\ \exists N\ \forall m,n\geq N,
\qquad d(x_m,x_n)<\varepsilon
\]
的序列 $(x_n)$, 都存在 $x\in X$ 使 $d(x_n,x)\to0$.
\end{definition}

条件中的 ``收敛于 $X$ 中某点'' 必不可少. 在更大的空间中找到极限, 不能直接证明原空间完备. 每个收敛序列都是 Cauchy 序列: 若 $x_n\to x$, 当 $m,n$ 足够大时,
\[
d(x_m,x_n)\leq d(x_m,x)+d(x,x_n)<\varepsilon.
\]
反向推理必须用到空间的完备性.

\begin{example}
已知 $\mathbb R$ 在通常距离下完备. 证明 $\mathbb Q$ 和 $(0,1)$ 在通常距离下都不完备.
\end{example}
\begin{solution}
取有理数 $q_n$ 使 $|q_n-\sqrt2|<1/n$, 则 $(q_n)$ 在 $\mathbb R$ 中收敛, 因而在 $\mathbb Q$ 的限制距离下是 Cauchy 序列. 若它在 $\mathbb Q$ 中收敛于 $q$, 就会在 $\mathbb R$ 中同时收敛于 $q$ 和 $\sqrt2$. 极限唯一性给出 $q=\sqrt2$, 与 $q\in\mathbb Q$ 矛盾.

区间 $(0,1)$ 也不完备, 因为 $x_n=1/(n+1)$ 是其中的 Cauchy 序列, 在 $\mathbb R$ 中的唯一极限 $0$ 却不属于该区间. 缺失的极限不一定是无理数; 一个被删去的端点也足以破坏完备性.
\end{solution}

\begin{example}
证明离散度量下的任意集合都完备.
\end{example}
\begin{solution}
若集合为空, 没有需要检验的序列. 否则对 Cauchy 序列使用 $\varepsilon=1/2$, 可知从某一项开始任意两项的距离都小于 $1/2$. 离散距离只有 $0$ 和 $1$, 所以这些项必须相同, 序列最终为常值并收敛.
\end{solution}

\begin{proposition}\label{prop:complete-subsequence}
若 Cauchy 序列 $(x_n)$ 有一个收敛子列 $x_{n_j}\to x$, 则整个序列也收敛于 $x$.
\end{proposition}
\begin{proof}
给定 $\varepsilon>0$, 由 Cauchy 条件取 $N$ 使 $m,n\geq N$ 时 $d(x_m,x_n)<\varepsilon/2$. 再选一个 $j$ 使 $n_j\geq N$ 且 $d(x_{n_j},x)<\varepsilon/2$. 于是对所有 $n\geq N$,
\[
d(x_n,x)\leq d(x_n,x_{n_j})+d(x_{n_j},x)<\varepsilon.
\]
这里用于比较的子列项只需选定一个; Cauchy 条件随后控制所有足够靠后的项.
\end{proof}

因此紧度量空间是完备的: 先用紧性找到收敛子列, 再用上面的命题收回整个 Cauchy 序列. 完备性本身并不保证紧性, 例如 $\mathbb R$ 完备但不紧. 完备性约束 Cauchy 序列, 紧性则还要求空间不能无限分散.

\originalsection{子空间的完备性}

\begin{theorem}\label{thm:closed-complete}
设 $A$ 是完备度量空间 $X$ 的子集, 并赋予限制距离. 则 $A$ 完备当且仅当 $A$ 在 $X$ 中闭.
\end{theorem}
\begin{proof}
先设 $A$ 闭. 若 $(a_n)$ 是 $A$ 中的 Cauchy 序列, 它也是 $X$ 中的 Cauchy 序列. $X$ 完备给出 $a_n\to x\in X$. 由于 $A$ 闭, $x\in A$, 故序列在 $A$ 中收敛.

反过来, 设 $A$ 完备. 取 $x\in\overline A$. 每个球 $B(x,1/n)$ 都与 $A$ 相交, 所以可取 $a_n\in A$ 使 $d(a_n,x)<1/n$. 序列 $(a_n)$ 在 $X$ 中收敛于 $x$, 因而是 Cauchy 序列. $A$ 完备给出某个 $a\in A$ 使 $a_n\to a$. 同一个限制距离下的收敛也就是在 $X$ 中的收敛, 极限唯一性给出 $x=a\in A$. 因此 $\overline A=A$, 即 $A$ 闭.
\end{proof}

这段证明还说明: 完备子空间在任意度量母空间中都闭, 反方向才需要母空间完备. 若母空间不完备, 闭子集也可能不完备, 因为母空间自身就是一个闭子集.

\begin{corollary}\label{cor:dense-complete}
若 $A$ 是度量空间 $X$ 的稠密完备子空间, 则 $A=X$.
\end{corollary}
\begin{proof}
由刚才证明的方向, $A$ 在 $X$ 中闭. 稠密性又给出 $\overline A=X$, 所以 $A=\overline A=X$.
\end{proof}

\begin{remark}
稠密性不能保证完备性. $\mathbb Q$ 在 $\mathbb R$ 中稠密, 却不完备. 母空间也可能是完备空间的扩张而仍然不完备: $\{0\}$ 完备, 把它扩张成 $[0,1)$ 后, 序列 $1-1/n$ 缺少极限. 这两种现象的区别在于, 完备性关乎限制距离下所有 Cauchy 序列的极限, 不是集合的大小或点的多少.
\end{remark}

\begin{theorem}[嵌套闭集判据]\label{thm:complete-nested}
非空度量空间 $X$ 完备, 当且仅当下述性质成立: 对每个非空闭集序列
\[
F_1\supset F_2\supset\cdots,
\qquad \operatorname{diam}F_n\longrightarrow0,
\]
交集 $\bigcap_{n=1}^{\infty}F_n$ 恰有一个点. 这里
$\operatorname{diam}F=\sup\{d(x,y):x,y\in F\}$, 允许前面有限项的直径为无穷大.
\end{theorem}
\begin{proof}
若 $X$ 完备, 从每个 $F_n$ 取 $x_n$. 对 $m,n\geq N$, 嵌套性保证 $x_m,x_n\in F_N$, 因而
$d(x_m,x_n)\leq\operatorname{diam}F_N$. 直径趋于 $0$ 表明 $(x_n)$ 是 Cauchy 序列, 所以 $x_n\to x\in X$. 固定 $N$, 后段全部位于闭集 $F_N$, 故 $x\in F_N$. 这对每个 $N$ 成立, 得到 $x\in\bigcap_NF_N$. 若 $y$ 也在交集中, 则对每个 $N$ 都有 $d(x,y)\leq\operatorname{diam}F_N$, 令 $N\to\infty$ 得 $x=y$.

反过来, 假设嵌套闭集性质成立. 对任意 Cauchy 序列 $(x_n)$, 令
\[
F_N=\overline{\{x_n:n\geq N\}}.
\]
这些集非空、闭且嵌套. 给定 $\varepsilon>0$, 取 $N$ 使 $m,n\geq N$ 时 $d(x_m,x_n)<\varepsilon/2$. 对任意 $u,v\in F_N$, 从序列后段分别取 $u_j\to u$, $v_j\to v$. 三角不等式给出距离函数的连续性, 所以
$d(u,v)=\lim_jd(u_j,v_j)\leq\varepsilon/2$. 因此 $\operatorname{diam}F_N\leq\varepsilon/2$, 直径趋于 $0$. 假设给出 $x\in\bigcap_NF_N$. 当 $n\geq N$ 时 $x_n,x\in F_N$, 故
$d(x_n,x)\leq\operatorname{diam}F_N\to0$. 于是 $(x_n)$ 在 $X$ 中收敛.
\end{proof}

直径趋于 $0$ 不是可删去的附带条件. 在完备空间 $\mathbb R$ 中, 非空闭集 $F_n=[n,\infty)$ 嵌套, 交集却为空. 它们向无穷远移动, 并没有集中到一个点.

\originalsection{上确界距离与函数空间的完备性}

在函数空间中, 一个点是一整个函数. 若要用极限寻找函数方程的解, 就必须先选定函数之间的距离. 本节的距离同时控制所有自变量, 因而能把函数列的收敛转化为一致收敛.

\begin{definition}\label{def:cb-space}
对非空度量空间 $X$, 定义
\[
C_b(X)=\{f:X\to\mathbb R:f\text{ 连续且有界}\},
\qquad \|f\|_\infty=\sup_{x\in X}|f(x)|.
\]
函数 $f,g\in C_b(X)$ 的\textbf{上确界距离}为
$d_\infty(f,g)=\|f-g\|_\infty$.
\end{definition}

有界性保证上确界是有限实数. 这里不要求最大值在某点取得: 例如 $f(x)=x$ 在 $X=(0,1)$ 上有上确界 $1$, 却没有最大值. 原课将 $C_b(X)$ 记为 $C_\infty(X)$; 本书用下标 $b$ 强调有界, 避免与光滑函数空间 $C^\infty$ 混淆.

先核对这是一个距离. 非负性和对称性来自绝对值. 若 $\|f-g\|_\infty=0$, 则每个 $x$ 上都有 $|f(x)-g(x)|=0$, 所以 $f=g$. 对每个 $x$,
\[
|f(x)-h(x)|\leq |f(x)-g(x)|+|g(x)-h(x)|
\leq\|f-g\|_\infty+\|g-h\|_\infty.
\]
对左端取上确界, 即得到三角不等式. 这个证明没有使用最大值点, 因而适用于非紧空间 $X$.

\begin{definition}
函数列 $(f_n)$ \textbf{一致收敛}于函数 $f$, 如果
\[
\forall\varepsilon>0\ \exists N\ \forall n\geq N\ \forall x\in X,
\qquad |f_n(x)-f(x)|<\varepsilon.
\]
相比点态收敛, 一致收敛要求同一个 $N$ 对所有 $x$ 有效.
\end{definition}

在有界函数空间中, $d_\infty(f_n,f)\to0$ 等价于一致收敛. 一个方向直接由
$|f_n(x)-f(x)|\leq\|f_n-f\|_\infty$ 得到; 另一个方向对一致收敛的不等式取上确界, 用 $\varepsilon/2$ 代替 $\varepsilon$ 即可.

\begin{theorem}\label{thm:cb-complete}
对任意非空度量空间 $X$, $(C_b(X),d_\infty)$ 是完备空间.
\end{theorem}
\begin{proof}
设 $(f_n)$ 是上确界距离下的 Cauchy 序列. 目标是构造一个仍然属于 $C_b(X)$ 的极限. 构造分为点态极限、一致控制、有界性和连续性四个相互衔接的步骤.

固定 $x\in X$. 由于
\[
|f_n(x)-f_m(x)|\leq\|f_n-f_m\|_\infty,
\]
实数列 $(f_n(x))$ 是 Cauchy 序列. $\mathbb R$ 完备, 因而可定义
$f(x)=\lim_{n\to\infty}f_n(x)$. 这样在每一点定义出一个函数 $f:X\to\mathbb R$, 但此刻尚未知道它连续或有界.

给定 $\varepsilon>0$, 取 $N$ 使 $m,n\geq N$ 时
$\|f_n-f_m\|_\infty<\varepsilon/2$. 固定任意 $n\geq N$ 和任意 $x$, 在
$|f_n(x)-f_m(x)|<\varepsilon/2$ 中令 $m\to\infty$, 得
\[
|f_n(x)-f(x)|\leq\varepsilon/2.
\]
$N$ 不依赖 $x$, 所以取上确界可得
$\|f_n-f\|_\infty\leq\varepsilon/2<\varepsilon$. 这就是一致收敛. 注意极限可能把严格不等式变成非严格不等式, 提前使用 $\varepsilon/2$ 可保留最终所需的严格估计.

对上述估计取一次 $\varepsilon=2$, 则某个 $N$ 满足所有 $x$ 上
$|f(x)-f_N(x)|\leq1$. 于是
\[
|f(x)|\leq |f_N(x)|+1\leq\|f_N\|_\infty+1.
\]
因此 $f$ 有界. 无须假设所有 $f_n$ 的界一开始就是同一个常数; 一个固定函数 $f_N$ 和一致误差已经足够.

最后证明连续性. 固定 $x_0\in X$ 和 $\varepsilon>0$. 一致收敛允许先选一个 $n$ 使
$\|f-f_n\|_\infty<\varepsilon/3$. $f_n$ 在 $x_0$ 连续, 再选 $\delta>0$ 使
$d(x,x_0)<\delta$ 时 $|f_n(x)-f_n(x_0)|<\varepsilon/3$. 于是
\begin{align*}
|f(x)-f(x_0)|
&\leq |f(x)-f_n(x)|+|f_n(x)-f_n(x_0)|
     +|f_n(x_0)-f(x_0)|\\
&<\varepsilon.
\end{align*}
故 $f$ 连续. 至此 $f\in C_b(X)$, 且 $d_\infty(f_n,f)\to0$, 完备性得证.
\end{proof}

\begin{corollary}\label{cor:c-interval-complete}
设 $a<b$. 实连续函数空间 $C([a,b])=C^0([a,b])$ 在上确界距离下完备.
\end{corollary}
\begin{proof}
连续函数在紧区间 $[a,b]$ 上有界, 所以
$C([a,b])=C_b([a,b])$. 应用定理~\ref{thm:cb-complete} 即得. 具体证明仍是: 由一致 Cauchy 得每一点的实数极限, 由同一个尾部估计得一致收敛, 再以三段误差证明极限连续.
\end{proof}

\begin{remark}
点态极限本身不能保留连续性. $f_n(x)=x^n$ 在 $[0,1]$ 上连续, 点态极限为
\[
f(x)=\begin{cases}0,&0\leq x<1,\\1,&x=1,\end{cases}
\]
它在 $1$ 不连续. 因此这些 $f_n$ 不可能在上确界距离下形成 Cauchy 序列: 若是, 上面的完整性证明会迫使其点态极限连续. 这也解释了为什么函数方程中的距离必须认真选择.
\end{remark}

\originalsection{等距嵌入与完备化}

对不完备空间, 我们希望补上 Cauchy 序列缺少的极限, 又不改变原有点之间的距离. 仅说 ``把所有极限添进去'' 并不是构造, 因为在构造完成之前, 这些极限还没有一个共同的母空间. 原课提供了一个巧妙办法: 先把空间等距放入已经证明完备的函数空间, 再取闭包.

\begin{definition}\label{def:completion}
度量空间 $X$ 的一个\textbf{完备化}是一个完备度量空间 $\widehat X$, 连同等距嵌入 $i:X\to\widehat X$, 满足 $i(X)$ 在 $\widehat X$ 中稠密. 等距嵌入的意思是
\[
\widehat d(i(x),i(y))=d(x,y)\qquad(x,y\in X).
\]
等距性自动保证 $i$ 单射. 我们常把 $x$ 和 $i(x)$ 视为同一点, 但定义中保留嵌入能准确表达完备化的唯一性.
\end{definition}

\begin{theorem}\label{thm:completion-exists}
每个度量空间都有完备化.
\end{theorem}
\begin{proof}
空空间已经完备, 取自身即可. 以下设 $X\ne\varnothing$, 并固定一个基点 $o\in X$. 对每个 $x\in X$ 定义
\[
g_x(p)=d(p,x)-d(p,o)\qquad(p\in X).
\]
为什么减去第二项? 若空间无界, $p\mapsto d(p,x)$ 可能无界, 不能放进 $C_b(X)$. 两个距离的差却由反三角不等式满足
\[
|g_x(p)|=|d(p,x)-d(p,o)|\leq d(x,o).
\]
因此对固定 $x$, 函数 $g_x$ 有界. 又由三角不等式,
\begin{align*}
|g_x(p)-g_x(q)|
&\leq |d(p,x)-d(q,x)|+|d(p,o)-d(q,o)|\\
&\leq 2d(p,q),
\end{align*}
所以 $g_x$ 连续. 于是 $g_x\in C_b(X)$.

定义 $i(x)=g_x$. 对任意 $x,y\in X$,
\[
|g_x(p)-g_y(p)|=|d(p,x)-d(p,y)|\leq d(x,y),
\]
取上确界得到 $\|g_x-g_y\|_\infty\leq d(x,y)$. 为了得到反向不等式, 只需找一个取等的测试点. 取 $p=x$, 就有
\[
|g_x(x)-g_y(x)|=|d(x,x)-d(x,y)|=d(x,y).
\]
故 $\|g_x-g_y\|_\infty=d(x,y)$, $i$ 是等距嵌入. 这里映射只需是到它的像 $i(X)$ 的双射, 无须满射到整个 $C_b(X)$.

现在令 $\widehat X=\overline{i(X)}$, 闭包在 $C_b(X)$ 中取得, 并赋予上确界距离的限制. 定理~\ref{thm:cb-complete} 给出 $C_b(X)$ 完备, 定理~\ref{thm:closed-complete} 给出闭子空间 $\widehat X$ 完备. 由闭包的定义, $i(X)$ 在 $\widehat X$ 中稠密. 这正好满足完备化的全部要求.
\end{proof}

这个构造的用途超出证明存在性: 原空间里的几何距离成为两个函数之间的上确界距离. 基点 $o$ 只用于让函数有界, 不改变最后恢复出来的距离. 不同基点可能给出不同的函数模型, 下节说明这些模型仍表示同一个完备化.

\originalsection{完备化的唯一性}

\begin{theorem}\label{thm:completion-unique}
若 $(\widehat X_1,i_1)$ 和 $(\widehat X_2,i_2)$ 都是 $X$ 的完备化, 则存在唯一的满射等距映射
$U:\widehat X_1\to\widehat X_2$, 满足 $U\circ i_1=i_2$.
\end{theorem}
\begin{proof}
对 $z\in\widehat X_1$, 由 $i_1(X)$ 稠密, 选 $x_n\in X$ 使
$i_1(x_n)\to z$. 序列 $(i_1(x_n))$ 是 Cauchy 序列, 等距性给出
\[
\widehat d_2(i_2(x_n),i_2(x_m))
=d(x_n,x_m)
=\widehat d_1(i_1(x_n),i_1(x_m)).
\]
所以 $(i_2(x_n))$ 也是 Cauchy 序列. 由 $\widehat X_2$ 完备, 可以定义
$U(z)=\lim_ni_2(x_n)$.

先核对定义与近似序列无关. 若 $i_1(y_n)\to z$, 则
\[
d(x_n,y_n)\leq\widehat d_1(i_1(x_n),z)
                 +\widehat d_1(z,i_1(y_n))\longrightarrow0.
\]
等距性让 $\widehat d_2(i_2(x_n),i_2(y_n))\to0$, 因而两个极限相同. 取常值近似序列可得 $U(i_1(x))=i_2(x)$.

对 $z,w\in\widehat X_1$, 分别选近似序列 $i_1(x_n)\to z$, $i_1(y_n)\to w$. 距离对两个变量连续, 所以
\begin{align*}
\widehat d_2(U(z),U(w))
&=\lim_n\widehat d_2(i_2(x_n),i_2(y_n))\\
&=\lim_n d(x_n,y_n)=\widehat d_1(z,w).
\end{align*}
因此 $U$ 等距, 特别地单射.

为证明满射, 任取 $v\in\widehat X_2$, 选 $x_n\in X$ 使 $i_2(x_n)\to v$. 等距性表明 $(i_1(x_n))$ 是 Cauchy 序列. 由 $\widehat X_1$ 完备, 它收敛于某个 $z$, 而定义给出 $U(z)=v$.

最后, 若 $V$ 是另一满足 $V\circ i_1=i_2$ 的等距映射, 则它连续. 对任意 $i_1(x_n)\to z$,
\[
V(z)=\lim_nV(i_1(x_n))=\lim_ni_2(x_n)=U(z).
\]
所以 $V=U$, 唯一性得证.
\end{proof}

\begin{remark}
``唯一'' 指的是固定原空间嵌入后, 延伸该嵌入的等距同构唯一. 它不表示两个构造所得的点集在字面上相等, 也不表示完备空间不能有其他等距对称. 例如 $\mathbb R$ 有等距映射 $x\mapsto-x$, 但它不在原嵌入的每一点上保持恒等.
\end{remark}

\originalsection{Cauchy序列模型}

\emph{本节为编者补充, 将原课提出但未展开的 Cauchy 序列思路写成完整构造. 它与上一节的函数模型互补: 新点就是原空间中近似过程的等价类.}

设 $\mathcal C(X)$ 是 $X$ 中所有 Cauchy 序列组成的集合. 对 $u=(u_n)$ 和 $v=(v_n)$, 定义
\[
u\sim v\quad\Longleftrightarrow\quad d(u_n,v_n)\longrightarrow0.
\]
自反性和对称性直接来自距离; 若 $u\sim v$ 且 $v\sim w$, 则
$d(u_n,w_n)\leq d(u_n,v_n)+d(v_n,w_n)\to0$, 所以关系传递. 因此可以取商集 $\mathcal C(X)/\!\sim$.

\begin{proposition}\label{prop:cauchy-completion-distance}
在等价类上定义
\[
\rho([u],[v])=\lim_{n\to\infty}d(u_n,v_n).
\]
该极限存在, 与代表序列无关, 且 $\rho$ 是商集上的距离.
\end{proposition}
\begin{proof}
由三角不等式,
\[
|d(u_n,v_n)-d(u_m,v_m)|\leq d(u_n,u_m)+d(v_n,v_m).
\]
右端由两个 Cauchy 条件趋于 $0$, 所以 $(d(u_n,v_n))$ 是实数 Cauchy 序列, 极限存在且有限. 若 $u\sim u'$、$v\sim v'$, 同样的估计给出
\[
|d(u_n,v_n)-d(u'_n,v'_n)|
\leq d(u_n,u'_n)+d(v_n,v'_n)\longrightarrow0,
\]
故极限与代表元无关.

非负性和对称性对每一项成立, 取极限后仍成立. $\rho([u],[v])=0$ 恰好等价于 $u\sim v$, 也就是 $[u]=[v]$. 最后从
$d(u_n,w_n)\leq d(u_n,v_n)+d(v_n,w_n)$ 取极限, 得到三角不等式.
\end{proof}

把 $x\in X$ 对应到常值序列的类, 记为 $j(x)=[(x,x,\ldots)]$. 则
$\rho(j(x),j(y))=d(x,y)$, 所以 $j$ 等距.

\begin{theorem}\label{thm:cauchy-completion}
赋予上述距离的 $\mathcal C(X)/\!\sim$, 连同嵌入 $j$, 是 $X$ 的完备化.
\end{theorem}
\begin{proof}
先证明 $j(X)$ 稠密. 任取 $A=[(u_n)]$ 和 $\varepsilon>0$, 由 $(u_n)$ 的 Cauchy 性取 $N$, 使 $m,n\geq N$ 时 $d(u_n,u_m)<\varepsilon/2$. 当 $n\geq N$ 时,
\[
\rho(j(u_n),A)=\lim_{m\to\infty}d(u_n,u_m)
\leq\varepsilon/2<\varepsilon.
\]
因此 $j(u_n)\to A$.

再证明商空间完备. 设 $(A_k)$ 是商空间中的 Cauchy 序列. 利用刚证明的稠密性, 对每个 $k\geq1$ 选 $x_k\in X$ 使
$\rho(A_k,j(x_k))<2^{-k}$. 对 $k,l$,
\[
d(x_k,x_l)=\rho(j(x_k),j(x_l))
\leq 2^{-k}+\rho(A_k,A_l)+2^{-l}.
\]
当 $k,l$ 足够大时右端任意小, 所以 $(x_k)$ 是 $X$ 中的 Cauchy 序列. 令 $A=[(x_k)]$. 第一段稠密性证明对这个序列给出 $j(x_k)\to A$. 因此
\[
\rho(A_k,A)\leq\rho(A_k,j(x_k))+\rho(j(x_k),A)
\longrightarrow0.
\]
于是每个 Cauchy 序列 $(A_k)$ 都在商空间中收敛, 完备性得证.
\end{proof}

两个构造之间的对应也可以直接写出: 把 $[(x_n)]$ 送到 $C_b(X)$ 中的一致极限 $\lim_ng_{x_n}$. 等距性保证 $(g_{x_n})$ 一致 Cauchy, 等价序列给出同一个极限. 唯一性定理说明这正是两个完备化之间延伸原嵌入的唯一等距同构. 对 $\mathbb Q$ 的通常距离而言, 这个构造恢复 $\mathbb R$; 换一个距离, 得到的完备化也可能改变.

\originalsection{练习与解答}

\emph{以下题目均为编者练习, 用于检查完备性条件、函数空间距离与完备化构造, 不冒充 MIT 原题.}

\begin{exercise}\label{ex:complete-topology-metric}
在 $\mathbb R$ 上定义 $d_*(x,y)=|\arctan x-\arctan y|$. 证明 $d_*$ 是距离, 它与通常距离诱导同一拓扑, 但 $(\mathbb R,d_*)$ 不完备. 描述它的一个完备化.
\end{exercise}
\begin{solution}
$\arctan$ 单射, 所以 $d_*(x,y)=0$ 当且仅当 $x=y$. 非负性、对称性和三角不等式由实数绝对值立即得到. 映射
$\arctan:\mathbb R\to(-\pi/2,\pi/2)$ 在通常拓扑下连续, 逆映射 $\tan$ 也连续, 因而是同胚. $d_*$ 又恰好是这个映射拉回的通常距离, 所以两种距离给出同一拓扑.

取 $x_n=n$. 则 $\arctan n\to\pi/2$, 所以 $(n)$ 对 $d_*$ 是 Cauchy 序列. 如果它对 $d_*$ 收敛于 $x\in\mathbb R$, 就有 $\arctan n\to\arctan x$, 迫使 $\arctan x=\pi/2$, 不可能. 因此该度量不完备. 闭区间 $[-\pi/2,\pi/2]$ 连同嵌入 $i(x)=\arctan x$ 是其完备化: 区间完备, 映射等距, 像为稠密的开区间. 这个例子说明完备性属于距离结构, 同一拓扑不能单独决定它.
\end{solution}

\begin{exercise}\label{ex:c1-complete}
在 $C^1([a,b])$ 上取距离
\[
d_{C^1}(f,g)=\|f-g\|_\infty+\|f'-g'\|_\infty.
\]
证明它是完备距离. 若只使用 $\|f-g\|_\infty$, $C^1([-1,1])$ 是否完备?
\end{exercise}
\begin{solution}
距离公理由两项上确界估计得到; 第一项已保证正定性. 若 $(f_n)$ 对 $d_{C^1}$ 是 Cauchy 序列, 则 $(f_n)$ 和 $(f_n')$ 分别在 $C([a,b])$ 的上确界距离下是 Cauchy 序列. 由推论~\ref{cor:c-interval-complete}, 存在连续函数 $f,h$ 使 $f_n\to f$, $f_n'\to h$, 两者均为一致收敛.

对每个 $x\in[a,b]$, 微积分基本定理给出
\[
f_n(x)=f_n(a)+\int_a^x f_n'(t)\,\mathrm dt.
\]
而
\[
\left|\int_a^x(f_n'(t)-h(t))\,\mathrm dt\right|
\leq(b-a)\|f_n'-h\|_\infty\longrightarrow0.
\]
令 $n\to\infty$ 得 $f(x)=f(a)+\int_a^xh(t)\,\mathrm dt$. 因 $h$ 连续, 再用微积分基本定理得 $f\in C^1([a,b])$ 且 $f'=h$. 因而
$d_{C^1}(f_n,f)=\|f_n-f\|_\infty+\|f_n'-h\|_\infty\to0$.

若只控制函数值, 取 $f_n(x)=\sqrt{x^2+1/n}$, 它们属于 $C^1([-1,1])$, 且
\[
0\leq f_n(x)-|x|
=\frac{1/n}{\sqrt{x^2+1/n}+|x|}\leq\frac1{\sqrt n}.
\]
故 $f_n$ 一致收敛于 $|x|$, 从而是上确界 Cauchy 序列. $|x|$ 在 $0$ 不可微, 不属于 $C^1([-1,1])$. 任何可能的上确界极限必为它, 所以这一距离下不完备.
\end{solution}

\begin{exercise}\label{ex:l1-incomplete-continuous}
在 $C([0,1])$ 上取 $d_1(f,g)=\int_0^1|f(x)-g(x)|\,\mathrm dx$. 证明这一空间不完备.
\end{exercise}
\begin{solution}
对 $n\geq3$, 令 $f_n$ 在 $[0,1/2-1/n]$ 上为 $0$, 在 $[1/2+1/n,1]$ 上为 $1$, 中间用直线连接. 设 $s(x)=0$ 当 $x<1/2$, $s(x)=1$ 当 $x\geq1/2$. 两者只在长度 $2/n$ 的区间上可能不同, 且差的绝对值不超过 $1$, 所以
$\int_0^1|f_n-s|\leq2/n$. 于是
$d_1(f_n,f_m)\leq2/n+2/m$, $(f_n)$ 是 Cauchy 序列.

若它在该空间中收敛于连续函数 $f$, 则积分三角不等式给出
\[
\int_0^1|f-s|\leq d_1(f,f_n)+\int_0^1|f_n-s|\longrightarrow0.
\]
因此左端为 $0$. 在区间 $(0,1/2)$ 上, 连续非负函数 $|f|$ 的积分为 $0$, 故每一点都为 $0$: 若某点为正, 连续性会使一段正长度区间上的值有正下界, 积分便为正. 同理在 $(1/2,1)$ 上 $f=1$. 这与 $f$ 在 $1/2$ 连续矛盾. 因而不存在连续极限, 该距离下不完备.
\end{solution}

\begin{exercise}\label{ex:completion-cauchy-limits}
设 $(\widehat X,i)$ 是 $X$ 的完备化. 证明 $\widehat X$ 中每个点都是 $X$ 中某个 Cauchy 序列在嵌入后的极限. 若 $X$ 已经完备, 证明 $i$ 满射.
\end{exercise}
\begin{solution}
任取 $z\in\widehat X$. 稠密性允许取 $x_n\in X$ 使
$\widehat d(i(x_n),z)<1/n$. 故 $i(x_n)\to z$, 从而是 Cauchy 序列. 等距性让 $(x_n)$ 在 $X$ 中也是 Cauchy 序列.

若 $X$ 完备, 则 $x_n\to x\in X$. 等距性给出 $i(x_n)\to i(x)$, 极限唯一性迫使 $z=i(x)$. 每个 $z$ 都在像中, 所以 $i$ 满射. 已完备的空间没有需要补入的新极限.
\end{solution}

\originalchapter{一致连续与压缩映射}\label{ch:fixed}

\emph{来源: 本章重编自 Paige Bright, MIT OCW 18.S190, IAP 2023, Lecture 5, pp.\ 1--4. 一致连续、参数积分、Banach 不动点定理和积分算子来自原课; 误差界、一致连续延拓、局部 Picard--Lindel\"of 定理及章末练习为编者补充. 原课积分式的积分变量已订正, 仅连续核不保证可微性的限制在正文说明. 来源与许可详见书末说明.}

\originalsection{连续性与Lipschitz条件}

连续性说, 在固定点附近足够小的输入变化造成足够小的输出变化. 但 ``足够小'' 可以随点的位置改变. 在用一整个函数作为未知量时, 我们往往需要对所有点同时有效的控制, 这就引出一致连续性和 Lipschitz 条件.

\begin{definition}\label{def:uniform-continuity}
设 $(X,d_X)$、$(Y,d_Y)$ 为度量空间. 函数 $f:X\to Y$ 称为\textbf{一致连续}, 如果
\[
\forall\varepsilon>0\ \exists\delta>0\ \forall x,y\in X,
\qquad d_X(x,y)<\delta\ \Longrightarrow\ d_Y(f(x),f(y))<\varepsilon.
\]
同一个 $\delta$ 必须控制 $X$ 中所有点对.
\end{definition}

连续性的量词顺序是先固定 $x$, 再对给定 $\varepsilon$ 选择 $\delta$, 允许 $\delta$ 依赖 $x$. 一致连续性先选 $\delta$, 再检查所有 $x,y$. 把一个点的连续性证明逐点做完, 并不自动给出一个可同时使用的 $\delta$.

\begin{definition}\label{def:lipschitz}
若存在常数 $K\geq0$, 使对所有 $x,y\in X$ 都有
\[
d_Y(f(x),f(y))\leq Kd_X(x,y),
\]
则称 $f$ 为 \textbf{Lipschitz 映射}, 或 $K$-Lipschitz 映射.
\end{definition}

\begin{proposition}\label{prop:lipschitz-uniform}
Lipschitz 映射一致连续, 一致连续映射连续.
\end{proposition}
\begin{proof}
若 $K>0$, 对给定的 $\varepsilon>0$ 取 $\delta=\varepsilon/K$. 当 $d_X(x,y)<\delta$ 时,
\[
d_Y(f(x),f(y))\leq Kd_X(x,y)<K\delta=\varepsilon.
\]
这个 $\delta$ 不依赖 $x,y$, 所以 $f$ 一致连续. 若 $K=0$, 则任意两点的像的距离都为 $0$, 函数为常值, 可取任何 $\delta>0$. 最后, 一致连续的定义对每个固定 $x$ 都有效, 因而给出该点的连续性.
\end{proof}

\begin{example}
证明 $f(x)=\sqrt{x}$ 在 $[0,1]$ 上一致连续但不满足 Lipschitz 条件, 而 $x\mapsto1/x$ 在 $(0,1)$ 上连续却不一致连续.
\end{example}
\begin{solution}
对 $x\geq y\geq0$,
\[
(\sqrt{x}-\sqrt{y})^2
\leq(\sqrt{x}-\sqrt{y})(\sqrt{x}+\sqrt{y})=x-y.
\]
所以 $|\sqrt{x}-\sqrt{y}|\leq\sqrt{|x-y|}$, 取 $\delta=\varepsilon^2$ 就得到一致连续. 若存在 Lipschitz 常数 $K$, 令 $y=0$ 将给出 $\sqrt{x}\leq Kx$, 即 $1/\sqrt{x}\leq K$ 对所有 $x>0$ 成立, 不可能.

$f(x)=1/x$ 在 $(0,1)$ 上连续却不一致连续. 取 $x_n=1/(n+1)$, $y_n=1/(n+2)$. 则 $|x_n-y_n|\to0$, 但 $|f(x_n)-f(y_n)|=1$. 对 $\varepsilon=1/2$, 无论给出多小的 $\delta$, 最终这些点对都违反一致连续的要求.
\end{solution}

这些例子区分了三个层次. Lipschitz 条件给出线性的全局估计; 一致连续允许非线性的变化尺度; 普通连续只给出每一点附近的局部控制.

\originalsection{紧空间上的一致连续性与参数积分}

\begin{theorem}\label{thm:heine-cantor}
若 $X$ 是紧度量空间, $f:X\to Y$ 连续, 则 $f$ 一致连续.
\end{theorem}
\begin{proof}
给定 $\varepsilon>0$. 对每个 $c\in X$, 连续性允许选 $r_c>0$, 使
\[
d_X(z,c)<r_c\quad\Longrightarrow\quad
d_Y(f(z),f(c))<\varepsilon/2.
\]
所有 $B(c,r_c)$ 构成 $X$ 的开覆盖. 引理~\ref{ch:compact:lebesgue} 给出 $\delta>0$, 使每个球 $B(x,\delta)$ 都包含在该覆盖的某一个成员中. 若 $d_X(x,y)<\delta$, 就可取 $c$ 使
$B(x,\delta)\subset B(c,r_c)$. $x,y$ 都属于 $B(c,r_c)$, 所以
\[
d_Y(f(x),f(y))\leq d_Y(f(x),f(c))+d_Y(f(c),f(y))
<\varepsilon.
\]
同一个 $\delta$ 对所有点对有效, 故为一致连续.
\end{proof}

证明中的紧性将各点的局部控制汇合成一个全局尺度. 不必要求值域 $Y$ 紧, 也不必要求 $Y$ 完备; 这里用的是定义域的紧性和 $f$ 的连续性.

\begin{proposition}\label{prop:parameter-integral}
设 $a<b$, $c<d$, 且 $F:[a,b]\times[c,d]\to\mathbb R$ 连续. 则
\[
G(y)=\int_a^bF(x,y)\,\mathrm dx
\]
是 $[c,d]$ 上的连续函数.
\end{proposition}
\begin{proof}
赋予矩形上确界距离, 它是紧度量空间. 由定理~\ref{thm:heine-cantor}, $F$ 一致连续. 对给定的 $\varepsilon>0$, 取 $\delta>0$, 使矩形中两点距离小于 $\delta$ 时, 函数值的差小于 $\varepsilon/(2(b-a))$. 若 $|y-y'|<\delta$, 则对所有 $x\in[a,b]$,
\[
|F(x,y)-F(x,y')|<\frac{\varepsilon}{2(b-a)}.
\]
因此
\begin{align*}
|G(y)-G(y')|
&\leq\int_a^b|F(x,y)-F(x,y')|\,\mathrm dx\\
&\leq(b-a)\sup_{x\in[a,b]}|F(x,y)-F(x,y')|\\
&\leq\varepsilon/2<\varepsilon.
\end{align*}
这个直接估计已证明一致连续, 因而也证明连续. 若 $a=b$, 积分恒为 $0$, 结论仍成立.
\end{proof}

此处关键在于对所有积分变量 $x$ 使用同一个控制. 单独知道每个 $x$ 上 $F(x,y')\to F(x,y)$, 并不足以直接交换积分和极限; 上式展示了本例交换的实际依据.

\originalsection{一致连续映射的延拓}

\emph{本节为编者补充, 说明上一章的完备化如何与一致连续性配合.}

\begin{lemma}\label{lem:uniform-cauchy}
一致连续映射把 Cauchy 序列送到 Cauchy 序列.
\end{lemma}
\begin{proof}
设 $f:X\to Y$ 一致连续, $(x_n)$ 是 $X$ 中的 Cauchy 序列. 给定 $\varepsilon>0$, 先取一致连续性中的 $\delta>0$, 再取 $N$ 使 $m,n\geq N$ 时 $d_X(x_m,x_n)<\delta$. 于是 $d_Y(f(x_m),f(x_n))<\varepsilon$. 这就是像序列的 Cauchy 条件.
\end{proof}

\begin{theorem}\label{thm:uniform-extension}
设 $A$ 在度量空间 $X$ 中稠密, $Y$ 完备. 每个一致连续映射 $f:A\to Y$ 都有唯一的一致连续延拓 $\widetilde f:X\to Y$. 若 $f$ 是 $K$-Lipschitz, 延拓也可取同一个 Lipschitz 常数 $K$.
\end{theorem}
\begin{proof}
对 $x\in X$, 稠密性允许选择 $a_n\in A$ 使 $a_n\to x$. 它是 Cauchy 序列, 由引理~\ref{lem:uniform-cauchy}, $(f(a_n))$ 也是 Cauchy 序列. $Y$ 完备保证极限存在, 定义
$\widetilde f(x)=\lim_nf(a_n)$.

若另有 $b_n\in A$ 趋于 $x$, 则
$d_X(a_n,b_n)\leq d_X(a_n,x)+d_X(x,b_n)\to0$. 一致连续性给出 $d_Y(f(a_n),f(b_n))\to0$, 所以两个极限相同. 定义不依赖所选序列. 对 $x\in A$ 取常值序列即可得到 $\widetilde f(x)=f(x)$.

证明延拓的一致连续性时, 需要照顾取极限后的非严格不等式. 给定 $\varepsilon>0$, 对原映射用 $\varepsilon/2$ 选择 $\eta>0$: 当 $a,b\in A$ 且 $d_X(a,b)<\eta$ 时, $d_Y(f(a),f(b))<\varepsilon/2$. 取 $\delta=\eta/3$. 对 $d_X(x,y)<\delta$ 的任意 $x,y\in X$, 选 $a_n\to x$, $b_n\to y$. 对足够大的 $n$,
\[
d_X(a_n,b_n)\leq d_X(a_n,x)+d_X(x,y)+d_X(y,b_n)<\eta.
\]
故令 $n\to\infty$ 得
$d_Y(\widetilde f(x),\widetilde f(y))\leq\varepsilon/2<\varepsilon$. $\delta$ 不依赖 $x,y$, 延拓一致连续.

若 $H:X\to Y$ 是另一连续延拓, 对 $a_n\to x$ 有
$H(x)=\lim_nH(a_n)=\lim_nf(a_n)=\widetilde f(x)$, 所以延拓唯一. 最后, 若原映射满足 Lipschitz 估计, 则
\[
d_Y(\widetilde f(x),\widetilde f(y))
=\lim_nd_Y(f(a_n),f(b_n))
\leq K\lim_nd_X(a_n,b_n)=Kd_X(x,y).
\]
因此常数 $K$ 被保留.
\end{proof}

连续性本身不能代替一致连续性. 例如 $1/x$ 在稠密子空间 $(0,1)$ 上连续, 却不能连续延拓到 $[0,1]$. 延拓构造需要把趋于边界的 Cauchy 序列送到仍有极限的序列.

\originalsection{Banach 不动点定理}

现在把 Lipschitz 常数降到 $1$ 以下. 每次应用映射后, 两个输入的距离都会按固定比例缩小. 若反复应用映射, 相邻迭代之间的距离形成几何级数, 完备性便能承接这个过程的极限.

\begin{definition}\label{def:contraction}
度量空间 $X$ 上的自映射 $T:X\to X$ 称为\textbf{压缩映射}, 如果存在 $q\in[0,1)$ 使
\[
d(Tx,Ty)\leq qd(x,y)\qquad(x,y\in X).
\]
满足 $Tx_*=x_*$ 的点 $x_*$ 称为 $T$ 的\textbf{不动点}.
\end{definition}

\begin{theorem}[Banach 不动点定理]\label{thm:banach-fixed}
设 $X$ 是非空完备度量空间, $T:X\to X$ 是压缩映射, 压缩常数为 $q<1$. 则 $T$ 恰有一个不动点 $x_*$. 对任意 $x_0\in X$, 递归定义 $x_{n+1}=Tx_n$, 都有 $x_n\to x_*$.
\end{theorem}
\begin{proof}
先设 $0<q<1$. 由压缩条件逐次得到
\[
d(x_{n+1},x_n)\leq qd(x_n,x_{n-1})
\leq\cdots\leq q^nd(x_1,x_0).
\]
仅仅证明相邻距离趋于 $0$ 不够, 因为小步长仍可能累积出无穷路程. 这里几何级数的可求和性给出所有后段项之间的估计. 当 $m>n$ 时,
\begin{align*}
d(x_m,x_n)
&\leq\sum_{j=n}^{m-1}d(x_{j+1},x_j)\\
&\leq d(x_1,x_0)\sum_{j=n}^{m-1}q^j
\leq\frac{q^n}{1-q}d(x_1,x_0).
\end{align*}
右端与 $m$ 无关并趋于 $0$, 因而 $(x_n)$ 是 Cauchy 序列. 完备性保证存在 $x_*\in X$ 使 $x_n\to x_*$. 压缩映射是 Lipschitz 映射, 因而连续. 所以
\[
Tx_*=\lim_{n\to\infty}Tx_n
=\lim_{n\to\infty}x_{n+1}=x_*.
\]
这证明存在性. 若 $y_*$ 也为不动点, 则
\[
d(x_*,y_*)=d(Tx_*,Ty_*)\leq qd(x_*,y_*).
\]
因 $1-q>0$, 得 $d(x_*,y_*)=0$, 即 $x_*=y_*$.

若 $q=0$, 则 $T$ 的所有像相同, 记为 $c\in X$. 因 $Tc=c$, 它是唯一不动点, 任意迭代从第一项起就恒为 $c$. 这一情况无需几何级数估计.
\end{proof}

证明同时给出了寻找解的算法. 完备性用于获得极限, 自映射条件用于保证每一步和极限都仍在空间内, $q<1$ 同时用于几何级数求和与唯一性. 在应用之前, 这三个条件都必须分别核对.

\begin{proposition}[迭代误差界]\label{prop:banach-error}
在定理~\ref{thm:banach-fixed} 的条件下, 当 $0<q<1$ 时, 对 $n\geq0$ 有先验误差界
\[
d(x_n,x_*)\leq\frac{q^n}{1-q}d(x_1,x_0).
\]
对 $n\geq1$ 有后验误差界
\[
d(x_n,x_*)\leq\frac{q}{1-q}d(x_n,x_{n-1}).
\]
此外任意 $x\in X$ 都满足残差估计
\[
d(x,x_*)\leq\frac{d(x,Tx)}{1-q}.
\]
\end{proposition}
\begin{proof}
在上一定理的 $d(x_m,x_n)$ 估计中令 $m\to\infty$, 得到先验界. 对后验界, 以当前一步距离为起点, 压缩条件给出
\[
d(x_{n+j+1},x_{n+j})\leq q^{j+1}d(x_n,x_{n-1})
\qquad(j\geq0).
\]
把这些距离相加再令后端趋于极限, 得
\[
d(x_n,x_*)\leq\sum_{j=0}^{\infty}q^{j+1}d(x_n,x_{n-1})
=\frac{q}{1-q}d(x_n,x_{n-1}).
\]
最后, 对任意 $x$ 使用三角不等式和压缩条件,
\[
d(x,x_*)\leq d(x,Tx)+d(Tx,Tx_*)
\leq d(x,Tx)+qd(x,x_*),
\]
移项即得残差估计. 当 $q=0$ 时, 残差界仍成立, 且 $x_n=x_*$ 对 $n\geq1$ 成立.
\end{proof}

先验界只用初始一步和迭代次数, 能在计算前估计需要多少步. 后验界用已经算出的相邻两项, 能在计算中决定何时停止. 残差界则可以检验任何候选点, 无须它一定来自同一条迭代.

\originalsection{积分方程的连续解}

\begin{theorem}\label{thm:fredholm-contraction}
设 $a<b$, $g\in C([a,b])$, $k\in C([a,b]\times[a,b])$, $\lambda\in\mathbb R$. 令
\[
M=\max_{(x,y)\in[a,b]^2}|k(x,y)|.
\]
若 $q=|\lambda|M(b-a)<1$, 则积分方程
\[
f(x)=g(x)+\lambda\int_a^b k(x,y)f(y)\,\mathrm dy
\]
在 $C([a,b])$ 中恰有一个解. 从任意 $f_0\in C([a,b])$ 出发的迭代
\[
f_{n+1}(x)=g(x)+\lambda\int_a^b k(x,y)f_n(y)\,\mathrm dy
\]
一致收敛于该解, 并满足命题~\ref{prop:banach-error} 的上确界误差估计.
\end{theorem}
\begin{proof}
在 $C([a,b])$ 上定义
\[
(Tf)(x)=g(x)+\lambda\int_a^b k(x,y)f(y)\,\mathrm dy.
\]
先核对自映射条件. 对固定的连续 $f$, 二元函数 $k(x,y)f(y)$ 连续; 命题~\ref{prop:parameter-integral} 在交换两个坐标的角色后表明其关于 $y$ 的积分是 $x$ 的连续函数. 加上连续的 $g$, 得 $Tf\in C([a,b])$.

再核对压缩性. 对 $f,h\in C([a,b])$, 每个 $x$ 都满足
\begin{align*}
|(Tf)(x)-(Th)(x)|
&\leq|\lambda|\int_a^b|k(x,y)|\,|f(y)-h(y)|\,\mathrm dy\\
&\leq|\lambda|M(b-a)\|f-h\|_\infty.
\end{align*}
取上确界得 $\|Tf-Th\|_\infty\leq q\|f-h\|_\infty$, 且 $q<1$. 最后, 推论~\ref{cor:c-interval-complete} 表明该函数空间完备, 并且零函数保证它非空. Banach 不动点定理给出唯一的不动点, 正是积分方程的唯一连续解; 同时给出迭代收敛及误差估计.
\end{proof}

\begin{remark}
条件 $|\lambda|M(b-a)<1$ 是本次上确界估计得到的充分条件, 不是所有积分方程有解的必要条件. 有些核可以得到更小的算子估计, 有些方程即使不满足这个条件也有唯一解; 那些结论需要另行证明.

仅连续的核不能保证解属于 $C^1$. 例如在 $[-1,1]$ 上取 $g=0$, $k(x,y)=|x|$, $\lambda=1$. 连续函数 $f(x)=|x|$ 满足
$f(x)=\int_{-1}^1|x|f(y)\,\mathrm dy$, 却在 $0$ 不可微. 即使在压缩范围内也不能加强为 $C^1$: 取 $g(x)=1$, $k(x,y)=|x|$, $\lambda=1/4$, 则压缩常数为 $1/2$, 唯一解为 $f(x)=1+(2/3)|x|$, 仍不可微. 确实, 若 $I=\int_{-1}^1f(y)\,\mathrm dy$, 方程给出 $f(x)=1+(I/4)|x|$, 从而 $I=2+I/4$, 即 $I=8/3$.
\end{remark}

这一限制并不削弱不动点结论, 而是区分了两种不同的推理: 完备性和压缩性给出连续函数空间中的解, 可微性必须来自积分式的额外结构. 下一节的变上限积分恰好提供这种结构.

\originalsection{初值问题与 Picard 迭代}

考虑初值问题
\[
y'(t)=F(t,y(t)),\qquad y(t_0)=y_0.
\]
如果 $y$ 是 $C^1$ 解, 微积分基本定理给出
\[
y(t)=y_0+\int_{t_0}^{t}F(s,y(s))\,\mathrm ds.
\]
反过来, 若一个连续函数满足这个积分式, 且 $F$ 在函数图像上连续, 那么被积函数连续, 微积分基本定理使 $y$ 可微, 导数正是 $F(t,y(t))$. 因此我们可以先在连续函数空间中寻找积分算子的不动点, 再恢复微分方程.

\begin{theorem}[局部 Picard--Lindel\"of 定理]\label{thm:picard-local}
设 $\alpha,\beta>0$, 函数 $F$ 在闭矩形
\[
R=[t_0-\alpha,t_0+\alpha]\times[y_0-\beta,y_0+\beta]
\]
上连续, 并存在 $L\geq0$ 使同一 $t$ 下的所有 $u,v\in[y_0-\beta,y_0+\beta]$ 都满足
\[
|F(t,u)-F(t,v)|\leq L|u-v|.
\]
令 $M=\max_R|F|$. 选 $h>0$ 满足
\[
h\leq\alpha,\qquad hM\leq\beta,\qquad hL<1.
\]
则在 $I=[t_0-h,t_0+h]$ 上, 初值问题有唯一一个图像位于 $R$ 的 $C^1$ 解. 若另一解只在 $t_0$ 的某个邻域定义, 并具有相同初值, 则它在某个更小的共同邻域上与这个解相同.
\end{theorem}
\begin{proof}
所需的 $h$ 总能取到: $M,L$ 是有限非负数, $\alpha,\beta$ 为正, 所以取充分小的正数即可. 若 $M=0$ 或 $L=0$, 对应条件不施加额外限制.

在完备空间 $C(I)$ 中考虑闭球
\[
\mathcal B=\{u\in C(I):\|u-y_0\|_\infty\leq\beta\}.
\]
它非空, 因为含常值函数 $y_0$. 它是闭集: 若 $u_n\to u$ 一致且 $\|u_n-y_0\|_\infty\leq\beta$, 则逐点取极限得 $|u(t)-y_0|\leq\beta$ 对每个 $t$ 成立. 因而由定理~\ref{thm:closed-complete}, $\mathcal B$ 完备.

定义 Picard 算子
\[
(Tu)(t)=y_0+\int_{t_0}^{t}F(s,u(s))\,\mathrm ds.
\]
若 $u\in\mathcal B$, 则 $(s,u(s))\in R$ 对所有 $s\in I$ 成立, 所以积分有定义, 被积函数连续, $Tu$ 也连续. 对 $t<t_0$ 将积分方向反转再估计, 对 $t\geq t_0$ 直接估计; 两种情形统一为
\[
|(Tu)(t)-y_0|\leq|t-t_0|M\leq hM\leq\beta.
\]
因此 $T$ 确实把 $\mathcal B$ 映入自身. 这一步解释了保球条件 $hM\leq\beta$ 的用途; 仅有压缩估计不足以保证不动点定理可用.

对 $u,v\in\mathcal B$, Lipschitz 条件给出
\begin{align*}
|(Tu)(t)-(Tv)(t)|
&\leq\left|\int_{t_0}^{t}|F(s,u(s))-F(s,v(s))|\,\mathrm ds\right|\\
&\leq |t-t_0|L\|u-v\|_\infty
\leq hL\|u-v\|_\infty.
\end{align*}
第一行右端表示在两个端点之间对非负函数积分的绝对值, 因而同时适用于正反两个方向. 取上确界后得到压缩常数 $q=hL<1$. Banach 不动点定理给出唯一 $y\in\mathcal B$ 满足 $Ty=y$. 它连续, 而 $F(t,y(t))$ 连续, 所以积分式给出
$y'(t)=F(t,y(t))$ 和 $y(t_0)=y_0$. 端点的导数按单侧导数理解; 在 $I$ 的内部这是通常的 $C^1$ 解.

若 $z$ 是另一个在 $I$ 上、图像位于 $R$ 的 $C^1$ 解, 将其微分方程积分便得 $Tz=z$, 且 $z\in\mathcal B$, 不动点唯一性给出 $z=y$.

最后设 $z$ 只在包含 $t_0$ 的开区间上定义, 具有相同初值, 并解同一方程. 由 $z(t_0)=y_0$ 和连续性, 可取 $0<\eta\leq h$, 使 $J=[t_0-\eta,t_0+\eta]$ 包含在它的定义区间中, 且 $|z(t)-y_0|\leq\beta$ 对所有 $t\in J$ 成立. $y|_J$ 和 $z|_J$ 都属于 $C(J)$ 中同一半径为 $\beta$ 的闭球. 限制后的 Picard 算子仍是自映射, 压缩常数为 $\eta L\leq hL<1$. 两者都为其不动点, 因而在 $J$ 上相同. 这正是所陈述的局部唯一性.
\end{proof}

定理没有宣称解能在任意长的时间区间上存在. 矩形限定了可用的 $F$ 与 Lipschitz 条件, 保球估计限定了函数值的范围, 压缩估计限定了本次直接证明的时间长度. 若要继续延拓解, 必须在新的邻域重新核对条件.

\begin{example}
用 Picard 迭代求初值问题 $y'=y$, $y(0)=1$ 的局部解, 并核对迭代可用的区间.
\end{example}
\begin{solution}
取 $R=[-1,1]\times[0,2]$. 此时 $M=2$, $L=1$. 选 $h=1/2$, 则 $hM=1\leq\beta=1$, $hL=1/2<1$. 在 $[-1/2,1/2]$ 上, 从 $y^{(0)}(t)=1$ 开始迭代,
\[
y^{(n+1)}(t)=1+\int_0^t y^{(n)}(s)\,\mathrm ds.
\]
前几项为 $1+t$, $1+t+t^2/2$, $1+t+t^2/2+t^3/6$. 归纳可得
$y^{(n)}(t)=\sum_{j=0}^n t^j/j!$. 不动点定理证明这些多项式在上述区间上一致收敛于唯一解; 已知指数函数满足同一方程和初值, 所以极限为 $e^t$. 本次论证给出的是一个受控局部区间, 而不是从形式级数直接假定解存在.
\end{solution}

\originalsection{练习与解答}

\emph{以下均为编者练习. 题目分别检查一致控制、压缩定理的假设、积分方程的求解与初值问题的唯一性.}

\begin{exercise}\label{ex:uniform-x-square}
证明 $x\mapsto x^2$ 在每个有界闭区间上 Lipschitz, 但在 $\mathbb R$ 上不一致连续.
\end{exercise}
\begin{solution}
若 $x,y\in[-A,A]$, 则
$|x^2-y^2|=|x-y||x+y|\leq2A|x-y|$, 所以常数 $2A$ 有效; 任意有界闭区间都包含在这样的区间中. 在整个 $\mathbb R$ 上, 取 $x_n=n$, $y_n=n+1/n$. 输入距离为 $1/n\to0$, 输出距离为
\[
y_n^2-x_n^2=2+1/n^2\geq2.
\]
对 $\varepsilon=1$ 不存在统一的 $\delta$, 故不一致连续.
\end{solution}

\begin{exercise}\label{ex:banach-assumptions}
分别给出例子, 说明 Banach 不动点定理中 ``完备''、``自映射'' 和 ``$q<1$'' 三个条件在当前表述中不能直接删除.
\end{exercise}
\begin{solution}
在不完备空间 $X=(0,1)$ 上取 $T(x)=x/2$. 它是常数 $q=1/2$ 的压缩自映射, 却无不动点, 因为唯一可能的解 $x=0$ 不在 $X$ 中.

在完备空间 $X=[0,1]$ 上取 $T:X\to\mathbb R$, $T(x)=2+x/2$. 它满足同样的压缩距离估计, 但像不在 $X$ 中, 而不动点方程的解 $x=4$ 也不在 $X$. 因此只有距离估计、没有自映射条件并不够.

在完备空间 $\mathbb R$ 上取 $T(x)=x+1$. 它是 $1$-Lipschitz 自映射, 却无不动点. 同一空间上的恒等映射也是 $1$-Lipschitz, 却有无穷多个不动点. 因而将 $q<1$ 放宽为 $q\leq1$ 后, 存在性和唯一性都不能保证.
\end{solution}

\begin{exercise}\label{ex:scalar-contraction-error}
在 $\mathbb R$ 上解 $x=\cos x$. 证明恰有一个解, 并给出从 $x_0=0$ 开始迭代时保证误差小于 $\varepsilon$ 的先验条件.
\end{exercise}
\begin{solution}
任意解必在 $[-1,1]$ 中. $T(x)=\cos x$ 把闭区间 $[-1,1]$ 映入自身. 对 $x,y\in[-1,1]$, 中值定理给出
\[
|\cos x-\cos y|\leq\sin(1)|x-y|,
\]
因为 $|\sin t|\leq\sin(1)<1$ 对 $|t|\leq1$ 成立. 区间完备, 所以 Banach 定理给出唯一不动点 $x_*$. 区间之外不可能有解, 因此这也是 $\mathbb R$ 中的唯一解.

令 $q=\sin(1)$, $x_0=0$, 则 $x_1=1$, 先验误差界为
$|x_n-x_*|\leq q^n/(1-q)$. 取足够大的整数 $n$ 使
$q^n<\varepsilon(1-q)$, 即保证误差小于 $\varepsilon$. 无须预先知道 $x_*$ 的数值.
\end{solution}

\begin{exercise}\label{ex:constant-kernel}
设 $g\in C([0,1])$, $|\lambda|<1$. 求积分方程
\[
f(x)=g(x)+\lambda\int_0^1 f(y)\,\mathrm dy
\]
的唯一连续解. 当 $\lambda=1$ 时, 完整讨论是否有解及是否唯一.
\end{exercise}
\begin{solution}
记 $A=\int_0^1f(y)\,\mathrm dy$, $G=\int_0^1g(y)\,\mathrm dy$. 对方程两侧积分, 得 $A=G+\lambda A$. 当 $\lambda\ne1$ 时,
$A=G/(1-\lambda)$, 从而
\[
f(x)=g(x)+\frac{\lambda G}{1-\lambda}.
\]
直接代回可验证这是解, 推导也证明它唯一. 在 $|\lambda|<1$ 的范围, 这一结论与压缩定理一致; 公式还说明压缩条件不是必要条件.

若 $\lambda=1$, 关系 $A=G+A$ 要求 $G=0$. 若 $G\ne0$, 无解. 若 $G=0$, 对任意常数 $A\in\mathbb R$, 函数 $f=g+A$ 都满足积分为 $A$, 因而都是解. 所有解也必须是这种形式, 所以有无穷多个解.
\end{solution}

\begin{exercise}\label{ex:ode-nonunique}
考察初值问题 $y'=\sqrt{|y|}$, $y(0)=0$. 证明它有不止一个局部 $C^1$ 解, 并指出为何不与局部 Picard--Lindel\"of 定理矛盾.
\end{exercise}
\begin{solution}
零函数是一个解. 对任意 $c\geq0$, 定义
\[
y_c(t)=\begin{cases}
0,&t\leq c,\\
(t-c)^2/4,&t>c.
\end{cases}
\]
函数在连接点两侧的导数都趋于 $0$, 所以是 $C^1$ 函数. 当 $t\leq c$ 时两侧都为 $0$; 当 $t>c$ 时,
$y_c'(t)=(t-c)/2=\sqrt{|y_c(t)|}$. 因 $c\geq0$, 还满足 $y_c(0)=0$. 特别地 $y_0$ 与零函数在 $0$ 的任何邻域都不同, 因而局部唯一性失败.

$F(t,y)=\sqrt{|y|}$ 连续, 但在 $y=0$ 的任何邻域都不能对 $y$ 满足统一 Lipschitz 条件. 对 $u>0$, $v=0$, 其差商为
\[
\frac{|F(t,u)-F(t,0)|}{|u-0|}=\frac1{\sqrt u}\longrightarrow\infty.
\]
定理要求的 Lipschitz 假设不成立, 所以没有矛盾.
\end{solution}

\originalchapter{拓扑空间与连通性}\label{ch:topology}
\emph{本章以 18.S190 第6讲第1--2页的拓扑空间定义和第2讲第4页的连通性定义为起点. 基、分离性、积与商、连通性的完整论证为编者补充.}

\originalsection{拓扑公理与基}
度量同时规定距离的数值和开集的结构. 连续函数的逆像判据只涉及开集, 因而可以将开集本身作为基本对象. 这样做保留了连续性、紧致性和连通性, 但不再自动保留柯西序列、距离误差或收缩系数. 后面每次从度量空间推广到拓扑空间时, 都要检查证明是否曾使用半径或距离.

\begin{definition}
设 $X$ 是集合. 一个子集族 $\mathcal T\subseteq\mathcal P(X)$ 称为 $X$ 上的拓扑, 如果 $\varnothing,X\in\mathcal T$, 任意多个 $\mathcal T$ 中集合的并仍属 $\mathcal T$, 有限多个的交仍属 $\mathcal T$. $(X,\mathcal T)$ 称为拓扑空间. $\mathcal T$ 中的集合称为开集, 补集为开集的集合称为闭集. 点 $x$ 的邻域指包含一个含 $x$ 的开集的集合; 开邻域则还要求自身开.
\end{definition}
闭包 $\overline A$ 仍定义为所有包含 $A$ 的闭集的交. 等价地, $x\in\overline A$ 当且仅当 $x$ 的每个开邻域与 $A$ 相交: 若有不交的开邻域 $U$, 闭集 $X\setminus U$ 包含 $A$ 却不包含 $x$; 若 $x$ 不在某个闭超集 $F$ 内, 开邻域 $X\setminus F$ 就与 $A$ 不交. 内部是所有包含于 $A$ 的开集的并, 边界仍是 $\overline A\setminus A^\circ$. 这些定义只使用开闭结构, 因而适用于一般拓扑.

这里允许空并和空交, 分别是 $\varnothing$ 和 $X$. “任意并”与“有限交”不能交换: 在通常实数拓扑中, $\bigcap_{n\ge1}(-1/n,1/n)=\{0\}$ 不是开集.

\begin{example}
$\mathcal P(X)$ 是离散拓扑, $\{\varnothing,X\}$ 是平凡拓扑. 验证它们确为拓扑, 并说明度量空间的开集形成拓扑, 而至少含两点的平凡拓扑不能由度量产生.
\end{example}
\begin{solution}
离散拓扑包含所有子集, 而平凡拓扑只包含空集与全集, 两者均直接满足三个公理. 度量空间的开集形成拓扑 $\mathcal T_d$, 因为开集的任意并和有限交仍开. 若 $X$ 至少有两点, 平凡拓扑不能由度量产生: 度量中的半径小于两点距离的开球能够包含其中一点而排除另一点.
\end{solution}

\begin{definition}
一个开集族 $\mathcal B$ 称为拓扑 $\mathcal T$ 的基, 如果每个开集都是某些 $\mathcal B$ 中集合的并. 等价地, 对每个 $x\in U\in\mathcal T$, 都存在 $B\in\mathcal B$ 满足 $x\in B\subseteq U$.
\end{definition}
开球是度量拓扑的一组基. 基使检验连续性和写出开集更有效: 不必处理所有开集, 只须处理足够小的一批基本开集.

\begin{proposition}\label{thm:basis-criterion}
子集族 $\mathcal B$ 能作为某个拓扑的基, 当且仅当它覆盖 $X$, 且若 $x\in B_1\cap B_2$, 则有 $B_3\in\mathcal B$ 满足 $x\in B_3\subseteq B_1\cap B_2$. 此时由它生成的拓扑就是它的所有并组成的集合族.
\end{proposition}
\begin{proof}
若已有拓扑且 $\mathcal B$ 是基, 则把基的局部包含性质分别用于 $X$ 和开集 $B_1\cap B_2$, 即得两个条件. 反过来, 令 $\mathcal T$ 为基元素的所有并. 覆盖条件保证 $X\in\mathcal T$, 空并保证 $\varnothing\in\mathcal T$, 任意并的性质由并集结合律直接得到. 若 $U,V\in\mathcal T$ 且 $x\in U\cap V$, 可选 $B_1,B_2$ 使 $x\in B_1\subseteq U$ 和 $x\in B_2\subseteq V$. 再选 $B_3$ 使 $x\in B_3\subseteq U\cap V$. 将对所有 $x$ 选出的这些 $B_3$ 取并就得到 $U\cap V$, 故其属于 $\mathcal T$. 两两交闭合加上 $X$ 可推出所有有限交闭合.
\end{proof}

实线上端点为有理数的开区间也是一组基. 给定 $x\in(a,b)$, 有理数的稠密性使我们可取 $a<q<x<r<b$. 因而 $x\in(q,r)\subseteq(a,b)$. 这说明一组基不必包含所有开球, 只须能够在每一点局部缩小.

\begin{definition}
若对任意不同的点 $x,y\in X$, 都存在互不相交的开邻域 $U\ni x$ 和 $V\ni y$, 则称 $X$ 为 Hausdorff 空间. 若对每个点 $x$ 存在可数个开邻域 $U_1,U_2,\ldots$, 使任意开邻域 $U\ni x$ 都包含某个 $U_n$, 则称 $X$ 满足第一可数性.
\end{definition}
度量空间是 Hausdorff 空间: 取半径 $d(x,y)/3$ 的两个球, 三角不等式使两球不相交. 它也满足第一可数性: $B(x,1/n)$ 构成点 $x$ 的可数邻域基.

\originalsection{序列与连续映射}
在拓扑空间中定义 $x_n\to x$, 是指每个 $x$ 的开邻域都包含序列的所有充分靠后的项. 度量空间中这与 $\varepsilon$ 定义等价, 因为开球构成基. 在 Hausdorff 空间中极限唯一: 若序列同时收敛到不同的 $x,y$, 两个不相交的开邻域都应包含充分靠后的项, 产生矛盾. 一般拓扑空间中这个结论会失效; 平凡拓扑中每个序列收敛到每个点.

\begin{definition}
拓扑空间间的映射 $f:X\to Y$ 称为连续, 如果对每个 $Y$ 的开集 $V$, $f^{-1}(V)$ 都是 $X$ 的开集. 若 $f$ 是双射, 且 $f$ 与 $f^{-1}$ 均连续, 则称 $f$ 为同胚.
\end{definition}
复合连续映射仍连续, 因为 $(g\circ f)^{-1}(W)=f^{-1}(g^{-1}(W))$. 逆像保持任意并, 所以检查连续性时只需检查目标空间的基元素. 对闭集的逆像仍闭也是等价判据, 因为 $f^{-1}(Y\setminus F)=X\setminus f^{-1}(F)$.

\begin{proposition}\label{thm:first-countable-sequence}
连续映射保持序列收敛. 若 $X$ 满足第一可数性, 则反过来也成立: 若 $x_n\to x$ 总能推出 $f(x_n)\to f(x)$, 则 $f$ 连续. 同样, 在第一可数空间中, $x\in\overline A$ 当且仅当存在 $A$ 中序列收敛到 $x$.
\end{proposition}
\begin{proof}
若 $f$ 连续, 则 $f(x)$ 的每个开邻域 $V$ 的逆像都是 $x$ 的开邻域; $x_n\to x$ 保证 $x_n$ 最终在这个逆像中. 反过来, 若 $f$ 不连续, 存在 $Y$ 中开集 $V$, 使 $f^{-1}(V)$ 非开. 取 $x\in f^{-1}(V)$ 且没有包含于此逆像的 $x$ 的开邻域. 将 $x$ 的可数邻域基换成 $W_n=U_1\cap\cdots\cap U_n$, 可使其递减. 按 $x$ 的选法, 每个 $x$ 的开邻域都有点的像落在 $V$ 外. 从每个 $W_n$ 选 $x_n$ 满足 $f(x_n)\notin V$. 邻域基的性质保证 $x_n\to x$, 而像序列不收敛到 $f(x)$, 矛盾. 若 $x\in\overline A$, 每个 $W_n$ 都与 $A$ 相交, 选 $a_n\in A\cap W_n$ 即得 $a_n\to x$. 反向由收敛定义直接得到所有开邻域与 $A$ 相交.
\end{proof}
上述反向结论中, 第一可数性承担了从所有邻域中选出一条序列的作用. 不能把度量空间的序列判据不加条件地搬到一般拓扑空间.

\begin{example}\label{ex:cocountable}
在不可数集合 $X$ 上定义余可数拓扑: 非空开集是补集可数的集合. 验证这是拓扑, 并说明闭包和连续性的序列判据在这个空间中失效.
\end{example}
\begin{solution}
这是拓扑, 因为若一个并集非空, 可选其中一个非空成员 $U$, 并集的补集包含于可数集 $X\setminus U$; 空并或成员全为空的并仍为空集. 有限交只需在成员均非空时检查, 其补集是有限个可数集的并; 若有空成员则交为空. 在此拓扑中, 收敛序列必须最终恒等于极限. 若 $x_n\to x$ 但有无穷多项不同于 $x$, 令 $C=\{x_n:x_n\ne x\}$, 则 $X\setminus C$ 是 $x$ 的开邻域, 这些无穷多项都不在其中, 矛盾. 然而任意不可数真子集 $A$ 都满足 $\overline A=X$, 因为不可数的 $A$ 不可能包含在可数的开集补集中. 因此 $x\notin A$ 属于 $\overline A$, 却没有 $A$ 中序列收敛到它.

这也是一个非 Hausdorff 空间: 两个非空开集的交非空, 因为不可数集合删去两个可数集仍非空. 从该空间到同一集合的离散拓扑的恒等映射保持所有序列极限, 却不连续. 此例同时指出闭包序列判据和序列连续判据的适用边界.
\end{solution}

\originalsection{子空间、积空间与商空间}
\begin{definition}
若 $A\subseteq X$, 则 $\mathcal T_A=\{A\cap U:U\in\mathcal T_X\}$ 称为子空间拓扑. 对拓扑空间 $X,Y$, 由所有矩形 $U\times V$ 生成的拓扑称为 $X\times Y$ 的积拓扑, 其中 $U,V$ 分别开.
\end{definition}
子空间的开集不必在原空间中开. 例如 $[0,1]$ 中的 $[0,1/2)$ 是相对开集, 因为它等于 $[0,1]\cap(-1,1/2)$. 相对开闭必须注明所处空间.
矩形基满足命题\ref{thm:basis-criterion}, 因为矩形交仍是矩形. 若 $X,Y$ 是度量空间, 则
\[
 d((x,y),(x',y'))=\max\{d_X(x,x'),d_Y(y,y')\}
\]
是度量, 其半径 $r$ 的球正是 $B_X(x,r)\times B_Y(y,r)$. 反过来, 任一点在开矩形中时, 可取两个分量允许半径的最小值, 找到其中的积度量球. 因而这个度量恰生成积拓扑. 有限多个空间同理. 这里不把有限积结论误当成任意无限积的度量公式.

\begin{proposition}
映射 $h:Z\to X\times Y$ 连续, 当且仅当两个坐标映射 $\pi_X\circ h$ 和 $\pi_Y\circ h$ 连续.
\end{proposition}
\begin{proof}
投影的开集逆像是开矩形 $U\times Y$ 或 $X\times V$, 故投影连续. 若 $h$ 连续, 两个复合当然连续. 反过来, 开矩形的逆像等于两个坐标逆像的交, 因而开; 用矩形基检查连续性即可.
\end{proof}

\begin{definition}
设 $\sim$ 是 $X$ 上的等价关系, $Q=X/{\sim}$ 是等价类集合, $q:X\to Q$ 是自然满射. 定义 $V\subseteq Q$ 开, 当且仅当 $q^{-1}(V)$ 在 $X$ 中开. 这称为商拓扑. 更一般地, 满射 $q:X\to Q$ 若满足此开集等价判据, 则称为商映射.
\end{definition}
这个定义确实给出拓扑, 因为逆像保持任意并、有限交以及全集和空集. “像的开集”与“逆像的开集”作用不同: 商映射本身连续, 但连续满射未必是商映射.

\begin{proposition}\label{thm:quotient-universal}
若 $q:X\to Q$ 是商映射, $h:Q\to Y$ 是映射, 则 $h$ 连续当且仅当 $h\circ q$ 连续. 因而一个在每个等价类上常值的连续映射 $f:X\to Y$ 唯一诱导连续映射 $\widetilde f:Q\to Y$.
\end{proposition}
\begin{proof}
正向用复合连续性. 反向取 $Y$ 中开集 $V$, 由复合连续可知 $q^{-1}(h^{-1}(V))$ 开, 商拓扑定义给出 $h^{-1}(V)$ 开. 若 $f$ 在等价类上常值, 定义 $\widetilde f([x])=f(x)$ 是良定义的, 并满足 $\widetilde f\circ q=f$. 满射性使这一定义唯一.
\end{proof}

\begin{lemma}\label{thm:compact-hausdorff}
Hausdorff 空间的紧子集闭. 从紧空间到 Hausdorff 空间的连续双射是同胚.
\end{lemma}
\begin{proof}
设 $K$ 紧, $y\notin K$. 对每个 $x\in K$, 取互不相交的开邻域 $U_x\ni x$ 和 $V_x\ni y$. 有限个 $U_{x_i}$ 覆盖 $K$, 则 $\bigcap_i V_{x_i}$ 是 $y$ 的开邻域且与 $K$ 不交. 故 $X\setminus K$ 开. 空 $K$ 的结论也成立. 对连续双射 $f:X\to Y$, $X$ 的闭集 $F$ 紧: 给定 $F$ 的相对开覆盖, 先把各成员写成 $F\cap O_i$, 其中 $O_i$ 在 $X$ 中开. 集合族 $\{O_i\}$ 加上 $X\setminus F$ 是 $X$ 的开覆盖, 从其有限子覆盖中即可得到 $F$ 的有限子覆盖. 连续像 $f(F)$ 紧, 因而在 $Y$ 中闭. 于是逆映射的闭集逆像都闭, 故逆映射连续. 紧性与连续像论证只使用开覆盖, 在一般拓扑空间中同样成立.
\end{proof}

\begin{example}
将 $[0,1]$ 的两个端点视为等价, 其余点仅与自身等价. 证明商空间同胚于单位圆 $S^1$.
\end{example}
\begin{solution}
映射 $f(t)=(\cos 2\pi t,\sin 2\pi t)$ 连续, 并且其纤维正好是这些等价类, 故由命题\ref{thm:quotient-universal}诱导连续双射 $\widetilde f:[0,1]/{\sim}\to S^1$. 商空间是紧区间的连续像, 因而紧; $S^1$ 是欧氏空间子空间, 因而 Hausdorff. 引理\ref{thm:compact-hausdorff}保证 $\widetilde f$ 为同胚. 这说明粘合端点的几何直觉怎样变成可验证的拓扑论证.
\end{solution}

\originalsection{连通性与道路连通性}
\begin{definition}
若空间 $X$ 不能写成两个互不相交、非空开集的并, 则称 $X$ 连通. 这样的表示称为一个分离. 子集的连通性总相对于子空间拓扑定义. 若任意 $x,y\in X$ 都可用连续映射 $\gamma:[0,1]\to X$ 连接, 即 $\gamma(0)=x$, $\gamma(1)=y$, 则称 $X$ 道路连通.
\end{definition}
连通不意味着只有一个点. 它意味着没有开闭的方式把空间整体分成两部分. 空空间也按定义连通, 但许多连通像和并集论证会另外要求有关集合非空.

\begin{proposition}
空间 $X$ 连通当且仅当其同时开闭的子集只有 $\varnothing$ 和 $X$. 连通空间的连续像连通.
\end{proposition}
\begin{proof}
非空真子集 $A$ 同时开闭时, $A$ 和 $X\setminus A$ 给出分离. 反之, 分离的两部分互为补集, 所以都同时开闭. 若连续像 $f(X)$ 有分离 $U,V$, 它们在 $X$ 中的逆像是互不相交的非空开集并覆盖 $X$, 与连通性矛盾. 这里用到 $f:X\to f(X)$ 的连续性, 它由子空间开集逆像判据保证.
\end{proof}

\begin{theorem}\label{thm:interval-connected}
实数集 $I$ 连通当且仅当它是区间, 即对 $a,b\in I$ 且 $a<t<b$, 必有 $t\in I$. 空集和单点集包括在此约定中.
\end{theorem}
\begin{proof}
若 $a<t<b$, $a,b\in I$ 而 $t\notin I$, 则 $I\cap(-\infty,t)$ 与 $I\cap(t,\infty)$ 给出分离, 因而连通集必须是区间. 反向设 $I$ 是区间但有分离 $U,V$. 取 $a\in U,b\in V$, 必要时交换名称, 使 $a<b$. 令 $s=\sup(U\cap[a,b])$, 则 $s\in[a,b]\subseteq I$. 若 $s\in U$, 相对开放性给出 $\varepsilon>0$ 使 $I\cap(s-\varepsilon,s+\varepsilon)\subseteq U$. 又 $b\in V$, 故 $s<b$; 从 $s$ 右侧取一个仍在 $[a,b]$ 中的点会违反上确界. 若 $s\in V$, 相对开放性给出 $\varepsilon>0$ 使 $I\cap(s-\varepsilon,s+\varepsilon)\subseteq V$. 由于 $a\in U$, 有 $s>a$; 上确界的定义保证这小段里又有 $U$ 的点, 矛盾. 两种情形均不可能, 故区间连通.
\end{proof}
若 $f:I\to\mathbb R$ 连续, 其像是连通集, 因而是区间. 这立即给出介值定理, 并说明其背后的条件是定义域的连通性, 而不只是函数有局部极限.

\begin{proposition}
道路连通空间连通. 连通子集族若有公共点, 则其并连通.
\end{proposition}
\begin{proof}
若道路连通空间有分离, 取两个分量中的端点并用道路连接. 道路的像是连通区间的连续像, 不可能同时与分离的两部分相交, 矛盾. 若连通集合 $A_i$ 有公共点 $p$ 而并集有分离, 设 $p$ 属于其中一部分. 每个 $A_i$ 都因连通而必须完全在该部分, 另一部分因此空, 矛盾.
\end{proof}
欧氏空间中凸集道路连通, 因为 $\gamma(t)=(1-t)x+ty$ 处处留在凸集内. 更一般地, 星形集可先把 $x$ 连接到星形中心, 再连接到 $y$.

\begin{example}\label{ex:sine-curve}
连通不必道路连通. 令
\[
 G=\{(x,\sin(1/x)):0<x\le1\},\qquad S=G\cup(\{0\}\times[-1,1]).
\]
证明 $S$ 连通而非道路连通.
\end{example}
\begin{solution}
$G$ 是区间 $(0,1]$ 的连续像, 因而连通. 它在平面中的闭包正是 $S$: 对任意 $a\in[-1,1]$, 选 $\theta$ 使 $\sin\theta=a$, 则充分大的 $n$ 所对应的 $x_n=1/(\theta+2\pi n)$ 在 $(0,1]$ 中且趋于 $0$, 点 $(x_n,\sin(1/x_n))$ 收敛到 $(0,a)$; 其他新极限点不存在, 因为 $y$ 坐标总在该闭区间中, 而 $x>0$ 时由连续性仍在图像上. 连通集的闭包连通: 若闭包有分离, 连通的原集只能在一部分, 另一非空相对开部分应与稠密原集相交, 矛盾.

但是 $S$ 中不存在从竖直线段到 $G$ 的道路. 假设有道路 $\gamma(t)=(x(t),y(t))$, 从 $x(0)=0$ 到 $x(1)>0$. 连续函数 $x(t)$ 的零点集合是非空闭集, 因为它包含 $0$, 且是紧区间的闭子集. 因此它有最大值 $t_0$. 因为 $x(1)>0$, 有 $t_0<1$; 按最大值定义和 $S$ 中横坐标非负的条件, $x(t)>0$ 对 $t\in(t_0,1]$ 成立. 对任意小的 $\delta>0$, 从 $x(t_0)=0$ 到 $x(t_0+\delta)>0$ 的介值性保证该小时间段内 $x(t)$ 取得所有充分小的正值. 其中可选 $1/x=\pi/2+2k\pi$ 以及 $1/x=3\pi/2+2k\pi$, 所以 $y(t)$ 在任意这样的时间段内分别取得 $1$ 与 $-1$. 这与 $y$ 在 $t_0$ 连续矛盾. 因而 $S$ 连通而非道路连通.
\end{solution}

\originalsection{练习与解答}
以下练习由编者编写, 不作为官方试题.
\begin{exercise}
证明有限积 $X\times Y$ 在 $X,Y$ 连通时连通. 同时证明道路连通空间的有限积道路连通.
\end{exercise}
\begin{solution}
若某一因子为空, 积为空, 结论按定义成立. 否则固定 $(x_0,y_0)$. 对每个 $x\in X$, 子集 $\{x\}\times Y$ 同胚于 $Y$, 而 $X\times\{y_0\}$ 同胚于 $X$. 两者都连通且交于 $(x,y_0)$, 因而 $A_x=(\{x\}\times Y)\cup(X\times\{y_0\})$ 连通. 所有 $A_x$ 都含 $(x_0,y_0)$, 故其并即整个积空间连通. 若两因子道路连通, 分别选从 $x_1$ 到 $x_2$ 的道路 $\alpha$ 和从 $y_1$ 到 $y_2$ 的道路 $\beta$. 坐标连续性保证 $t\mapsto(\alpha(t),\beta(t))$ 为积空间的道路.
\end{solution}
\begin{exercise}
设 $X$ 是紧空间, $Y$ 是 Hausdorff 空间, $f:X\to Y$ 是连续满射. 证明 $f$ 是商映射. 为什么去掉紧性后, 连续双射不一定是同胚?
\end{exercise}
\begin{solution}
若 $f^{-1}(V)$ 开, 则 $F=X\setminus f^{-1}(V)$ 闭而紧. 由于 $f$ 满射, $f(F)=Y\setminus V$; 它紧, 在 $Y$ 中闭, 故 $V$ 开. 反向由连续性保证, 所以 $f$ 是商映射. 对反例, 取离散拓扑的 $\mathbb R$ 到通常拓扑的 $\mathbb R$ 的恒等映射. 它连续且双射, 但逆映射不连续, 因为单点集在离散拓扑中开而在通常拓扑中不开. 离散的无限实数空间不是紧空间, 其单点开覆盖无有限子覆盖.
\end{solution}
\begin{exercise}
设 $X$ 是连通空间, $f:X\to\{0,1\}$ 连续, 目标赋予离散拓扑. 证明 $f$ 是常值. 用这个判据证明任意至少有两点的离散空间不连通.
\end{exercise}
\begin{solution}
像是连通子集, 而离散二点空间唯一的非空连通子集是两个单点, 所以当 $X$ 非空时 $f$ 常值; 空定义域按通常约定也无非恒定映射. 若离散空间有不同点 $a,b$, 令 $f(a)=0$, 其余点值为 $1$. 所有逆像均开, 所以 $f$ 连续且非常值, 该空间不连通.
\end{solution}

\originalchapter{函数空间与分析中的拓扑}\label{ch:function-spaces}
\emph{本章补齐 18.S190 第6讲第3--5页关于 Banach 空间、线性泛函与 Dirichlet 问题的导论. Baire 定理和一致有界原理参考 MIT OCW 18.102 Spring 2021 第3讲 PDF 第3--4页(正文第16--17页); 函数族紧性与具体反例为编者补充, 不列为已核实的必考细项.}

\originalsection{有界线性算子与对偶空间}
范数空间既有度量又有线性运算. 距离 $d(u,v)=\lVert u-v\rVert$ 保证加法和数乘连续: 例如 $u_n\to u,v_n\to v$ 时,
$\lVert(u_n+v_n)-(u+v)\rVert\le\lVert u_n-u\rVert+\lVert v_n-v\rVert\to0$.
若 $\lambda_n\to\lambda$ 且 $u_n\to u$, 序列 $u_n$ 有界, 估计
$\lVert\lambda_nu_n-\lambda u\rVert\le |\lambda_n-\lambda|\lVert u_n\rVert+|\lambda|\lVert u_n-u\rVert$ 给出数乘连续. 这些运算在纯度量空间中未必有定义.

\begin{definition}
实或复范数空间 $E,F$ 间的线性映射 $T:E\to F$ 称为有界线性算子, 如果存在 $C\ge0$ 使 $\lVert Tx\rVert_F\le C\lVert x\rVert_E$ 对所有 $x\in E$ 成立. 算子范数定义为
\[
 \lVert T\rVert=\sup_{\lVert x\rVert_E\le1}\lVert Tx\rVert_F.
\]
所有有界线性算子构成空间 $\mathcal L(E,F)$. 当 $F$ 是底域 $\mathbb R$ 或 $\mathbb C$ 时, 它们称为连续线性泛函, 其空间记为 $E^*$. “线性泛函”也常泛指任意线性标量映射, 故本书谈对偶范数时明确要求连续或有界.
\end{definition}
若 $E\ne\{0\}$, 上确界也可在单位球面上取, 因为线性性允许缩放; 若 $E=\{0\}$, 用单位闭球的定义仍给出零范数, 避免空单位球面的约定问题. $\lVert T\rVert$ 是不等式中最小可用常数, 上确界无需取到.

\begin{proposition}
线性算子 $T:E\to F$ 有界, 当且仅当它在 $0$ 连续, 也当且仅当它处处连续. 有界算子满足 $\lVert Tx-Ty\rVert\le\lVert T\rVert\lVert x-y\rVert$.
\end{proposition}
\begin{proof}
有界性和线性性给出最后的不等式, 因而连续. 若在 $0$ 连续, 可取 $\delta>0$ 使 $\lVert z\rVert<\delta$ 时 $\lVert Tz\rVert<1$. 对 $x\ne0$, 令 $z=\delta x/(2\lVert x\rVert)$, 则 $\lVert Tx\rVert<2\lVert x\rVert/\delta$. 对 $x=0$ 不等式成立, 所以 $T$ 有界. 处处连续当然包含在 $0$ 连续.
\end{proof}

\begin{example}
在非退化区间 $[a,b]$ 的连续函数空间中, 求点值泛函与积分泛函的范数, 并比较不同定义域范数下微分算子的连续性.
\end{example}
\begin{solution}
在 $C([a,b])$ 中, 点值泛函 $f\mapsto f(t_0)$ 的范数是 $1$, 因为 $|f(t_0)|\le\lVert f\rVert_\infty$, 且常函数 $1$ 取等. 当 $a<b$ 时, 积分泛函 $I(f)=\int_a^b f(t)\,dt$ 的范数是 $b-a$, 由积分估计及常函数取等得到. 导数算子 $D:C^1([a,b])\to C([a,b])$ 若定义域使用 $\lVert f\rVert_{C^1}=\lVert f\rVert_\infty+\lVert f'\rVert_\infty$ 则有界, 但若只使用 $\lVert f\rVert_\infty$ 就不有界. 例如 $f_n(t)=\sin(n(t-a))/n$ 在上确界范数下趋于 $0$, 而 $\lVert f_n'\rVert_\infty=1$. 同一个代数算子的连续性取决于选用的范数.
\end{solution}

\begin{theorem}\label{thm:operator-banach}
若 $F$ 是 Banach 空间, 则 $\mathcal L(E,F)$ 在算子范数下也是 Banach 空间. 特别地, 任意范数空间 $E$ 的连续对偶 $E^*$ 都是 Banach 空间, 即使 $E$ 本身不完备.
\end{theorem}
\begin{proof}
算子范数的正定性、齐次性和三角不等式分别来自标量缩放及逐点估计, 所以它是范数. 设 $(T_n)$ 是算子范数下的柯西序列. 对每个 $x\in E$,
$\lVert T_nx-T_mx\rVert\le\lVert T_n-T_m\rVert\lVert x\rVert$, 因而 $(T_nx)$ 在 $F$ 中有极限, 定义 $Tx=\lim_n T_nx$. 将极限通过有限线性组合, 得到 $T$ 线性. 柯西序列 $(T_n)$ 有统一范数上界 $M$, 于是 $\lVert Tx\rVert=\lim_n\lVert T_nx\rVert\le M\lVert x\rVert$, 所以 $T$ 有界. 给定 $\varepsilon>0$, 当 $n,m\ge N$ 时 $\lVert T_n-T_m\rVert<\varepsilon/2$. 固定 $n\ge N$, 令 $m\to\infty$ 得 $\lVert(T_n-T)x\rVert\le(\varepsilon/2)\lVert x\rVert$, 再对单位闭球取上确界, 得 $\lVert T_n-T\rVert\le\varepsilon/2<\varepsilon$. 因此收敛在算子范数中成立.
\end{proof}
这里两个层次都不可省: 先借助目标空间的完备性定义逐点极限, 再用对所有单位向量统一的柯西估计证明算子范数收敛. 仅有逐点收敛不足以给出算子范数收敛.

\originalsection{Baire 定理与一致有界原理}
全有界性与完备性一起产生紧致性. 但无限维 Banach 空间的单位闭球通常不紧; 完备性单独还能控制什么? Baire 定理给出的控制是: 不能用可数多个处处稀疏的闭集覆盖一个非空完备空间.

\begin{definition}
度量空间 $X$ 中的子集 $A$ 称为无处稠密, 如果 $\overline A$ 的内部为空. 可数个无处稠密集的并称为第一纲集.
\end{definition}
闭集无处稠密, 等价于不包含任何非空开球. 无处稠密集不必为空; 例如实数中的整数集闭且内部空. 第一纲集则仍可能稠密. 有理数是可数个单点的并, 因而在 $\mathbb R$ 中是第一纲集, 同时又稠密.

\begin{theorem}[Baire 定理]\label{thm:baire}
在完备度量空间 $X$ 中, 可数个稠密开集 $G_n$ 的交仍稠密. 特别地, 非空完备度量空间不能是可数个无处稠密集的并; 若非空 $X=\bigcup_n F_n$ 且各 $F_n$ 闭, 则至少一个 $F_n$ 有非空内部.
\end{theorem}
\begin{proof}
若 $X$ 为空, 稠密交的结论按定义成立. 否则要证明稠密, 只须在每个非空开集 $U$ 中找到交集的点. 因 $G_1$ 稠密且开, $U\cap G_1$ 非空开. 取中心 $x_1$ 与 $0<r_1<1/2$, 使闭球 $\overline B_d(x_1,r_1)=\{x:d(x,x_1)\le r_1\}$ 包含在 $U\cap G_1$ 内. 这可做到: 从中心的一个已包含开球取更小半径. 这里闭球记号表示距离不大于半径的集合, 不假定它等于同半径开球的闭包.

归纳地, $B(x_n,r_n)\cap G_{n+1}$ 非空开, 取 $x_{n+1}$ 与 $0<r_{n+1}<2^{-n-1}$, 使
\[
 \overline B_d(x_{n+1},r_{n+1})\subseteq B(x_n,r_n)\cap G_{n+1}.
\]
这些非空闭球嵌套且直径趋于 $0$. 第5章的完备空间嵌套闭集定理给出唯一公共点 $x$. 若不调用该定理, 也可直接看出: 对 $m>n$, 有 $x_m\in\overline B_d(x_n,r_n)$, 故中心序列柯西; 由完备性有极限, 各闭球包含相应尾部所以也包含极限. 于是 $x\in U$ 且 $x\in G_n$ 对所有 $n$ 成立, 证明稠密性.

若 $A_n$ 无处稠密, 其闭包补集 $G_n=X\setminus\overline A_n$ 稠密且开. 非空 $X$ 中稠密交非空, 其中的点不在任何 $A_n$ 内, 故这些集合不能覆盖 $X$. 对闭集覆盖若每个 $F_n$ 内部均空, 就得到这样的不可能覆盖.
\end{proof}
反例的环境同样重要: $\mathbb Q$ 在通常度量下是可数个单点的并, 每个单点在 $\mathbb Q$ 内闭且内部空, 但 $\mathbb Q$ 不完备. 这不违反 Baire 定理.

\begin{theorem}[一致有界原理]\label{thm:uniform-boundedness}
设 $E$ 是 Banach 空间, $F$ 是范数空间, $T_n\in\mathcal L(E,F)$. 如果对每个 $x\in E$ 都有 $\sup_n\lVert T_nx\rVert<\infty$, 则 $\sup_n\lVert T_n\rVert<\infty$.
\end{theorem}
\begin{proof}
令 $A_k=\{x\in E:\lVert T_nx\rVert\le k\text{ 对所有 }n\}$, $k\ge1$. 每个 $A_k$ 是闭集, 因为它是闭集 $T_n^{-1}(\overline B_F(0,k))$ 的交. 逐点有界保证 $E=\bigcup_k A_k$. 若 $E=\{0\}$, 所有算子范数为零; 否则 Baire 定理仍对整个非空完备空间 $E$ 适用, 给出 $B(x_0,r)\subseteq A_k$, $r>0$. 对 $\lVert h\rVert<r$, $x_0+h$ 与 $x_0$ 均属 $A_k$, 故
$\lVert T_nh\rVert\le\lVert T_n(x_0+h)\rVert+\lVert T_nx_0\rVert\le2k$.
取 $h=(r/2)x$ 且 $\lVert x\rVert\le1$, 得 $\lVert T_nx\rVert\le4k/r$. 这个界同时适用于所有 $n,x$, 因而 $\lVert T_n\rVert\le4k/r$.
\end{proof}
本证明在整个 $E$ 中使用 Baire 定理. 若只在单位闭球中得到相对开球, 不能直接假定它包含范数空间中的整球; 必须额外处理边界. 将闭集 $A_k$ 定义在全空间中可以避开这个不必要的障碍.

\originalsection{连续函数族的紧性}
函数序列逐点有界并不足以有一致收敛子列. 在 $[0,1]$ 上, $f_n(x)=x^n$ 全部满足 $\lVert f_n\rVert_\infty=1$, 但逐点极限在 $1$ 不连续; 所以任何子列也不可能一致收敛到连续函数. 缺失的是所有函数共同的局部变化控制.

\begin{definition}
设 $K$ 为非空紧度量空间. 函数族 $\mathcal F\subseteq C(K,\mathbb R)$ 称为一致等度连续, 若对每个 $\varepsilon>0$ 存在 $\delta>0$, 使所有 $f\in\mathcal F$ 和所有 $x,y\in K$ 都满足
$d(x,y)<\delta\Rightarrow |f(x)-f(y)|<\varepsilon$.
称它一致有界, 若存在 $M$ 使 $|f(x)|\le M$ 对所有 $f,x$ 成立.
\end{definition}
紧空间上常用的逐点等度连续定义与本定义等价. 若每个 $x$ 都有适用于全族的半径 $r_x$ 使 $d(x,y)<r_x$ 时 $|f(x)-f(y)|<\varepsilon/2$, 则球 $B(x,r_x/2)$ 覆盖 $K$. 取有限子覆盖 $B(x_i,r_i/2)$, 并令 $\delta<\min_i r_i/2$. 若 $y$ 在第 $i$ 个小球内且 $d(y,z)<\delta$, 则 $y,z$ 都在 $B(x_i,r_i)$ 内, 用三角不等式给出 $|f(y)-f(z)|<\varepsilon$. 这样将逐点半径变为全域统一半径.

\begin{theorem}[Arzel\`a--Ascoli 判据]\label{thm:ascoli}
设 $K$ 为非空紧度量空间. $\mathcal F\subseteq C(K,\mathbb R)$ 的闭包在上确界范数下紧, 当且仅当 $\mathcal F$ 一致有界且一致等度连续.
\end{theorem}
\begin{proof}
若闭包 $H$ 紧, 则它在范数度量下有界, 即有统一上确界界. 给定 $\varepsilon>0$, 用有限个半径 $\varepsilon/4$ 的范数球覆盖 $H$, 中心为 $g_1,\ldots,g_m\in H$. 每个 $g_i$ 在紧 $K$ 上一致连续, 故存在统一的 $\delta>0$ 使 $d(x,y)<\delta$ 时 $|g_i(x)-g_i(y)|<\varepsilon/2$ 对所有中心成立. 任意 $f\in H$ 选相距小于 $\varepsilon/4$ 的中心, 得
$|f(x)-f(y)|<\varepsilon/4+\varepsilon/2+\varepsilon/4=\varepsilon$. 因而 $\mathcal F$ 一致等度连续.

反过来, 要利用第4章的“完备且全有界”等价判据. $C(K)$ 完备, 故 $\overline{\mathcal F}$ 完备. 只须证明 $\mathcal F$ 全有界, 因为全有界集的闭包仍全有界: 先用半径 $\eta/2$ 的有限球覆盖原集, 再用半径 $\eta$ 的同心球覆盖闭包.

固定 $\varepsilon>0$, 取等度连续半径 $\delta$ 使近点函数差小于 $\varepsilon/4$. 用有限个半径 $\delta$ 的球覆盖 $K$, 中心为 $x_1,\ldots,x_N$. 在有界区间 $[-M,M]$ 中分出有限个长度小于 $\varepsilon/4$ 的小区间, 把每个 $f\in\mathcal F$ 按 $N$ 个值 $f(x_i)$ 各属哪一个小区间分类; 端点可以用半开分割避免多重归属. 类别仅有限个. 对每个非空类别选一个代表 $g$. 若 $f$ 与 $g$ 同类, 则对任意 $x\in K$ 选 $i$ 满足 $d(x,x_i)<\delta$, 有
\[
 |f(x)-g(x)|\le |f(x)-f(x_i)|+|f(x_i)-g(x_i)|+|g(x_i)-g(x)|<3\varepsilon/4.
\]
所以这些有限代表的半径 $\varepsilon$ 球覆盖 $\mathcal F$. 若函数族为空, 结论直接成立. 故闭包全有界且完备, 从而紧.
\end{proof}

在 $[0,1]$ 上, 族 $\{f:\lVert f\rVert_\infty\le M,\ |f(x)-f(y)|\le L|x-y|\}$ 一致有界且等度连续. 两个不等式在一致极限下仍成立, 所以这个族还闭, 因而紧. 相比之下, $C([0,1])$ 的整个单位闭球不紧. 这是有限维“闭且有界即紧”和无限维函数空间之间最常见的分界.

\originalsection{流形与 Dirichlet 问题的动机}
18.S190 的终点还包括流形和偏微分方程的学习方向. 这里保留其与拓扑的关系, 不把导论提升为已证明的流形理论或椭圆方程存在性理论.

\begin{definition}
一个 $n$ 维拓扑流形通常定义为 Hausdorff、第二可数且局部同胚于 $\mathbb R^n$ 的空间. 第二可数指存在可数的全局开集基. “局部同胚”指每点有开邻域 $U$, 以及从 $U$ 到 $\mathbb R^n$ 某开子集的同胚. 光滑流形进一步要求所选坐标图之间的转移映射光滑.
\end{definition}
本定义并不直接把流形定义为嵌在欧氏空间中的光滑曲面. 例如圆周在每点附近同胚于开区间, 但圆周整体不与开区间同胚, 因为前者紧而后者不紧. 开圆弧上可用角度作局部坐标. 度量不是这里写出的定义项; 在常用的 Hausdorff、第二可数约定下有可度量化定理, 该定理超出本书证明范围. 局部坐标、全局拓扑和所选度量是不同层次.

平面区域 $\Omega$ 上的经典 Dirichlet 问题要求寻找 $u\in C^2(\Omega)\cap C(\overline\Omega)$, 使
\[
 \Delta u=\frac{\partial^2u}{\partial x^2}+\frac{\partial^2u}{\partial y^2}=0\quad\text{在 }\Omega,
 \qquad u|_{\partial\Omega}=f.
\]
存在性还取决于区域边界和边界数据的条件, 不能只凭 $f$ 连续就对任意区域宣布存在. 一个自然的研究方向是 Dirichlet 能量
$E(u)=\frac12\int_\Omega|\nabla u|^2\,dA$. 在积分有意义、且允许光滑紧支撑变分的函数类中, 若光滑 $u$ 满足 $E(u)<\infty$ 且已是能量极小元, 则对每个 $\phi\in C_c^\infty(\Omega)$,
\[
 0=\left.\frac{d}{dt}E(u+t\phi)\right|_{t=0}
   =\int_\Omega\nabla u\cdot\nabla\phi\,dA
   =-\int_\Omega(\Delta u)\phi\,dA.
\]
最后一步用紧支撑消除边界项. 由于 $\Delta u$ 连续, 若它在某点为正或负, 可选择支撑在 $\Delta u$ 保持同号的小邻域内的非零非负测试函数, 使最后的积分非零. 因此 $\Delta u=0$. 这个论证证明的是“光滑极小元存在时满足方程”, 并未证明极小元存在.

\begin{example}\label{ex:dirichlet-bumps}
构造零边界条件的光滑函数序列, 使其Dirichlet能量趋于零, 而一阶导数上确界范数趋于无穷.
\end{example}
\begin{solution}
小能量不能控制 $C^1$ 范数, 所以不能直接用 $C^1$ 紧性为上述极小元存在作保证. 取 $x_0\in\Omega$ 及一个非恒定的光滑函数 $\phi$, 其支撑紧含于单位圆盘内. 这样的函数可由 $\exp(-1/(1-|x|^2))$ 在 $|x|<1$ 定义、在外部置零并适当缩放构造; 所有阶导数在边界趋零是由指数衰减快于任意幂得到的. 选 $r_n\downarrow0$ 使 $B(x_0,r_n)\subseteq\Omega$, 并定义
\[
 u_n(x)=r_n^{1/2}\phi((x-x_0)/r_n).
\]
它们在边界附近为零. 对这些有紧支撑的函数, 本例使用 $\lVert u\rVert_{C^1}=\lVert u\rVert_\infty+\lVert\nabla u\rVert_\infty$, 各项均有限. 变量代换给出
\[
 E(u_n)=\frac{r_n}{2}\int_{\mathbb R^2}|\nabla\phi(y)|^2\,dy\longrightarrow0,
 \qquad \lVert\nabla u_n\rVert_\infty=r_n^{-1/2}\lVert\nabla\phi\rVert_\infty\longrightarrow\infty.
\]
因此零边界问题的极小能量是 $0$, 这些 $u_n$ 是趋近该能量的序列, 却没有 $C^1$ 收敛子列. 能量趋近下确界与适当拓扑中的紧性是两件不同的事. 为解决一般存在性问题需要选择更合适的函数空间与收敛方式, 而不是把能量有界直接当成一致导数有界.
\end{solution}

\originalsection{练习与解答}
以下练习由编者编写, 不作为官方试题.
\begin{exercise}
在 $C([0,1])$ 中, 令 $\mathcal F=\{f:f(0)=0,\ |f(x)-f(y)|\le|x-y|\}$. 证明 $\mathcal F$ 紧. 若去掉 Lipschitz 条件而只留下 $\lVert f\rVert_\infty\le1$ 与 $f(0)=0$, 结论是否仍成立?
\end{exercise}
\begin{solution}
从 $f(0)=0$ 得 $|f(x)|\le x\le1$, 所以一致有界. Lipschitz 条件保证一致等度连续. 这些条件在一致极限下保持, 故族闭; Arzel\`a--Ascoli 判据给出紧性. 去掉 Lipschitz 条件后, $f_n(x)=x^n$ 都属于新族. 若有一致收敛子列, 其连续极限也应等于所有子列的逐点极限, 即在 $[0,1)$ 为 $0$ 而在 $1$ 为 $1$, 不连续, 矛盾. 因而新族不紧.
\end{solution}
\begin{exercise}
设 $c_{00}$ 是只有有限个非零项的实数序列空间, 赋予上确界范数. 定义 $T_n(x)=nx_n$. 证明每个 $T_n$ 是有界线性泛函, 对每个固定 $x$ 有 $\sup_n|T_nx|<\infty$, 而 $\sup_n\lVert T_n\rVert=\infty$. 指出这为什么不反驳一致有界原理.
\end{exercise}
\begin{solution}
$|T_nx|\le n\lVert x\rVert_\infty$, 取第 $n$ 项为 $1$ 其余为零的序列即可取等, 所以 $\lVert T_n\rVert=n$. 对固定有限支撑序列, $T_nx$ 只有有限多个非零值, 故逐点有界. 但 $c_{00}$ 不完备: 序列 $x^{(m)}=(1,1/2,\ldots,1/m,0,\ldots)$ 是上确界范数下的柯西序列, 其任何可能极限的第 $k$ 项必为 $1/k$, 这个无限支撑序列不在 $c_{00}$ 中. 一致有界原理所需定义域完备性失效.
\end{solution}
\begin{exercise}
证明一个非零且完备的实或复范数空间, 若作为向量空间有一组可数无限 Hamel 基 $e_1,e_2,\ldots$, 将产生矛盾. 可调用第3章的有限维范数等价, 但不能把可数稠密集当成 Hamel 基.
\end{exercise}
\begin{solution}
令 $E_n=\operatorname{span}(e_1,\ldots,e_n)$. 有限维范数等价保证 $E_n$ 完备, 因而在完备空间中闭. 它是一个真线性子空间, 所以内部空: 若包含某个球 $B(x,r)$, 则 $x\in E_n$, 减去 $x$ 得到 $B(0,r)\subseteq E_n$; 再用缩放使任意向量落入该球, 可得 $E_n=E$, 矛盾. Hamel 基的定义要求每个向量是有限线性组合, 故 $E=\bigcup_n E_n$. 这与 Baire 定理矛盾. 可分空间可以有可数稠密集, 但稠密集允许取极限, 而 Hamel 基只允许有限线性组合, 二者不能混淆.
\end{solution}

\originalchapter{可数性与分离公理}\label{ch:separation}
\emph{本章沿18.901 Fall 2004的可数性、分离公理与可度量化教学安排展开. Notes D提供可数性比较的直接依据; 完整证明、反例及练习由编者补充. 不复制该课商业教材.}

\originalsection{闭包与连续性}
第7章已将连续性表述为开集逆像的性质. 下面重新检查依赖距离的概念, 使后续构造能够在一般拓扑空间中使用.
\begin{definition}
对拓扑空间$X$中的子集$A$, 内部$A^\circ$是包含于$A$的所有开集的并, 闭包$\overline A$是包含$A$的所有闭集的交. 边界为$\partial A=\overline A\setminus A^\circ$. 若$\overline A=X$, 称$A$在$X$中稠密.
\end{definition}
这些并与交分别开、闭, 因而内部是最大的开子集, 闭包是最小的闭超集. 点$x$属于$\overline A$当且仅当$x$的每个开邻域都与$A$相交: 若有开邻域$U$与$A$不交, 闭集$X\setminus U$包含$A$却不含$x$; 反之, 若$x$不在某个包含$A$的闭集$F$中, 则$X\setminus F$就是这样的邻域. 这个判据没有使用序列, 所以无需第一可数性.

\begin{proposition}\label{thm:closure-continuity}
映射$f:X\to Y$连续当且仅当对每个$A\subseteq X$有$f(\overline A)\subseteq\overline{f(A)}$.
\end{proposition}
\begin{proof}
若$f$连续, $x\in\overline A$, 则$f(x)$的每个开邻域$V$的逆像都是$x$的开邻域, 与$A$相交, 故$V$与$f(A)$相交. 反之, 设所述包含关系成立. 对闭集$F\subseteq Y$, 取$A=f^{-1}(F)$. 因为$\overline{f(A)}\subseteq F$, 有$\overline A\subseteq f^{-1}(F)=A$, 因而$f^{-1}(F)$闭, 得到连续性.
\end{proof}

\begin{lemma}[粘合引理]\label{lem:pasting}
设$X=\bigcup_i A_i$, 且各$f_i:A_i\to Y$连续并在重叠处相同. 若$\{A_i\}$是开覆盖, 或是有限闭覆盖, 则逐片定义的$f:X\to Y$连续.
\end{lemma}
\begin{proof}
对开覆盖, 开集$V\subseteq Y$的逆像是各相对开集$f_i^{-1}(V)$的并; 因为$A_i$本身开, 这些集合在$X$中也开. 对有限闭覆盖, 改用闭集$F\subseteq Y$: 各$f_i^{-1}(F)$在闭子空间$A_i$中闭, 故在$X$中闭, 有限并仍闭. 重叠处的一致性保证$f$良定义.
\end{proof}
有限这一条件不可任意删除. 将$\R$按单点分成无限闭覆盖时, 任意函数在每片上都连续, 并不能推出全局连续. 后面拼接道路时使用的是有限个闭区间, 正好符合引理条件.

\originalsection{三种可数性条件}
\begin{definition}
空间$X$若有可数的开集基, 称为第二可数. 若有可数稠密子集, 称为可分. 若每个开覆盖都有可数子覆盖, 称为Lindel\"of空间. 第一可数性只要求每点各自具有可数邻域基, 并不要求这些基的总数可数.
\end{definition}
这些条件分别控制全局开集、稠密点集和开覆盖. $\R^n$中有理中心、有理半径的球构成可数基: 对$x\in U$先取$B(x,\varepsilon)\subseteq U$, 再选有理中心$q$足够靠近$x$, 以及满足$|x-q|<r<\varepsilon-|x-q|$的有理半径$r$, 即有$x\in B(q,r)\subseteq U$.

\begin{proposition}\label{thm:second-countable}
第二可数空间第一可数、可分且Lindel\"of. 第二可数性传给子空间.
\end{proposition}
\begin{proof}
设$\{B_n\}$为可数基. 含点$x$的基元素给出其邻域基. 从每个非空$B_n$选一点$x_n$, 则任意非空开集包含某个非空$B_n$, 因而与$\{x_n\}$相交, 得到可分性. 对任意开覆盖$\mathcal U$, 将那些包含在某个$U\in\mathcal U$中的$B_n$列出, 对每个这样的$n$选择一个$U_n\supseteq B_n$. 任意$x$先属于某个$U$, 再属于其中某个基元素, 故被某个$U_n$覆盖. 最后, $\{A\cap B_n\}$是子空间$A$的可数基. 若$X$为空, 各结论直接成立.
\end{proof}

\begin{theorem}\label{thm:metric-countability}
度量空间中, 第二可数、可分、Lindel\"of三者等价.
\end{theorem}
\begin{proof}
只须证明另外两者各自推出第二可数. 若$D$可数稠密, 则以$D$中点为中心、正有理数为半径的球构成可数基, 证明与上面的欧氏空间论证相同, 只用三角不等式. 若空间Lindel\"of, 对每个正整数$n$从全体半径$1/n$的球覆盖中取可数子覆盖, 记其中心集为$D_n$. 可数集$D=\bigcup_nD_n$稠密: 给定非空开球$B(x,\varepsilon)$, 取$1/n<\varepsilon$, 覆盖$x$的球中心属于$D\cap B(x,\varepsilon)$. 再用已经证明的可分情形.
\end{proof}

\begin{example}
在$\R$上以$[a,b)$为基定义下限拓扑. 证明此空间第一可数且可分, 但不是第二可数.
\end{example}
\begin{solution}
点$x$的$[x,x+1/n)$形成可数邻域基. 每个非空基本开集都含有理数, 故$\mathbb Q$稠密. 假设有可数基$\mathcal B$. 对每个$x$选$B_x\in\mathcal B$满足$x\in B_x\subseteq[x,x+1)$. 若$x<y$, 则$B_y\subseteq[y,y+1)$不含$x$, 而$B_x$含$x$, 所以$B_x\ne B_y$. 这要求基含有不可数多个不同成员, 矛盾. 因而度量空间中的等价不能照搬到一般拓扑空间.
\end{solution}

\originalsection{正则性与正规性}
\begin{definition}
$T_1$指每个单点集闭. Hausdorff也称$T_2$. 本书称一个$T_1$空间为正则空间, 若点$x$与不含它的闭集$F$有互不相交的开邻域; 称为正规空间, 若任意两不相交闭集有互不相交的开邻域. 这里将$T_1$纳入约定, 以避免不同教材的术语差异.
\end{definition}
Hausdorff蕴含$T_1$: 对固定$y$, 每个$x\ne y$都有不含$y$的开邻域, 这些邻域的并为$X\setminus\{y\}$. 正则也蕴含Hausdorff, 因为第二个点构成闭集. 正规蕴含正则, 因为点本身闭.

\begin{lemma}\label{lem:shrink}
$T_1$空间正则当且仅当对$x\in U$且$U$开, 存在开集$V$满足$x\in V\subseteq\overline V\subseteq U$. $T_1$空间正规当且仅当对闭集$A\subseteq U$且$U$开, 存在开集$V$满足$A\subseteq V\subseteq\overline V\subseteq U$.
\end{lemma}
\begin{proof}
若有分离性质, 将点或闭集$A$与$X\setminus U$分离为开集$V,W$. 因为$W\cap V=\varnothing$且$W$开, 有$\overline V\subseteq X\setminus W\subseteq U$. 反之, 将$U$取为另一个闭集的补集, 用$V$与$X\setminus\overline V$作所需的互不相交开邻域.
\end{proof}

\begin{proposition}\label{thm:normal-examples}
度量空间正规. 紧Hausdorff空间正规. 正则Lindel\"of空间正规.
\end{proposition}
\begin{proof}
度量情形, 对非空不交闭集$A,B$, 函数$d(x,A)=\inf_{a\in A}d(x,a)$是$1$-Lipschitz, 且其零点集正是$A$. 开集$\{d(x,A)<d(x,B)\}$与$\{d(x,B)<d(x,A)\}$不交并分别包含$A,B$. 空闭集的情形取空邻域即可.

紧Hausdorff情形先用有限覆盖证明正则: 给定$x\notin F$, 对每个$y\in F$选互不相交开邻域$U_y\ni x$, $V_y\ni y$, 用有限个$V_y$覆盖闭而紧的$F$, 再将相应$U_y$取交. 给定闭集$A\subseteq U$, 用正则性对每个$a\in A$取开邻域$V_a$使$\overline{V_a}\subseteq U$. 有限个$V_a$覆盖紧集$A$, 其并的闭包仍包含于$U$, 用引理\ref{lem:shrink}得到正规性.

最后设$X$正则Lindel\"of, $A,B$不交且闭; 任一为空时用空集及全集分离即可. 闭子空间$A$也Lindel\"of: 将相对开覆盖写为$A\cap O_i$, 其中$O_i$在$X$中开, 用$\{O_i\}\cup\{X\setminus A\}$覆盖$X$, 取可数子覆盖再限制到$A$. 因而可以找到可数开族$U_n,V_n$分别覆盖$A,B$, 满足$\overline{U_n}\cap B=\varnothing$, $\overline{V_n}\cap A=\varnothing$. 必要时重复成员补成序列. 令
\[
 U=\bigcup_n\left(U_n\setminus\bigcup_{i\le n}\overline{V_i}\right),\qquad
 V=\bigcup_n\left(V_n\setminus\bigcup_{i\le n}\overline{U_i}\right).
\]
各括号内集合开, 且$U,V$分别包含$A,B$. 若一点同时属于第$n$个$U$片和第$m$个$V$片, 当$m\le n$时第一片已删去$V_m$的闭包, 当$n\le m$时第二片已删去$U_n$的闭包, 两种情形都矛盾.
\end{proof}

\originalsection{Urysohn引理与延拓}
用连续函数分离闭集比用开集分离更强地连接了拓扑与分析. 距离公式只适用于已知的度量, 一般正规空间则要用一串嵌套开集构造函数.
\begin{theorem}[Urysohn引理]\label{thm:urysohn}
设$X$正规, $A,B$是不交闭集. 存在连续$f:X\to[0,1]$满足$f|_A=0$, $f|_B=1$.
\end{theorem}
\begin{proof}
记$D$为$[0,1]$中的二进有理数. 先取$U_1=X\setminus B$, 再由引理\ref{lem:shrink}取$A\subseteq U_0\subseteq\overline{U_0}\subseteq U_1$. 每一层相邻已定义的$r<s$之间, 取$U_{(r+s)/2}$使$\overline{U_r}\subseteq U_{(r+s)/2}\subseteq\overline{U_{(r+s)/2}}\subseteq U_s$. 逐层进行后, 对所有$r<s$均有$\overline{U_r}\subseteq U_s$.

定义$f(x)=\inf\{r\in D:x\in U_r\}$, 集合为空时令$f(x)=1$. $A\subseteq U_0$给出$f(A)=0$, 而$B$不属于任何$U_r$, 给出$f(B)=1$. 对$0<a\le1$, 有$\{f<a\}=\bigcup_{r<a}U_r$: 若下确界小于$a$, 必有一个候选$r<a$; 反向立即成立. 对$0\le a<1$, 有
$\{f>a\}=\bigcup_{r>a}(X\setminus\overline{U_r})$.
右侧一点不在$\overline{U_r}$, 就不在任何$U_s$($s\le r$), 因而$f\ge r>a$. 反之, 若$f(x)>a$, 选$a<r<s<f(x)$的二进有理数; 由于$\overline{U_r}\subseteq U_s$且$x\notin U_s$, 有$x\notin\overline{U_r}$. 这些逆像全开, 超出$[0,1]$的$a$则给出空集或全集. 开区间是两个这样的射线的交, 故$f$连续.
\end{proof}

\begin{theorem}[有界Tietze延拓]\label{thm:tietze}
设$X$正规, $A\subseteq X$闭, $f:A\to[-M,M]$连续, $M\ge0$. 则存在连续$F:X\to[-M,M]$使$F|_A=f$.
\end{theorem}
\begin{proof}
$M=0$时取零函数. 先对$r>0$及连续$h:A\to[-r,r]$构造误差缩减项. 集合$C=\{h\le-r/3\}$和$D=\{h\ge r/3\}$在$A$中闭, 因$A$闭而在$X$中闭, 且不交. 用Urysohn引理得到$g:X\to[-r/3,r/3]$, 在$C,D$上分别取$-r/3,r/3$. 在$C$上有$g=-r/3$, $h\in[-r,-r/3]$, 故$|h-g|\le2r/3$; 在$D$上有$g=r/3$, $h\in[r/3,r]$, 同界成立; 其余点满足$|h|<r/3$, $|g|\le r/3$, 仍得$|h-g|\le2r/3$.

对$h_0=f$, $r_0=M$做此构造, 再令$h_{n+1}=h_n-g_n|_A$, $r_{n+1}=2r_n/3$. 每步都有$\lVert g_n\rVert_\infty\le r_n/3$, $\lVert h_{n+1}\rVert_\infty\le r_{n+1}$. 级数$F=\sum_{n\ge0}g_n$一致收敛, 因为$\sum r_n/3=M$. 其值在$[-M,M]$中, 在$A$上与$f$的误差不超过$r_{n+1}\to0$. 一致极限在任意拓扑定义域上连续: 给定$x$和$\varepsilon$, 先使一致尾差小于$\varepsilon/3$, 再取有限部分和在$x$的邻域使变化小于$\varepsilon/3$, 三角不等式控制总变化. 因而$F$为所需延拓.
\end{proof}
这里证明的是实值有界版本. 它不表示从闭子集到任意目标空间的映射都能延拓; 例如圆周恒等映射不能延拓为从圆盘到圆周的连续映射, 第\ref{ch:fundamental}章将给出阻碍.

\originalsection{可度量化}
这里先约定可数立方体$[0,1]^{\mathbb N}$的积拓扑: 基本开集只对有限个坐标指定相对开区间, 其余坐标任意. 坐标逆像的有限交构成这个基, 因而映入立方体连续当且仅当各坐标连续. 第\ref{ch:products}章再讨论任意积.
\begin{theorem}[第二可数正则空间的可度量化]\label{thm:metrization}
每个第二可数正则空间都同胚于可数立方体$[0,1]^{\mathbb N}$的一个子空间, 因而可度量化.
\end{theorem}
\begin{proof}
第二可数蕴含Lindel\"of, 故空间正规. 设$\{B_n\}$为可数基. 对所有满足$\overline{B_i}\subseteq B_j$的有序对$(i,j)$, 用Urysohn引理选连续$f_{ij}:X\to[0,1]$, 在$\overline{B_i}$上取$1$, 在$X\setminus B_j$上取$0$. 这些函数可数, 列成$f_1,f_2,\ldots$, 有限时重复, 空间为空时结论直接成立.

对任意$x\in U$且$U$开, 先选$x\in B_j\subseteq U$, 再由正则性选$x\in V\subseteq\overline V\subseteq B_j$, 最后选$x\in B_i\subseteq V$. 相应$f_{ij}(x)=1$, 且$\{f_{ij}>1/2\}\subseteq B_j\subseteq U$. 因而这些函数既分离不同点, 又能检测全部开邻域. 映射$\Phi(x)=(f_n(x))_n$连续且单射, 其像中每个原开邻域都含一个坐标开集的逆像, 所以逆映射也连续.

可数立方体上定义$\rho(u,v)=\sum_{n\ge1}2^{-n}|u_n-v_n|$. 正定性与三角不等式逐项成立. 小$\rho$距离强制任意固定有限组坐标接近; 反之, 先令尾和小于$\varepsilon/2$, 再约束前有限组坐标使其贡献小于$\varepsilon/2$, 可在每个点找到包含于$\rho$球的基本积邻域. 故该度量生成积拓扑. 将其通过$\Phi$拉回即可.
\end{proof}
第二可数只是充分条件, 不是可度量化的必要条件: 任意不可数集合的离散度量生成离散拓扑, 而它没有可数基. 第\ref{ch:products}章会进一步区分可数积与不可数积.

\originalsection{练习与解答}
以下题目与解答均为编者补充.
\begin{exercise}
证明正则性传给子空间. 指出上述正规性的证明为何不能同样直接搬到任意子空间.
\end{exercise}
\begin{solution}
在$A\subseteq X$中给定$x\in A\cap U$, 其中$U$在$X$中开. 在$X$中取$x\in V\subseteq\overline V^{\,X}\subseteq U$. 则$A\cap V$是所需相对开邻域, 且$\overline{A\cap V}^{\,A}\subseteq A\cap\overline V^{\,X}\subseteq A\cap U$. $T_1$也传给子空间. 正规性要收缩的是整个闭集而非一个点; $A$中的闭集一般不在$X$中闭, 不能直接调用$X$的正规性. 若$A$闭, 其闭集确实在$X$中闭, 这时正规性传给$A$.
\end{solution}
\begin{exercise}
设$X$第二可数, $A$为其子集. 证明$A$可分. 为什么仅知$X$可分不能用同一证明?
\end{exercise}
\begin{solution}
子空间$A$第二可数, 从其非空基元素中各选一点可得可数稠密集. 若只知$D$在$X$中稠密, $D\cap A$未必在$A$中稠密: 在$X=\R$, $D=\mathbb Q$, $A=\R\setminus\mathbb Q$时交为空. 因而不能靠截取原稠密集证明任意子空间可分.
\end{solution}
\begin{exercise}
设$A,B$是度量空间中不交非空闭集. 验证$u(x)=d(x,A)/(d(x,A)+d(x,B))$连续且分别在$A,B$上取$0,1$. 分母是否存在全空间统一的正下界?
\end{exercise}
\begin{solution}
两个距离函数连续且非负, 同时为零会使$x\in A\cap B$, 故分母逐点为正, 商连续. 在$A$上第一项为零而第二项正, 在$B$上相反. 统一正下界未必存在: 在$\R^2$中取$A=\{(t,0):t\in\R\}$及$B=\{(t,e^{-t}):t\in\R\}$. 两者闭且不交, 而在$(t,0)$处分母不超过$e^{-t}$, 当$t\to\infty$趋零. 连续性只需逐点非零.
\end{solution}

\originalchapter{积空间与商空间}\label{ch:products}
\emph{本章沿18.901的任意积、商拓扑、局部紧性与Tychonoff定理安排补充完整论证. 商空间分离性参照Notes G; 超滤子的证明、模型计算和练习为编者补充.}

\originalsection{任意积与箱拓扑}
给定一族空间$\{X_i\}_{i\in I}$, 集合$P=\prod_{i\in I}X_i$由每个坐标取一个元素的函数组成. 讨论非空积时, 默认在通常含选择公理的集合论框架内工作. 积上的开集不能简单地定义为任意开矩形, 因为那会对无限多个坐标同时施加约束.
\begin{definition}
积拓扑的基本开集为$\prod_iU_i$, 其中$U_i\subseteq X_i$开, 且除有限多个$i$外均有$U_i=X_i$. 不要求这个有限性条件所得的是箱拓扑. 坐标投影记为$\pi_i:P\to X_i$.
\end{definition}
基本矩形的有限交仍有上述形式, 所以构成基. 积拓扑是使所有投影连续的最粗拓扑: 每个基本矩形是有限个$\pi_i^{-1}(U_i)$的交, 任一使投影连续的拓扑必包含它们. 因而$f:Z\to P$连续当且仅当所有$\pi_i\circ f$连续; 两方向分别用复合连续性和有限交逆像判据证明.

\begin{example}
在$\R^{\mathbb N}$的积拓扑中, 证明序列收敛当且仅当逐坐标收敛. 说明映射$t\mapsto(t,t,\ldots)$在积拓扑连续, 在箱拓扑不连续.
\end{example}
\begin{solution}
积中收敛由连续投影推出各坐标收敛. 反向对只约束有限组坐标的基本邻域, 在每个坐标上取充分靠后的指标, 再取其最大值即可. 对所述映射, 每个坐标都是恒等函数, 因而积拓扑下连续. 箱开集$\prod_{n\ge1}(-1/n,1/n)$的逆像是$\{0\}$, 在$\R$中不开, 因而箱拓扑下不连续.
\end{solution}

\begin{proposition}\label{thm:countable-product-metric}
若$(X_n,d_n)$可度量化, 则可数积的积拓扑由
$d(x,y)=\sum_{n\ge1}2^{-n}\min\{1,d_n(x_n,y_n)\}$生成. 非空空间的任意积Hausdorff当且仅当各因子Hausdorff. 可数个第二可数空间的积第二可数.
\end{proposition}
\begin{proof}
$\min(1,a+b)\le\min(1,a)+\min(1,b)$保证每项的三角不等式, 级数正定且收敛. 给定有限坐标约束, 取足够小$d$球使对应每一项都小于所需阈值, 可保证球在基本矩形中. 反向先截尾使$\sum_{n>N}2^{-n}<\varepsilon/2$, 再用前$N$个坐标的小球使头部总和小于$\varepsilon/2$, 得到包含于$d$球的基本矩形.

若各因子Hausdorff, 两个不同积点必在某一坐标不同, 该坐标的分离邻域逆像将两点分离. 反向固定其他坐标, 单个因子作为积中的一个子空间嵌入, 而Hausdorff性传给子空间. 对第二可数性, 将每个因子的可数基与有限坐标集组合; 可数集的有限子集可数, 每组有限选择也可数, 所得基可数.
\end{proof}
不可数积通常不满足第一可数性. 例如$\{0,1\}^I$的$0$点若有可数邻域基, 将每个邻域缩成基本矩形, 所有受约束坐标的并仍可数. 取不在其中的$j$, 没有这些基本邻域能包含于$\{x:x_j=0\}$, 矛盾. 这也说明不能用一条序列代替任意积中的全部邻域信息.

\originalsection{积空间的紧致性}
设$X,Y$均紧. 若一因子为空, 积为空而紧, 以下设非空. 有限积的证明可以直接从开覆盖开始. 对$X\times Y$的开覆盖, 固定$x$, 用有限个开矩形覆盖紧切片$\{x\}\times Y$, 每个矩形含于某个覆盖成员. 将这些矩形的$X$分量取交, 得$x$的邻域$U_x$, 使$U_x\times Y$由同一有限组成员覆盖. 再从$\{U_x\}$选有限个覆盖紧$X$, 即得全积的有限子覆盖. 这个论证亦称管状邻域论证.

任意积需要一个能同时处理所有坐标的工具. 下面从有限交性质构造超滤子, 不将度量序列的对角法误用于不可数积.
\begin{definition}
集合$X$上的滤子$\mathcal F$是一个子集族: $X\in\mathcal F$, $\varnothing\notin\mathcal F$, 对有限交封闭, 且$A\in\mathcal F$, $A\subseteq B$时$B\in\mathcal F$. 在真滤子中极大的滤子称为超滤子. 给定$x\in X$, 若$x$的每个邻域都属于$\mathcal F$, 称滤子收敛于$x$.
\end{definition}
一个具有有限交性质的集合族生成真滤子: 取包含它的有限交的所有超集. 记这个生成滤子为$\mathcal F_0$. 对所有包含$\mathcal F_0$的真滤子按包含排序, 该集合非空; 任一非空链的并仍是真滤子且包含$\mathcal F_0$, 空链也有$\mathcal F_0$作上界. 用Zorn引理得到极大元, 它也是所有真滤子中的极大元, 因为任何更大真滤子仍包含$\mathcal F_0$. 这是此处调用选择公理的具体位置. 超滤子$\mathcal U$对任意$A\subseteq X$满足$A\in\mathcal U$或$X\setminus A\in\mathcal U$, 且两者不能同时属于它. 若$A\notin\mathcal U$, 在$\mathcal U$中加入$A$不能生成真滤子, 故存在$F\in\mathcal U$与$A$不交, 上封闭性给出$X\setminus A\in\mathcal U$.

\begin{lemma}\label{lem:ultrafilter-compact}
空间紧致当且仅当每个超滤子都至少收敛于一点. 这里开覆盖与闭集有限交性质的等价, 由取补集直接得到, 对任意拓扑空间成立.
\end{lemma}
\begin{proof}
若$X$紧, 超滤子中各集合的闭包有有限交性质, 因为原集合的有限交非空. 由紧性选$x\in\bigcap_{F\in\mathcal U}\overline F$. 若$x$的开邻域$V$不在$\mathcal U$中, 则$X\setminus V\in\mathcal U$, 其闭包就是自身, 与$x$的选取矛盾. 因而$\mathcal U$收敛于$x$.

反之, 对有有限交性质的闭集族扩充为超滤子, 取其极限$x$. 每个原闭集$F$必须含$x$, 否则开邻域$X\setminus F$也在滤子中, 与$F$的交为空. 因而闭集有限交性质推出共同交非空, 等价于紧性. 空空间没有真滤子, 结论也成立.
\end{proof}

\begin{theorem}[Tychonoff定理]\label{thm:tychonoff}
任意一族紧空间的积, 在积拓扑中仍紧.
\end{theorem}
\begin{proof}
若积为空则紧. 否则设$\mathcal U$为积上的超滤子. 对每个$i$, 令$\mathcal U_i=\{A\subseteq X_i:\pi_i^{-1}(A)\in\mathcal U\}$. 逆像保持交与补集, 故$\mathcal U_i$是超滤子. 紧性使其收敛于某个$x_i$. 选取这些极限得到$x=(x_i)$. $x$的基本邻域由有限个坐标邻域逆像相交得到, 各逆像均在$\mathcal U$中, 因而基本邻域及其超集均在$\mathcal U$中. 所以$\mathcal U$收敛于$x$, 用引理\ref{lem:ultrafilter-compact}完成证明.
\end{proof}
Hausdorff不是积紧致定理的条件, 它的作用是极限唯一和紧集闭等性质. 箱拓扑则会使结论失败: $\{0,1\}^{\mathbb N}$在箱拓扑是无限离散空间, 单点开覆盖无有限子覆盖, 而在积拓扑紧致.

\originalsection{商映射与粘合}
第7章给出的商泛性质是构造映射的主要工具. 这里进一步处理如何验证商映射和商空间的分离性.
\begin{proposition}\label{thm:quotient-open-closed}
连续满射若为开映射或闭映射, 则是商映射. 商映射的复合仍为商映射. 若$q:X\to Q$为商映射, $W\subseteq Q$开, 则限制$q:q^{-1}(W)\to W$仍是商映射.
\end{proposition}
\begin{proof}
开映射情形, 若$q^{-1}(V)$开, 则$V=q(q^{-1}(V))$开. 闭映射情形将同一论证用于补集. 对复合, 连续性和反向开集判据依次使用即可. 最后若$V\subseteq W$且其逆像在$q^{-1}(W)$相对开, 因$q^{-1}(W)$开, 该逆像在$X$中开, 故$V$在$Q$中开, 特别在$W$中开; 反向由连续性.
\end{proof}
限制到任意子集时不能无条件使用最后结论. 同样, 两个商映射的积也不能无条件宣布是商映射; 后面的道路商构造只使用已验证的紧到Hausdorff情形或开映射情形.

\begin{theorem}\label{thm:closed-equivalence}
设$X$紧Hausdorff, $R\subseteq X\times X$是闭的等价关系. 则商空间$Q=X/R$紧Hausdorff, 且商映射$q$闭.
\end{theorem}
\begin{proof}
对闭集$F\subseteq X$, 饱和集$q^{-1}(q(F))$是紧集$R\cap(X\times F)$在第一坐标的投影, 故紧而闭. 商拓扑的闭集判据于是说明$q(F)$闭. 商空间紧是连续像定理.

不同等价类$C,D$各为闭紧子集, 因为类$C$是$R$与一条坐标切片的交所给出的闭集. 在紧Hausdorff空间的正规性下, 选不交开集$U\supseteq C$, $V\supseteq D$. 集合$U'=X\setminus q^{-1}(q(X\setminus U))$开、饱和, 包含$C$且包含于$U$; 同样构造$V'$. 商空间中的$q(U'),q(V')$开, 因为逆像分别是$U',V'$, 它们不交并分别含两类. 因而$Q$为Hausdorff.
\end{proof}

\begin{example}
证明$[0,1]^2$将相对边按相同参数粘合的商空间同胚于环面$S^1\times S^1$. 证明$S^n$的对径点商$\mathbb RP^n$是紧Hausdorff空间.
\end{example}
\begin{solution}
映射$(s,t)\mapsto(e^{2\pi is},e^{2\pi it})$的纤维正是相对边粘合所生成的等价类, 诱导从紧商空间到Hausdorff环面的连续双射, 由第7章同胚判据得到结论. 对对径点商, 等价关系是对角线与映射$x\mapsto-x$的图像之并. 两者分别为连续函数$(x,y)\mapsto|y-x|$与$(x,y)\mapsto|y+x|$的零点集, 所以在$S^n\times S^n$中闭. 有限并闭, 用定理\ref{thm:closed-equivalence}即可.
\end{solution}
反例也说明闭关系条件的意义. 在$\R$中将所有非零点粘为一类, 保留零点一类. 非零类在商中开, 但零点类没有与其不交的开邻域, 商不Hausdorff. 等价关系不闭: $(1/n,1)$属于关系, 却趋于不属于关系的$(0,1)$.

\originalsection{局部紧性与一点紧化}
\begin{definition}
Hausdorff空间若每点都有一个紧邻域, 称为局部紧空间. 邻域不一定开, 但须包含该点的开邻域. 一点紧化是在$X$外加入一点$\infty$, 使其邻域通过紧集的补集表达.
\end{definition}
\begin{lemma}\label{lem:local-compact-shrink}
局部紧Hausdorff空间中, 对$x\in U$且$U$开, 存在开$V$满足$x\in V\subseteq\overline V\subseteq U$, 且$\overline V$紧. 因而这种空间正则.
\end{lemma}
\begin{proof}
取紧邻域$K$及开$O$使$x\in O\subseteq K$. 紧$K$在Hausdorff空间中闭, 且作为子空间正规. 在$K$中将$x$的开邻域$K\cap O\cap U$收缩为相对开$W$, 使$\overline W^{\,K}\subseteq O\cap U$. 因$W\subseteq O$且$O$在$X$中开, $W$也是$X$中开. $K$闭使$\overline W^{\,X}=\overline W^{\,K}$, 后者为紧集, 且包含于$U$. 单点闭, 正则性由引理\ref{lem:shrink}得到.
\end{proof}

\begin{theorem}\label{thm:one-point}
若$X$局部紧Hausdorff且不紧, 在$X^*=X\cup\{\infty\}$中保留$X$的开集, 并把$\{\infty\}\cup(X\setminus K)$($K\subseteq X$紧)作为$\infty$的开邻域基. 则$X^*$紧Hausdorff, $X$为其开稠密子空间.
\end{theorem}
\begin{proof}
紧集在$X$中闭, 两个所述无穷点邻域的交对应紧集的有限并; 它们与$X$开集的交仍开. 因而给出拓扑基, 且诱导原子空间拓扑. 对$X$中的两个点, 原Hausdorff邻域仍可用. 对$x$与$\infty$, 用引理\ref{lem:local-compact-shrink}取紧闭包的开邻域$V$; $V$与$X^*\setminus\overline V$分离二者. 任意开覆盖取一个含$\infty$的成员, 它的补集包含于某个紧$K$并闭, 因而紧; 从剩余成员取有限子覆盖. 故$X^*$紧. 若$\infty$有不与$X$相交的开邻域, 将导致$X$本身紧, 与假设矛盾. 所以$X$稠密.
\end{proof}
对$n\ge1$, $\R^n$的一点紧化是$S^n$. 将北极删去后, 球面点$(u,t)\in\R^n\times\R$经立体投影变为$u/(1-t)$; 逆映射为
$x\mapsto(2x/(1+|x|^2),(|x|^2-1)/(1+|x|^2))$.
两式连续互逆, 当$|x|\to\infty$时逆像趋北极. 更准确地, 北极邻域的补集在其余球面紧, 映到$\R^n$中紧, 所以延伸为一点紧化的连续双射; 紧到Hausdorff判据使其为同胚.

\originalsection{练习与解答}
\begin{exercise}
证明积投影是开映射, 但未必是闭映射.
\end{exercise}
\begin{solution}
非空因子的积中, 基本矩形投影到所选坐标是对应开集, 或为空; 一般开集是基本矩形的并, 故投影开. $\R^2$中的闭集$\{(x,y):xy=1\}$投影为$\R\setminus\{0\}$, 不闭. 因而开投影的商映射性质不能误说成它总是闭映射.
\end{solution}
\begin{exercise}
证明可数个紧度量空间的积是紧度量空间, 并给出一条不用超滤子的序列证明.
\end{exercise}
\begin{solution}
度量由命题\ref{thm:countable-product-metric}给出. 对积中的序列, 先在第一坐标取收敛子列, 再在其内部取第二坐标收敛子列, 依次进行. 从第$n$层子列中选一个原指标, 要求这些原指标严格递增, 并保留其属于前$n$层的性质. 所得对角子列每个固定坐标均收敛, 因而在积中收敛. 积序列紧, 用度量空间中序列紧与紧的等价得紧性. 此证明不能直接处理不可数坐标.
\end{solution}
\begin{exercise}
设$X$紧Hausdorff, $A$为非空闭集. 将整个$A$缩成一点、其余点各自保留. 证明商空间紧Hausdorff.
\end{exercise}
\begin{solution}
关系$R=\Delta_X\cup(A\times A)$闭, 因为Hausdorff空间的对角线闭: 对不同点取不交开邻域, 其积避开对角线, 故对角线的补集开. $A\times A$也闭. 用定理\ref{thm:closed-equivalence}. 同样的定义若用于非闭$A$, 不能靠这个结论保证Hausdorff性.
\end{solution}

\originalchapter{同伦与基本群}\label{ch:fundamental}
\emph{同伦、星形集和形变收缩取材于18.905第5讲, 绕数参照18.900第4讲. 基本群路线与18.904 Spring 2011的Lecture Summaries相衔接, 该课仅提供讲课要求; 本章群结构、圆周计算及应用的完整证明为编者补充.}

\originalsection{同伦与形变收缩}
同胚要求两空间的点与开集完全对应, 而圆周和穿孔平面虽不同胚, 却有相同的环绕障碍. 为表达连续变形, 必须把时间变量也纳入连续性的定义, 不能只要求每个时刻分别连续.
\begin{definition}
两个连续映射$f_0,f_1:X\to Y$间的同伦是连续映射$H:X\times[0,1]\to Y$, 满足$H(x,0)=f_0(x)$, $H(x,1)=f_1(x)$. 若对$A\subseteq X$总有$H(a,s)=f_0(a)=f_1(a)$, 称同伦相对于$A$. 两空间同伦等价, 指存在$f:X\to Y$, $g:Y\to X$使$gf$同伦于$\operatorname{id}_X$, $fg$同伦于$\operatorname{id}_Y$.
\end{definition}
同伦是等价关系. 恒定时间给出自反性, $s\mapsto1-s$给出对称性; 将两个同伦各压缩至半段时间再粘合给出传递性, 连续性由有限闭覆盖的粘合引理保证. 与连续映射前后复合仍是同伦, 因而同伦等价亦可复合.

\begin{definition}
若$A\subseteq X$且连续$r:X\to A$满足$r|_A=\operatorname{id}_A$, 称$r$为收缩映射. 若存在$H:X\times[0,1]\to X$从恒等映射变到包含映射与$r$的复合, 且对$a\in A$恒有$H(a,s)=a$, 称$A$为$X$的强形变收缩. 本书说“形变收缩”时采用这个保持$A$的版本. 非空空间的恒等映射若同伦于常值映射, 称空间可缩.
\end{definition}
收缩映射与第6章具有Lipschitz常数小于$1$的收缩算子是不同概念. 此处只表示把空间映回一个子空间, 不涉及距离缩小.

\begin{example}
证明$\R^n\setminus\{0\}$形变收缩到$S^{n-1}$, 而星形区域可缩. 说明同伦等价不保持紧致性.
\end{example}
\begin{solution}
令$r(x)=x/|x|$, 并取$H(x,s)=((1-s)+s/|x|)x$. 系数始终正, 所以不经过零; $x$在单位球面上时始终不动, 最后到达$r(x)$. 星形区域以$b$为中心时, $H(x,s)=(1-s)x+sb$留在区域内, 给出可缩性. 球面紧而穿孔空间不紧, 两者却因上述形变收缩而同伦等价.
\end{solution}

\originalsection{道路的运算}
记$I=[0,1]$. 从$x$到$y$的道路是连续$\alpha:I\to X$满足$\alpha(0)=x$, $\alpha(1)=y$. 道路同伦须固定两个端点: $H(0,s)=x$, $H(1,s)=y$. 若$\alpha$从$x$到$y$, $\beta$从$y$到$z$, 定义先走$\alpha$再走$\beta$的拼接
\[
 (\alpha*\beta)(t)=
 \begin{cases}\alpha(2t),&0\le t\le1/2,\\ \beta(2t-1),&1/2\le t\le1.
 \end{cases}
\]
相交点处两值一致, 粘合引理保证连续. 逆道路为$\overline\alpha(t)=\alpha(1-t)$, 常值道路记$c_x$.

\begin{lemma}\label{lem:path-reparameter}
若$\varphi:I\to I$连续且$\varphi(0)=0$, $\varphi(1)=1$, 则$\alpha\circ\varphi$与$\alpha$固定端点同伦. 拼接在道路同伦类上良定义, 满足结合律、常值单位律以及$[\alpha*\overline\alpha]=[c_x]$.
\end{lemma}
\begin{proof}
同伦为$H(t,s)=\alpha((1-s)\varphi(t)+st)$. 若两段道路分别作固定端点同伦, 在两个半段上拼接这些同伦即可, 中间端点始终一致. 结合律的两种三段拼接只改变走各段的时间分配. 以每段占$1/3$的道路作参照, 将两种分配的折点对应作分段线性重参数, 用刚给出的同伦比较. 拼上常值段也由这种重参数同伦移去等待时间.

最后$\alpha*\overline\alpha$等于$\alpha(u(t))$, 其中$u(t)=\min(2t,2-2t)$. 取$H(t,s)=\alpha((1-s)u(t))$, 两端恒为$x$, 最后全为$x$, 得到逆道路消去律.
\end{proof}

\begin{definition}
以$x_0$为基点的环路是起点终点均为$x_0$的道路. 其固定端点同伦类的集合记$\pi_1(X,x_0)$, 乘法为$[\alpha][\beta]=[\alpha*\beta]$, 称为基本群.
\end{definition}
前一引理已证明群公理: 常值环路是单位, 反向环路给逆元. 基本群一般不交换. 两个圈粘在一点的空间中, “先走左圈再走右圈”与相反顺序并非总能固定基点变形, 第\ref{ch:vankampen}章将严格计算.

\originalsection{映射与基点}
\begin{proposition}\label{thm:pi1-functor}
连续映射$f:(X,x_0)\to(Y,y_0)$诱导群同态$f_*[\alpha]=[f\circ\alpha]$. 有$(gf)_*=g_*f_*$, $(\operatorname{id}_X)_*=\operatorname{id}$. 若$f_0,f_1$间的同伦保持基点$x_0$的像为$y_0$, 则$(f_0)_*=(f_1)_*$.
\end{proposition}
\begin{proof}
把道路同伦与$f$复合仍固定端点, 故定义与代表无关. 拼接与复合逐段相容, 故是同态. 复合和恒等式直接检查. 对保持基点的同伦$H$, 映射$(t,s)\mapsto H(\alpha(t),s)$就是像环路的固定端点同伦.
\end{proof}

\begin{proposition}\label{thm:basepoint-change}
从$x_0$到$x_1$的道路$u$给出同构
$c_u:\pi_1(X,x_1)\to\pi_1(X,x_0)$,
$c_u[\alpha]=[u*\alpha*\overline u]$.
若$H$为$f_0$到$f_1$的不保持基点同伦, $u(s)=H(x_0,s)$, 则$(f_0)_*=c_u\circ(f_1)_*$, 两边目标均为$\pi_1(Y,f_0(x_0))$.
\end{proposition}
\begin{proof}
相邻的$\overline u*u$可按引理\ref{lem:path-reparameter}消去, 因而$c_u$保持乘法, 逆同构为$c_{\overline u}$. 对第二式, 将正方形$(t,s)\mapsto H(\alpha(t),s)$的边界读作$f_0\alpha*u*\overline{f_1\alpha}*\overline u$. 这条边界环路可缩: 将正方形边界以其起点为中心在凸正方形内作直线收缩, 再与$H(\alpha(t),s)$复合. 因而边界代表单位, 整理得到所述公式.
\end{proof}
道路连通空间的不同基点给出同构群, 但同构依赖道路, 并非不带选择的相同群. 若两基点道路不同, 相应同构可能相差内共轭. 可缩且道路连通的空间基本群平凡: 常值映射诱导零同态, 上式表明恒等同态与它经基点改变相同, 故所有类均为单位. 对非空凸集还可直接以基点为中心作环路的直线收缩.

\begin{proposition}\label{thm:pi1-product}
有限积的基本群为各因子基本群的积. 形变收缩$A\subseteq X$对每个$a_0\in A$诱导$\pi_1(A,a_0)\cong\pi_1(X,a_0)$.
\end{proposition}
\begin{proof}
积中的环路和固定端点同伦分别等价于各坐标的环路与同伦, 拼接也逐坐标进行, 所以投影给出双射且保持群运算. 对形变收缩, 包含$i$与$r$均保持$a_0$, 满足$ri=\operatorname{id}_A$, $ir$保持基点同伦于恒等, 因而$i_*$与$r_*$互逆.
\end{proof}

\originalsection{圆周与绕数}
将平面与复数域对应, 记$p:\R\to S^1$, $p(t)=e^{2\pi it}$. 任意长度小于$1$且避开端点识别的角区间上, $p$都是到一段开圆弧的同胚. 开圆弧的全部逆像是这些角区间的整数平移, 彼此不交. 这使“展开角度”有严格含义.
\begin{lemma}\label{lem:circle-lift}
对道路$\alpha:I\to S^1$及$a\in\R$满足$p(a)=\alpha(0)$, 存在唯一连续$\widetilde\alpha:I\to\R$使$p\widetilde\alpha=\alpha$, $\widetilde\alpha(0)=a$.
\end{lemma}
\begin{proof}
以开圆弧覆盖$S^1$, 其道路逆像为$I$的开覆盖. 紧区间的Lebesgue数引理给出有限分割, 使每个闭小段的像包含于一段开圆弧. 在第一小段使用含$a$的角区间逆支. 在后续小段选含前一段末值的整数平移逆支. 段与段端点相同, 故粘成连续提升. 若有两个提升, 差值逐点为整数且连续, 由$I$连通而恒定; 初值差为零, 故完全相同.
\end{proof}

对基于$1$的环路, 取初值$0$的提升, 其末值为整数, 称为绕数$w(\alpha)$. 更一般的闭曲线可取任意初始角, 末角减初角仍为整数且与初值整数平移无关.
\begin{lemma}\label{lem:winding-homotopy}
绕数在固定基点同伦下不变, 且$w(\alpha*\beta)=w(\alpha)+w(\beta)$.
\end{lemma}
\begin{proof}
拼接的提升先走$\widetilde\alpha$, 再走$w(\alpha)+\widetilde\beta$, 因而末值相加. 为证明同伦不变, 先比较足够接近的两个环路. 若$|\alpha(t)-\beta(t)|<1$对全部$t$成立, 比值$\beta(t)/\alpha(t)$落在含$1$且不含$-1$的圆弧中, 可唯一写作$e^{2\pi i\theta(t)}$, 其中$\theta(t)\in(-1/2,1/2)$连续. 基点条件给$\theta(0)=\theta(1)=0$. 因而$\widetilde\alpha+\theta$是$\beta$初值为零的提升, 末值相同.

若$H:I^2\to S^1$为同伦, 由紧正方形上一致连续性, 可将时间$s$分成有限小段, 使相邻两时刻的环路在全部$t$上都距离小于$1$. 逐次使用前述比较, 得首尾绕数相同.
\end{proof}

\begin{theorem}\label{thm:pi1-circle}
绕数给出同构$\pi_1(S^1,1)\cong\mathbb Z$. 因而穿孔平面的基本群为$\mathbb Z$, 环面的基本群为$\mathbb Z^2$.
\end{theorem}
\begin{proof}
良定义和同态性质已证. 环路$t\mapsto e^{2\pi int}$的绕数为$n$, 故满射. 若两个环路绕数相同为$n$, 其提升都从$0$走到$n$. 在$\R$中取提升的直线同伦, 再与$p$复合, 两端始终映到$1$, 得固定基点同伦. 因而绕数也单射. 穿孔平面用径向形变收缩, 环面用积的基本群计算.
\end{proof}
绕数不只描述圆周内的环路. 对避开$a\in\R^2$的闭曲线$\gamma$, 将$(\gamma-a)/|\gamma-a|$投到圆周即可定义绕$a$的数. 若曲线能在穿孔平面内缩成一点, 此数必须为零. 零绕数对避开一个点的空间足以判定可缩, 对避开多个点则未必足够, 因为基本群可能不交换.

\originalsection{圆盘的拓扑阻碍}
\begin{theorem}\label{thm:no-retract-disk}
不存在连续收缩$r:D^2\to S^1$, 其中$D^2=\{x\in\R^2:|x|\le1\}$.
\end{theorem}
\begin{proof}
若有, 包含$i:S^1\hookrightarrow D^2$满足$ri=\operatorname{id}_{S^1}$, 因而$r_*i_*$为$\pi_1(S^1,1)$的恒等同态. 圆盘凸, 基本群平凡, 复合必为零同态, 与$\pi_1(S^1,1)=\mathbb Z$矛盾.
\end{proof}

\begin{theorem}[二维Brouwer不动点定理]
每个连续$f:D^2\to D^2$至少有一个不动点.
\end{theorem}
\begin{proof}
若处处$f(x)\ne x$, 令$v=x-f(x)\ne0$. 从$f(x)$沿$x$的方向继续走, 在圆周处停止. 为验证这一几何构造的连续性, 解$|x+\lambda v|^2=1$并取非负根:
\[
 \lambda(x)=\frac{-x\cdot v+\sqrt{(x\cdot v)^2+(1-|x|^2)|v|^2}}{|v|^2},
 \qquad r(x)=x+\lambda(x)v.
\]
根号内非负, 分母非零, 各运算连续, 且代回二次方程得$|r(x)|=1$. 若$|x|=1$, 因$f(x)\in D^2$且$f(x)\ne x$, 有$x\cdot v=1-x\cdot f(x)>0$, 所以$\lambda(x)=0$, $r(x)=x$. 因而$r$为不可能存在的收缩, 矛盾.
\end{proof}

\originalsection{练习与解答}
\begin{exercise}
证明闭环带$\{x\in\R^2:1\le|x|\le2\}$的基本群为$\mathbb Z$, 但它与圆周不同胚.
\end{exercise}
\begin{solution}
径向映射将它形变收缩到内圆周, 故基本群相同. 删去圆周任一点后得到开弧, 仍连通; 环带删去一个内点后也连通, 所以这一粗判据不足. 改看删去两点: 圆周删去两个不同点不连通, 环带删去任意两点仍道路连通. 后者用极坐标具体连接. 设删点集为$A$, 选$r\in(1,2)$不等于任何删点的半径. 对保留点$x$, 先在半径$|x|$的圆周上沿足够短且避开$A$的弧, 走到不等于任何删点角度的角度$\theta_x$; 有限性及$x\notin A$保证可选. 再沿$\theta_x$的径向段走到半径$r$, 这段不经过删点. 对另一点同样处理, 再沿半径$r$的圆周连接中间点. 该圆周没有删点, 全部路径留在闭环带内. 因同胚保持删点后的连通性, 二者不同胚.
\end{solution}
\begin{exercise}
对整数$n$, 计算$f_n:S^1\to S^1$, $f_n(z)=z^n$所诱导的基本群同态. 推出$n\ne\pm1$时$f_n$不能是同伦等价.
\end{exercise}
\begin{solution}
绕数为$k$的提升$\widetilde\alpha$在复合后变为$n\widetilde\alpha$, 末角为$nk$, 故同态为$k\mapsto nk$. 同伦等价诱导基本群同构, 即使同伦逆不保持基点, 基点改变也仅加入同构; $\mathbb Z$的乘$n$同态只有$n=\pm1$才可逆. $n=0$为常值映射, 同样不可逆.
\end{solution}
\begin{exercise}
给出一条绕数为零但并非常值的圆周环路, 并显式给出固定基点收缩.
\end{exercise}
\begin{solution}
取$\alpha(t)=\exp(2\pi i\sin(2\pi t))$. 提升$\sin(2\pi t)$两端均为零, 绕数为零. 同伦$H(t,s)=\exp(2\pi i(1-s)\sin(2\pi t))$连续, 两端恒为$1$, 最后为常值. 零绕数描述同伦类, 不要求路径本身不移动.
\end{solution}

\originalchapter{覆盖空间}\label{ch:coverings}
\emph{本章与18.904的覆盖、提升及基本群教学要求对应. 官方摘要并非完整证明文本, 以下局部构造、提升定理、分类论证及练习均为编者补充.}

\originalsection{局部多层结构}
圆周的角度提升将一圈展开为直线. 覆盖空间把这个构造推广为每个小邻域上若干互不相交的副本, 使全局环路能够通过提升端点被检测.
\begin{definition}
连续满射$p:E\to X$称为覆盖映射, 若每个$x\in X$有开邻域$U$, 使$p^{-1}(U)$是互不相交开集$V_\lambda$的并, 且每个$p|_{V_\lambda}:V_\lambda\to U$为同胚. 称$U$被均匀覆盖, $V_\lambda$为其上的片. 集合$p^{-1}(x)$称为纤维.
\end{definition}
覆盖映射局部同胚, 因而为开映射: 任意开集与各片相交后映为开集, 再取并. 纤维为离散子空间, 但总空间未必离散. 例如$\R\to S^1$的每个纤维是$a+\mathbb Z$, 各片是短角区间; $z\mapsto z^n$($n\ge1$)是$n$层圆周覆盖. 相反, $z\mapsto z^n$作为$\mathbb C\to\mathbb C$在零点有分支, $n>1$时不是覆盖.

\begin{lemma}\label{lem:lift-unique}
若$p:E\to X$为覆盖, $Y$连通, 两个连续$f,g:Y\to E$满足$pf=pg$并在一点相同, 则处处相同.
\end{lemma}
\begin{proof}
在任意$y$附近, 选$pf(y)$的均匀覆盖邻域$U$, 再缩小$Y$的邻域, 使$f,g$的像各留在其所在片中. 若$f(y)=g(y)$, 两者在同一片, 由该片上$p$单射而在整个邻域相同. 若不同, 两者所在片不同, 在整个邻域不同. 因而相等集及其补集均开; 连通性与非空相等集给出结论.
\end{proof}

\originalsection{道路与同伦提升}
\begin{theorem}\label{thm:cover-path-lift}
给定道路$\alpha:I\to X$及$e_0\in p^{-1}(\alpha(0))$, 存在唯一道路$\widetilde\alpha:I\to E$使$p\widetilde\alpha=\alpha$, $\widetilde\alpha(0)=e_0$.
\end{theorem}
\begin{proof}
均匀覆盖邻域的逆像覆盖紧区间$I$, 用Lebesgue数取有限分割, 使每个小闭区间的像落在一个均匀覆盖邻域中. 在第一段用含$e_0$的片上逆映射, 后续各段用含前一末值的片. 端点一致, 有限粘合给出连续提升. 唯一性由引理\ref{lem:lift-unique}.
\end{proof}

\begin{theorem}[环路同伦提升]\label{thm:cover-homotopy-lift}
设$H:I^2\to X$连续, 已给定底边$H(t,0)$的提升. 则存在唯一$\widetilde H:I^2\to E$提升$H$并与该底边提升一致. 若$H$是固定端点的道路同伦, 提升的两个端点在同伦时间中也不动.
\end{theorem}
\begin{proof}
均匀覆盖邻域在$H$下的逆像覆盖紧正方形, 用Lebesgue数取足够细的矩形网格, 使每个小闭格的像位于一个均匀覆盖邻域. 从底层由左向右, 再逐行向上构造. 在第一格, 已知底边提升落在一个片中: 底边连通, 不可能跨越互不相交的片. 用该片的逆映射提升整格. 后续格与已构造部分相交于下边、左边或两者, 这些边组成连通集, 已有提升均落在同一片中, 选此片的逆映射. 因而重叠处一致, 有限闭格粘合给出连续映射. 唯一性用引理\ref{lem:lift-unique}.

若端点在$X$中不动, 其提升随$s$连续且落在一个离散纤维中, 连通区间到离散空间的像为单点, 故提升端点不动.
\end{proof}
这条定理保证固定基点同伦的环路有相同的提升终点. 仅知每条道路可以提升并不足以得到这个结论, 必须用到同伦提升.

\begin{theorem}\label{thm:cover-injective}
覆盖诱导的$p_*:\pi_1(E,e_0)\to\pi_1(X,x_0)$单射. 其像正是从$e_0$提升后仍回到$e_0$的环路类.
\end{theorem}
\begin{proof}
若$\widetilde\alpha$是$E$中的环路, 则它本身就是$p\widetilde\alpha$从$e_0$出发的提升, 故像的类有闭提升. 反向, 有闭提升的$\alpha$就是该提升环路的像. 若$p\widetilde\alpha$固定端点同伦于常值环路, 提升该同伦. 最后一条道路提升是从$e_0$出发的常值道路, 两端始终固定, 故$\widetilde\alpha$在$E$中可缩. 因而核平凡, $p_*$单射.
\end{proof}

\originalsection{映射提升的判据}
\begin{definition}
空间局部道路连通, 指每点有由道路连通开集组成的邻域基. 空间半局部单连通, 指每个$x$有邻域$U$使$U$中基于$x$的环路在全空间$X$中可缩. 单连通指非空、道路连通且基本群平凡.
\end{definition}
局部道路连通空间的道路分支开: 每点的小道路连通邻域留在其道路分支内. 因而连通且局部道路连通的空间道路连通. 若每点有由单连通开集组成的邻域基, 称空间局部单连通. 局部单连通蕴含半局部单连通, 反向不能把“在$X$中可缩”误改为“在$U$中可缩”.

\begin{theorem}[提升判据]\label{thm:map-lifting}
设$Y$道路连通且局部道路连通, $f:(Y,y_0)\to(X,x_0)$连续, $p:(E,e_0)\to(X,x_0)$为覆盖. 则存在保持基点的连续提升$\widetilde f:Y\to E$, 当且仅当$f_*\pi_1(Y,y_0)\subseteq p_*\pi_1(E,e_0)$. 存在时唯一.
\end{theorem}
\begin{proof}
有提升时$f_*=p_*\widetilde f_*$, 得必要性. 反之, 对$y$选从$y_0$到$y$的道路$\gamma$, 定义$\widetilde f(y)$为$f\gamma$从$e_0$出发提升的终点. 若改选$\delta$, 则$f(\gamma*\overline\delta)$的类属于$p_*\pi_1(E,e_0)$, 其提升闭. 先提升$f\gamma$, 再提升$f\overline\delta$返回$e_0$, 将后一段反向, 由道路提升唯一性可知两种终点相同. 故定义不依赖选择.

为证连续, 给定$y$, 取$f(y)$的均匀覆盖邻域$U$, 再取道路连通开邻域$W\ni y$满足$f(W)\subseteq U$. 从$y$到$z\in W$的道路留在$W$内, 其像的提升留在含$\widetilde f(y)$的片$V$中. 于是对所有$z\in W$, $\widetilde f(z)=(p|_V)^{-1}f(z)$, 给出连续性. 唯一性由连通性和引理\ref{lem:lift-unique}.
\end{proof}

\begin{example}
对于$n\ge1$, 求$f:S^1\to S^1$何时能通过覆盖$p_n(z)=z^n$提升, 两者基点均取$1$.
\end{example}
\begin{solution}
$p_{n*}$是$\mathbb Z$上乘$n$, 像为$n\mathbb Z$. 若$f_*$是乘$d$, 提升判据给出$n\mid d$. 特别$f(z)=z^d$可提升为$z^{d/n}$. 一般连续$f$也按相同条件提升, 无需它具有幂函数形式. 提升的基点若未指定, 会有$n$种初始选择.
\end{solution}

\originalsection{普适覆盖的构造}
\begin{theorem}\label{thm:universal-cover}
若$X$道路连通、局部道路连通且半局部单连通, 则存在道路连通、单连通的覆盖空间$\widetilde X\to X$, 称为普适覆盖.
\end{theorem}
\begin{proof}
固定$x_0$. 将从$x_0$出发的道路按固定端点同伦分类, 类记$[\alpha]$, 终点投影$p[\alpha]=\alpha(1)$. 称道路连通开集$U$为合适邻域, 若其中环路在$X$中都可缩. 假设给出这种邻域基: 先在半局部单连通邻域中缩到道路连通开集, 对别的基点用$U$内道路改变基点, 可缩性不变.

对$\alpha(1)\in U$, 令$B([\alpha],U)=\{[\alpha*\beta]:\beta\text{是从}\alpha(1)\text{出发且留在}U\text{的道路}\}$. 任意$U$中点都有这样的$\beta$, 而同终点的两种$\beta$相差$U$内环路, 在$X$中可缩, 故$p$将$B([\alpha],U)$双射到$U$. 若两个这样的集合在$[\gamma]$处相交, 取终点的小合适邻域$W$包含于两原邻域的交, 则$B([\gamma],W)$包含于二者. 所以这些集合构成拓扑基.

在此拓扑中, $p$在每个$B([\alpha],U)$上为同胚: 若$[\gamma]$在基片$B([\alpha],U)$中, 且合适邻域$W$包含其终点并含于$U$, 则限制投影的逆像$(p|_{B([\alpha],U)})^{-1}(W)=B([\gamma],W)$; 这里是限制投影, 全局逆像通常有多片. 同时投影把基片中的小基片映成开邻域. 对固定$u\in U$, 用各$[\alpha]$终点为$u$的基片覆盖$p^{-1}(U)$. 两片若相交, 通过$U$内道路走回$u$, 得原两类相同, 故两片相等; 其余片互不相交. 因而$p$为覆盖.

道路$\alpha$有自然提升$t\mapsto[\alpha|_{[0,t]}]$, 各限制段重参数到$I$, 在$t=0$取常值类. 局部在一个合适$U$内变化时它落在同一基片, 并由片上的逆映射连续, 故整体连续. 它连接常值类与$[\alpha]$, 所以$\widetilde X$道路连通. 若$\widetilde X$中有基于常值类的环路, 其投影$\alpha$的自然提升终点为$[\alpha]$, 但也是常值类. 因而$\alpha$在$X$中可缩; 由$p_*$单射, 原环路也可缩, 得到单连通性.
\end{proof}
半局部单连通的用途出现在基片投影的单射性: 小邻域内不同路径只有在全空间可缩的前提下才被识别. 例如平面子空间$\bigcup_{n\ge1}C_n$, 其中$C_n$圆心为$(1/n,0)$、半径为$1/n$, 各圈共用原点, 称为Hawaiian earring. 原点每个邻域都包含足够小的整个圈. 固定$n$, 保留$C_n$而把其他圈压到原点给连续收缩: 在原点的连续性由所有足够小圈都留在小邻域保证, 在其他点可取不碰其他圈的小邻域. 因而该圈的一次环绕在全空间也不可缩. 半局部单连通在原点失效. 这特指平面子空间拓扑, 不能与抽象可数楔和的拓扑混淆.

\originalsection{子群与覆盖分类}
\begin{theorem}\label{thm:cover-classification}
在定理\ref{thm:universal-cover}的条件下, 基于$x_0$的道路连通覆盖, 按保持覆盖投影和基点的同胚分类, 与$G=\pi_1(X,x_0)$的子群$H$一一对应, 对应关系为$H=p_*\pi_1(E,e_0)$. 层数为指数$[G:H]$. 若不指定覆盖基点, 对应子群的共轭类.
\end{theorem}
\begin{proof}
给定$H$, 将同终点的道路$\alpha,\beta$视为等价, 当且仅当$[\alpha*\overline\beta]\in H$. 自反、对称、传递由子群和道路消去律得到, 其等价类记为$[\alpha]_H$. 沿普适覆盖证明中相同的基片构造拓扑, 因合适邻域的环路类为单位, 基片仍双射于$U$, 得覆盖$E_H\to X$. 自然道路提升使$E_H$道路连通. 基于$x_0$的环路提升闭恰好是其类属于$H$, 故对应子群正是$H$. $x_0$上纤维的类由$g h^{-1}\in H$判定相同, 即右陪集$Hg$, 其数目为$[G:H]$.

反过来, 对道路连通覆盖$E$及其子群$H$, 定义$F:E_H\to E$将$[\alpha]_H$送到从$e_0$提升$\alpha$的终点. 同终点的$\alpha,\beta$提升终点相同, 当且仅当$\alpha*\overline\beta$从$e_0$的提升闭, 即其类在$H$中. 因而$F$良定义且单射. 道路连通性使任意$e\in E$能从$e_0$沿道路到达, 投影后其提升终点为$e$, 故$F$满射. 为核对局部连续性, 取一个同时合适且被给定覆盖$E\to X$均匀覆盖的道路连通开邻域$U$: 先取均匀覆盖邻域, 再在其与半局部单连通邻域的交内缩小, 局部道路连通性保证可选. 均匀覆盖邻域限制到开子集仍均匀覆盖. 令$q:E_H\to X$, 基片$B$与给定覆盖的片$V$包含对应点, 道路提升唯一性给$F|_B=(p|_V)^{-1}\circ(q|_B)$. 这是两片间的同胚, 所以$F$连续且开, 结合双射得覆盖同胚. 最后在同一纤维改变基点, 以总空间中道路连接并投影为环路$u$, 由基点改变公式得到子群的共轭; 所有这种改变也都由提升环路实现.
\end{proof}

\begin{definition}
覆盖的甲板变换是同胚$T:E\to E$满足$pT=p$. 它们组成群$\operatorname{Deck}(p)$.
\end{definition}
连通覆盖的甲板变换由一点的像唯一确定, 因为它与另一变换是同一投影的两个提升, 可用引理\ref{lem:lift-unique}. 对正规子群$H\triangleleft G$, 构造$T_u[\alpha]_H=[u*\alpha]_H$. 正规性保证换代表后的共轭$u(\alpha\overline\beta)\overline u$仍在$H$中, 故良定义. 它逐片为同胚, 逆为$T_{\overline u}$. $u$在$H$中时作用恒等, 反之常值类的像改变; 由纤维上的传递作用及一点唯一性, 得$\operatorname{Deck}(E_H)=G/H$.

\begin{example}
分类所有连通圆周覆盖, 并计算$n$层覆盖的甲板群.
\end{example}
\begin{solution}
$\mathbb Z$的子群为$n\mathbb Z$($n\ge1$)及$\{0\}$. 前者由$z\mapsto z^n$实现, 后者由$\R\to S^1$实现, 分类定理保证没有遗漏. 前者甲板群为$\mathbb Z/n\mathbb Z$, 具体是乘$n$次单位根; 后者为$\mathbb Z$, 具体是整数平移. 这项分类指覆盖同胚类型, 不要求任意给定投影已经写成幂函数.
\end{solution}

\originalsection{练习与解答}
\begin{exercise}
证明$n$层覆盖中底空间每个点的纤维点数相同, 若底空间道路连通. 说明证明是否需要总空间连通.
\end{exercise}
\begin{solution}
选从$x$到$y$的道路, 将每个$x$纤维点提升到一个$y$纤维点. 反向道路的提升给出逆对应, 因道路提升唯一性, 两次复合回到原点. 所以纤维之间双射, 无需总空间连通. 无限纤维也有相同基数.
\end{solution}
\begin{exercise}
设$p:E\to X$覆盖且$X$单连通、局部道路连通. 证明每个道路连通的$E$必由$p$同胚到$X$.
\end{exercise}
\begin{solution}
$G$平凡, 唯一子群指数为$1$. 由分类或闭提升判据, 每个纤维只有一点: 若同纤维两点由道路连接, 投影为可缩环路, 同伦提升使终点回到起点. 因而$p$双射, 覆盖映射开且连续, 故为同胚. 没有总空间连通条件时可为多份$X$的不交并.
\end{solution}
\begin{exercise}
考虑$p:\R^2\to S^1\times S^1$, $p(s,t)=(e^{2\pi is},e^{2\pi it})$. 证明它是普适覆盖并求甲板变换.
\end{exercise}
\begin{solution}
小开弧的积逆像为互不相交角矩形, 投影逐片为同胚, 故是覆盖. $\R^2$凸且道路连通, 所以单连通. 平移$(s,t)\mapsto(s+m,t+n)$($m,n\in\mathbb Z$)都是甲板变换. 任意甲板变换对每点的位移为整数向量, 位移函数连续, 连通性使其恒定, 所以恰为这些变换, 群为$\mathbb Z^2$.
\end{solution}

\originalchapter{空间的拼接与基本群}\label{ch:vankampen}
\emph{van Kampen定理及曲面基本群与18.904的教学要求对应; 胞腔粘合的模型参照18.905第14讲. 本章自由积、网格论证、附加关系和完整题解为编者补充.}

\originalsection{自由群与自由积}
一个环路在不同开集之间穿行时, 可以分段记录每段的基本群元素. 分段记录首先形成一个词, 再由交集中的相容关系把不同记录识别.
\begin{definition}
集合$S$上的自由群$F(S)$由符号$s,s^{-1}$组成的有限约化词构成, 约化指没有相邻的$s s^{-1}$或$s^{-1}s$. 空词为单位, 乘法为拼接后消去相邻逆符号. 两群$A,B$的自由积$A*B$由来自两群的非单位元素交替组成的词构成, 同一因子的相邻元素合并为其乘积, 乘积为单位时删去.
\end{definition}
约化结果唯一可用逐字扫描说明. 用一串已经约化的符号作栈, 新符号若与栈顶相消就删去栈顶, 否则加入. 对自由积, 同因子栈顶与新元素先在因子中相乘, 单位即弹出, 非单位即替换. 在任意位置先合并或删去一个允许的相邻对, 再扫描, 与原词扫描结果一致; 这一事实对自由积只用因子内部的结合律, 对自由群只用逆符号消去. 因而任何消去次序都得相同的约化词. 特别$R(R(u)v)=R(uv)=R(uR(v))$, 其中$R$表示约化, 故乘法结合. 反向词并逐字取逆给逆元.

\begin{proposition}\label{thm:free-universal}
任意$S\to G$唯一扩充为$F(S)\to G$的同态. 给定同态$A\to G$, $B\to G$, 唯一得到$A*B\to G$的同态.
\end{proposition}
\begin{proof}
把词的每个符号映到指定元素, 按次序相乘. 相邻逆元的消去或同因子乘积的合并不改变结果, 故在约化词上良定义. 拼接对应乘法, 故为同态. 每个词由生成元素相乘构成, 唯一性随之成立.
\end{proof}
给群$G$的一组元素$R$, 记$\langle\!\langle R\rangle\!\rangle$为包含它们的最小正规子群, 即所有共轭$g r^{\pm1}g^{-1}$的有限乘积. 商群$G/\langle\!\langle R\rangle\!\rangle$将这些元素强制等于单位. 记号$\langle S\mid R\rangle$表示从$F(S)$添加关系$R$所得的群. 例如$\langle a,b\mid aba^{-1}b^{-1}\rangle\cong\mathbb Z^2$: 关系允许$a,b$交换, 每个词化成$a^m b^n$, 映到$(m,n)$表明不会产生别的相等关系.

\originalsection{van Kampen定理}
\begin{theorem}\label{thm:vankampen}
设$X=U\cup V$, $U,V$开且道路连通, $W=U\cap V$非空且道路连通, $x_0\in W$. 记包含诱导映射$i_U:\pi_1(W,x_0)\to\pi_1(U,x_0)$, $i_V:\pi_1(W,x_0)\to\pi_1(V,x_0)$. 则自然同态给出
\[
 \pi_1(X,x_0)\cong
 \frac{\pi_1(U,x_0)*\pi_1(V,x_0)}
 {\langle\!\langle i_U(w)i_V(w)^{-1}:w\in\pi_1(W,x_0)\rangle\!\rangle}.
\]
不要求$i_U,i_V$单射.
\end{theorem}
\begin{proof}
记右侧为$G$, 两因子到$G$的同态为$j_U,j_V$. 它们对来自$W$的环路给出相同元素. 包含到$X$的同态也有这种相容性, 因而自由积泛性质给出$\Phi:G\to\pi_1(X,x_0)$.

要构造逆映射, 对每个$z\in X$固定一条$c_z$从$x_0$到$z$: 当$z\in W$时选在$W$中, 当$z$只在$U$或只在$V$时选在相应开集中, 并令$c_{x_0}$常值. 若一段道路$\gamma$的像全在$U$, 给它赋值$j_U[c_{\gamma(0)}*\gamma*\overline{c_{\gamma(1)}}]$; 全在$V$时用$j_V$. 若同时在二者, 整个环路在$W$内, 两种赋值相同. 反向道路赋逆元素.

对环路$\alpha$, 将$I$细分使每段像落在$U$或$V$, 定义$\Theta(\alpha)$为这些赋值按次序的乘积. 在同一区域内加一个分点, 中间出现的$\overline c*c$可消去, 所以乘积不变. 两组分割有共同细分; 同一小段若被不同区域记录, 像必在$W$, 赋值仍相同. 因而$\Theta$与分割无关.

还须证明它与同伦无关, 这正是关系完备性的关键. 对固定基点同伦$H:I^2\to X$, 用$H^{-1}(U),H^{-1}(V)$的Lebesgue数将正方形分成有限小格, 每格像落在一个区域. 赋值于各有向格边. 单个格子的边界在其区域内可缩, 因为先将格边界在凸格内缩到一角, 再与$H$复合; 加上的$c_z$与逆道路互相消去. 因而每格都满足“下边乘右边等于左边乘上边”的群等式. 逐行逐格将一条网格道路移过小格, 每次使用这个等式, 得整个正方形下边乘右边等于左边乘上边. 两侧边在$x_0$恒定, 赋值为单位. 所以底环路与顶环路赋值相同.

由此$\Theta:\pi_1(X,x_0)\to G$良定义, 拼接分割表明它为同态. 对全在$U$或$V$内的环路, 它就是$j_U$或$j_V$. 因而$\Theta\Phi$在两个生成因子上为恒等, 在$G$上为恒等. 另一方面,$\Phi\Theta[\alpha]$把各段连同$c_z$拼接回$X$, 相邻逆道路消去后就是$[\alpha]$. 所以两映射互逆.
\end{proof}
开集条件保证有限细分和网格均存在, 交集道路连通保证可以在重叠处选择共同连接道路. 若交集不连通, 该版本不能直接应用, 即使图上两个区域看起来都很简单.

\originalsection{图与球面}
\begin{proposition}\label{thm:graph-pi1}
$k$个圆周粘在一点的楔和的基本群是$k$个生成元的自由群. 有限连通图有$V$个顶点、$E$条边时, 基本群是秩$E-V+1$的自由群.
\end{proposition}
\begin{proof}
对楔和, 取$U$为前$k-1$个圆周及第$k$圈在粘合点附近的短开弧, $V$为第$k$圈及其余圈在粘合点附近的短弧. 二者相对开, 分别形变收缩到前$k-1$圈与第$k$圈, 交集是一棵星状树, 可缩. van Kampen给出自由积, 由$\pi_1(S^1)=\mathbb Z$归纳.

一般连通图先选生成树$T$. 它有$V-1$条边: 从一个顶点起, 每次加入通向新顶点的边, 不造回路, 最后遍历全部顶点. 树可缩, 由不断将叶边线性压向其另一端得到. 再逐条加入其余边. 加入边的两端在$T$中由唯一简单道路连接. 对新图取$U$为去掉新边中点的空间, 它沿新边两侧开尾形变收缩到旧图; 取$V$为新边与该树中道路的一个小开邻域, 它形变收缩到所构成的圆周. 邻域在道路以外只含各岔边的短开尾, 这些尾均可线性压到所附顶点. $U\cap V$为删掉新边中点后的上述圆周连同短开尾, 可缩且道路连通. 故每加一条非树边, 基本群自由积上一个$\mathbb Z$. 一共$E-(V-1)$条, 得所述秩. 环边和重复边按独立边处理, 论证仍成立.
\end{proof}

\begin{proposition}\label{thm:sphere-simply-connected}
对$n\ge2$, $S^n$单连通. 对$n\ge2$, $\pi_1(\mathbb RP^n)=\mathbb Z/2\mathbb Z$.
\end{proposition}
\begin{proof}
以去掉南极、去掉北极的两个开集覆盖$S^n$, 各自经立体投影同胚于$\R^n$, 故单连通. 交集同胚于$S^{n-1}\times\R$, 在$n\ge2$时道路连通. van Kampen右侧为平凡群的自由积再取商, 仍平凡, 所以$S^n$单连通.

对径商$p:S^n\to\mathbb RP^n$为两层覆盖. 对点$x$取足够小开球帽$U$使$U\cap(-U)=\varnothing$, 饱和开集$U\cup(-U)$投影为开集, 且两个球帽各同胚到它; 此处用商判据及对径作用的开性. $S^n$单连通, 因而这是普适覆盖, 两层使底空间基本群只有两个元素, 即$\mathbb Z/2\mathbb Z$. 或用甲板群: 只有恒等与对径变换.
\end{proof}
当$n=1$时$\mathbb RP^1$本身为圆周, 基本群为$\mathbb Z$, 不能删掉$n\ge2$条件.

\originalsection{附加一个二维胞腔}
\begin{theorem}\label{thm:attach-cell-pi1}
设$A$为道路连通的紧Hausdorff空间, 沿连续$f:S^1\to A$附加一个圆盘, 得$X=A\cup_f D^2$. 以$f(1)$为基点, 则
$\pi_1(X)\cong\pi_1(A)/\langle\!\langle[f]\rangle\!\rangle$.
有限个圆盘依次附加时, 对全部附加环路取正规闭包.
\end{theorem}
\begin{proof}
在圆盘上以$r=|z|$为径向参数. 取$U$为$A$连同圆盘的$r>1/3$部分, $V$为圆盘内部$r<2/3$部分. 它们在商中开: 在不交并$A\sqcup D^2$中的饱和逆像开, 可用商判据. $U$将环带径向推到边界后形变收缩到$A$, $V$为开圆盘而可缩, 交集为开环带. 对$U$的形变在$A$上恒等、圆盘上沿半径移动, 并在边界与$f$一致. 为严格验证时间参数中的连续性, 注意$A\sqcup D^2$紧Hausdorff, 附加的等价关系闭: 它由对角线、边界与$f$的图像及其转置、满足$f(z)=f(z')$的边界点对组成, 各部分均闭. 定理\ref{thm:closed-equivalence}使$X$紧Hausdorff. 因而商映射与$I$之积也是从紧空间到Hausdorff空间的连续满射, 为商映射; 限制到开集$U\times I$仍为商映射. 径向同伦在不交并的相应开子集与$I$之积上连续并在纤维中常值, 可由商泛性质下降到$U\times I$.

交集的基本群由绕中心一圈生成, 包含到$U$并收缩到$A$后正是附加环路$f$, 包含到$V$则为单位. 用van Kampen得到所述正规闭包商. 基点若先取在交集中, 用径向道路改变到$f(1)$, 只改变附加环路的共轭代表, 正规闭包不变. 多个圆盘逐次重复此步骤.
\end{proof}
这个定理用于有限图和有限胞腔的附加, 下面的具体空间均满足所用的商邻域条件. 对更一般粘合, 不能将任意闭子空间的并直接当作本定理中的开覆盖.

\begin{example}
计算由边界词$aba^{-1}b^{-1}$粘合正方形所得环面, 以及由$a^2$附加圆盘所得实射影平面的基本群.
\end{example}
\begin{solution}
环面的一维骨架是两个圆周的楔和, 基本群$F(a,b)$. 填入面使边界词成为单位, 得$\langle a,b\mid aba^{-1}b^{-1}\rangle=\mathbb Z^2$, 与积空间计算一致. 射影平面的圆盘边界将对径点识别, 边界映到射影直线时绕行两次, 故附加词为$a^2$. 得$\langle a\mid a^2\rangle=\mathbb Z/2\mathbb Z$, 与覆盖计算一致.
\end{solution}

\originalsection{练习与解答}
\begin{exercise}
证明$\R^2\setminus\{(-1,0),(1,0)\}$的基本群是两个生成元的自由群. 给出绕两个点的绕数都为零但不可缩的环路类.
\end{exercise}
\begin{solution}
取$U=\{x<1/2\}\setminus\{(-1,0)\}$, $V=\{x>-1/2\}\setminus\{(1,0)\}$. 每个穿孔半平面同胚于穿孔平面, 基本群为$\mathbb Z$; 交集为竖直开条带而凸. van Kampen给出$F(a,b)$, $a,b$分别绕两点. 交换子$aba^{-1}b^{-1}$约化后仍是非空词, 所以不可缩. 它对每个点的总绕数均为零, 因为对应同态分别只记录$a,b$的指数和. 因而两个普通绕数不能检测非交换信息.
\end{solution}
\begin{exercise}
空间由一个圆周及沿$z\mapsto z^m$, $z\mapsto z^n$附加的两个圆盘组成, 其中$m,n\ge1$. 求基本群.
\end{exercise}
\begin{solution}
得到$\langle a\mid a^m,a^n\rangle$. 原圆周群为$\mathbb Z$, 被添加关系生成的子群为$m\mathbb Z+n\mathbb Z=d\mathbb Z$, $d=\gcd(m,n)$. 一个方向所有组合均为$d$倍数; 反向B\'ezout等式$um+vn=d$使$d$也在生成子群中. 故基本群为$\mathbb Z/d\mathbb Z$; 特别互素时单连通, 但不能据此断言空间可缩.
\end{solution}
\begin{exercise}
证明Klein瓶的边界词$aba^{-1}b$给出群$\langle a,b\mid aba^{-1}=b^{-1}\rangle$, 并证明此群不是交换群.
\end{exercise}
\begin{solution}
由附加胞腔定理得到该表示. 在$\R^2$上定义$A(x,y)=(x+1,-y)$, $B(x,y)=(x,y+1)$, 则$ABA^{-1}=B^{-1}$. 因而泛性质给出此表示群到这些变换群的同态. 两变换$AB$和$BA$分别将$(0,0)$送到$(1,-1)$和$(1,1)$, 不同, 所以群中$a,b$也不交换. 这一步用具体作用检验非交换性, 而不是仅看关系的外形.
\end{solution}

\originalchapter{曲面与欧拉示性数}\label{ch:surfaces}
\emph{有限复形、边界矩阵、可定向性及角亏依据Paul Seidel的18.900第29--31、33、40讲; 胞腔模型参照18.905第14、18讲. 第40讲为公开补充讲义, Spring 2023课堂未讲. 下文补齐略证并纠正符号笔误; 多边形模型及练习为编者补充.}

\originalsection{有限复形与曲面的局部条件}
多边形粘合能制造曲面, 但任意粘合不一定仍局部像平面. 三角剖分把这个要求转化为顶点周围的三角形如何排列.
\begin{definition}
有限二维单纯复形$K$由有限顶点集、无向边和三角形组成, 每条边的端点、每个三角形的边和顶点都属于$K$, 同一顶点集只允许一个单形. 几何实现$|K|$可取标准基向量为顶点, 将每个单形实现为它们的凸包; 两个单形仅在共同面相交. 若每条边恰邻接两个三角形, 且每个顶点的链环为一个圈, 称为闭组合曲面.
\end{definition}
顶点$v$的链环由与$v$共同成边的邻点, 以及与$v$共同成三角形的邻点对组成. 要求链环为圈意味着围绕$v$的三角形形成一圈扇面, 而不是两圈扇面仅共用$v$. 在边内部, 两个半三角形拼成圆盘邻域; 在顶点处, 圈的锥也是圆盘; 面内部本来就是平面开集. 因而$|K|$是紧拓扑曲面. 反例是把两个三角剖分球面仅在一点粘合: 每条边仍邻接两面, 但粘合点的链环为两个不交圈, 不是流形点.

有边界的组合曲面允许边邻接一个或两个三角形, 顶点链环为区间或圈. 区间链环处得到半圆盘邻域, 构成曲面的边界. 本章讨论有限组合模型, 不声称证明了任意拓扑曲面的可三角剖分性.

\begin{definition}
有限二维复形有$V$个顶点、$E$条边、$F$个三角形时, 定义$\chi(K)=V-E+F$, 称为欧拉示性数. 更一般的有限胞腔分解按各维胞腔数取交错和; 对曲面的多边形分解仍为$V-E+F$.
\end{definition}
细分一条边时增加一个顶点和一条边, 若它属于三角形, 每个邻面又同时增加一条边与一个三角形, 所以$\chi$不变. 在三角形内部加一个顶点并连到三个角时, 变化为$\Delta V=1$, $\Delta E=3$, $\Delta F=2$, 也不改变$\chi$. 更一般, 在一个有$m$条边的面内加一中心并扇形细分, 变化$1-m+(m-1)=0$. 因而这里使用的边细分与面细分均保持$\chi$; 存在共同此类细分的两分解给出同一数.

欧拉示性数对所有有限CW空间的同伦不变性是18.905第18讲的更强结论, 其证明调用前面建立的同调理论. 本册下面证明有限链复形的Euler--Poincar\'e恒等式, 不把它单独冒称为完整的奇异同调不变性证明, 也不据细分计算宣布已证明所有曲面的分类.

\originalsection{边界矩阵与Betti数}
给顶点编号, 每条边任选方向, 每个三角形任选循环方向. 以顶点、边、三角形为基, 定义实向量空间$C_0,C_1,C_2$. 对有向边$[i,j]$令$\partial_1[i,j]=[j]-[i]$; 对有向面$[i,j,k]$令
$\partial_2[i,j,k]=[j,k]-[i,k]+[i,j]$,
反向边按负号处理. 矩阵分别记$D_1,D_2$. 边的方向与顺序在两矩阵中必须一致, 不然乘积没有所声称的意义.

\begin{proposition}\label{thm:boundary-square}
$\partial_1\partial_2=0$. 改变边、面的方向或排序不改变下述同调空间的维数.
\end{proposition}
\begin{proof}
对单个面展开, 得$([k]-[j])-([k]-[i])+([j]-[i])=0$, 线性性给出一般结论. 方向改变是对应基向量乘$-1$, 排序是基置换, 均为可逆基变换; 同一个边基变换同时改变$D_1$的列与$D_2$的行, 核、像和商空间维数不变.
\end{proof}

\begin{definition}
定义$H_0=C_0/\operatorname{im}\partial_1$, $H_1=\ker\partial_1/\operatorname{im}\partial_2$, $H_2=\ker\partial_2$. $b_i=\dim_\R H_i$称为实系数Betti数. 商$H_1$良定义, 正因$\operatorname{im}\partial_2\subseteq\ker\partial_1$.
\end{definition}
这里“闭链”指边界为零的链, “边界链”指某个更高维链的边界. 闭环的边向量是闭链, 因为各中间端点消去. 三角形的边界是可以在面内填起的闭链. $H_1$记录未被面边界消去的线性环绕信息, 会失去基本群的非交换次序.

\begin{theorem}[Euler--Poincar\'e恒等式]\label{thm:euler-poincare}
有限二维复形满足
$b_0=V-\operatorname{rank}D_1$,
$b_1=E-\operatorname{rank}D_1-\operatorname{rank}D_2$,
$b_2=F-\operatorname{rank}D_2$,
故$\chi=b_0-b_1+b_2$.
对任意有限链复形, 链空间维数的交错和等于同调维数的交错和.
\end{theorem}
\begin{proof}
秩零度定理给出$\dim\ker\partial_1=E-\operatorname{rank}D_1$, 再减去其子空间$\operatorname{im}\partial_2$的维数得到$b_1$. 另两式直接由商维数与秩零度定理得到, 相减后秩项消去. 一般有限链复形令$r_i=\operatorname{rank}\partial_i$, 首尾缺失边界的秩取零, 则$b_i=\dim C_i-r_i-r_{i+1}$. 乘$(-1)^i$求和时每个$r_i$恰消去, 得一般恒等式.
\end{proof}

\begin{proposition}
$b_0$等于有限复形的连通分支数.
\end{proposition}
\begin{proof}
每条边的边界将它的两个端点在$H_0$中视为相等. 同一图分支内各顶点沿边路相等, 不同分支之间没有边界关系. 给每个分支取一个顶点类, 它们生成$H_0$且独立: 对每个分支求该分支的顶点系数和, 所得线性泛函在每条边的边界上为零, 并分辨这些类. 面不会连接新的顶点分支, 每个图分支的几何实现道路连通且相互不交, 有限性使它们同时开闭, 所以就是空间的连通分支.
\end{proof}

\begin{example}
计算四面体边界的$\chi$与三个Betti数.
\end{example}
\begin{solution}
$V=4$, $E=6$, $F=4$, $\chi=2$. 边界$D_1$的像由$[2]-[1],[3]-[1],[4]-[1]$生成且独立, 秩为$3$. 将各面方向选成外向, 四个面边界的和为零. 若一组面系数的边界为零, 每条边两侧面系数相等, 四个面相互邻接, 因而全部相等. 所以$\ker D_2$一维, $\operatorname{rank}D_2=3$. 得$(b_0,b_1,b_2)=(1,0,1)$.
\end{solution}

\originalsection{可定向性}
\begin{definition}
若可给各三角形选择方向, 使共边上两面的诱导方向相反, 称闭组合曲面可定向. 否则称不可定向. 对带边界曲面, 条件只在邻接两面的边上检查.
\end{definition}
定向不是把所有平面画出的三角形一律画成逆时针: 图中被分开画出的配对边仍须按实际粘合检查. 一个圈走一周后局部方向可能翻转, 正是M\"obius带的现象.

\begin{lemma}\label{lem:dual-connected}
连通闭组合曲面的面邻接图连通, 其中两面共边时连边.
\end{lemma}
\begin{proof}
以共边可达性划分面. 同一顶点周围的面沿链环依次共边, 故都在同一类. 因而不同类既不共边也不共顶点. 每类的面连同其边、顶点形成子复形, 各类互不相交; 有限个闭子复形的并覆盖曲面, 每类补集也闭, 所以每类同时开闭. 连通性只允许一类.
\end{proof}

\begin{theorem}\label{thm:surface-b2}
连通闭组合曲面在实数系数下, 可定向时$b_2=1$, 不可定向时$b_2=0$.
\end{theorem}
\begin{proof}
选任意初始面方向. 闭$2$链$c=\sum_T a_T T$沿每条边的系数为零. 两侧面系数因此满足$a_{T'}=\pm a_T$, 正负号由面方向决定. 面邻接图连通, 所以从一个面的系数决定全部, 核维数至多$1$. 若曲面可定向, 按相容方向选基, 所有$a_T=1$给出非零闭链, 故核维数恰$1$.

反向若有非零闭链, 传播关系使每个$a_T$非零且绝对值相同. 把负系数面的方向反转, 使所有系数相同且正, 此时每条边的零边界条件迫使两侧诱导方向相反, 给出定向. 因而不可定向时核为零.
\end{proof}
这里底域条件有实质作用. 在$\mathbb F_2$中正负号相同, 所有面的和总是闭链, 连通闭曲面的$b_2$均为$1$. 因而$\mathbb RP^2$的实系数$b_2=0$与模$2$系数$b_2=1$不矛盾.

\begin{proposition}\label{thm:orientable-even}
闭可定向组合曲面的$\chi$为偶数. 连通时$b_1$也是偶数.
\end{proposition}
\begin{proof}
用相容方向把每个面的三条边看作有向“面侧”, 集合大小$3F=2E$. 置换$\tau$将每个面侧送到相邻面的同一边, 因方向相反而翻转, 是$E$个换位, 符号为$(-1)^E$. 置换$\sigma$在同一面中将面侧送到下一条, 为$F$个三循环, 符号为$1$.

$\sigma\tau$保持有向边的起点: 先过到对面的反向边, 再取该面下一边, 起点回到原起点. 顶点链环条件使它在每个顶点周围形成一个循环, 恰有$V$个循环. 一个有$N$个元素、$c$个循环的置换符号为$(-1)^{N-c}$, 故$\operatorname{sign}(\sigma\tau)=(-1)^{2E-V}=(-1)^V$. 同时它等于$(-1)^E$, 所以$V-E$为偶数. $3F=2E$还使$F$为偶数, 因而$\chi=V-E+F$为偶数. 连通可定向时$b_0=b_2=1$, Euler--Poincar\'e式给$b_1=2-\chi$, 亦为偶数.
\end{proof}

\originalsection{多边形曲面模型}
将一个多边形的同名边按箭头方向配对, 可得到一个有限胞腔曲面. 沿多边形边界逆时针读边, 同向记字母, 反向记逆字母. 这是一种胞腔分解, 允许同边两端识别, 不直接满足单纯复形“同顶点集只一个单形”的限制. 可以具体细分为单纯三角剖分. 先将每对边同步三等分, 使每个商边上有两个不同内点, 即使原边两端识别也不产生环边或重边. 在尚未粘合的多边形内, 取一圈不作任何边界识别的互异顶点$u_i$, 分别对应逐段边界顶点$b_i$. 将环带四边形$b_i b_{i+1}u_{i+1}u_i$沿$b_i u_{i+1}$分成两面, 再以一个内部中心$o$扇形细分内圈. 每个三角形有不同的三个顶点, 含有不被识别的内圈顶点, 不同面也不具有同一顶点集, 所以得到真正的单纯复形. 原边三等分不改变$V-E$. 若细分后边界有$p$段, 后续插入$p+1$个顶点及$4p$条内边, 用$3p$面替换原一面, 变化$\Delta\chi=(p+1)-4p+(3p-1)=0$. 这给出胞腔计数与三角剖分计数的一致性, 两种记账仍须区分.

\begin{center}
\begin{tikzpicture}[x=1cm,y=1cm,>=stealth]
\draw (0,0) rectangle (2,1.5);
\draw[->] (.55,0)--(1.45,0); \node[below] at (1,0) {$a$};
\draw[->] (.55,1.5)--(1.45,1.5); \node[above] at (1,1.5) {$a$};
\draw[->] (0,.35)--(0,1.15); \node[left] at (0,.75) {$b$};
\draw[->] (2,.35)--(2,1.15); \node[right] at (2,.75) {$b$};
\node at (1,-.8) {环面};
\draw (4,0) rectangle (6,1.5);
\draw[->] (4.55,0)--(5.45,0); \node[below] at (5,0) {$a$};
\draw[->] (4.55,1.5)--(5.45,1.5); \node[above] at (5,1.5) {$a$};
\draw[->] (4,.35)--(4,1.15); \node[left] at (4,.75) {$b$};
\draw[->] (6,1.15)--(6,.35); \node[right] at (6,.75) {$b$};
\node at (5,-.8) {Klein瓶};
\end{tikzpicture}
\end{center}
\noindent 图中只配对同名边, 箭头同向端点对应. 左图边界词为$aba^{-1}b^{-1}$, 右图为$ab^{-1}a^{-1}b^{-1}$; 将右图的$b$换名为$b^{-1}$, 即得到上一章的$aba^{-1}b$约定.

\begin{example}
计算球面、环面、实射影平面、Klein瓶的胞腔数、基本群及实系数Betti数.
\end{example}
\begin{solution}
球面可将圆盘的整条边界缩成一点, 得一个$0$胞腔和一个$2$胞腔, $\chi=2$. 基本群平凡, 三角模型的Betti数为$(1,0,1)$. 环面有一个顶点、两条边、一个面, $\chi=0$, 基本群$\mathbb Z^2$, 可定向, 因而$(b_0,b_1,b_2)=(1,2,1)$.

实射影平面用圆盘边界对径粘合, 有一个顶点、一条边、一个面, $\chi=1$, 基本群$\mathbb Z/2\mathbb Z$. 它不可定向: 沿射影直线一周, 对径识别使横向方向反转; 也可在圆盘的两条配对半边上检查, 同一面会要求相同方向的边互相消去, 无法选定相容方向. 实系数$b_2=0$, 得$(1,0,0)$. Klein瓶有一个顶点、两条边、一个面, $\chi=0$, 群为上一章的非交换群. 反向粘合的一对边同样翻转横向方向, 不可定向, 得$(1,1,0)$.

上述胞腔计数可由细分多边形验证: 在边上和面中添加的顶点、边、面按前述细分公式成对消去. 球面的两半球三角剖分、环面和Klein瓶的方格细分、射影圆盘边界细分均给同样数值. 有限扭元素不能由实系数Betti数看见, 所以$\mathbb RP^2$的$b_1=0$并不意味着基本群平凡.
\end{solution}

\begin{definition}
连通和$M\#N$通过在两闭曲面各删去一个小开圆盘, 再将新边界圆周配对得到. 定向曲面采用使边界方向相反的配对. 记$\Sigma_g$为$g$个环面的连通和, $g=0$时取球面; 记$N_k$为$k$个实射影平面的连通和, $k\ge1$.
\end{definition}
\begin{proposition}\label{thm:connected-sum-euler}
对上述有限组合模型,$\chi(M\#N)=\chi(M)+\chi(N)-2$. 因而$\chi(\Sigma_g)=2-2g$, $\chi(N_k)=2-k$. 相应基本群表示为
\[
 \pi_1(\Sigma_g)=\left\langle a_1,b_1,\ldots,a_g,b_g\ \middle|\
 \prod_{i=1}^g a_ib_ia_i^{-1}b_i^{-1}=1\right\rangle,
 \quad
 \pi_1(N_k)=\langle c_1,\ldots,c_k\mid c_1^2\cdots c_k^2=1\rangle.
\]
第一式$g=0$时群平凡.
\end{proposition}
\begin{proof}
先细分使删去的盘为一个面. 每删去一个开面使$F$减$1$, 边界仍保留. 两边界均细分为$m$条边后粘合, $V$与$E$各减少$m$, 所以粘合不改变两穿孔曲面的欧拉数之和, 得$-2$公式. 环面数值为零, 射影平面数值为一, 归纳得到两个公式.

将各穿孔多边形从孔边界连一条割线到原边界, 沿割线切开, 得原边界词与孔边界词组合的盘. 连通和粘合孔边界时, 两条孔词方向相反而消去, 保留两个原边界词的顺序拼接. 反复进行后, $\Sigma_g$用一个$4g$边形、边界词为各交换子的积; $N_k$用一个$2k$边形、边界词为各平方的积. 所有多边形角点在这些配对下成为一个顶点: 对每个交换子块四角依次识别, 对每个平方块两边首尾识别, 相邻块的角点连接同一类. 一维骨架于是为相应数目的圆周楔和, 用附加胞腔定理得到群表示.
\end{proof}
这些模型给出重要曲面族及其不变量, 不等于已经证明“每个紧连通曲面都在此表中”. 曲面分类还须证明一般多边形词可化为标准形, 并核对唯一性; 任意拓扑曲面还涉及三角剖分定理, 均不在本册证明范围.

\originalsection{角亏与离散Gauss--Bonnet}
给每条边指定正长度, 每个三角形的三边满足严格三角不等式, 就能按该长度实现为欧氏三角形并沿等长边粘合. 这里的内在几何不要求整个曲面嵌入$\R^3$, 相邻面可以在空间中没有指定的夹角.
\begin{definition}
闭组合曲面的顶点$v$处, 定义角亏$\kappa_v=2\pi-\sum_{T\ni v}\theta_{T,v}$, 其中$\theta_{T,v}$为面$T$在$v$处的内角. 正角亏表示总角小于一整周, 负角亏表示总角大于一整周.
\end{definition}
\begin{theorem}[离散Gauss--Bonnet]\label{thm:discrete-gauss-bonnet}
有限闭欧氏三角曲面满足$\sum_v\kappa_v=2\pi\chi$. 有边界时, 边界顶点改用$\kappa_v=\pi-\sum_{T\ni v}\theta_{T,v}$, 仍有$\sum_v\kappa_v=2\pi\chi$.
\end{theorem}
\begin{proof}
闭情形每面三角和为$\pi$, 故$\sum_v\kappa_v=2\pi V-\pi F$. 每条边邻接两面, 数面边关联得到$3F=2E$, 因而$2\pi(V-E+F)=2\pi V-\pi F$.

有边界时记边界顶点数与边数为$V_b,E_b$, 内部顶点数为$V_i$. 边界是有限个组合圆周, 故$V_b=E_b$; 关联计数给$3F=2E-E_b$. 于是
\[
 \sum_v\kappa_v=2\pi V_i+\pi V_b-\pi F
 =2\pi(V_i+V_b-E+F)=2\pi\chi.
\]
无需曲面可定向; 两种计数只依赖局部曲面条件.
\end{proof}

\begin{example}
求正二十面体各顶点的角亏, 并说明平坦环面为何可以处处零角亏.
\end{example}
\begin{solution}
每顶点邻接五个等边三角形, $\kappa=2\pi-5\pi/3=\pi/3$, 十二顶点总和$4\pi=2\pi\cdot2$. 平面等边三角形格按两个独立平移周期取商, 每点邻接六个$\pi/3$角, 角亏为零, 符合环面$\chi=0$. 若需严格单纯复形, 取足够大的周期格避免重复边和相同顶点面的识别; 小周期图可仅作为胞腔记账, 不能冒称单纯剖分.
\end{solution}
角亏总和为零不要求每一点都零: 可以有正负角亏互相抵消. 离散定理不等于光滑曲面的高斯曲率积分定理, 后者需要新的曲率定义与分析论证.

\originalsection{练习与解答}
\begin{exercise}
在有四个顶点的复形中, 取所有六条边, 但只填三个三角面$[1,2,3]$, $[1,2,4]$, $[1,3,4]$. 求Betti数, 并说明是否为闭曲面.
\end{exercise}
\begin{solution}
$\operatorname{rank}D_1=3$. 三面的边界独立, 因为边$[2,3],[2,4],[3,4]$各只属于其中一个面, 零边界迫使三个系数都零, 故$\operatorname{rank}D_2=3$. 得$(b_0,b_1,b_2)=(1,0,0)$, $\chi=1$. 它为四面体边界删去一个开面的圆盘, 三条外边只邻接一个面, 故不是闭曲面.
\end{solution}
\begin{exercise}
设闭三角曲面全部面为等边三角形, 每个顶点恰有六条邻边. 求$\chi$. 能否据此断言它为环面?
\end{exercise}
\begin{solution}
链环为圈使邻面数等于邻边数, 每顶点总角为$2\pi$, 角亏为零, 故$\chi=0$. 仅此不能断言环面, 因为Klein瓶也有这种平坦组合模型, 且不要求曲面连通. 即使另外假设连通可定向, 从$\chi=0$推出环面还要调用本册未证明的完整曲面分类定理.
\end{solution}
\begin{exercise}
对边闭链$c$及满足$D_2^{\mathsf T}u=0$的边向量$u$, 证明$u\cdot c$只依赖$c$在$H_1$中的类. 若$u=D_1^{\mathsf T}v$, 证明它对所有闭链均为零.
\end{exercise}
\begin{solution}
换代表$c'=c+D_2w$时, $u\cdot(c'-c)=(D_2^{\mathsf T}u)\cdot w=0$. 若$u=D_1^{\mathsf T}v$, 则$u\cdot c=v\cdot D_1c=0$. 这些线性绕数只看边的净次数, 不看排列次序, 所以不能代替基本群; 例如两圈楔和的交换子对所有此类实值测量都为零, 但在自由群中非平凡.
\end{solution}

\originalchapter{流形与长度几何}\label{ch:manifolds}
\emph{坐标图、光滑映射与逆函数论证依据Tomasz S. Mrowka的18.965 Fall 2004第1、2、4讲; 路径长度及内在距离参照Paul Seidel的18.950 Fall 2008第4章Lecture 36、39. 有限维局部构造、一般已给定黎曼度量后的距离拓扑证明及配套题解为编者补充. 本章不涵盖这些课程的全部流形、曲率或测地线理论.}

\originalsection{局部欧氏空间}
第8章将流形作为分析的发展方向. 现在可以将局部坐标、全局拓扑和距离的关系具体建立起来. 一个圆周可以用小开弧的角度描述, 但全局没有一张实区间坐标图; 坐标重叠的变换使多张局部描述能够相容.
\begin{definition}
一个$n$维拓扑流形是Hausdorff、第二可数且每点有开邻域同胚于$\R^n$开子集的空间. 坐标图为$(U,\varphi)$, 其中$U$开且$\varphi:U\to\varphi(U)\subseteq\R^n$为同胚, $\varphi(U)$须开. 若图族覆盖空间, 且转移映射$\psi\circ\varphi^{-1}$在重叠坐标域上均光滑, 称为光滑图册. 指定相容光滑图册的流形称为光滑流形.
\end{definition}
维数$n$在本册的定义中固定; 本册不证明维数不变定理. 有边界流形把坐标像改为半空间$\{x_n\ge0\}$的相对开集, 边界点映到$x_n=0$. 以下光滑微分讨论默认无边界, 带边界情形需额外约定延拓的光滑性.

\begin{proposition}\label{thm:manifold-metrizable}
拓扑流形局部紧、局部道路连通且可度量化. 连通流形道路连通.
\end{proposition}
\begin{proof}
在坐标像内以点为中心取一个闭球, 半径足够小使其留在坐标域. 逆像是紧邻域, 因闭球紧而图逆映射连续, 故局部紧. 更小的开球逆像道路连通, 并可包含于任意给定邻域, 故局部道路连通. 局部紧Hausdorff空间正则, 用引理\ref{lem:local-compact-shrink}; 第二可数正则空间可度量化, 用定理\ref{thm:metrization}. 道路分支在局部道路连通空间中开, 若连通只能有一个.
\end{proof}
Hausdorff和第二可数不能只因有局部坐标就省去. 例如两份实直线将全部非零点按相同坐标识别, 却保留两个零点, 形成“双原点直线”. 每点有实区间坐标, 但两个原点的任意邻域均相交, 不是Hausdorff, 因而不符合本书流形定义.

\begin{proposition}
每个光滑图册唯一包含于一个极大相容图册. 构造不需要另用Zorn引理.
\end{proposition}
\begin{proof}
把所有与原图册每张图都相容的图加入. 任取两张新图, 在它们重叠的每一点, 选一张原图覆盖该点; 两新图之间的转移映射在该点附近为通过原图的两个光滑转移映射的复合, 故光滑, 反向也一样. 因而新族仍为图册. 任何与原图册相容的图已经在新族内, 故极大且唯一.
\end{proof}

\originalsection{球面与射影空间的坐标}
\begin{example}
写出$S^n$与$\mathbb RP^n$的光滑图册, 并核对转移映射.
\end{example}
\begin{solution}
在$S^n\subseteq\R^n\times\R$中记点$(u,t)$. 去掉北极时取$\varphi_N(u,t)=u/(1-t)$, 去掉南极时取$\varphi_S(u,t)=u/(1+t)$. 逆映射分别为
\[
 x\longmapsto\left(\frac{2x}{1+|x|^2},\frac{|x|^2-1}{1+|x|^2}\right),\qquad
 x\longmapsto\left(\frac{2x}{1+|x|^2},\frac{1-|x|^2}{1+|x|^2}\right).
\]
重叠处$x\ne0$, 转移为$x\mapsto x/|x|^2$, 光滑且自逆. 球面为欧氏子空间, Hausdorff且第二可数, 两图覆盖, 所以确为光滑流形.

射影空间沿用第10章的球面对径商拓扑, 将一个对径类记为$[x_0:\cdots:x_n]$, 即非零向量所在的直线. 每条直线恰有两个单位向量, 所以这只是同一商空间的另一种记法. 开集$U_i=\{[x]:x_i\ne0\}$上取$\varphi_i([x])=(x_j/x_i)_{j\ne i}$. 比值在对径点上相同, 由商泛性质连续; 逆映射插入$x_i=1$, 再归一化并取对径类, 连续. 图间转移通过非零坐标相除, 因而在其开放定义域光滑. $\mathbb RP^n$已证紧Hausdorff, 有限张第二可数图的基之并为可数全局基, 故满足流形条件. 写比值时必须删去第$i$项, 否则坐标多出恒等于$1$的一项.
\end{solution}

\begin{definition}
光滑流形间映射$f:M\to N$称光滑, 若连续且所有局部坐标表达$\psi\circ f\circ\varphi^{-1}$光滑. 双射$f$及其逆均光滑时称为微分同胚.
\end{definition}
只需对覆盖定义域、目标域的一组图检查. 换图后表达式在两侧复合光滑转移映射, 由链式法则仍光滑. 同胚只约束连续结构, 微分同胚还约束指定光滑结构; 在$\R$上$x\mapsto x^3$为同胚, 但其逆$x^{1/3}$在零点不可微, 不是通常光滑结构下的微分同胚.

\originalsection{逆函数与正则水平集}
本节用多元微积分的导数、链式法则及矩阵运算作为先修. 关键存在性由第6章收缩映射原理证明, 使局部坐标的构造不只停在定理名称.
\begin{theorem}[有限维逆函数定理]\label{thm:inverse-function}
设$F:O\subseteq\R^n\to\R^n$为$C^1$映射, $O$开, $DF(a)$可逆. 则有$a$的开邻域$U$及$F(a)$的开邻域$V$, 使$F:U\to V$为$C^1$微分同胚, 且$D(F^{-1})(y)=DF(F^{-1}(y))^{-1}$. 若$F$光滑, 逆也光滑.
\end{theorem}
\begin{proof}
平移并复合可逆线性映射, 可令$a=0$, $F(0)=0$, $DF(0)=I$. 写$F(x)=x+G(x)$. 由导数连续性选$r>0$, 使闭球$\overline B_r\subseteq O$且其上$\lVert DG\rVert\le1/2$. 沿线段积分导数得$|G(x)-G(x')|\le|x-x'|/2$. 对$|y|<r/2$, 映射$T_y(x)=y-G(x)$将闭球映入自身, 因为$G(0)=0$且$|T_y(x)|\le|y|+r/2<r$. 它为收缩, 故有唯一固定点$\phi(y)$满足$F(\phi(y))=y$.

两个固定点方程相减, 得$|\phi(y)-\phi(y')|\le2|y-y'|$, 所以$\phi$连续. 还由$|\phi(y)|\le2|y|<r$得其值在开球内. 令$V=B_{r/2}$, $U=B_r\cap F^{-1}(V)$, 则$U$开, 每个$y\in V$恰对应一个$x\in U$, 所以$F:U\to V$与$\phi$互逆.

在闭球内$DF=I+DG$均可逆: 级数$\sum_{k\ge0}(-DG)^k$在算子范数下收敛, 与$I+DG$相乘的有限部分和余项趋零, 得逆矩阵且逆范数不超过$2$. 对$x=\phi(y)$, $h=\phi(y+k)-\phi(y)$, 有$k=DF(x)h+o(|h|)$, 而$|h|\le2|k|$. 乘逆矩阵得到$h=DF(x)^{-1}k+o(|k|)$, 故$D\phi(y)=DF(\phi(y))^{-1}$. 右侧连续, 因而逆为$C^1$. 若$F$光滑, 矩阵求逆在可逆矩阵集合中光滑, 可用伴随矩阵除以非零行列式核对; 由此公式和链式法则逐阶归纳, 得$\phi$光滑.
\end{proof}

\begin{theorem}[正则水平集]\label{thm:regular-level}
设$f:O\subseteq\R^n\to\R^m$光滑, $O$开, $m\le n$, 且对每个$x\in f^{-1}(c)$, $Df(x)$满射. 则非空水平集$f^{-1}(c)$按子空间拓扑为光滑$(n-m)$维流形.
\end{theorem}
\begin{proof}
在某点$a$, 满秩矩阵有一个$m\times m$可逆子块. 重排坐标为$(u,v)\in\R^{n-m}\times\R^m$, 使$D_vf(a)$可逆. 映射$\Phi(u,v)=(u,f(u,v))$的导数为块三角矩阵, 对角块为$I$与$D_vf(a)$, 故可逆. 由逆函数定理, $\Phi$在$a$附近为微分同胚. 水平集在此坐标下恰为第二组坐标等于$c$的切片, 第一组坐标$u$形成开集, 给局部图. 两图间转移是环境光滑变换在该切片上的限制, 因而光滑. 水平集作为欧氏子空间Hausdorff且第二可数, 具备全部流形条件.
\end{proof}
例如$f(x)=|x|^2$, $c=1$时$Df(x)h=2x\cdot h$满射, 得球面为光滑流形. 零水平集未满足正则条件; 它仍可能碰巧是其他维数的流形, 所以满秩是充分条件, 不能称为所有水平集成为流形的必要条件. 曲线$xy=0$在原点呈十字, 则确实不局部同胚于实区间: 删去原点的小邻域有四个分支, 区间删去一点仅两个.

\originalsection{切空间与微分}
\begin{definition}
在光滑流形点$p$, 将过$p$的光滑曲线按在一张局部图中有相同的零时刻导数分类, 得切向量. 这些类组成切空间$T_pM$. 光滑映射$f:M\to N$的微分$df_p:T_pM\to T_{f(p)}N$将曲线类$[\gamma]$送到$[f\circ\gamma]$.
\end{definition}
这里曲线定义在含$0$的小开区间, $\gamma(0)=p$. 换坐标时速度由可逆矩阵$D(\psi\circ\varphi^{-1})(\varphi(p))$变换, 所以曲线是否同类与图无关. 每个坐标速度$v\in\R^n$由曲线$\varphi^{-1}(\varphi(p)+tv)$实现, 因而一张图给出$T_pM\cong\R^n$. 加法和数乘在图中定义, 可逆线性坐标变化使之独立于图. 同理链式法则使$df_p$良定义且线性, 并给$d(gf)_p=dg_{f(p)}\circ df_p$.

\begin{proposition}\label{thm:tangent-level}
正则水平集$M=f^{-1}(c)\subseteq\R^n$的切空间为$T_pM=\ker Df(p)$, 其中环境速度用于识别切向量.
\end{proposition}
\begin{proof}
对留在$M$中的曲线,$f(\gamma(t))=c$, 求导得$Df(p)\gamma'(0)=0$. 反向用正则水平集证明中的坐标$\Phi=(u,f)$, 微分将$\ker Df(p)$映成第一坐标空间$\R^{n-m}\times\{0\}$. 在水平切片中取给定速度的直线, 经$\Phi^{-1}$得到$M$中的曲线, 环境速度正是原向量. 因而两个包含方向均成立.
\end{proof}
球面上$T_pS^n=\{v:p\cdot v=0\}$. 切空间是附在一点的线性空间, 不等于曲面附近的一块; 将它画为过$p$的仿射平面时, 实际画的是$p+T_pM$.

\originalsection{长度与内在距离}
\begin{definition}
黎曼度量是在每个$T_pM$上指定正定内积$g_p$, 在局部图中其矩阵系数$g_{ij}(x)$随$x$光滑. 分段$C^1$曲线$\gamma:[a,b]\to M$的长度为
$L_g(\gamma)=\int_a^b\sqrt{g_{\gamma(t)}(\gamma'(t),\gamma'(t))}\,\mathrm dt$,
其中对$C^1$曲线, $\gamma^{\prime}(t)$指坐标导数经$T_{\gamma(t)}M\cong\R^n$对应得到的切向量. 在分段点分别积分相加, 有限个端点的速度取值不影响积分.
\end{definition}
坐标变化使切向量和度量矩阵按相反方向变换, 内积值不变, 所以长度与图无关. 保端点的单调满射$C^1$重参数在可逆区间上由积分代换保持长度; 允许分段重参数或等待段时分段核对即可. 对嵌入欧氏空间的正则水平集, 环境内积限制在切空间给诱导度量. 参数化$F(u)$的坐标矩阵为$g_{ij}=\partial_iF\cdot\partial_jF$, 正定性来自参数导数的单射性.

\begin{theorem}\label{thm:riemannian-distance}
连通光滑流形给定黎曼度量后, $d_g(p,q)=\inf_\gamma L_g(\gamma)$是有限度量, 并生成原流形拓扑, 下确界取遍连接两点的分段$C^1$曲线. 本定理不声称下确界必被某条曲线取得.
\end{theorem}
\begin{proof}
坐标小凸球中的点可以用坐标直线连接, 对应光滑曲线. 按可用分段光滑曲线连通划分, 每类开, 连通性使只有一类, 所以下确界有限. 反向道路给对称性, 拼接并按任意小误差选近乎最短曲线给三角不等式, 常值曲线给$d_g(p,p)=0$.

须另外证明不同点距离正以及拓扑一致. 在$p$的一张坐标图中取中心为$p$的闭坐标球$\overline B_r$, 紧含于图域. 度量矩阵连续且正定, 在$\overline B_r\times S^{n-1}$上取最小、最大值, 得$0<c\le C<\infty$满足
$c|v|\le\sqrt{g_x(v,v)}\le C|v|$.
对留在此球内的曲线, 长度至少为$c$乘坐标端点距离. 对离开球的曲线, 到第一次碰到坐标球边界为止的长度至少$cr$; 首次碰边存在, 因曲线连续且闭球紧含于图域, 从中心离开开球必须先经过其边界. 因而对$q$在球内, 任意曲线长度至少$\min\{c|\varphi(q)-\varphi(p)|,cr\}$, 若$q$在球外则至少$cr$. 不同$q$的下界正. 球内直线又给$d_g(p,q)\le C|\varphi(q)-\varphi(p)|$.

上界说明坐标点趋于$p$时$d_g$趋零; 下界说明足够小$d_g$球包含于任意给定的小坐标球. 两种邻域系相互包含, 故拓扑一致. $n=0$时连通非空流形只有一点, 结论直接成立.
\end{proof}
这是内在距离: 路径须留在流形内. 在单位圆周上, 两点的内在距离为较短圆弧长度, 环境距离为弦长, 两者生成相同拓扑却数值不同. 任意图册给出的可度量化也不唯一指定黎曼度量; 本节定理从已经给定的度量出发, 不代替一般流形上黎曼度量存在性的单位分解证明.

\begin{example}
计算圆柱$M=\{(\cos\theta,\sin\theta,z)\}$的内在距离, 并解释局部坐标为何不能直接消除全局绕圈.
\end{example}
\begin{solution}
参数导数分别为$(-\sin\theta,\cos\theta,0)$与$(0,0,1)$, 故局部度量为$\mathrm d\theta^2+\mathrm dz^2$. 对两点取角代表$\theta_1,\theta_2$. 记覆盖$P(\theta,z)=(\cos\theta,\sin\theta,z)$, 任意连接曲线取提升初值$(\theta_1,z_1)$, 终点为$(\theta_2+2\pi k,z_2)$, $k\in\mathbb Z$. 在小开圆弧与高度区间上, $P$的逆支是光滑角度坐标. 紧参数区间有有限分割使每小段留在一张这样的图内, 并可细化原来的$C^1$分段; 提升在各段等于光滑逆支与原曲线复合, 故仍分段$C^1$. 若提升写为$(\theta(t),z(t))$, 原曲线速度为$(-\sin\theta\,\theta^{\prime},\cos\theta\,\theta^{\prime},z^{\prime})$, 范数平方为$(\theta^{\prime})^2+(z^{\prime})^2$, 因而两曲线长度相等. 欧氏平面曲线长度至少为端点距离, 等号可由直线实现, 投影回圆柱得
\[
 d_M(p,q)=\min_{k\in\mathbb Z}
 \sqrt{(\theta_2-\theta_1+2\pi k)^2+(z_2-z_1)^2}.
\]
整数最小值存在, 因表达式随$|k|\to\infty$趋无穷. 单一角坐标只能覆盖开圆弧, 整圈的整数差须由覆盖提升记录. 这是拓扑与长度计算共同作用的一个具体例子.
\end{solution}

\originalsection{练习与解答}
\begin{exercise}
证明$S^1\times S^1$为光滑二维流形, 写出标准乘积度量并计算两个角坐标点的距离.
\end{exercise}
\begin{solution}
两开圆弧图的积覆盖环面, 转移在各重叠连通分支上为角度加常数$2\pi k$, 因而光滑; Hausdorff和第二可数由有限积结论. 乘积度量为$\mathrm d\theta^2+\mathrm d\phi^2$. 由圆柱例题的逆支分段论证, 覆盖$P(\theta,\phi)=(e^{i\theta},e^{i\phi})$将连接曲线提升为$\R^2$中的分段$C^1$曲线, 且逐段速度平方为$(\theta^{\prime})^2+(\phi^{\prime})^2$, 长度保持. 固定提升初值后, 两终点差为$(\Delta\theta+2\pi m,\Delta\phi+2\pi n)$, 所以距离是这些欧氏范数的最小值, $m,n\in\mathbb Z$. 直线投影实现最小值. 这是平坦环面的内在几何, 不等于常见空间旋转环面从$\R^3$诱导的度量.
\end{solution}
\begin{exercise}
对$f(x,y,z)=x^2+y^2-z^2$, 证明每个非零水平集是光滑曲面; 分析零水平集在原点附近的拓扑.
\end{exercise}
\begin{solution}
$Df=(2x,2y,-2z)$仅在原点为零, 非零水平集不含原点, 故正则且二维. 零水平集为双圆锥. 删掉原点后分成$z>0$与$z<0$两部分, 两部分在任意原点邻域内都非空且不能在删点圆锥中连接. 若原点有二维圆盘坐标邻域, 缩小为坐标圆盘后删中心仍道路连通, 与圆锥的这两分支矛盾. 因而原点不是二维流形点; 其余点仍正则.
\end{solution}
\begin{exercise}
在开单位圆盘内采用欧氏诱导度量. 证明它的内在距离就是欧氏距离, 却不是完备空间. 说明这与局部欧氏性是否冲突.
\end{exercise}
\begin{solution}
圆盘凸, 线段连接给长度等于弦长, 任何曲线长度不小于弦长, 故两距离相同. 点列$(1-1/n,0)$($n\ge2$)为Cauchy列, 在平面极限为$(1,0)$但不属于开圆盘, 因而在圆盘内不收敛. 完备性是选定距离的全局性质, 局部坐标只描述邻域结构, 所以没有冲突. 本例也不包含关于完备黎曼流形上最短曲线存在的Hopf--Rinow定理.
\end{solution}

\originalchapter{度量空间原课程习题}\label{ch:official-metric}
\emph{题面来自 Paige Bright, MIT OCW 18.S190 Introduction to Metric Spaces, IAP 2023 的三份 Problem Set. 每份原件为2页题目及1页OCW署名页. 以下保留24道原编号大题, 包括全部选做题和33个显式末级子问; 解答均为中文编者补充, 不是MIT发布的标准答案.}

前15章的58道编者练习不计作原课程作业. 有些原题已在正文中完整证明, 此处给出中文题面、准确的证明链接及所需应用, 避免再抄一遍同一论证. 原件中的错误与记号歧义另作说明, 不把修正版冒充原题.

\originalsection{第一份题单}
\begin{exercise}\label{ps190:1:1}
\textbf{PS1第1题.} 在$\R^2$上定义英国铁路距离: 若$x,y,0$共线, 令$d(x,y)=\norm{x-y}_2$; 否则令$d(x,y)=\norm{x}_2+\norm{y}_2$. 证明这是度量, 并解释名称. 原提示是画图并使用欧氏距离.
\end{exercise}
\begin{solution}
距离有限非负、对称, 且$d(x,y)\ge\norm{x-y}_2$, 因而正定. 对三角不等式, 若$x,z,0$共线, 则$d(x,z)=\norm{x-z}_2\le\norm{x-y}_2+\norm{y-z}_2\le d(x,y)+d(y,z)$.

若$x,z,0$不共线, 则$x,z$都非零. 若$y=0$, 三角不等式取等号. 若$y\ne0$, 它不可能同时与$x,0$及$z,0$共线. 两侧都不共线时, $d(x,y)+d(y,z)=\norm{x}_2+2\norm{y}_2+\norm{z}_2\ge d(x,z)$. 只有$x,y,0$共线时, 由$\norm{x}_2\le\norm{x-y}_2+\norm{y}_2$, 得同一不等式; 只有$y,z,0$共线时交换端点. 同一线路上的两点直接行走, 不同线路之间须经过原点这个枢纽, 由此得名.
\end{solution}

\begin{exercise}\label{ps190:1:2}
\textbf{PS1第2题.} $d(f,g)=\sup_{x\in[0,1]}|f'(x)-g'(x)|$是否为$C^1([0,1])$上的度量? 若不是, 说明满足哪些公理、违反哪条公理.
\end{exercise}
\begin{solution}
非负、有限、对称, $d(f,f)=0$, 且由逐点三角不等式再取上确界得$d(f,h)\le d(f,g)+d(g,h)$. 但$f=0,g=1$不同而距离为零, 故正定性失败. 更精确地, $d(f,g)=0$当且仅当$f-g$为常数, 这是中值定理的直接推论. 习题\ref{ch:metric:ex-c1}已证明加上$f(0)=0$的条件后, 导数上确界才给出范数.
\end{solution}

\begin{exercise}\label{ps190:1:3}
\textbf{PS1第3题.} 证明$d(x,y)=|x-y|/(1+|x-y|)$是$\R$上的度量.
\end{exercise}
\begin{solution}
这正是习题\ref{ch:metric:ex-bounded}对通常实数距离的应用. 该处已逐项核对公理, 关键为递增函数$h(t)=t/(1+t)$满足$h(a+b)\le h(a)+h(b)$; 因而$h(|x-z|)\le h(|x-y|)+h(|y-z|)$.
\end{solution}

\begin{exercise}\label{ps190:1:4}
\textbf{PS1第4题.} 原课称半度量为满足非负性、对称性、三角不等式以及$d'(x,x)=0$的函数, 但允许不同点间距离为零. 证明一个度量$d$与一个半度量$d'$之和仍是度量.
\end{exercise}
\begin{solution}
和有限非负且对称, 对角线上为零. 若$x\ne y$, 则$d(x,y)+d'(x,y)\ge d(x,y)>0$. 两个三角不等式相加, 即得$(d+d')(x,z)\le(d+d')(x,y)+(d+d')(y,z)$. 本题的半度量也常称伪度量; 不采用其他文献中不同的同名定义.
\end{solution}

\begin{exercise}\label{ps190:1:5}
\textbf{PS1第5题.} 原题写作: 证明$I_t:C^0([a,b])\to C^1([a,b])$连续, 其中$I_t(f)=\int_a^t f(x)\dd x$, $t\in[a,b]$. 原提示说证明与课堂例子相似, 需要调整$\varepsilon$和$\delta$.
\end{exercise}
\begin{solution}
固定$t$时右端是一个实数, 与所写值域不符. 按积分算子的含义, 应理解为$(If)(t)=\int_a^t f(x)\dd x$, 定义域取上确界度量, 值域取$C^1$度量. 习题\ref{ch:metric:ex-integration}已完整证明$If\in C^1$及$\norm{If-Ig}_{C^1}\le(1+b-a)\norm{f-g}_\infty$. 对给定$\varepsilon>0$, 取$\delta=\varepsilon/(1+b-a)$即得连续性. 若确实固定$t$并把值域改为$\R$, 则$|I_t(f)-I_t(g)|\le(b-a)\norm{f-g}_\infty$, 也连续. 两种解释的对象与值域须分别说明.
\end{solution}

\begin{exercise}\label{ps190:1:6}
\textbf{PS1第6题(选做).} 设$a_k,b_k\in\R$, $1\le k\le n$.
\begin{enumerate}[label=(\alph*)]
\item 当$1<p<\infty$, $1/p+1/q=1$时, 证明H\"older不等式$\sum_k|a_kb_k|\le(\sum_k|a_k|^p)^{1/p}(\sum_k|b_k|^q)^{1/q}$. 原提示是先证$A^tB^{1-t}\le tA+(1-t)B$.
\item 当$1\le p<\infty$时, 原件实际把左端印为括号内含比较符号的下式:
\[
 \left(\sum_k|a_k+b_k|^p\le\sum_k|a_k+b_k|^p\right)^{1/p}
 \le\left(\sum_k|a_k|^p\right)^{1/p}
   +\left(\sum_k|b_k|^p\right)^{1/p}.
\]
原提示为将$\sum_k|a_k+b_k|^p$控制为$\sum_k|a_k||a_k+b_k|^{p-1}+\sum_k|b_k||a_k+b_k|^{p-1}$, 再用(a).
\end{enumerate}
\end{exercise}
\begin{solution}
(a) 固定$A,B>0$及$0<t<1$, 函数$F(x)=tx+(1-t)B-x^tB^{1-t}$的导数为$t[1-(B/x)^{1-t}]$. 在$x<B$时为负, $x>B$时为正, 故$F(x)\ge F(B)=0$. 取$x=A$即得提示中的不等式; 其中一个数为零时由非负性直接成立.

令$A_0=(\sum|a_k|^p)^{1/p}$, $B_0=(\sum|b_k|^q)^{1/q}$. 若一个为零, 相应全体坐标为零, 所需式成立. 否则对$u_k=|a_k|/A_0$, $v_k=|b_k|/B_0$, 用刚证的式子取$t=1/p$, $A=u_k^p$, $B=v_k^q$, 得$u_kv_k\le u_k^p/p+v_k^q/q$. 求和得$\sum u_kv_k\le1/p+1/q=1$, 乘回$A_0B_0$即得H\"older不等式.

(b) 原式括号内多了一段重复求和与比较符号, 左端并非正常的数值表达式, 属于排印错误. 若按抽取文字把它读作$\sum|a_k+b_k|^p\le(\sum|a_k+b_k|^p)^{1/p}$, 也一般不成立: $n=1,p=2,a_1=2,b_1=0$会给$4\le2$. 正确的Minkowski式应只保留左端的$p$次根, 再与右端两范数之和比较. $p=1$时直接对$|a_k+b_k|\le|a_k|+|b_k|$求和. $p>1$时令$S=(\sum|a_k+b_k|^p)^{1/p}$. 若$S=0$, 结论成立. 否则由原提示和(a), 注意$(p-1)q=p$, 有
\[
 S^p\le
 \left[\left(\sum|a_k|^p\right)^{1/p}
 +\left(\sum|b_k|^p\right)^{1/p}\right]S^{p-1}.
\]
除以$S^{p-1}$即得正确结论. 正定性和齐次性直接由求和式得到, 所以$\ell^p$范数及其诱导距离确实成立; 这补上了第1章明确暂未证明的一般$p$情形.
\end{solution}

\begin{exercise}\label{ps190:1:7}
\textbf{PS1第7题(选做).} 对$f,g\in C^\infty([a,b])$, 定义$d_n(f,g)=\sup_{x\in[a,b]}|f^{(n)}(x)-g^{(n)}(x)|$, 其中$f^{(0)}=f$. $n\ge1$时$d_n$是半度量, $d_0$是度量. 证明$d(f,g)=\sum_{n=0}^\infty 2^{-n}d_n(f,g)/(1+d_n(f,g))$是度量. 原课指出这与Fr\'echet空间有关.
\end{exercise}
\begin{solution}
每个导数在闭区间上连续有界, 所以各$d_n$有限. 每项非负且小于$2^{-n}$, 故级数收敛并不超过$2$. 对称性、对角线上为零显然. 若$f\ne g$, 则$d_0(f,g)>0$, 第零项正, 从而$d(f,g)>0$. 习题\ref{ch:metric:ex-bounded}的$h(t)=t/(1+t)$证明只用非负性、单调性和三角不等式, 因而对第三个函数$u$, 有$h(d_n(f,u))\le h(d_n(f,g))+h(d_n(g,u))$, 对半度量也成立. 乘$2^{-n}$, 对有限部分和求和再令截断指标趋于无穷, 得$d$的三角不等式. 此处不进一步声称已证明整个Fr\'echet空间理论.
\end{solution}

\originalsection{第二份题单}
本份题单统一设$(X,d)$为度量空间. 原预备说明的$C^0([a,b])$距离为$\max|f-g|$, 范数具有正定性、齐次性与三角不等式, $K\Subset X$表示$K$为紧子集. 若非空实数集$A$有有限上确界$a$, 则对每个$n\ge1$可取$x_n\in A$使$a-1/n<x_n\le a$.

\begin{exercise}\label{ps190:2:1}
\textbf{PS2第1题.} (1) 若$x_n\to x$, $y_n\to y$, 证明$d(x_n,y_n)\to d(x,y)$. (2) 若$(x_n)$与$(y_n)$都是$X$中的Cauchy序列, 证明实数列$d(x_n,y_n)$收敛, 不得假设两点列已在$X$中收敛. 原提示误写为$|d(a,b)+d(a',b')|\le d(a,a')+d(b,b')$.
\end{exercise}
\begin{solution}
(1) 定理\ref{ch:seq:basic-limits}已证明, 所需估计为$|d(x_n,y_n)-d(x,y)|\le d(x_n,x)+d(y_n,y)\to0$. (2) 命题\ref{ch:seq:cauchy-distance}的完整证明使用$|d(x_n,y_n)-d(x_m,y_m)|\le d(x_n,x_m)+d(y_n,y_m)$: 右端随$m,n$充分大而任意小, 实数完备性给出收敛. 这不要求$X$完备. 原提示的加号应改为减号; 例如$a=a'\ne b=b'$时原左端正而右端为零, 原式不成立.
\end{solution}

\begin{exercise}\label{ps190:2:2}
\textbf{PS2第2题.} 证明$A\subseteq X$闭当且仅当$A$中每个在$X$中收敛的序列, 其极限仍属于$A$.
\end{exercise}
\begin{solution}
这是定理\ref{ch:seq:closure-sequences}的闭集判据, 该处已证明双向. 其中逆向的关键是对$x\in\overline A$逐次选$x_n\in A\cap B(x,1/n)$; 点列趋于$x$, 题设给$x\in A$. “在$X$中收敛”须保留, 本题并未断言$A$中的所有序列都收敛.
\end{solution}

\begin{exercise}\label{ps190:2:3}
\textbf{PS2第3题.} 设$(f_n)$是$C^0([0,1])$上确界距离下的Cauchy序列.
\begin{enumerate}[label=(\alph*)]
\item 固定$x_0\in[0,1]$, 证明$\lim_n f_n(x_0)$存在.
\item 定义$f(x)=\lim_n f_n(x)$. 证明对每个$\varepsilon>0$, 存在同一个$N$, 使所有$x\in[0,1]$及$n\ge N$满足$|f_n(x)-f(x)|\le\varepsilon$.
\item 证明$f$在$[0,1]$连续. 原提示插入$f_n(x)$与$f_n(x_0)$, 将差拆为三项.
\item 用(a)--(c)说明$f_n$作为函数空间中的序列收敛于$f$.
\end{enumerate}
\end{exercise}
\begin{solution}
这四步已包含在定理\ref{thm:cb-complete}及推论\ref{cor:c-interval-complete}中, 逐问对应如下. (a) $|f_n(x_0)-f_m(x_0)|\le\norm{f_n-f_m}_\infty$, 故实数列Cauchy, 由$\R$完备得极限. (b) 先取$N$使$m,n\ge N$时$\norm{f_n-f_m}_\infty<\varepsilon$. 固定$n\ge N,x$, 令$m\to\infty$, 得$|f_n(x)-f(x)|\le\varepsilon$; 同一$N$适合全部$x$.

(c) 给定$x_0$与$\varepsilon>0$, 用(b)取某个$n$使$\norm{f-f_n}_\infty\le\varepsilon/4$. 由$f_n$在$x_0$连续, 当$|x-x_0|$足够小时$|f_n(x)-f_n(x_0)|<\varepsilon/4$. 三角不等式给$|f(x)-f(x_0)|<3\varepsilon/4<\varepsilon$. 端点采用相对邻域. (d) 已有$f\in C^0([0,1])$, 对(b)取上确界并令$\varepsilon$任意, 得$\norm{f_n-f}_\infty\to0$. 极限在原空间中, 因而该空间完备.
\end{solution}

\begin{exercise}\label{ps190:2:4}
\textbf{PS2第4题.} 设$V$为实赋范空间, $d(x,y)=\norm{x-y}$.
\begin{enumerate}[label=(\alph*)]
\item 证明$d(\lambda x,\lambda y)=|\lambda|d(x,y)$.
\item 证明平移不变性$d(x+z,y+z)=d(x,y)$.
\item 证明$d$为$V$上的度量, 称为范数诱导的度量.
\end{enumerate}
\end{exercise}
\begin{solution}
(a) $\norm{\lambda x-\lambda y}=\norm{\lambda(x-y)}=|\lambda|\norm{x-y}$. (b) $(x+z)-(y+z)=x-y$, 取范数即得. (c) 完整公理证明见命题\ref{ch:metric:norm-metric}; 正定性来自$x-y=0$当且仅当$x=y$, 对称性用$\norm{-v}=\norm v$, 三角不等式用$x-z=(x-y)+(y-z)$.
\end{solution}

\begin{exercise}\label{ps190:2:5}
\textbf{PS2第5题.} 证明度量空间中的开集$U$可表示为任意多个开球的并.
\end{exercise}
\begin{solution}
命题\ref{ch:seq:union-balls}已完整证明此结论: 对每个$x\in U$选$r_x>0$使$B(x,r_x)\subseteq U$, 则$U=\bigcup_{x\in U}B(x,r_x)$. 原文的“任意多个”表示允许无限指标族, 不表示任意预先指定的球族都能覆盖$U$. 空集对应空并.
\end{solution}

\begin{exercise}\label{ps190:2:6}
\textbf{PS2第6题(选做).} (a) 原题: 若$K\Subset\R$, 证明$K$有最大值与最小值. (b) 将$\R$上的Heine--Borel定理推广到$\R^n$. 原提示说此证明与课堂证明相似.
\end{exercise}
\begin{solution}
(a) 原题须增加$K\ne\varnothing$: 空集紧却没有最大值或最小值. 对非空紧集, 定理\ref{ch:compact:extreme-values}应用于恒等实函数即得两极值. 也可由紧集闭且有界, 取$s=\sup K$, 用$x_n\in K$且$s-1/n<x_n\le s$, 得$x_n\to s$; 闭性给$s\in K$. 下确界同理. (b) 定理\ref{ch:eucompact:heine-borel}已完整证明$K\subseteq\R^n$紧当且仅当在$\R^n$中闭且有界. 必要性见定理\ref{ch:eucompact:compact-closed-bounded}; 充分性把有界集放进紧闭方盒, 再用命题\ref{ch:eucompact:closed-subset}. 不把这一结论推广到任意度量空间.
\end{solution}

\begin{exercise}\label{ps190:2:7}
\textbf{PS2第7题(选做).} 对度量空间之间的映射$f:X\to Y$, 证明$f$连续当且仅当每个$Y$中开集$U$的逆像$f^{-1}(U)$在$X$中开.
\end{exercise}
\begin{solution}
这是定理\ref{ch:seq:preimage}, 该处已给双向完整证明. 正向在$x\in f^{-1}(U)$处取像点周围的小球, 再用连续性取$x$周围映入它的小球; 逆向对$B_Y(f(x),\varepsilon)$的开逆像取包含$x$的小球, 即得$\varepsilon$--$\delta$条件.
\end{solution}

\begin{exercise}\label{ps190:2:8}
\textbf{PS2第8题(选做).} 两范数等价指存在常数$C_1>0,C_2$使$C_1\norm{x}_1\le\norm{x}_2\le C_2\norm{x}_1$. 在$\R^n$证明
\[
 \norm{x}_\infty\le\norm{x}_2\le\norm{x}_1
 \le\sqrt n\norm{x}_2\le n\norm{x}_\infty,
\]
并说明为何三个范数两两等价. 此处下标$1,2$也分别表示$\ell^1,\ell^2$范数.
\end{exercise}
\begin{solution}
命题\ref{ch:metric:three-norms}已证明前三段以及$\norm{x}_1\le\sqrt n\norm{x}_2$. 最后一段补用$\norm{x}_2^2=\sum_i|x_i|^2\le n\norm{x}_\infty^2$, 开平方并乘$\sqrt n$即可. 特别地, $\norm{x}_\infty\le\norm{x}_2\le\sqrt n\norm{x}_\infty$, $\norm{x}_\infty\le\norm{x}_1\le n\norm{x}_\infty$, $\norm{x}_2\le\norm{x}_1\le\sqrt n\norm{x}_2$. 每一对均有正的有限双边常数, 因而两两等价. 对非零向量, 原定义中的双边式也自动迫使$C_2>0$.
\end{solution}

\begin{exercise}\label{ps190:2:9}
\textbf{PS2第9题(选做).} 对任意映射$f:X\to Y$, 若$A\subseteq B\subseteq X$, 证明$f(A)\subseteq f(B)$. 原题末尾另链一个站外讨论, 不是本题的附加题单.
\end{exercise}
\begin{solution}
若$y\in f(A)$, 则存在$a\in A$使$y=f(a)$. 因$a\in B$, 有$y\in f(B)$. 这不需要连续性、单射性或满射性.
\end{solution}

\originalsection{第三份题单}
原预备说明定义$A+B=\{a+b:a\in A,b\in B\}$, 以及$\ell^p=\{a=(a_k)_{k\ge1}:\sum_k|a_k|^p<\infty\}$, $1\le p<\infty$, 配范数$\norm{a}_p=(\sum_k|a_k|^p)^{1/p}$. 将习题\ref{ps190:1:6}(b)应用于有限截断再取极限, 得$\norm{a+b}_p\le\norm a_p+\norm b_p$; 这也保证加法留在$\ell^p$中. 正定性、齐次性由级数直接得到, 所以确为赋范空间.

\begin{exercise}\label{ps190:3:1}
\textbf{PS3第1题.} 若$A\subseteq\R^n$闭, $B\subseteq\R^n$紧, 证明$A+B$闭. 原提示用序列判据; 原课另指出可推广到拓扑向量空间, 此处不要求该推广.
\end{exercise}
\begin{solution}
若有一个集合为空, 和为空, 已闭. 否则设$z_m=a_m+b_m\to z$, $a_m\in A,b_m\in B$. 紧性给子列$b_{m_j}\to b\in B$, 于是$a_{m_j}=z_{m_j}-b_{m_j}\to z-b$. $A$闭给$z-b\in A$, 所以$z=(z-b)+b\in A+B$. 闭集序列判据给结论. 抽子列必须与$z_m$使用同一组下标.
\end{solution}

\begin{exercise}\label{ps190:3:2}
\textbf{PS3第2题.} 设$S\subseteq X$, $x\in S$. 称$x$是孤立点, 如果某个$\varepsilon>0$满足$B(x,\varepsilon)\cap S=\{x\}$. 证明$x$孤立当且仅当每个在$S$中趋于$x$的序列$(a_n)$最终恒等于$x$.
\end{exercise}
\begin{solution}
若$x$孤立, 收敛使$a_n$最终落入上述球中, 因而只能等于$x$. 若$x$不孤立, 对每个$n$取$a_n\in B(x,1/n)\cap(S\setminus\{x\})$. 则$a_n\to x$, 却没有一项等于$x$, 与最终恒等条件矛盾. 这是正文聚点序列判据的直接补充, 不是断言一般拓扑空间也有同一序列刻画.
\end{solution}

\begin{exercise}\label{ps190:3:3}
\textbf{PS3第3题.} 证明度量空间$X$中有限多个紧子集的并紧.
\end{exercise}
\begin{solution}
这与习题\ref{ch:eucompact:ex-finite-union}的第一问完全相同, 该处有限子覆盖的证明直接适用: 分别给每个紧子集取有限子覆盖, 再合并有限多个有限族. 无需再增加闭性或欧氏空间条件.
\end{solution}

\begin{exercise}\label{ps190:3:4}
\textbf{PS3第4题.} 对$1\le p<\infty$, 令$S=\{a\in\ell^p:\norm a_p\le1\}$. 说明$S$闭且有界, 并证明它不紧. 原提示取标准单位向量$e_n=(\delta_{kn})_{k\ge1}$, 说明它没有收敛子列, 再解释与紧性的关系.
\end{exercise}
\begin{solution}
反三角不等式$|\norm a_p-\norm b_p|\le\norm{a-b}_p$说明范数连续, 故$S$为闭区间$[0,1]$的逆像, 是闭集. 每个点到$0$的距离不超过$1$, 所以有界. 又$\norm{e_n}_p=1$, 且$m\ne n$时$\norm{e_n-e_m}_p=2^{1/p}$. 因而任何子列都不Cauchy, 不能收敛. 定理\ref{ch:compact:compact-to-sequential}说明紧度量空间中的每个序列有收敛子列, 故$S$不紧. 第4章的$\ell^2$反例是$p=2$的特例, 不足以单独代替本题全部$p$的证明.
\end{solution}

\begin{exercise}\label{ps190:3:5}
\textbf{PS3第5题.} 在上确界度量下, 证明$S=\{f\in C^0([0,1]):\norm f_\infty\le1\}$不紧. 原提示取$f_n(x)=x^n$.
\end{exercise}
\begin{solution}
该反例及完整证明已用于\hyperlink{exercise.8.1}{习题8.1的第二问}; 那里额外要求$f(0)=0$, 这里的同一序列仍全部属于$S$. 若$S$紧, 则$(x^n)$有在上确界距离下收敛的子列$f_{n_j}\to f\in C^0([0,1])$. 逐点极限只能是$0$在$[0,1)$、$1$在$x=1$的函数, 它不连续, 矛盾. 因而$S$不紧. 两题的集合不同, 此处复用的是同一反例与论证.
\end{solution}

\begin{exercise}\label{ps190:3:6}
\textbf{PS3第6题.} 对$0<k<1,b\in\R$, 设$f(x)=kx+b$. 证明$f:\R\to\R$是压缩映射, 求不动点并直接证明唯一性.
\end{exercise}
\begin{solution}
$|f(x)-f(y)|=k|x-y|$, 故压缩常数为$k<1$. 方程$x=kx+b$等价于$(1-k)x=b$, 所以唯一解为$x_*=b/(1-k)$. 代入确为不动点, 任何不动点都必须满足同一线性方程, 因而唯一; 不需要借助不动点定理代替本题要求的直接论证.
\end{solution}

\begin{exercise}\label{ps190:3:7}
\textbf{PS3第7题(选做, 保留原命题).} 在$\ell^2$中, 原题要求证明$A=\{a=(a_k):|a_k|<k^{-3}\text{ 对每个 }k\ge1\}$是紧子集.
\end{exercise}
\begin{solution}
\textbf{原命题不成立.} 对$m\ge2$, 取$a^{(m)}=(1-1/m)e_1$. 每个坐标均满足严格不等式, 故$a^{(m)}\in A$, 但$\norm{a^{(m)}-e_1}_2=1/m\to0$, 而$e_1\notin A$. 因而$A$在$\ell^2$中不闭, 由紧集必闭知它不紧.

\textbf{编者修正版.} 把每个$<$改为$\le$, 令$\widehat A=\{a\in\ell^2:|a_k|\le k^{-3}\text{ 对每个 }k\ge1\}$, 则$\widehat A$紧. 为完整证明, 取其中任意序列$(a^{(m)})$. 每个固定坐标都在紧区间$[-k^{-3},k^{-3}]$中, 逐坐标抽子列, 再取对角子列$a^{(m_j)}$, 使每个坐标趋于$a_k$, 仍有$|a_k|\le k^{-3}$. 因$\sum k^{-6}<\infty$, 有$a\in\ell^2$且$a\in\widehat A$.

给定$\varepsilon>0$, 先取$N$使$4\sum_{k>N}k^{-6}<\varepsilon^2/2$. 对充分大的$j$, 前$N$个坐标的差平方和小于$\varepsilon^2/2$. 尾部利用$|a_k^{(m_j)}-a_k|\le2k^{-3}$, 得$\norm{a^{(m_j)}-a}_2^2<\varepsilon^2$. 所以对角子列在$\ell^2$中收敛, $\widehat A$序列紧, 由定理\ref{ch:compact:sequential-to-compact}得紧. 只给各坐标有界不能完成论证, 统一尾部控制是关键.
\end{solution}

\begin{exercise}\label{ps190:3:8}
\textbf{PS3第8题(选做, 保留原命题).} 考虑形如$\sum_n a_ne^{inx}$的函数, 其中$|a_n|<(1+|n|)^{-2}$. (a) 证明每个函数属于$C^0([0,2\pi])$. (b) 原题要求证明这个函数集在$C^0([0,2\pi])$中紧.
\end{exercise}
\begin{solution}
原件未标求和指标的范围, 又使用复指数. 以下明确采用$n\in\mathbb Z$和复值连续函数的通常解释, 配上确界距离. 若只取$n\ge0$, 同一证明仍适用; 若要求实值, 需加$a_{-n}=\overline{a_n}$, 以下反例及修正也仍适用. 这是解释题面所需的记号说明, 不是将它悄悄改成另一题.

(a) 令$M_n=(1+|n|)^{-2}$. 因$\sum_{n\in\mathbb Z}M_n<\infty$, 对称有限部分和的任意尾部的上确界不超过相应$\sum M_n$的尾部. 因而级数一致收敛: 逐点在完备的$\mathbb C$中取得极限后, 同一尾部估计给一致收敛. 有限和连续, 一致极限仍连续, 可分别对实部、虚部应用第5章的一致极限证明. 所得函数还以$2\pi$为周期.

(b) \textbf{严格不等式下不紧.} 常数函数$f_m=1-1/m$, $m\ge2$, 取$a_0=1-1/m$, 其余系数为零, 全部满足原条件. 它们一致趋于常数$1$, 但$1$不在原函数集. 为避免假定级数表示天然唯一, 对任一满足系数界的表示, 利用一致收敛逐项积分得$a_0=(2\pi)^{-1}\int_0^{2\pi}f(x)\dd x$, 因非零整频率的积分为零. 若$f=1$, 必有$a_0=1$, 违反$|a_0|<1$. 所以原集不闭, 从而不紧.

\textbf{编者修正版.} 令$\widehat S$为满足$|a_n|\le M_n$的同类级数全体. 每个级数由(a)连续. 给定$\widehat S$中的函数序列$f_m$, 为每个选取一组满足界的系数$a_n^{(m)}$. 枚举可数集$\mathbb Z$, 在每个系数的紧复圆盘中逐次抽子列并取对角子列, 得$a_n^{(m_j)}\to a_n$对每个$n$成立, 且$|a_n|\le M_n$. 令$g=\sum_n a_ne^{inx}\in\widehat S$.

给定$\varepsilon>0$, 先取$N$使$2\sum_{|n|>N}M_n<\varepsilon/2$. 再取$j$充分大, 使有限和$\sum_{|n|\le N}|a_n^{(m_j)}-a_n|<\varepsilon/2$. 则$\norm{f_{m_j}-g}_\infty\le\sum_{|n|\le N}|a_n^{(m_j)}-a_n|+2\sum_{|n|>N}M_n<\varepsilon$. 因而$\widehat S$序列紧, 再用度量空间序列紧与紧的等价得紧. 实值情形的共轭对称条件在逐坐标极限下保持, 所以同样成立.
\end{solution}

% 官方CQ题面逐题译录; 图以自包含TikZ路径重绘.
\originalchapter{官方理解题与解答}\label{ch:official-comprehension}
本章题面译自MIT OpenCourseWare, Paul Seidel的18.900 Geometry and Topology in the Plane, Spring 2023所选第4、29、30、31、33、40讲的Comprehension Questions. 保留官方题号, 共22题. 下列解答均为中文编者补充, 不是MIT官方答案. 前文的编者习题与这些官方题分属不同来源.

第40讲的材料由OCW公开, 但官方课程说明将它列为Spring 2023实际教学中省略的讲次; 本章将其作为公开补充材料收录. 涉及Betti数时采用实数系数. 除第40讲说明的平移曲面记账外, 复形不允许两条不同边具有相同端点.

\originalsection{第4讲: 绕数}
\begin{exercise}\label{cq:4.1}
\textbf{官方题4.1.} 图中的黑点在内部还是外部? 这是一条多边形闭路, 无须检查这一点.
\begin{center}
\begin{tikzpicture}[x=1cm,y=1cm,line width=.35pt]
\draw (4.0000,0.8889) --(3.4074,1.0370) --(4.0000,1.1852) --(2.6667,2.6667) --(2.9630,3.1111) --(3.7037,2.5185) --(3.5556,2.3704) --(2.9630,2.8148) --(2.9630,2.6667) --(3.7037,1.7778) --(4.0000,1.9259) --(3.7037,2.0741) --(4.1481,2.8148) --(3.8519,3.5556) --(3.2593,3.5556) --(3.2593,3.1111) --(3.7037,2.9630) --(3.5556,3.1111) --(3.4074,3.2593) --(3.5556,3.4074) --(3.8519,2.9630) --(3.7037,2.6667) --(2.8148,3.4074) --(2.8148,3.2593) --(2.6667,3.5556) --(3.4074,3.8519) --(4.1481,3.5556) --(4.1481,3.8519) --(2.9630,4.1481) --(1.7778,2.5185) --(1.9259,2.3704) --(2.3704,2.0741) --(2.0741,2.5185) --(2.3704,3.1111) --(2.6667,3.1111) --(2.5185,2.8148) --(2.2222,2.5185) --(2.8148,2.2222) --(1.6296,1.3333) --(2.2222,0.8889) --(1.9259,1.3333) --(2.8148,2.0741) --(3.7037,1.1852) --(2.9630,1.0370) --(3.4074,0.5926) --(1.6296,0.4444) --(1.6296,0.5926) --(2.5185,0.5926) --(2.8148,0.7407) --(2.3704,1.1852) --(2.6667,1.6296) --(2.8148,1.4815) --(2.6667,1.4815) --(2.8148,1.1852) --(3.1111,1.4815) --(2.8148,1.9259) --(2.0741,1.3333) --(2.6667,0.7407) --(1.9259,0.7407) --(1.3333,1.0370) --(1.6296,1.6296) --(2.0741,1.7778) --(1.6296,2.2222) --(2.2222,1.9259) --(1.6296,2.5185) --(0.2963,2.3704) --(0.1481,3.4074) --(0.8889,4.2963) --(1.7778,3.7037) --(2.2222,4.1481) --(2.5185,4.0000) --(1.6296,2.8148) --(1.1852,2.8148) --(1.9259,3.5556) --(1.1852,3.7037) --(0.8889,4.0000) --(0.4444,3.4074) --(0.8889,3.7037) --(1.2237,3.4163) --(1.5704,3.4963) --(1.0370,2.9630) --(0.8889,3.4815) --(0.8889,2.8148) --(0.4444,3.1111) --(0.4444,2.5185) --(1.7778,2.6667) --(2.9630,4.4444) --(4.2963,3.8519) --(4.5926,2.3704) --(4.0000,1.6296) --(4.1481,0.4444) --(1.9259,-0.0000) --(0.0000,0.7407) --(0.2963,2.0741) --(1.3333,2.3704) --(1.6296,1.7778) --(0.7407,1.0370) --(1.4815,0.7407) --(1.3333,0.4444) --(0.4000,0.8444) --(0.5733,1.6844) --(0.5907,1.2351) --(0.7257,1.7235) --(0.8726,1.3723) --(0.8889,1.7778) --(1.4815,1.7778) --(1.1852,2.0741) --(0.4444,1.9259) --(0.1481,0.7407) --(1.3333,0.2963) --(2.3704,0.1481) --(3.4074,0.4444) --(4.0000,0.8889);
\fill (2.5185,2.2222) circle (0.0741);
\end{tikzpicture}
\end{center}
\end{exercise}
\begin{solution}
在内部. 从黑点向右略向上作射线, 例如取斜率$0.01$, 与边界相交5次且避开顶点. 每穿过一次边界就在内部与外部之间切换; 射线终端在外部, 交点数为奇数, 所以起点在内部. 这正是第4讲公式(4.10)的奇偶判别.
\end{solution}

\begin{exercise}\label{cq:4.2}
\textbf{官方题4.2.} 计算按箭头方向行进时围绕黑点的绕数.
\begin{center}
\begin{tikzpicture}[x=1cm,y=1cm,line width=.35pt]
\fill (4.1333,1.8667) circle (0.0667);
\draw (6.8000,1.8667) --(6.2667,4.0000) --(4.0000,1.2000) --(2.6667,2.8000) --(2.5333,2.1333) --(4.4000,2.4000) --(4.9333,0.4000) --(2.9333,0.2667) --(4.2667,3.2000) --(6.4000,2.1333) --(2.8000,1.2000) --(3.6000,0.0000) --(5.6000,1.3333) --(2.6667,1.6000) --(2.9333,3.6000) --(4.8000,4.2667) --(5.3333,0.5333) --(6.4000,1.2000) --(3.4667,2.8000) --(1.3333,1.7333) --(1.8667,0.1333) --(4.4000,1.0667) --(2.2667,1.4667) --(2.0000,2.8000) --(4.0000,3.6000) --(4.6667,2.5333) --(5.4667,3.7333) --(7.2000,3.2000) --(5.3333,2.0000) --(2.1333,3.6000) --(0.0000,2.4000) --(0.6667,1.3333) --(3.4667,0.9333) --(1.0667,2.4000) --(2.6667,4.0000) --(5.8667,0.0000) --(7.0667,0.6667) --(6.2667,1.6000) --(6.8000,1.8667);
\draw (6.5333,2.9333) --(6.6667,2.4000);
\draw (6.6855,2.6452) --(6.6667,2.4000) --(6.5347,2.6075);
\end{tikzpicture}
\end{center}
\end{exercise}
\begin{solution}
绕数为$3$. 从黑点向右略向上作射线, 避开顶点和交叉点. 从近到远, 七次穿越按第4讲公式(4.9)的约定依次贡献$+1,+1,+1,+1,-1,+1,-1$, 总和为$3$. 这里正号表示曲线在交点沿射线的左法向穿过. 必须沿题图箭头取向; 反向绕行会得到$-3$.
\end{solution}

\begin{exercise}\label{cq:4.3}
\textbf{官方题4.3.} 仿照第4讲例4.7, 对下图逐步使用消除自交点的方法, 并记录每一步绕数怎样改变.
\begin{center}
\begin{tikzpicture}[x=.7cm,y=.55cm,line width=.4pt,>=stealth]
\draw (0.000,0.000) -- (0.000,-7.000) -- (6.000,-7.000) -- (6.000,-1.000) -- (1.000,-1.000) -- (1.000,-6.000) -- (5.000,-6.000) -- (5.000,-2.000) -- (2.000,-2.000) -- (2.000,-5.000) -- (4.000,-5.000) -- (4.000,-3.000) -- (3.000,-3.000) -- (3.000,-4.000) -- (7.000,-4.000) -- (7.000,0.000) -- cycle;
\draw[->] (4,0)--(2.5,0);
\foreach \i in {1,...,4}{\node at (\i-.5,-3.5) {$\i$};}
\node[left] at (0,-3.5) {$0$};
\foreach \x/\i in {6/1,5/2,4/3}{\node[below right] at (\x,-4) {$X_{\i}$};}
\end{tikzpicture}
\end{center}
\end{exercise}
\begin{solution}
图中外部绕数为$0$, 四个有界区域的初始绕数依次为$1,2,3,4$. 记它们为$R_1,R_2,R_3,R_4$, 将三个交叉点从外到内标作$X_1,X_2,X_3$. 图中的这些数字及$X_i$为编者增加的标记.

依次处理$X_1,X_2,X_3$. 每次将交叉处分成的两段闭路中较内的一段反向, 再在一个不含其他交点的小圆盘内连接邻近的两端, 使结果仍为一条闭路. 下图给出依次得到的闭路; 顶部箭头保持原方向.
\begin{center}
\begin{tikzpicture}[x=.4cm,y=.34cm,line width=.4pt,>=stealth]
\begin{scope}[xshift=0.0cm]
\draw (0.000,0.000) -- (0.000,-7.000) -- (6.000,-7.000) -- (6.000,-4.150) -- (5.850,-4.000) -- (3.000,-4.000) -- (3.000,-3.000) -- (4.000,-3.000) -- (4.000,-5.000) -- (2.000,-5.000) -- (2.000,-2.000) -- (5.000,-2.000) -- (5.000,-6.000) -- (1.000,-6.000) -- (1.000,-1.000) -- (6.000,-1.000) -- (6.000,-3.850) -- (6.150,-4.000) -- (7.000,-4.000) -- (7.000,0.000) -- cycle;
\draw[->] (4,0)--(2.5,0);
\node at (3.5,-8.1) {第1步};
\end{scope}
\begin{scope}[xshift=3.3cm]
\draw (0.000,0.000) -- (0.000,-7.000) -- (6.000,-7.000) -- (6.000,-4.150) -- (5.850,-4.000) -- (5.150,-4.000) -- (5.000,-3.850) -- (5.000,-2.000) -- (2.000,-2.000) -- (2.000,-5.000) -- (4.000,-5.000) -- (4.000,-3.000) -- (3.000,-3.000) -- (3.000,-4.000) -- (4.850,-4.000) -- (5.000,-4.150) -- (5.000,-6.000) -- (1.000,-6.000) -- (1.000,-1.000) -- (6.000,-1.000) -- (6.000,-3.850) -- (6.150,-4.000) -- (7.000,-4.000) -- (7.000,0.000) -- cycle;
\draw[->] (4,0)--(2.5,0);
\node at (3.5,-8.1) {第2步};
\end{scope}
\begin{scope}[xshift=6.6cm]
\draw (0.000,0.000) -- (0.000,-7.000) -- (6.000,-7.000) -- (6.000,-4.150) -- (5.850,-4.000) -- (5.150,-4.000) -- (5.000,-3.850) -- (5.000,-2.000) -- (2.000,-2.000) -- (2.000,-5.000) -- (4.000,-5.000) -- (4.000,-4.150) -- (3.850,-4.000) -- (3.000,-4.000) -- (3.000,-3.000) -- (4.000,-3.000) -- (4.000,-3.850) -- (4.150,-4.000) -- (4.850,-4.000) -- (5.000,-4.150) -- (5.000,-6.000) -- (1.000,-6.000) -- (1.000,-1.000) -- (6.000,-1.000) -- (6.000,-3.850) -- (6.150,-4.000) -- (7.000,-4.000) -- (7.000,0.000) -- cycle;
\draw[->] (4,0)--(2.5,0);
\node at (3.5,-8.1) {第3步};
\end{scope}
\end{tikzpicture}
\end{center}
在远离三个小圆盘的原区域内各固定一点, 绕数的变化如下. 表中$R_0$表示原外部; 同值区域有时会因消除交点而合并.
\begin{center}
\begin{tabular}{c|rrrrr}
 & $R_0$ & $R_1$ & $R_2$ & $R_3$ & $R_4$ \\\hline
初始 & $0$ & $1$ & $2$ & $3$ & $4$ \\
消除$X_1$后 & $0$ & $1$ & $0$ & $-1$ & $-2$ \\
再消除$X_2$后 & $0$ & $1$ & $0$ & $1$ & $2$ \\
再消除$X_3$后 & $0$ & $1$ & $0$ & $1$ & $0$
\end{tabular}
\end{center}
为核对这些数, 令$L_1,L_2,L_3,L_4$是把原交叉处拆开后得到的四条逆时针嵌套简单闭路, 从外到内编号. 初始绕数是四个$\operatorname{wind}(L_j,q)$之和; 三次反向操作后的系数依次为$(1,-1,-1,-1)$、$(1,-1,1,1)$、$(1,-1,1,-1)$. 在$R_j$内恰有前$j$条闭路的绕数等于$1$, 求部分和即得全表. 局部改接不影响圆盘外固定点的绕数, 因为连续变形过程中没有经过这些点.
\end{solution}

\begin{exercise}\label{cq:4.4}
\textbf{官方题4.4.} 一条多边形闭路只有简单自交点, 且恰有4个这样的自交点. 它能有的最大绕数是多少? 解释, 并画出达到最大绕数的例子.
\end{exercise}
\begin{solution}
最大为$5$. 一般地, 在自交点把有向闭路拆成两条有向闭路时, 原绕数等于两条闭路绕数之和. 持续拆开各闭路自身的交点, 最后至多得到$N+1$条简单闭路, 因为每次拆开增加一条闭路, 并消去至少一个原交点. 每条简单闭路在给定点的绕数属于$\{-1,0,1\}$, 所以$|\operatorname{wind}|\le N+1$. 两条拆开后的闭路之间即使还有相交, 也不影响各自为简单闭路的判断及绕数的可加性.

下面的五层方形螺旋恰有4个简单自交点. 黑点$q$位于五个逆时针圈共同包围的区域, 向右射线的五次穿越都贡献$+1$, 因而绕数为$5$.
\begin{center}
\begin{tikzpicture}[x=.43cm,y=.33cm,line width=.4pt,>=stealth]
\draw (0.000,0.000) -- (0.000,-10.000) -- (9.000,-10.000) -- (9.000,-1.000) -- (1.000,-1.000) -- (1.000,-9.000) -- (8.000,-9.000) -- (8.000,-2.000) -- (2.000,-2.000) -- (2.000,-8.000) -- (7.000,-8.000) -- (7.000,-3.000) -- (3.000,-3.000) -- (3.000,-7.000) -- (6.000,-7.000) -- (6.000,-4.000) -- (4.000,-4.000) -- (4.000,-6.000) -- (10.000,-6.000) -- (10.000,0.000) -- cycle;
\draw[->] (6,0)--(4,0);
\fill (5,-5) circle (.085);
\node[right] at (5,-5) {$q$};
\end{tikzpicture}
\end{center}
\end{solution}

\originalsection{第29讲: 边界矩阵}
本组题的官方说明要求只用截至第29讲已介绍的材料; 不能预先使用下一讲关于平面复形Betti数的结论. 本节解答遵守这一限制.

\begin{exercise}\label{cq:29.1}
\textbf{官方题29.1.} 哪些图是平面复形, 哪些不是? 与讲义相同, 属于复形的顶点均画成粗黑点, 属于复形的三角形均涂灰. 四幅图从左到右依次编号为(a)--(d).
\begin{center}
\begin{tikzpicture}[x=1cm,y=1cm,line width=.35pt]
\draw[fill=black!26] (1.3333,0.9167) --(0.3333,0.4167) --(0.5000,1.5833) --(1.3333,0.9167);
\draw[fill=black!26] (5.3333,0.7500) --(5.3333,2.0833) --(6.6667,1.4167) --(5.3333,0.7500);
\draw[fill=black!26] (6.6667,0.0833) --(5.3333,0.7500) --(6.6667,1.4167) --(6.6667,0.0833);
\draw[fill=black!17] (2.6667,0.2500) --(1.3333,0.9167) --(2.6667,0.9167) --(2.6667,0.2500);
\fill (1.3333,0.9167) circle (0.0833);
\fill (2.6667,0.2500) circle (0.0833);
\fill (2.6667,0.9167) circle (0.0833);
\fill (5.3333,2.0833) circle (0.0833);
\fill (6.6667,1.4167) circle (0.0833);
\fill (5.3333,0.7500) circle (0.0833);
\fill (6.6667,0.0833) circle (0.0833);
\draw (5.3333,2.0833) --(6.6667,0.0833);
\draw[fill=black!26] (8.1667,0.2500) --(7.5000,0.5833) --(8.0000,0.9167) --(8.1667,0.2500);
\draw[fill=black!26] (8.6667,0.7500) --(8.0000,0.9167) --(8.5000,1.2500) --(8.6667,0.7500);
\draw[fill=black!26] (8.8333,0.0833) --(8.1667,0.2500) --(8.5000,0.5833) --(8.8333,0.0833);
\draw[fill=black!26] (8.8333,0.0833) --(9.0000,0.7500) --(9.5000,0.4167) --(8.8333,0.0833);
\draw[fill=black!26] (9.0000,0.7500) --(9.3333,1.4167) --(8.8333,1.2500) --(9.0000,0.7500);
\draw[fill=black!26] (8.8333,1.2500) --(8.5000,1.2500) --(8.6667,1.7500) --(8.8333,1.2500);
\fill (7.5000,0.5833) circle (0.0833);
\fill (8.0000,0.9167) circle (0.0833);
\fill (8.1667,0.2500) circle (0.0833);
\fill (8.5000,0.5833) circle (0.0833);
\fill (8.6667,0.7500) circle (0.0833);
\fill (8.5000,1.2500) circle (0.0833);
\fill (8.6667,1.7500) circle (0.0833);
\fill (8.8333,1.2500) circle (0.0833);
\fill (9.3333,1.4167) circle (0.0833);
\fill (9.0000,0.7500) circle (0.0833);
\fill (9.5000,0.4167) circle (0.0833);
\fill (8.8333,0.0833) circle (0.0833);
\draw (4.5000,1.4167) --(4.0000,1.9167) --(3.5000,1.2500) --(3.6667,0.4167) --(4.3333,0.4167) --(4.1667,0.9167) --(4.5000,1.4167);
\fill (4.0000,1.9167) circle (0.0833);
\fill (3.5000,1.2500) circle (0.0833);
\fill (4.0000,1.4167) circle (0.0833);
\fill (3.8333,0.7500) circle (0.0833);
\fill (3.6667,0.4167) circle (0.0833);
\fill (4.3333,0.4167) circle (0.0833);
\fill (4.1667,0.9167) circle (0.0833);
\fill (4.5000,1.4167) circle (0.0833);
\fill (0.5000,1.5833) circle (0.0833);
\fill (0.3333,0.4167) circle (0.0833);
\draw (2.6667,0.2500) .. controls (3.1667,1.0833) and (2.6667,1.4167) ..(2.1667,1.6111) .. controls (1.6667,1.8055) and (1.1667,1.8611) ..(0.8055,1.8611) .. controls (0.4445,1.8611) and (0.2222,1.8055) ..(0.1111,1.6111) .. controls (0.0000,1.4167) and (0.0000,1.0833) ..(0.3333,0.4167);
\draw (0.3333,0.4167) --(2.6667,0.2500);
\end{tikzpicture}
\end{center}
\end{exercise}
\begin{solution}
(b)、(d)是平面复形, (a)、(c)不是. (a)中的外侧曲边不是线段, 违反平面复形的直边要求; 内部还画出了相同两端点之间的直边. (c)中的两条对角线在没有黑点处相交, 另画的对角线穿过灰三角形的内部. (b)允许有孤立顶点及未填充的多边形区域, 所以内部的黑点无妨. (d)中各三角形只共用整条边或顶点, 且每个三角形的边、每条边的端点都在图中, 满足全部条件.
\end{solution}

\begin{exercise}\label{cq:29.2}
\textbf{官方题29.2.} 给定下列边界矩阵, 画出复形:
\[
D_1=\begin{pmatrix}
-1&-1&-1&-1&0&0&0\\
1&0&0&0&-1&-1&0\\
0&1&0&0&1&0&-1\\
0&0&1&0&0&1&1\\
0&0&0&1&0&0&0
\end{pmatrix},\qquad
D_2=\begin{pmatrix}
1&0&0\\-1&1&0\\0&-1&0\\0&0&0\\1&0&1\\0&0&-1\\0&1&1
\end{pmatrix}.
\]
\end{exercise}
\begin{solution}
$D_1$的列依次给出边$12,13,14,15,23,24,34$. 按这个边顺序读取$D_2$的三列, 得面$123,134,234$. 取$3$为三角形$124$的内部顶点, 将大三角形分成三面, 再从$1$接一条伸向外部顶点$5$的边, 即得所求图.
\begin{center}
\begin{tikzpicture}[x=1cm,y=1cm,line width=.45pt]
\coordinate (v1) at (0,0); \coordinate (v2) at (3,0);
\coordinate (v3) at (1.5,.8); \coordinate (v4) at (1.5,2.2);
\coordinate (v5) at (-1.2,1.2);
\fill[black!15] (v1)--(v2)--(v4)--cycle;
\draw (v1)--(v2)--(v4)--cycle;
\draw (v1)--(v3)--(v2); \draw (v3)--(v4); \draw (v1)--(v5);
\foreach \i in {1,2,3,4,5}{\fill (v\i) circle (.045);}
\node[below left] at (v1) {$1$}; \node[below right] at (v2) {$2$};
\node[below] at (v3) {$3$}; \node[above] at (v4) {$4$}; \node[above left] at (v5) {$5$};
\end{tikzpicture}
\end{center}
\end{solution}

\begin{exercise}\label{cq:29.3}
\textbf{官方题29.3.} 哪些整数能成为平面复形的欧拉示性数? 对能出现的每个值给出例子, 对不能出现的值给出解释.
\end{exercise}
\begin{solution}
每个整数都能出现. 对$m\ge1$, 取$m$个孤立顶点, 得$\chi=m$. 对$m\le0$, 令$r=1-m$, 在互不重叠的角形扇区内画$r$个三角形的边界, 所有边界仅共用一个顶点, 不填面. 此时$V=1+2r$, $E=3r$, $F=0$, 所以$\chi=1-r=m$. 每条边都是线段, 不同三角形只在公共顶点相交, 确实是平面复形. 这些计数不需要下一讲的孔洞公式.
\end{solution}

\begin{exercise}\label{cq:29.4}
\textbf{官方题29.4.} 计算一个三角形的Betti数. 这里三角形包含填充的内部, 即$n_0=3,n_1=3,n_2=1$.
\end{exercise}
\begin{solution}
三条边的边界生成二维空间$\{(x_1,x_2,x_3):x_1+x_2+x_3=0\}$, 所以$\operatorname{rank}D_1=2$. 唯一面的边界非零, 所以$\operatorname{rank}D_2=1$. 代入定理\ref{thm:euler-poincare}, 得$(b_0,b_1,b_2)=(1,0,0)$.
\end{solution}

\originalsection{第30讲: Betti数}
\begin{exercise}\label{cq:30.1}
\textbf{官方题30.1.} 假设下图的灰色部分已分解成三角形, 使整体成为复形. 计算它的Betti数.
\begin{center}
\begin{tikzpicture}[x=1cm,y=1cm,line width=.35pt]
\draw[fill=black!17] (0.0000,1.4286) --(0.2857,1.7143) --(0.5714,1.7143) --(0.5714,0.0000) --(0.2857,0.0000) --(0.2857,1.4286) --(0.0000,1.1429) --(0.0000,1.4286);
\draw[fill=black!17] (1.1429,1.7143) --(0.8571,1.4286) --(0.8571,1.1429) --(1.1429,0.8571) --(0.8571,0.5714) --(0.8571,0.2857) --(1.1429,0.0000) --(1.4286,0.0000) --(1.7143,0.2857) --(1.7143,0.5714) --(1.4286,0.8571) --(1.7143,1.1429) --(1.7143,1.4286) --(1.4286,1.7143) --(1.1429,1.7143);
\draw[fill=black!0] (1.1429,1.1429) rectangle (1.4286,1.4286);
\draw[fill=black!0] (1.1429,0.2857) rectangle (1.4286,0.5714);
\draw[fill=black!17] (2.0000,0.2857) --(2.0000,0.0000) --(2.2857,0.0000) --(2.2857,0.2857) --(2.0000,0.2857);
\draw[fill=black!17] (2.5714,1.4286) --(2.8571,1.7143) --(3.1429,1.7143) --(3.4286,1.4286) --(3.4286,0.2857) --(3.1429,0.0000) --(2.8571,0.0000) --(2.5714,0.2857) --(2.5714,0.5714) --(2.8571,0.5714) --(2.8571,0.2857) --(3.1429,0.2857) --(3.1429,0.8571) --(2.8571,0.8571) --(2.5714,1.1429) --(2.5714,1.4286);
\draw[fill=black!0] (2.8571,1.1429) rectangle (3.1429,1.4286);
\draw[fill=black!17] (3.7143,1.4286) --(4.0000,1.7143) --(4.2857,1.7143) --(4.5714,1.4286) --(4.5714,0.2857) --(4.2857,0.0000) --(4.0000,0.0000) --(3.7143,0.2857) --(3.7143,1.4286);
\draw[fill=black!0] (4.0000,0.2857) rectangle (4.2857,1.4286);
\draw (2.0000,0.2857) --(2.0000,0.2857);
\draw[fill=black!17] (4.8571,1.4286) --(5.1429,1.7143) --(5.4286,1.7143) --(5.7143,1.4286) --(5.7143,0.2857) --(5.4286,0.0000) --(5.1429,0.0000) --(4.8571,0.2857) --(4.8571,1.4286);
\draw[fill=black!0] (5.1429,0.2857) rectangle (5.4286,1.4286);
\end{tikzpicture}
\end{center}
\end{exercise}
\begin{solution}
$(b_0,b_1,b_2)=(6,5,0)$. 字形$1,8,\text{小数点},9,0,0$是六个分支. $8$有两个孔洞, $9$及两个$0$各有一个, 合计五个; $1$与小数点没有孔洞. 用第30讲定理30.2及推论30.5, 平面复形的$b_2=0$, $b_1$等于有界补集分支数, 就得到这些值. 小数点是填充的小方块, 不能忽略其连通分支.
\end{solution}

\begin{exercise}\label{cq:30.2}
\textbf{官方题30.2.} 画一个满足$b_0=3,b_1=3$的平面复形.
\end{exercise}
\begin{solution}
取三个相互分离、内部不填充的三角形边界:
\begin{center}
\begin{tikzpicture}[x=1cm,y=1cm,line width=.45pt]
\foreach \x in {0,2.5,5}{
\draw (\x,0)--(\x+1.5,0)--(\x+.75,1.2)--cycle;
\fill (\x,0) circle (.04); \fill (\x+1.5,0) circle (.04); \fill (\x+.75,1.2) circle (.04);}
\end{tikzpicture}
\end{center}
每个分支有$V=E=3$, $F=0$, $\operatorname{rank}D_1=2$, 所以每个的Betti数为$(1,1,0)$. 三个分支的矩阵互不混合, 总Betti数为$(3,3,0)$.
\end{solution}

\begin{exercise}\label{cq:30.3}
\textbf{官方题30.3.} 两个抽象复形的不交并, 是将它们各自的顶点、边和三角形直接合在一起, 并将两套对象视为不同对象. 若$K=K_1\sqcup K_2$, 它们的Betti数有何关系? 为什么?
\end{exercise}
\begin{solution}
对$i=0,1,2$, 有$b_i(K)=b_i(K_1)+b_i(K_2)$. 先列$K_1$的基, 再列$K_2$的基, 两个边界矩阵就都成为对应矩阵的块对角和. 块对角矩阵的秩等于两块的秩之和, 各维单形数也相加; 将这些等式代入定理\ref{thm:euler-poincare}的三个公式即得结论. 等价地, 同调向量空间本身也是两个同调空间的直和.
\end{solution}

\begin{exercise}\label{cq:30.4}
\textbf{官方题30.4.} 将M\"obius带构造成一个抽象复形, 并计算它的Betti数. 注意按本课程的定义, 不允许不同边具有相同的端点. 允许用计算机辅助计算矩阵的秩.
\end{exercise}
\begin{solution}
将一个长方形分成三个小方格, 每格从左上角连到右下角, 再将左右两边反向粘合. 六个不同顶点记作$a_0,a_1,a_2,b_0,b_1,b_2$, 并约定$a_3=b_0,b_3=a_0$. 六个面的顶点集为
\[
\{a_i,b_i,b_{i+1}\},\quad \{a_i,a_{i+1},b_{i+1}\},\qquad i=0,1,2.
\]
边取这些面的全部二元子集. 图中相同字母表示同一顶点, 两侧箭头表示实际粘合方向.
\begin{center}
\begin{tikzpicture}[x=1.3cm,y=1cm,line width=.45pt,>=stealth]
\draw (0,0) rectangle (3,1.4);
\foreach \i in {0,1,2}{\draw (\i,1.4)--(\i+1,0);}
\foreach \i in {1,2}{\draw (\i,0)--(\i,1.4);}
\foreach \i/\a/\b in {0/a_0/b_0,1/a_1/b_1,2/a_2/b_2,3/b_0/a_0}{
\fill (\i,1.4) circle (.035); \fill (\i,0) circle (.035);
\node[above] at (\i,1.4) {$\a$}; \node[below] at (\i,0) {$\b$};}
\draw[->] (-.25,.3)--(-.25,1.1); \draw[->] (3.25,1.1)--(3.25,.3);
\end{tikzpicture}
\end{center}
边恰有三条$a_ib_i$、三条$a_ib_{i+1}$及六条边界边$a_ia_{i+1},b_ib_{i+1}$, 共12条; 代入端点约定检查可知没有重复的无序端点对. 面的顶点集也两两不同. 因而这确为抽象复形, 几何实现正是上述反向粘合得到的M\"obius带.

它连通, 故$\operatorname{rank}D_1=V-1=5$. 每个三角形各有一条只属于该面的边界边, 所以$D_2t=0$时沿这些边逐一读出该面系数为零; 六列独立, $\operatorname{rank}D_2=6$. 因此$V=6,E=12,F=6$给出$(b_0,b_1,b_2)=(1,1,0)$. 这也与$\chi=0$一致.
\end{solution}

\begin{exercise}\label{cq:30.5}
\textbf{官方题30.5.} 平面上有$v_1=(0,0),v_2=(1,0),v_3=(0,1),v_4=(1,1)$. 尺度$\sigma=0.7,1.1,1.5$时的Vietoris--Rips复形各是什么?
\end{exercise}
\begin{solution}
四条正方形边长为$1$, 两条对角线长为$\sqrt2$. 按第30讲的定义, 两点距离小于$\sigma$就加边, 三点两两距离小于$\sigma$就填三角形.
\begin{enumerate}[label=(\arabic*)]
\item $\sigma=0.7$: 只有四个孤立顶点, 没有边或面.
\item $\sigma=1.1$: 有边$12,13,24,34$, 即正方形边界, 没有面; 任意三个顶点都包含一对对角点.
\item $\sigma=1.5$: 有全部六条边和四个三角面$123,124,134,234$, 是抽象四面体的边界.
\end{enumerate}
本课程的复形定义只保留到二维, 所以最后一个模型没有三维四面体内部. 若采用通常的全维Vietoris--Rips定义, 最后一项还应包含一个三维单形; 这不是本题采用的约定. 前述三个二维模型的Betti数分别为$(4,0,0)$、$(1,1,0)$、$(1,0,1)$, 最后一项也见\hyperlink{example.14.1}{例题14.1}.
\end{solution}

\originalsection{第31讲: 组合曲面}
\begin{exercise}\label{cq:31.1}
\textbf{官方题31.1.} 将两个四面体在一个顶点粘在一起. 结果是曲面吗? 简短解释.
\begin{center}
\begin{tikzpicture}[x=1cm,y=1cm,line width=.45pt]
\coordinate (v) at (0,0);
\coordinate (a) at (-.8,.9); \coordinate (b) at (-1.3,.25); \coordinate (c) at (-.8,-.5);
\coordinate (d) at (.8,.9); \coordinate (e) at (1.3,.25); \coordinate (f) at (.8,-.5);
\fill[black!15] (a)--(b)--(c)--(v)--cycle;
\fill[black!15] (d)--(e)--(f)--(v)--cycle;
\draw (a)--(b)--(c)--(v)--(a)--(c); \draw[dashed] (b)--(v);
\draw (d)--(e)--(f)--(v)--(d)--(f); \draw[dashed] (e)--(v);
\foreach \p in {v,a,b,c,d,e,f}{\fill (\p) circle (.04);}
\node[below] at (v) {$v$};
\end{tikzpicture}
\end{center}
\end{exercise}
\begin{solution}
不是. 粘合点$v$的链环是两个互不相交的三角形圆周, 而非一个圆周, 违反第14章的顶点局部条件. 等价地, 其足够小邻域删去$v$后分成两部分, 而圆盘邻域删去中心仍连通. 每条边邻接两面这一条件单独不能保证是曲面.
\end{solution}

\begin{exercise}\label{cq:31.2}
\textbf{官方题31.2.} 明确证明八面体可定向.
\end{exercise}
\begin{solution}
将上下两个顶点记为$N,S$, 赤道四个顶点依次记为$v_0,v_1,v_2,v_3$, 下标模$4$. 给四个上半面的方向$[N,v_i,v_{i+1}]$, 四个下半面的方向$[S,v_{i+1},v_i]$.

在赤道边$v_iv_{i+1}$上, 上面沿$v_i\to v_{i+1}$, 下面沿反方向. 在边$Nv_i$上, 相邻的两个上半面分别沿$N\to v_i$和$v_i\to N$; $Sv_i$同理. 这些是全部边, 故每条共边两侧的诱导方向相反, 所列八个方向组成定向.
\end{solution}

\begin{exercise}\label{cq:31.3}
\textbf{官方题31.3.} 讲义指出曲面的三角形个数必为偶数. 为什么?
\end{exercise}
\begin{solution}
按本讲的闭组合曲面定义, 每个三角形有3条边, 每条边恰邻接2个三角形. 对面边关联计数, 得$3F=2E$, 因而$F$为偶数. 这里不能直接套到有边界曲面, 其边界边只邻接一个三角形.
\end{solution}

\begin{exercise}\label{cq:31.4}
\textbf{官方题31.4.} 是否存在欧拉示性数为$3$的曲面? 允许任意多个连通分支. 若存在, 举例; 若不存在, 解释.
\end{exercise}
\begin{solution}
存在: 取球面与实射影平面的不交并$S^2\sqcup\mathbb RP^2$. 球面的四面体边界模型有$\chi=4-6+4=2$; 射影平面的组合模型有$\chi=1$, 见\hyperlink{example.14.2}{例题14.2}. 两者先采用互不相交的顶点集, 不交并的顶点、边、面数分别相加, 故总欧拉数为$2+1=3$. 该曲面有一个不可定向分支, 与命题\ref{thm:orientable-even}不冲突.
\end{solution}

\begin{exercise}\label{cq:31.5}
\textbf{官方题31.5.} 第6讲的三角形铺砌(6.9)、(6.11)是周期的, 所以各自都能给出一个环面的组合版本. 画出这两个曲面. 注意按本课程的定义, 不允许不同边具有相同端点.
\end{exercise}
\begin{solution}
下面画的是两个足够大的周期平行四边形, 同向箭头所标的两对边分别按平移粘合. 顶点与边界上的中点也依同一平移配对. 图中的黑点是三角剖分顶点; 四个角都记作$v_{00}$, 其余重复的边界顶点按下述周期规则识别.
\begin{center}
\begin{tikzpicture}[x=.75cm,y=.75cm,line width=.35pt,>=stealth]
\draw (0,0)--(0,4);
\draw (0,0)--(4,0);
\draw (1,0)--(1,4);
\draw (0,1)--(4,1);
\draw (2,0)--(2,4);
\draw (0,2)--(4,2);
\draw (3,0)--(3,4);
\draw (0,3)--(4,3);
\draw (4,0)--(4,4);
\draw (0,4)--(4,4);
\draw (0.00000,0.00000)--(1.00000,1.00000);
\draw (1.00000,1.00000)--(0.00000,2.00000);
\draw (0.00000,2.00000)--(1.00000,3.00000);
\draw (1.00000,3.00000)--(0.00000,4.00000);
\draw (2.00000,0.00000)--(1.00000,1.00000);
\draw (1.00000,1.00000)--(2.00000,2.00000);
\draw (2.00000,2.00000)--(1.00000,3.00000);
\draw (1.00000,3.00000)--(2.00000,4.00000);
\draw (2.00000,0.00000)--(3.00000,1.00000);
\draw (3.00000,1.00000)--(2.00000,2.00000);
\draw (2.00000,2.00000)--(3.00000,3.00000);
\draw (3.00000,3.00000)--(2.00000,4.00000);
\draw (4.00000,0.00000)--(3.00000,1.00000);
\draw (3.00000,1.00000)--(4.00000,2.00000);
\draw (4.00000,2.00000)--(3.00000,3.00000);
\draw (3.00000,3.00000)--(4.00000,4.00000);
\fill (0,0) circle (.035);
\fill (0,1) circle (.035);
\fill (0,2) circle (.035);
\fill (0,3) circle (.035);
\fill (0,4) circle (.035);
\fill (1,0) circle (.035);
\fill (1,1) circle (.035);
\fill (1,2) circle (.035);
\fill (1,3) circle (.035);
\fill (1,4) circle (.035);
\fill (2,0) circle (.035);
\fill (2,1) circle (.035);
\fill (2,2) circle (.035);
\fill (2,3) circle (.035);
\fill (2,4) circle (.035);
\fill (3,0) circle (.035);
\fill (3,1) circle (.035);
\fill (3,2) circle (.035);
\fill (3,3) circle (.035);
\fill (3,4) circle (.035);
\fill (4,0) circle (.035);
\fill (4,1) circle (.035);
\fill (4,2) circle (.035);
\fill (4,3) circle (.035);
\fill (4,4) circle (.035);
\draw[->] (1,-.3)--(3,-.3); \draw[->] (1,4.3)--(3,4.3);
\draw[->] (-.3,1)--(-.3,3); \draw[->] (4.3,1)--(4.3,3);
\node at (2,-.85) {(6.9) 的周期商};
\node[below left] at (0,0) {$v_{00}$};
\node[below right] at (4,0) {$v_{00}$};
\node[above left] at (0,4) {$v_{00}$};
\node[above right] at (4,4) {$v_{00}$};
\begin{scope}[xshift=5.3cm]
\draw (0.00000,0.00000)--(1.00000,0.00000);
\draw (1.00000,0.00000)--(0.50000,0.86603);
\draw (0.50000,0.86603)--(0.00000,0.00000);
\draw (0.00000,0.00000)--(0.75000,0.43301);
\draw (0.50000,0.86603)--(1.50000,0.86603);
\draw (1.50000,0.86603)--(1.00000,0.00000);
\draw (1.50000,0.86603)--(0.75000,0.43301);
\draw (1.50000,0.86603)--(1.00000,1.73205);
\draw (1.00000,1.73205)--(0.50000,0.86603);
\draw (1.50000,0.86603)--(0.75000,1.29904);
\draw (1.00000,1.73205)--(2.00000,1.73205);
\draw (2.00000,1.73205)--(1.50000,0.86603);
\draw (1.50000,0.86603)--(1.50000,1.73205);
\draw (2.00000,1.73205)--(1.50000,2.59808);
\draw (1.50000,2.59808)--(1.00000,1.73205);
\draw (1.50000,2.59808)--(1.50000,1.73205);
\draw (1.50000,2.59808)--(2.50000,2.59808);
\draw (2.50000,2.59808)--(2.00000,1.73205);
\draw (1.50000,2.59808)--(2.25000,2.16506);
\draw (1.00000,0.00000)--(2.00000,0.00000);
\draw (2.00000,0.00000)--(1.50000,0.86603);
\draw (1.50000,0.86603)--(1.50000,0.00000);
\draw (1.50000,0.86603)--(2.50000,0.86603);
\draw (2.50000,0.86603)--(2.00000,0.00000);
\draw (1.50000,0.86603)--(2.25000,0.43301);
\draw (2.50000,0.86603)--(2.00000,1.73205);
\draw (1.50000,0.86603)--(2.25000,1.29904);
\draw (2.00000,1.73205)--(3.00000,1.73205);
\draw (3.00000,1.73205)--(2.50000,0.86603);
\draw (3.00000,1.73205)--(2.25000,1.29904);
\draw (3.00000,1.73205)--(2.50000,2.59808);
\draw (3.00000,1.73205)--(2.25000,2.16506);
\draw (2.50000,2.59808)--(3.50000,2.59808);
\draw (3.50000,2.59808)--(3.00000,1.73205);
\draw (3.00000,1.73205)--(3.00000,2.59808);
\draw (2.00000,0.00000)--(3.00000,0.00000);
\draw (3.00000,0.00000)--(2.50000,0.86603);
\draw (3.00000,0.00000)--(2.25000,0.43301);
\draw (2.50000,0.86603)--(3.50000,0.86603);
\draw (3.50000,0.86603)--(3.00000,0.00000);
\draw (3.00000,0.00000)--(3.00000,0.86603);
\draw (3.50000,0.86603)--(3.00000,1.73205);
\draw (3.00000,1.73205)--(3.00000,0.86603);
\draw (3.00000,1.73205)--(4.00000,1.73205);
\draw (4.00000,1.73205)--(3.50000,0.86603);
\draw (3.00000,1.73205)--(3.75000,1.29904);
\draw (4.00000,1.73205)--(3.50000,2.59808);
\draw (3.00000,1.73205)--(3.75000,2.16506);
\draw (3.50000,2.59808)--(4.50000,2.59808);
\draw (4.50000,2.59808)--(4.00000,1.73205);
\draw (4.50000,2.59808)--(3.75000,2.16506);
\fill (0.00000,0.00000) circle (.035);
\fill (0.50000,0.86603) circle (.035);
\fill (1.00000,1.73205) circle (.035);
\fill (1.50000,2.59808) circle (.035);
\fill (1.00000,0.00000) circle (.035);
\fill (1.50000,0.86603) circle (.035);
\fill (2.00000,1.73205) circle (.035);
\fill (2.50000,2.59808) circle (.035);
\fill (2.00000,0.00000) circle (.035);
\fill (2.50000,0.86603) circle (.035);
\fill (3.00000,1.73205) circle (.035);
\fill (3.50000,2.59808) circle (.035);
\fill (3.00000,0.00000) circle (.035);
\fill (3.50000,0.86603) circle (.035);
\fill (4.00000,1.73205) circle (.035);
\fill (4.50000,2.59808) circle (.035);
\fill (2.25000,0.43301) circle (.03);
\fill (2.25000,2.16506) circle (.03);
\fill (0.75000,0.43301) circle (.03);
\fill (2.25000,1.29904) circle (.03);
\fill (1.50000,1.73205) circle (.03);
\fill (0.75000,1.29904) circle (.03);
\fill (3.00000,0.86603) circle (.03);
\fill (3.75000,2.16506) circle (.03);
\fill (3.75000,1.29904) circle (.03);
\fill (1.50000,0.00000) circle (.03);
\fill (3.00000,2.59808) circle (.03);
\draw[->] (.5,-.3)--(2.5,-.3);
\draw[->] (2,2.9)--(4,2.9);
\draw[->] (-.2,.3)--(.8,2.03);
\draw[->] (3.2,.3)--(4.2,2.03);
\node at (2.2,-.85) {(6.11) 的周期商};
\node[below left] at (0,0) {$v_{00}$};
\node[below right] at (3,0) {$v_{00}$};
\node[above left] at (1.5,2.59808) {$v_{00}$};
\node[above right] at (4.5,2.59808) {$v_{00}$};
\end{scope}
\end{tikzpicture}
\end{center}
对于(6.9), 取$4\times4$方格, 每格沿一条对角线分成两面, 相邻格的对角线方向交替. 顶点是$v_{ij}$, $i,j\in\mathbb Z/4\mathbb Z$, 水平边、竖直边分别连接$v_{ij}$与$v_{i+1,j},v_{i,j+1}$. 当$i+j$为偶数时对角边为$v_{ij}v_{i+1,j+1}$, 奇数时为$v_{i+1,j}v_{i,j+1}$. 这给出$V=16,E=48,F=32$.

对于(6.11), 取$u=(1,0),w=(1/2,\sqrt3/2)$张成的等边三角形格. 格点$p_{ij}=iu+jw$的颜色取$i+2j\pmod3$. 每个等边三角形恰有一个颜色$0$的顶点; 从该顶点连到对边中点, 分成两个$30^\circ/60^\circ/90^\circ$三角形. 这正是(6.11)的反射铺砌模式: 每个大三角形由颜色$0$的角向对边作一条高. 按周期$3u,3w$取商, 格点成为$v_{ij}$, $i,j\in\mathbb Z/3\mathbb Z$. 18个大三角形各分成两面; 颜色$1$、$2$之间的9条边各增加一个中点, 所以$V=18,E=54,F=36$.

两幅图都对平面周期格取商, 因而每个顶点的链环仍为一个圆周, 每条边邻接两面, 几何实现是环面. 还须检查端点限制: 第一图中非零边位移只为水平、竖直及相应对角位移, 在模$4$下这些无序端点对不重复. 第二图原格边的位移为$u,w,w-u$, 模$3$下不同方向及不同起点也不产生重复无序端点对; 新中点按所在旧边独立编号, 连接它的两个端点及两个对侧顶点互不相同. 因而细分后也无环边、重边或重复面. 这避免把过小基本域的胞腔模型误当成题目要求的抽象复形.
\end{solution}

\originalsection{第33讲: 组合绕数}
\begin{exercise}\label{cq:33.1}
\textbf{官方题33.1.} 考虑至少有一条边的抽象复形, 即$n_1>0$. 方程$D_2^{\mathsf T}c=0$是否可能只有零解$c=0$? 若可能, 举例; 若不可能, 解释.
\end{exercise}
\begin{solution}
不可能. 选一条边$ij$, 取顶点向量$x$在$i$处为$1$, 其余为$0$. 边向量$c=D_1^{\mathsf T}x$在边$ij$上的分量为$\pm1$, 因而非零. 又由命题\ref{thm:boundary-square}, $D_2^{\mathsf T}D_1^{\mathsf T}=(D_1D_2)^{\mathsf T}=0$, 所以$c$是所需的非零解.

这个解是顶点势的梯度, 对所有闭链的测量都为零, 正如\hyperlink{exercise.14.3}{习题14.3}所证. 因此“有非零解”不能直接解释成“有非零的有效绕数”或$b_1>0$.
\end{solution}

\originalsection{第40讲: 组合曲面的几何}
\begin{exercise}\label{cq:40.1}
\textbf{官方题40.1.} 在由等边三角形组成的八面体上, 找一条周期测地线.
\end{exercise}
\begin{solution}
取八面体顶点$a,b,c,d,e,f$, 其中$a,e$、$b,d$、$c,f$各为一对相对顶点. 按顺序连接边$ab,ac,cd,de,ef,fb$的中点, 最后回到$ab$中点. 这六段分别位于面$abc,acd,cde,def,efb,fba$内.

将这六个等边三角形沿依次经过的边展开, 得下列三角形带. $a',b'$分别是$a,b$的另一份平面表示, 左端边$ab$与右端边$a'b'$粘合. 箭头所示水平线穿过全部六面的相应边中点.
\begin{center}
\begin{tikzpicture}[x=1.7cm,y=1.7cm,line width=.45pt,>=stealth]
\coordinate (a) at (0,.8660254); \coordinate (b) at (-.5,0);
\coordinate (c) at (.5,0); \coordinate (d) at (1,.8660254);
\coordinate (e) at (1.5,0); \coordinate (f) at (2,.8660254);
\coordinate (bp) at (2.5,0); \coordinate (ap) at (3,.8660254);
\draw (b)--(a)--(d)--(f)--(ap)--(bp)--(e)--(c)--cycle;
\draw (a)--(c)--(d)--(e)--(f)--(bp);
\draw[->,line width=.7pt] (-.25,.4330127)--(2.75,.4330127);
\foreach \p/\lab in {a/a,b/b,c/c,d/d,e/e,f/f,bp/b',ap/a'}{\fill (\p) circle (.025);}
\node[above] at (a) {$a$}; \node[above] at (d) {$d$}; \node[above] at (f) {$f$}; \node[above] at (ap) {$a'$};
\node[below] at (b) {$b$}; \node[below] at (c) {$c$}; \node[below] at (e) {$e$}; \node[below] at (bp) {$b'$};
\end{tikzpicture}
\end{center}
展开后的路线是一条直线, 每次穿越共边都按相同速度和相同角度继续, 且不经过顶点. 最后一条边与第一条边的识别在图中是平移$(3,0)$, 所以终点和出发点相同, 切向方向也相同, 可周期延续. 若八面体边长为$1$, 一周期长为$6\cdot(1/2)=3$.
\end{solution}

\begin{exercise}\label{cq:40.2}
\textbf{官方题40.2.} 平移曲面能有正的欧拉示性数吗?
\end{exercise}
\begin{solution}
不能. 第40讲中, 原三角形的一个角为$\pi a/b$, $a,b$互素正整数时, 每个对应商顶点由$2b$份该角围成, 总角为$2\pi a$, 角亏为$2\pi(1-a)\le0$. 所有顶点角亏均非正, 由定理\ref{thm:discrete-gauss-bonnet}有$2\pi\chi=\sum_v\kappa_v\le0$.

更一般的平移粘合模型也有整倍$2\pi$的锥角, 结论相同. 第40讲允许平移曲面的记账出现相同端点间的不同边, 并将边另行编号; 此时$3F=2E$及角亏计数仍成立, 所以不影响上述证明. 如需严格单纯复形, 可细分后使用同一欧拉数.
\end{solution}

\begin{exercise}\label{cq:40.3}
\textbf{官方题40.3.} 一个三角形的角为$(\pi/7,2\pi/7,4\pi/7)$. 它对应的平移曲面由多少个三角形组成? 欧拉示性数是多少?
\end{exercise}
\begin{solution}
共14个三角形, $\chi=-4$. 两条交于$\pi/7$角的边的反射复合生成$2\pi/7$旋转, 共有7种旋转. 另外两角产生的旋转已在其中, 再加上7种反射后的方向, 一共14种标记三角形的方向. 不能把仅经旋转或反射重合的三角形当成平移副本删除.

每个角的既约分母都是$7$, 对应一个商顶点需要14份该角, 而各类型角恰各有14份. 不同类型的角在边粘合时不会互相识别, 因而共有3个顶点. 闭曲面每条边邻接两面, $E=3F/2=21$, 得$\chi=3-21+14=-4$. 也可用三个顶点的角亏$0,-2\pi,-6\pi$核对: 总和为$-8\pi=2\pi\chi$.
\end{solution}

\originalchapter{18.901 官方作业}\label{ch:official-topology-assignment}

本章依据 MIT OCW 的 18.901 \emph{Introduction to Topology}, Fall 2004, James Munkres 课程作业页及其 Problem Set 5 原件整理. PS5 原件共两页: 第一页说明作业规则, 第二页是一张九行、十四列的空表. 前六行是具体空间或空间类别, 后三行是操作; 总计 $9\times14=126$ 个判定格. 以下中文题面依据原件翻译, 解答均为编者补充, 不是 MIT 发布的官方答案. 集合论采用 ZFC; 特别是任意积及基数比较按含选择公理的通常约定理解.

\originalsection{Problem Set 5 题面}
\begin{exercise}\label{ex:official-ps5}
本作业为选做, 用来补救任何一次成绩低于 $70$ 分的必做作业. 从下表选一行或多行, 在所选行的每个格内填 $+$ 表示“是”, 或填 $-$ 表示“否”. 每完整答对一行, 教师在当前最低分的那份作业上加 $6$ 分, 该份作业最多加至 $70$ 分. 同往常一样, 不允许合作. 可以在期末考试之前或考试时提交; 评判结果为最终结果. 原件以“No box tops needed to enter!”的玩笑结束这段说明, 意为不必凭商品包装上的兑奖标志参加.

前六行问:“给定空间是否满足该性质?” 后三行问:“给定操作是否保持该性质?” 为使原件中的课程记号可以独立阅读, 六行空间及三行操作说明如下.
\begin{enumerate}
\item $S_\Omega$: 所有可数序数构成的集合 $[0,\omega_1)$, 配以序拓扑. $\Omega=\omega_1$ 是第一个不可数序数, 不属于这个空间; $[0,\omega_1]$ 则是另一个包含最大元的空间.
\item 有序方形: 集合 $[0,1]\times[0,1]$ 按字典序排列, 即先比较第一坐标, 第一坐标相同才比较第二坐标, 再取这个线性序的序拓扑. 它与通常的欧氏方形拓扑不同.
\item $\R_\ell\times\R_\ell$: 每个 $\R_\ell$ 的下限拓扑以 $[a,b)$ ($a<b$) 为基, 两因子取积拓扑. 此空间通常称为 Sorgenfrey 平面.
\item $\R^\omega$ 的一致拓扑: $\omega$ 表示可数坐标集, 一致度量为 $\rho(x,y)=\sup_n\min\{1,|x_n-y_n|\}$. 不限制各序列本身有界.
\item $\R^I$ 的积拓扑: $I$ 是指标集, 基本开集只对有限个坐标施加开区间条件. 原件未限定 $|I|$, 所以答案必须区分空集、有限集、可数无限集及不可数集. 此行不是箱拓扑: 箱拓扑允许同时限制所有坐标. 第\ref{ch:products}章已说明二者的差别.
\item 任意度量空间: 问某性质是否由“是度量空间”这一条件保证, 并非挑选一个特定度量空间.
\item 取可数积: 各因子都有该性质时, 它们的可数积是否仍有该性质? 这里取积拓扑.
\item 取闭子空间: 空间有该性质时, 每个闭子空间是否都有该性质?
\item 取开子空间: 空间有该性质时, 每个开子空间是否都有该性质?
\end{enumerate}

十四列依原件顺序为以下性质. 原表没有零维或有限维列.
\begin{enumerate}
\item 连通: 不存在两个不交非空开集将空间分成两部分.
\item 道路连通: 任意两点之间都有连续道路 $\gamma:[0,1]\to X$.
\item 紧致: 每个开覆盖都有有限子覆盖.
\item 局部紧 Hausdorff: 空间 Hausdorff, 且每点有一个紧邻域. 在 Hausdorff 情形, 这等价于每点有开邻域, 其闭包紧.
\item Hausdorff: 两个不同点有互不相交的开邻域. 这是原件中独立的一列.
\item 正则: 按第\ref{ch:separation}章约定包含 $T_1$, 且点与不含它的闭集可以用不交开邻域分离.
\item 正规: 包含 $T_1$, 且不交闭集可以用不交开邻域分离.
\item 第一可数: 每点有可数邻域基.
\item 第二可数: 有可数开集基.
\item Lindel\"of: 每个开覆盖都有可数子覆盖.
\item 有可数稠密子集: 即可分.
\item 局部可度量化: 每点有一个开邻域, 该邻域的子空间拓扑可由某个度量生成. 此列不另加全空间 Hausdorff 条件.
\item 可度量化: 全空间拓扑可由某个度量生成.
\item 完全正则: 包含 $T_1$, 且对每个闭集 $F$ 和 $x\notin F$, 存在连续函数 $f:X\to[0,1]$ 满足 $f(x)=1$ 和 $f|_F=0$.
\end{enumerate}
后三行的肯定项是对所有满足所讨论性质的空间作出的断言; 不为其他列额外假定 Hausdorff、正则或非空. 空空间视为连通、道路连通且紧致, 这些约定不改变表中的反例.
\end{exercise}

\originalsection{判定与证明}
\begin{solution}
下面两张表按原顺序各列七项, 合起来覆盖全部 $126$ 格. 对 $\R^I$ 行, $\mathrm E$ 表示“当且仅当 $I=\varnothing$”, $\mathrm F$ 表示“当且仅当 $I$ 有限”, $\mathrm C$ 表示“当且仅当 $I$ 至多可数”, $\mathrm D$ 表示“当且仅当 $|I|\le\mathfrak c$”, 其中 $\mathfrak c=|\R|$. 这些条件是相应格中 $+$ 的准确条件, 条件不成立即填 $-$. 其他行的 $+$ 与 $-$ 均无附加条件.

\begin{center}
\setlength{\tabcolsep}{4pt}
\begin{tabular}{lccccccc}
\hline
空间或操作 & 连通 & \shortstack{道路\\连通} & 紧致 & \shortstack{局部紧\\Hausdorff} & Hausdorff & 正则 & 正规\\
\hline
$S_\Omega$ & $-$ & $-$ & $-$ & $+$ & $+$ & $+$ & $+$\\
有序方形 & $+$ & $-$ & $+$ & $+$ & $+$ & $+$ & $+$\\
$\R_\ell\times\R_\ell$ & $-$ & $-$ & $-$ & $-$ & $+$ & $+$ & $-$\\
$\R^\omega$ 一致拓扑 & $-$ & $-$ & $-$ & $-$ & $+$ & $+$ & $+$\\
$\R^I$ 积拓扑 & $+$ & $+$ & $\mathrm E$ & $\mathrm F$ & $+$ & $+$ & $\mathrm C$\\
任意度量空间 & $-$ & $-$ & $-$ & $-$ & $+$ & $+$ & $+$\\
取可数积 & $+$ & $+$ & $+$ & $-$ & $+$ & $+$ & $-$\\
取闭子空间 & $-$ & $-$ & $+$ & $+$ & $+$ & $+$ & $+$\\
取开子空间 & $-$ & $-$ & $-$ & $+$ & $+$ & $+$ & $-$\\
\hline
\end{tabular}
\end{center}

\begin{center}
\setlength{\tabcolsep}{4pt}
\begin{tabular}{lccccccc}
\hline
空间或操作 & \shortstack{第一\\可数} & \shortstack{第二\\可数} & Lindel\"of & 可分 & \shortstack{局部\\可度量} & 可度量 & \shortstack{完全\\正则}\\
\hline
$S_\Omega$ & $+$ & $-$ & $-$ & $-$ & $+$ & $-$ & $+$\\
有序方形 & $+$ & $-$ & $+$ & $-$ & $-$ & $-$ & $+$\\
$\R_\ell\times\R_\ell$ & $+$ & $-$ & $-$ & $+$ & $-$ & $-$ & $+$\\
$\R^\omega$ 一致拓扑 & $+$ & $-$ & $-$ & $-$ & $+$ & $+$ & $+$\\
$\R^I$ 积拓扑 & $\mathrm C$ & $\mathrm C$ & $\mathrm C$ & $\mathrm D$ & $\mathrm C$ & $\mathrm C$ & $+$\\
任意度量空间 & $+$ & $-$ & $-$ & $-$ & $+$ & $+$ & $+$\\
取可数积 & $+$ & $+$ & $-$ & $+$ & $-$ & $+$ & $+$\\
取闭子空间 & $+$ & $+$ & $+$ & $-$ & $+$ & $+$ & $+$\\
取开子空间 & $+$ & $+$ & $-$ & $+$ & $+$ & $+$ & $+$\\
\hline
\end{tabular}
\end{center}

先说明会重复使用的判据. 命题\ref{thm:second-countable}给出第二可数蕴含第一可数、可分和 Lindel\"of; 定理\ref{thm:metric-countability}给出度量空间中第二可数、可分与 Lindel\"of 等价. 命题\ref{thm:normal-examples}给出度量空间正规、紧 Hausdorff 空间正规, 以及正则 Lindel\"of 空间正规. 正规空间由 Urysohn 引理\ref{thm:urysohn}完全正则. Hausdorff 空间中紧集闭, 见引理\ref{thm:compact-hausdorff}.

完全正则性传给子空间及任意积的细节也在这里固定. 若 $A\subseteq X$, $F\subseteq A$ 相对闭, 将 $F$ 写成 $A\cap H$, 其中 $H$ 在 $X$ 中闭; 对 $x\in A\setminus F$ 用 $X$ 上分离 $x,H$ 的函数再限制到 $A$. 积空间中, 对 $x\notin F$ 取基本矩形 $\bigcap_{i\in J}\pi_i^{-1}(U_i)$ 包含 $x$ 并避开 $F$, 其中 $J$ 有限. 在每个因子取 $f_i(x_i)=1$, $f_i=0$ 于 $X_i\setminus U_i$. 函数 $\prod_{i\in J}f_i\circ\pi_i$ 连续, 在 $x$ 处为 $1$, 在 $F$ 上为 $0$. $T_1$ 性同样传给子空间与积. 完全正则蕴含正则: 上述函数的开集 $\{f>2/3\}$ 与 $\{f<1/3\}$ 分别包含该点和该闭集并互不相交. 下述证明依次对应原件九行.

\begin{enumerate}
\item \textbf{$S_\Omega$.}
每个初始区间 $[0,\alpha]$ 都紧. 事实上, 对它的任一开覆盖, 用超限归纳证明每个 $[0,\beta]$ 都有有限子覆盖: 后继步骤加一个覆盖后继点的开集; 极限步骤中, 覆盖 $\beta$ 的开集包含某个 $(\gamma,\beta]$, 再覆盖 $[0,\gamma]$ 即可. 这也证明 $[0,\omega_1]$ 紧. 序拓扑 Hausdorff; 每个 $\alpha$ 的开邻域 $[0,\alpha+1)$ 就是紧开集 $[0,\alpha]$, 所以局部紧 Hausdorff.

对每个 $\alpha<\omega_1$, $[0,\alpha]$ 可数, 其序拓扑有可数基, 并且紧 Hausdorff 而正则; 定理\ref{thm:metrization}使它可度量化. 上述开邻域证明局部可度量化, 从而第一可数. 集合 $[0,\alpha]$ 及其补集都开, 可用于分离任意两个不同序数, 所以空间不连通, 且任意连通子集至多一个点; 因道路像连通, 它也不道路连通.

开覆盖 $\{[0,\alpha):\alpha<\omega_1\}$ 没有可数子覆盖: 可数个可数序数的上确界仍小于 $\omega_1$. 同理每个可数点集有界而不稠密. 所以不紧、非 Lindel\"of、不可分且非第二可数.

初始区间 $[0,\alpha]$ 闭开, 其差集 $(\beta,\alpha]=[0,\alpha]\setminus[0,\beta]$, 连同含最小元的初始区间, 形成闭开邻域基, 从而得到正则性. 若两闭集 $A,B$ 都无界, 可交替选择严格递增的 $a_n\in A,b_n\in B$, 其可数上确界 $\lambda<\omega_1$ 同时是两条子序列的极限; 闭性给出 $\lambda\in A\cap B$. 因此不交闭集至少有一个有界. 有界闭集位于某个紧初始区间中, 因而紧. 正则性分离这个紧集的每一点与另一个闭集, 取有限个邻域并, 即将两闭集分离. 故正规并完全正则.

最后, $S_\Omega$ 可数紧: 任一可数开覆盖若无有限子覆盖, 可选 $x_n$ 避开前 $n$ 个成员; 所有 $x_n$ 落在某个紧 $[0,\beta]$ 中, 此区间上的有限子覆盖与选择矛盾. 可数紧的度量空间紧: 每个序列都有聚点, 否则取一个无聚点的无限值域, 其补集及分别只含一个值域点的开邻域组成没有有限子覆盖的可数开覆盖; 第一可数性再给收敛子序列, 用第\ref{ch:compact}章的紧性等价. 因本空间不紧, 它不可能可度量化.

\item \textbf{有序方形.}
字典序是稠密序: 同一竖直纤维内可取第二坐标的中间值; 不同纤维之间可取第一坐标的中间值. 它还有最小元 $(0,0)$、最大元 $(1,1)$, 每个非空子集有上确界. 具体地, 先取第一坐标的上确界 $a$; 若集合在纤维 $a$ 中有点, 再取那些点第二坐标的上确界, 否则取 $(a,0)$.

完备且有两端点的序拓扑紧. 对开覆盖令 $T$ 为所有使初始区间至该点有有限子覆盖的点; 覆盖最小元的成员保证 $T$ 非空. 取 $c=\sup T$. 覆盖 $c$ 的开邻域含其左侧一段; 与 $T$ 中足够靠近 $c$ 的点的有限子覆盖合并, 即覆盖至 $c$, 所以 $c\in T$. 若 $c$ 不是最大元, 同一开邻域还覆盖右侧一段, 与 $c=\sup T$ 矛盾. 因而全空间有有限子覆盖. 序拓扑 Hausdorff, 故空间紧、局部紧 Hausdorff、正规、正则、Lindel\"of 和完全正则.

稠密完备序的区间连通: 若 $x<y$ 分别属于不交开集构成的分离两侧, 考察左侧在 $[x,y]$ 内的上确界 $c$. 若 $c$ 在左侧, 其开邻域和序稠密性给出更大的左侧点; 若在右侧, 其左邻域排除逼近 $c$ 的左侧点. 端点情形因 $x,y$ 分居两侧也排除. 所以有序方形连通.

集合 $\{x\}\times(0,1)$ ($x\in[0,1]$) 是不可数族两两不交非空开集, 可数稠密集不可能逐一与之相交. 故不可分、非第二可数且不可度量化, 最后一项也可由紧度量空间第二可数推出. 每点仍有可数邻域基: 纤维内部只用第二坐标的可数区间; 在 $(x,0)$ 或 $(x,1)$ 处, 从相邻纤维取第一坐标逐步逼近 $x$ 的序列, 再配合同一纤维内的第二坐标序列即可. 两个全局端点只取存在的那一侧.

若一道路的端点第一坐标不同, 其连通像必须包含两端点之间的整个序区间, 因而遇到不可数多个竖直开区间. 但道路像可分, 因为 $[0,1]$ 中有理点的像稠密, 与这些不交开集矛盾. 因此不道路连通. 在 $(1/2,0)$ 的任意开邻域内, 都能找到一个包含不可数多个完整纤维的闭序区间. 此区间紧而不可分; 若该邻域可度量化, 此区间将是紧度量空间并可分, 矛盾. 故空间也不局部可度量化.

\item \textbf{Sorgenfrey 平面.}
每个 $[a,b)$ 在 $\R_\ell$ 中既开又闭, 基本矩形也闭开. 因而平面 Hausdorff、正则, 且完全正则: 对点与闭集选择避开闭集的闭开矩形, 用其示性函数分离. 闭开矩形分离不同点, 所以任意连通子集至多一点, 不连通也不道路连通. 点 $(x,y)$ 的 $[x,x+1/n)\times[y,y+1/n)$ 给出第一可数性. $\mathbb Q^2$ 稠密.

反对角线 $A=\{(t,-t):t\in\R\}$ 是闭离散子空间. 在其点 $(t,-t)$, 一个下限基本矩形只与 $A$ 相交于该点. 对不在 $A$ 上的点, 按坐标和为正或负取足够小的矩形避开 $A$, 故 $A$ 闭; 其每个子集也在平面中闭. 如果平面正规, Urysohn 引理会为每个 $B\subseteq A$ 提供一个连续 $[0,1]$ 值函数, 在 $B$ 上为 $0$, 在 $A\setminus B$ 上为 $1$. 这些函数互不相同, 至少有 $2^{\mathfrak c}$ 个. 然而连续实值函数由其在可数稠密集上的值唯一决定, 最多只有 $\mathfrak c^{\aleph_0}=\mathfrak c$ 个, 与 Cantor 定理矛盾. 因而平面不正规. 正则 Lindel\"of 空间正规, 故它非 Lindel\"of; 第二可数蕴含 Lindel\"of, 故也非第二可数. 由非正规还知不紧和不可度量化.

它不局部紧. 若某点有紧邻域 $K$, 则第一投影 $C=\pi_1(K)$ 是下限直线的紧集, 并包含一个非空通常开区间. 恒等映射 $C\to C$ 的通常实数子空间拓扑是连续双射, 紧到 Hausdorff 的判据使它为同胚. 但在该区间内部取 $c<d$, 下限开集 $[c,d)\cap C$ 不是通常相对开集, 矛盾.

它不局部可度量化. 任意非空基本矩形都包含一条相对闭的水平线段, 同胚于下限区间 $[a,b)$. 这个区间可分却非第二可数: 对每个 $t\in[a,b)$, 在一个假定的基中选择 $t\in B_t\subseteq[t,b)$, 不同 $t$ 所选 $B_t$ 不同. 由度量空间的可数性等价, 这个区间不能可度量化; 所以任何包含矩形的开邻域也不能可度量化.

\item \textbf{$\R^\omega$ 的一致拓扑.}
$\rho$ 是度量, 所以 Hausdorff、正规、正则、完全正则、第一可数且局部可度量化. 对每个 $x$, 集合 $x+\ell^\infty$ 既开又闭: 小于 $1$ 的一致球中坐标差一致有界, 不同的有界差等价类互不相交并构成开分划. 例如零序列与 $(n)_n$ 不在同一类. 因而空间不连通, 不道路连通. 同一类内的线性道路确实连续, 因为其坐标差有统一有限界.

$\{0,1\}^{\mathbb N}$ 中不同点距离为 $1$, 是不可数分离集. 半径 $1/3$ 的球两两不交, 所以没有可数稠密集; 由度量可数性等价也非第二可数且非 Lindel\"of. 它不局部紧: 任何点的任何紧邻域都包含某个球 $B_\rho(x,r)$, 取 $0<t<\min\{r,1\}$, 则其中的 $x+t e_n$ 两两距离为 $t$, 没有收敛子序列, 与紧度量空间的序列紧性矛盾. 因而也不紧.

\item \textbf{$\R^I$ 的积拓扑.}
任意两点的逐坐标线性道路 $\gamma(t)_i=(1-t)x_i+ty_i$ 连续, 因为每个坐标连续; 所以道路连通且连通. 完全正则性由前述乘积函数证明, 也得到正则和 Hausdorff. 若 $I=\varnothing$, 积是一个点, 全部十四项都成立. 若 $I\ne\varnothing$, 投影到一个 $\R$ 因子连续且满射, 所以不紧.

有限 $I$ 的积是欧氏空间, 局部紧 Hausdorff. 无限 $I$ 时, 每个基本邻域都留下一个不受限制的坐标. 若有紧邻域 $K$, 它包含一个基本邻域, 在这个坐标的投影就是 $\R$, 而紧集投影不能为 $\R$, 矛盾. 因而局部紧性恰好对应有限 $I$.

至多可数 $I$ 时, 命题\ref{thm:countable-product-metric}给出可度量化和第二可数, 于是第一可数、正规、Lindel\"of 且局部可度量化. 不可数 $I$ 时, 任一点都不第一可数: 假设有可数邻域基, 将其各缩成一个基本邻域, 全部受约束坐标的并仍可数; 取未受约束的 $j$, 这些邻域不能包含于 $\{y:|y_j-x_j|<1\}$. 所以非第二可数、不可度量化且不局部可度量化. 后一项使用局部可度量化必推出第一可数.

可分性的准确界是 $|I|\le\mathfrak c$. 必要性: 对可数稠密集 $D$, 各坐标在 $D$ 上的取值向量必须两两不同, 否则 $D$ 落在两个坐标相等的闭集内而不稠密. 因此 $I$ 注入 $\R^D$, 后者至多有 $\mathfrak c$ 个元素. 充分性: 为各 $i$ 选择不同的实数 $t_i$, 取可数集 $\{(p(t_i))_{i\in I}:p\in\mathbb Q[t]\}$. 给定有限坐标的开区间条件, 用 Lagrange 插值使一个实系数多项式在那些不同的 $t_i$ 处取指定的区间内部值, 再充分逼近其有限个系数为有理数, 仍满足全部条件. 故这个可数集稠密.

不可数积的非正规性需要真正的反例, 不能仅从不可度量化推出. 下面完整给出 Stone 反例的有限坐标证明. 在 $X=\mathbb N^{\omega_1}$ 中取 $\mathbb N=\{0,1,2,\ldots\}$ 为离散空间, 记 $[F,x]=\{y:y|_F=x|_F\}$, 其中 $F$ 有限. 令 $H_j$ ($j=0,1$) 由那些除数值 $j$ 外每个数值至多出现一次的函数组成. 若一函数不在 $H_j$, 用两个坐标同取某个 $m\ne j$ 的柱集就能见证, 所以 $H_j$ 闭. 若函数同时属于 $H_0,H_1$, 每个数值都至多出现一次, 成为 $\omega_1\to\mathbb N$ 的单射, 不可能. 因此两闭集不交.

任取开集 $U\supseteq H_0$, $V\supseteq H_1$. 递归构造有限集 $F_n$ 和相容单射 $q_n:F_n\to\{2,3,\ldots\}$. 从 $F_0=\varnothing$ 开始, 令 $f_n$ 在 $F_n$ 上等于 $q_n$, 外面为 $0$. 因 $f_n\in H_0$, 可取有限 $F_{n+1}\supseteq F_n$ 使 $[F_{n+1},f_n]\subseteq U$, 再把 $q_n$ 扩成单射 $q_{n+1}$. 令 $F=\bigcup_nF_n$, $q=\bigcup_nq_n$, 并令 $g$ 在 $F$ 上等于 $q$, 外面为 $1$; 则 $g\in H_1$. 取有限 $G$ 使 $[G,g]\subseteq V$, 再取 $n$ 使 $G\cap F\subseteq F_n$. 在 $G\cap F_{n+1}$ 上, $g$ 与 $f_n$ 一致, 因为这个交集落在 $F_n$ 内. 两组有限坐标条件相容, 所以 $[G,g]\cap[F_{n+1},f_n]\ne\varnothing$. 得 $U\cap V\ne\varnothing$, 证明 $X$ 非正规.

对不可数 $I$, 在 ZFC 中选择 $J\subseteq I$, $|J|=\aleph_1$. 闭子空间 $\mathbb N^J\times\{0\}^{I\setminus J}$ 同胚于上述 $X$; 正规性传给闭子空间, 所以 $\R^I$ 非正规. 它已是正则空间, 正则 Lindel\"of 空间正规, 故它也非 Lindel\"of. 这同时完成两格的不可数情形, 无须连续统假设. Stone 反例的来源见本章末尾; 上面的有限坐标论证已给出所需全部步骤.

\item \textbf{任意度量空间.}
度量空间 Hausdorff、正规和正则, 由距离函数或 Urysohn 引理完全正则; 球 $B(x,1/n)$ 给第一可数性. 定义本身给可度量化及局部可度量化. 两点离散空间反驳普遍的连通性与道路连通性; $\R$ 反驳普遍紧性; 上一行的一致度量空间反驳普遍局部紧性. 不可数离散度量空间不第二可数、不可分且非 Lindel\"of: 稠密集必须是全集, 单点开覆盖没有可数子覆盖.

\item \textbf{取可数积.}
先说明有限积连通: 两因子积中, 固定 $(a,b)$, 各集合 $(X\times\{b\})\cup(\{x\}\times Y)$ 是两连通切片的并, 共同含 $(a,b)$; 对 $x$ 取并得到全积, 再归纳到有限积. 可数积的连通性由有限坐标改变点的集合证明. 固定积点 $a$, 只改变有限坐标的各有限子积都连通, 共同含 $a$, 故其并连通; 该并在积中稠密. 连通集的闭包连通, 因为闭包若有分离, 两个非空相对开侧都会遇到原稠密连通集并将其分离. 所以全积连通. 因子或积为空时结论直接成立. 道路连通性由逐坐标道路及积的连续性判据得到. 紧性由 Tychonoff 定理\ref{thm:tychonoff}得到; Hausdorff 性按不同坐标分离. 正则性则在有限受约束坐标逐一取闭包包含于给定坐标开集的小邻域, 其基本矩形闭包仍包含在原矩形中. 第一可数的局部基由有限坐标的可数局部基组合而成; 第二可数和可度量化见命题\ref{thm:countable-product-metric}. 可分性可取各因子的可数稠密集, 固定其各一点, 只在有限坐标改取稠密集中的点; 所得点集可数且稠密. 完全正则性已由有限个分离函数的乘积证明.

局部紧 Hausdorff 性不保持: $\R^{\mathbb N}$ 即为反例. 正规性和 Lindel\"of 性不保持, 均可用 $\R_\ell\times\R_\ell$; 若要求写成可数无限积, 补上单点因子即可. 为确认因子性质, 这里证明 $\R_\ell$ 是 Lindel\"of. 给定开覆盖, 为每个 $x$ 取 $[x,b_x)$ 包含于某个覆盖成员. 通常开集 $O=\bigcup_x(x,b_x)$ 可从这些通常开区间中选可数子覆盖, 因通常实线第二可数. 对 $x\in\R\setminus O$, 区间 $(x,b_x)$ 两两不交: 若 $x<y$ 均在此补集, 必有 $b_x\le y$. 各区间选一个有理数, 知补集可数. 用覆盖 $O$ 的可数个区间对应的原覆盖成员, 再加补集各点对应的成员, 得到下限直线的可数子覆盖. 下限直线正则, 因而正规.

局部可度量化也不保持: $S_\Omega^{\mathbb N}$ 的每个因子局部可度量化, 但任何非空基本邻域都留下一个完整 $S_\Omega$ 因子. 固定其他坐标, 它含一个同胚于 $S_\Omega$ 的子空间, 因而不能可度量化. 若积中某点有可度量开邻域, 将其缩成基本邻域就得到矛盾.

\item \textbf{取闭子空间.}
$\{0,1\}\subseteq\R$ 是闭子空间, 但不连通也不道路连通. 紧性由闭子集紧得到; 正规性因为闭子空间中的闭集仍在母空间中闭而保持. Lindel\"of 性由给相对开覆盖添上闭子空间的补集得到母空间开覆盖, 再取可数子覆盖. Hausdorff、正则、第一可数、第二可数、可度量化、局部可度量化和完全正则均传给任意子空间: 限制分离邻域、邻域基、基、度量或分离函数即可; 正则性的闭包收缩论证见第\ref{ch:separation}章练习. 闭子空间的局部紧 Hausdorff 性, 用原紧邻域与闭子空间相交得到.

可分性不保持: Sorgenfrey 平面可分, 但其闭离散反对角线 $A$ 不可数, 不是可分空间. 此反例说明不能简单把母空间的可数稠密集与闭子空间取交.

\item \textbf{取开子空间.}
$\R\setminus\{0\}$ 不连通也不道路连通; $(0,1)\subseteq[0,1]$ 说明紧性不保持. 第一可数、第二可数、Hausdorff、正则、可度量化、局部可度量化及完全正则按上一项传给所有子空间. 可分性在开子空间保持: 若 $D$ 可数稠密且 $U$ 开, 每个非空相对开集也是母空间的非空开集, 因而与 $D\cap U$ 相交. 局部紧 Hausdorff 也保持: 给定 $x\in U$, 用局部紧 Hausdorff 的邻域收缩引理\ref{lem:local-compact-shrink}取得闭包紧且包含于 $U$ 的小开邻域.

Lindel\"of 性不保持: $[0,\omega_1)$ 是紧空间 $[0,\omega_1]$ 的开子空间, 已知非 Lindel\"of. 正规性不保持的反例是删除角点的 Tychonoff 板. 令 $P=[0,\omega_1]\times[0,\omega]$, 两因子取序拓扑. 它们紧 Hausdorff, 所以 $P$ 紧 Hausdorff 并正规; $Y=P\setminus\{(\omega_1,\omega)\}$ 开. 在 $Y$ 中, 集合 $A=\{\omega_1\}\times[0,\omega)$ 和 $B=[0,\omega_1)\times\{\omega\}$ 不交且闭, 因其在 $P$ 中的闭包只额外添入已删除的角点. 若开集 $U\subseteq Y$ 包含 $A$, 对每个有限序数 $n$ 可取 $\alpha_n<\omega_1$ 使 $(\alpha_n,\omega_1]\times\{n\}\subseteq U$. 令 $\alpha=\sup_n\alpha_n<\omega_1$, 取 $\beta=\alpha+1$. 点 $(\beta,\omega)\in B$ 的每个邻域都含 $(\beta,n)$, 其中 $n$ 足够大; 而这些点全在 $U$. 因此任何包含 $B$ 的开集都与 $U$ 相交, $Y$ 不正规.
\end{enumerate}
\end{solution}

\originalsection{官方作业目录与题源}
\label{sec:official-topology-directory}
官方作业页将 weekly exercises 与 problem sets 区分开: weekly exercises 不交评分, 用于备考, 鼓励合作; PS0 是集合论语言的诊断作业, 其分数只供选课建议; PS1--4 要求独立提交规范书面解答; PS5 原件明确标为 optional. 下列编号均指 James R. Munkres, \emph{Topology}, 2nd ed., Prentice-Hall, 1999 的节号和习题号. 官方页对 PS0--4 及 weekly exercises 仅给教材题号, \textbf{缺少可公开下载的完整题面}; 本章只列目录, 不据题号重构或转录版权教材题面. PS5 则有完整两页原件, 已在习题\ref{ex:official-ps5}翻译并解答.

\begin{description}
\item[PS0.] \S1: 2, 5; \S2: 4(c), 4(e). 均只给答案. \S1 第2题的 (j), (k), (l) 答“是”或“否”, 其余答所列逻辑或包含关系符号. 原件未在作业页重刊这道教材题的各子问.
\item[PS1.] \S3: 13; \S9: 8; \S13: 7; \S17: 16, 18. \S13 第7题及 \S17 第16、18题只给答案.
\item[PS2.] \S18: 13; \S20: 6(a)(b), 8(b)(c); \S24: 4. \S20 第8题假定 (a), 其中 (c) 只给答案.
\item[PS3.] \S26: 9, 12; \S28: 4; \S30: 5.
\item[PS4.] \S31: 7(a)(b)(d); \S33: 4; \S34: 4, 5; \S38: 9. \S34 第4、5题给证明或反例.
\item[PS5.] 官方完整 PDF, 九行十四列, 选做.
\end{description}

\begin{description}
\item[第1周: 逻辑与基础.] \S1: 3; \S2: 1, 2, 4, 5.
\item[第2周: 关系、基数、选择公理.] \S3: 11, 12, 15; \S5: 4, 5; \S6: 3, 6; \S7: 3, 4, 5.
\item[第3周: 拓扑、闭集.] \S7: 6; \S13: 2, 6, 8; \S16: 3, 4, 8, 10; \S17: 3, 4, 5, 6, 8, 9.
\item[第4周: 连续函数、任意积.] \S17: 10, 11, 12; \S18: 2, 3, 7, 8; \S19: 1, 2, 3, 6, 8.
\item[第5周: 度量拓扑.] \S19: 7; \S20: 2, 4, 5, 8(a); \S21: 3, 4, 6, 7.
\item[第6周: 商拓扑.] \S22: 2, 3, 6.
\item[第7周: 连通空间、紧空间.] \S23: 2, 3, 5, 7, 8; \S24: 1, 5, 8, 9; \S25: 1, 2, 3; \S26: 1, 4, 5, 6.
\item[第8周: 紧性续讲.] \S26: 7, 8; \S27: 1, 4; \S28: 1, 2, 3.
\item[第9周: 良序集、极大原理; 期中考试.] 官方页未列习题号.
\item[第10周: 可数性与分离公理.] \S10: 2, 3, 6; \S29: 5, 6, 7, 8; \S30: \texttt{1, 2, 4 7, 8, 9}. 官方页面在 $4$ 和 $7$ 之间未印分隔符, 此处照录; 不将这一处排印歧义视为已核实的独立题号分隔.
\item[第11周: Urysohn 引理、可度量化.] \S31: 1, 3, 6; \S32: 1, 2, 3, 6, 7.
\item[第12周: Tietze 定理.] \S33: 2, 6, 7, 8.
\item[第13周: Tychonoff 定理、Stone--\v{C}ech 紧化.] \S34: 1, 3, 7, 8; \S37: 2, 3; \S38: 5, 6.
\item[第14周: Baire 空间、维数理论.] \S38: 3, 4; \S36: 1, 5; \S48: 1, 2, 3, 5, 6. 保留官方页面的节号顺序.
\item[第15周: 嵌入欧氏空间.] 官方页未列习题号. 随后的期末考试行也未列习题号.
\end{description}

官方目录与 PS5 下载链接见
\href{https://ocw.mit.edu/courses/18-901-introduction-to-topology-fall-2004/pages/assignments/}{MIT OCW: 18.901 Fall 2004 Assignments}.
关于不可数积的非正规性, 原始出处是 A. H. Stone, ``Paracompactness and product spaces'', \emph{Bulletin of the American Mathematical Society} \textbf{54} (1948), 977--982. 闭集 $H_0,H_1$ 的定义及这一反例的不同证明也见 N. Noble, \href{https://arxiv.org/abs/2002.02483}{\emph{Poorly Separated Infinite Normal Products}}, \S6.1. 本章已展开有限坐标证明, 并未将这一结论视为前十五章已经证明的教材结论.

\begin{thebibliography}{99}
\bibitem{bright2023}
Paige Bright. \textit{18.S190 Introduction to Metric Spaces}. MIT OpenCourseWare, January IAP 2023. 六讲Lecture Notes及官方TeX, 另附三份Problem Sets. 主讲义PDF物理页数依次为8、8、8、6、7、6.\\
课程: \url{https://ocw.mit.edu/courses/18-s190-introduction-to-metric-spaces-january-iap-2023/}.\\
讲义: \url{https://ocw.mit.edu/courses/18-s190-introduction-to-metric-spaces-january-iap-2023/pages/lecture-notes/}.\\
三份原题: \href{https://ocw.mit.edu/courses/18-s190-introduction-to-metric-spaces-january-iap-2023/pages/problem-sets/}{Problem Sets}, 各3物理页, 包含2题面页与1OCW署名页; 本册第16章逐题译录与编者解答.

\bibitem{rodriguez2021}
Casey Rodriguez(授课), Andrew Lin(笔记). \textit{18.102/18.1021 Introduction to Functional Analysis}, Lecture 3: Quotient Spaces, the Baire Category Theorem and the Uniform Boundedness Theorem. MIT OpenCourseWare, Spring 2021. PDF第3--4页, 正文印刷第16--17页: Baire定理及一致有界原理. 本册重新给出证明, 修正原证明中相对开球与全空间开球的区分.\\
官方资源页: \href{https://ocw.mit.edu/courses/18-102-introduction-to-functional-analysis-spring-2021/resources/lecture-3-quotient-spaces-the-baire-category-theorem-and-the-uniform-boundedness-theorem/}{18.102第3讲}.

\bibitem{munkres2004}
James Munkres(授课). \textit{18.901 Introduction to Topology}, Supplementary Notes D: Countability Axioms. MIT OpenCourseWare, Fall 2004, 3页. 本册可数性讨论的选读参考; 不作为18.901整课或商业教材的全文翻译. Notes G第1--3页为商映射、邻接空间与正规性的参考; Notes C与K保留为选读归档, 未纳入新增核心.\\
\url{https://ocw.mit.edu/courses/18-901-introduction-to-topology-fall-2004/pages/lecture-notes/}.\\
作业: \href{https://ocw.mit.edu/courses/18-901-introduction-to-topology-fall-2004/pages/assignments/}{Assignments}; \href{https://ocw.mit.edu/courses/18-901-introduction-to-topology-fall-2004/d6c5d6d30bf26300b6f1b2c278791796_problemset_5.pdf}{PS5完整原件}, 2物理页(第二页印刷第8页). 本册第18章解答126格; PS0--4及周练只登记官方教材题号, 缺完整题面.

\bibitem{miller2016}
Haynes Miller(授课), Sanath Devalapurkar(课堂LaTeX记录), Xianglong Ni(插图). \textit{18.905 Algebraic Topology I}, MIT OpenCourseWare, Fall 2016. 第5讲物理页1--2: 同伦与形变收缩; 第14讲页1--3: 胞腔粘合; 第18讲页1: 欧拉示性数与同调交错和. 不把后续同调代数与上同调材料称为本册已覆盖.\\
\url{https://ocw.mit.edu/courses/18-905-algebraic-topology-i-fall-2016/pages/lecture-notes/}.

\bibitem{snowden2011}
Andrew Snowden. \textit{18.904 Seminar in Topology}, MIT OpenCourseWare, Spring 2011. Lecture Summaries: 基本群、van Kampen与覆盖的教学安排, 不是完整逐讲证明讲义; 本册相应证明为编者补充.\\
\url{https://ocw.mit.edu/courses/18-904-seminar-in-topology-spring-2011/pages/lecture-summaries/}.

\bibitem{seidel2023}
Paul Seidel. \textit{18.900 Geometry and Topology in the Plane}, MIT OpenCourseWare, Spring 2023. 第4讲印刷页30--33: 绕数; 第29--31讲印刷页219--222、225--228、232--235: 有限复形、边界矩阵与可定向性; 第33讲页246--249: 环路线性计数; 第40讲页292--293: 角亏与离散Gauss--Bonnet. 第40讲为官方公开补充, 当学期课堂未讲. 第30、31、33讲中的符号笔误已在正文修正, 原件保留.\\
\url{https://ocw.mit.edu/courses/18-900-geometry-and-topology-in-the-plane-spring-2023/}.\\
所选CQ原题: \href{https://ocw.mit.edu/courses/18-900-geometry-and-topology-in-the-plane-spring-2023/mit18_900s23_q4.pdf}{第4讲}, \href{https://ocw.mit.edu/courses/18-900-geometry-and-topology-in-the-plane-spring-2023/mit18_900s23_q29.pdf}{第29讲}, \href{https://ocw.mit.edu/courses/18-900-geometry-and-topology-in-the-plane-spring-2023/mit18_900s23_q30.pdf}{第30讲}, \href{https://ocw.mit.edu/courses/18-900-geometry-and-topology-in-the-plane-spring-2023/mit18_900s23_q31.pdf}{第31讲}, \href{https://ocw.mit.edu/courses/18-900-geometry-and-topology-in-the-plane-spring-2023/mit18_900s23_q33.pdf}{第33讲}, \href{https://ocw.mit.edu/courses/18-900-geometry-and-topology-in-the-plane-spring-2023/mit18_900s23_q40.pdf}{第40讲}, 各2物理页(1题面页与1署名页), 共22题. 第31讲第5题另需\href{https://ocw.mit.edu/courses/18-900-geometry-and-topology-in-the-plane-spring-2023/mit18_900s23_lec6.pdf}{第6讲}图(6.9)、(6.11), 原件7物理页, 图在物理第4页、印刷第51页.

\bibitem{mrowka2004}
Tomasz S. Mrowka. \textit{18.965 Geometry of Manifolds}, MIT OpenCourseWare, Fall 2004. 第1、2讲: 图册、光滑映射、基本病态; 第4讲印刷页9--12: 逆函数与隐函数论证, 本册限有限维并补齐条件.\\
\url{https://ocw.mit.edu/courses/18-965-geometry-of-manifolds-fall-2004/pages/lecture-notes/}.

\bibitem{seidel2008}
Paul Seidel. \textit{18.950 Differential Geometry}, MIT OpenCourseWare, Fall 2008. Chapter 4, Lectures 36、39: 长度与内在距离; 本册不译编该章其余测地线与曲率理论. 原件13个物理页, 包括来源页与章扉页.\\
\url{https://ocw.mit.edu/courses/18-950-differential-geometry-fall-2008/pages/lecture-notes/}.

\bibitem{tsinghua2027}
清华大学本科招生网. \textit{2027年丘成桐数学科学领军人才培养计划招生简章}. 2026年9月28日. 第六条第3项: 一试数学1包括分析、代数、几何与拓扑.\\
\url{https://www.admissions.tsinghua.edu.cn/info/1033/2240.htm}.

\bibitem{qzc-syllabus}
清华大学求真书院. \textit{数学领军计划之数学与物理科目考纲发布!}. 官方通知, 2026年8月21日. 本次核实其通知与详细考纲链接存在, 未获得所链原文; 具体考纲核实边界见前言和\texttt{research/syllabus-verification.md}.\\
\url{https://qzc.tsinghua.edu.cn/}. 既有仓库记录的资料题名为\textit{领军考试大纲(2).pdf}, 文件内日期2026年8月20日, 共6页; 本次没有将该摘要提升为重新逐页核验的证据.

\bibitem{cc-license}
Creative Commons. \textit{Attribution--NonCommercial--ShareAlike 4.0 International}.\\
许可: \url{https://creativecommons.org/licenses/by-nc-sa/4.0/}.\\
MIT OCW使用条款: \url{https://ocw.mit.edu/pages/privacy-and-terms-of-use/}.
\end{thebibliography}

\end{document}
